\documentclass[12pt,a4paper]{article}

\usepackage{mypreamble}
\usepackage{folland_sol}

\title{\textbf{\texttt{Real Analysis \\ Modern Techniques and Their Applications}} \\ \texttt{Solutions}}
\author{\texttt{Arkaraj Mukherjee}}
\date{\texttt{\today}}

\begin{document}
\maketitle
\newpage
\tableofcontents
\newpage
\stepcounter{section}
\section*{Chapter \thesection\ : Measures}
\addcontentsline{toc}{section}{Chapter \thesection\ : Measures}
\begin{problem}\label{prob:1.1}
A family of sets $\mathcal{R} \subset \mathcal{P}(X)$ is called a \textbf{ring} if it is closed under finite unions and differences (i.e., if $E_1, \dots, E_n \in \mathcal{R}$, then $\cup_{1}^{n} E_j \in \mathcal{R}$, and if $E, F \in \mathcal{R}$, then $E \setminus F \in \mathcal{R}$). A ring that is closed under countable unions is called a \textbf{$\sigma$-ring}.
\begin{enumerate}[label=\alph*.]
    \item Rings (resp. $\sigma$-rings) are closed under finite (resp. countable) intersections.
    \item If $\mathcal{R}$ is a ring (resp. $\sigma$-ring), then $\mathcal{R}$ is an algebra (resp. $\sigma$-algebra) iff $X \in \mathcal{R}$.
    \item If $\mathcal{R}$ is a $\sigma$-ring, then $\{E \subset X : E \in \mathcal{R} \text{ or } E^c \in \mathcal{R}\}$ is a $\sigma$-algebra.
    \item If $\mathcal{R}$ is a $\sigma$-ring, then $\{E \subset X : E \cap F \in \mathcal{R} \text{ for all } F \in \mathcal{R}\}$ is a $\sigma$-algebra.
\end{enumerate}
\end{problem}
\begin{solution}
\begin{enumerate}[label=\alph*.]
    \item If $A,B\in\mR$ then, $A\setminus B\in\mR$ and again, $A\setminus(A\setminus B)=A\cap B\in\mR.$ 
    We can induct on the number of sets and conclude that $\mR$ is closed under finite intersections. 
    If $\mR$ was a $\sigma$-ring and $\{A_n\}\subseteq\mR$ then $A_1\setminus A_1,A_1\setminus A_2,A_1\setminus A_3,\ldots\in\mR$ and $\cup (A_1\setminus A_n)=A_1\setminus(\cap A_n)\in\mR$ by definition. 
    Thus, $A_1\setminus(A_1\setminus \cap A_n)=\cap A_n\in\mR.$
    \item The existence of $X\in\mR$ implies that its closed under complements and we can conclude using part (a). 
    Now if it was a $\sigma$-algebra to begin with then, $X\in\mR$ follows from definition. 
    \item Let this set be $\mathcal S$.
    As, $\phi\in\mR\subseteq\mathcal S$ we must also have $X\in\mathcal S.$ 
    If $A$ is in this set then $A\in\mR$ or, $A^c\in\mR$, in the first case $X\setminus A=A^c\in\mR$ and in the second case $A^c\in\mR$ already, proving the closure of this set under complements. 
    Now let $\{E_n\}_{n=1}^{\infty} \subset \mathcal{S}$. 
    If $E_n \in \mathcal{R}$ for all $n$, then $\bigcup E_n \in \mathcal{R} \subset \mathcal{S}$ because $\mathcal{R}$ is a $\sigma$-ring. 
    If there exists at least one $k$ such that $E_k \notin \mathcal{R}$, then necessarily $E_k^c \in \mathcal{R}$. 
    We examine the complement of the union, 
    $$E^c_k\supseteq\left(\bigcup_{n=1}^\infty E_n\right)^c = \bigcap_{n=1}^\infty E_n^c = \bigcap_{n=1}^\infty (E_n^c \cap E_k^c)$$
    For any $n$, if $E_n \in \mathcal{R}$, then $E_n^c \cap E_k^c = E_k^c \setminus E_n \in \mathcal{R}$. If $E_n^c \in \mathcal{R}$, then $E_n^c \cap E_k^c \in \mathcal{R}$ by closure under intersections. 
    Thus, the sequence $\{E_n^c \cap E_k^c\}_n$ lies entirely in $\mathcal{R}$. Since $\sigma$-rings are closed under countable intersections by part (a), we conclude $(\bigcup E_n)^c \in \mathcal{R}$, which implies $\bigcup E_n \in \mathcal{S}$.
    \item Clearly, $\phi,X$ are included in this set. If $A$ is in this then $\forall F\in\mR,A\cap F\in\mR$ and by definition $F\setminus (F\cap A)=F\cap A^c\in\mR$ so $A^c$ is in it, thus its closed under complements. 
    If $A_1,A_2\ldots$ are in it then $\forall n\in\bN,\forall F\in\mR,A_n\cap F\in\mR$ and a $\mR$ is a $\sigma$-ring we must have $\cup(A_n\cap F)=F\cap(\cup A_n)\in\mR$ which is enough to show that its closed under coutnable unions, thus its a $\sigma$-algebra.
\end{enumerate}
\end{solution}
\begin{problem}\label{prob:1.2}
Complete the proof of Proposition 1.2.
\end{problem}
\begin{solution}
    Trivial.
\end{solution}
\begin{problem}\label{prob:1.3}
Let $\mathcal{M}$ be an infinite $\sigma$-algebra.
\begin{enumerate}[label=\alph*.]
    \item $\mathcal{M}$ contains an infinite sequence of disjoint sets.
    \item $\text{card}(\mathcal{M}) \ge \mathfrak{c}$.
\end{enumerate}
\end{problem}
\begin{solution}
    Let $\{E_n\}_{n\in\bN}$ be a countably infinite collection of distinct sets in $\mM$ which we can choose by the virtue of it being infinite. Consider the atoms formed by this collection, 
    $$\left\{\bigcap_{n\in\bN}E_n^{f(n)}\Big\vert f:\bN\to\{1,c\}\right\}$$
    This is a subset of $\mM$ as $\mM$ is a sigma algebra. 
    Clearly for any two different $f,g$ the intersections they produce are disjoint. Also, every set in our collection is an union of some sets from these atoms,
    $$E_n=\bigcup_{\substack{f\in\{1,c\}^\bN\\ f(n)=1}}\bigcap_{n\in\bN}E_n^{f(n)}$$
    So if there are finitely many nonempty atoms, say $n$ then its easy to show that there are atmose $2^n$ sets to begin with which neccesitates that there be infinite atoms. 
    Thus there exists some countably infinite sequence of pairwise disjoint nonempty sets in $\mM$(choose some atoms) proving part (a). 
    Let $\{F_n\}_{n\in\bN}$ be the collection of some pairwise disjoint sets in $\mM$. 
    Identifying $A^0$ with $\phi,$ consider the collection
    $$\mathcal G:=\left\{\bigcup_{n\in\bN}F_n^{f(n)}\Big\vert f:\bN\to\{0,1\}\right\}$$
    Clearly $\mathcal G\subseteq\mM$ and all the sets in this are distinct as $F_n$ are all pairwise disjoint. 
    Also this is in bijection to $\{0,1\}^\bN.$ 
    Thus, $|\mM|\geq|\mathcal G|=|\{0,1\}^\bN|=\mathfrak c.$
\end{solution}
\begin{problem}\label{prob:1.4}
An algebra $\mathcal{A}$ is a $\sigma$-algebra iff $\mathcal{A}$ is closed under countable increasing unions (i.e., if $\{E_j\}_{1}^{\infty} \subset \mathcal{A}$ and $E_1 \subset E_2 \subset \cdots$, then $\cup_{1}^{\infty} E_j \in \mathcal{A}$).
\end{problem}
\begin{solution}
    If its a $\sigma$-algebra then its tirvially an algebra closed under all sorts of intersections. 
    If its an algebra closed under increasing countable unions then for $A_1,\ldots\in\mathcal A$ consider $B_n:=\cup_{m\leq n}A_m\in\mathcal A$. 
    Clearly $B_n$ is increasing and thus, $\cup B_n=\cup A_n\in\mathcal A.$
\end{solution}
\begin{problem}\label{prob:1.5}
If $\mathcal{M}$ is the $\sigma$-algebra generated by $\mathcal{E}$, then $\mathcal{M}$ is the union of the $\sigma$-algebras generated by $\mathcal{F}$ as $\mathcal{F}$ ranges over all countable subsets of $\mathcal{E}$. (Hint: Show that the latter object is a $\sigma$-algebra.)
\end{problem}
\begin{solution}
    We wish to prove that,
    $$\mM(\mE)=\bigcup_{\substack{A\subseteq\mE\\ |A|\leq|\bN|}}\mM(A)$$
    Izt suffices to show that the set on the right is a $\sigma$-algebra as $\mE$ is already a subset and we can conclude using the definition of generated $\sigma$-algebra. 
    The set is clearly closed under complements and also $\phi,X$ are in it. 
    Now, say $\{A_n\}$ is subset, then we can find $\mE_n\subseteq\mE$ such that $A_n\in\mM(\mE_n)$. 
    Then, $|\cup\mE_n|\leq|\bN|$ and
    $$\bigcup_{n\in\bN}A_n\in\mM\left(\bigcup_{n\in\bN}\mM(\mE_n)\right)\subseteq\mM\left(\bigcup_{n\in\bN}\mE_n\right)\subseteq\bigcup_{\substack{A\subseteq\mE\\ |A|\leq|\bN|}}\mM(A)$$
    proving that its a sigma algebra.
\end{solution}
\begin{problem}\label{prob:1.6}
Complete the proof of Theorem 1.9.
\end{problem}
\begin{solution}
    First we show that $\overline{\mM}$ is a sigma algebra. It clearly has $\phi,X$ and is closed under countable unions, all thats left to show is that its closed under complements. Now, if $E\cup F\in\overline{\mM}$ with $E\in\mM$ and $F\subseteq N\in mN$, write it as $E\cup F=E\cup(E\setminus F)\in\overline{\mM}$ with $E\subseteq\mM$ and $F\setminus E\subseteq N\setminus E\in\mN$. 
    As, $E$ and $N\setminus E$ are disjoint we see that it suffices to assume $E\cap N=\phi$ from the start and we do so. 
    Write $E\cup F=(E\cup N)\cap(F\cup N^c)=(E\cup N)\cap(N\setminus F)^c$ and $(E\cup F)^c=(E\cup N)^c\cup(N\setminus F)$ with $E\cup N\in\mM$ and $N\setminus F\subseteq N\in\mN$ so $(E\cup F)\in\overline{\mM}$ proving that its a $\sigma$-algebra. 
    Set $\overline{\mu}(E\cup F):=\mu(E)$ when $E\in\mM,F\subseteq N\in\mN$, this is well defined because if $E_1\cup F_1=E_2\cup F_2$ with $E_1,E_2\in\mM$ and $F_1,F_2\subseteq N_1,N_2$ resp. in $\mN$. 
    Then, $\mu(E_1)=\overline{\mu}(E_1\cup F_1)=\overline{\mu}(E_1\cup N_1)=\mu(E_1\cup N_1)\geq\mu(E_2)$ as $E_2\subseteq E_1\cup F_1\subseteq E_1\cup N_1$. 
    Similarly we get that $\mu(E_1)\leq\mu(E_2)$, thus $\overline{\mu}(E_2\cup E_2)=\mu(E_2)=\mu(E_1)=\overline{\mu}(E_1\cup F_1)$. 
    $\overline{\mu}$ is measure as trivially, $\overline{\mu}(\phi)=0$ and for any disjoint countable collection of sets $\{E_n\cup F_n\}_{n\in\bN}\subseteq\overline{\mM}$ with $E_n\in\mM$ and $F_n\subseteq N_n\in\mN$ we have,
    $$\overline{\mu}\left(\bigcup_{n\in N}(E_n\cup F_n)\right)=\overline{\mu}\left(\left(\bigcup_{n\in\bN}E_n\right)\cup\left(\bigcup_{n\in\bN}F_n\right)\right)=\mu\left(\bigcup_{n\in N}E_n\right)=\sum_{n=1}^\infty\mu(E_n)=\sum_{n=1}^\infty \overline{\mu}(E_n\cup F_n)$$
    where $\cup E_n\in\mM$ and $\cup F_n\subseteq\cup N_n\in\mN$ by sigma additivity of $\mu.$
    Now we prove uniqueness, suppose $\mu_*$ is another such measure which extends $\mu$. 
    Then, by definition $\mu_*|_{\mM}=\mu$ so, for $E\cup F\in\overline{\mM}$ as before, 
    $$\overline{\mu}(E\cup F)=\mu(E)=\mu_*(E)\leq\mu_*(E\cup F)\leq\mu_*(E\cup N)=\mu(E\cup N)=\mu(E)+0=\overline{\mu}(E\cup F)$$
    by monotonicity, which proves that $\mu_*=\overline{\mu}$ and we are done.
\end{solution}
\begin{problem}\label{prob:1.7}
If $\mu_1, \dots, \mu_n$ are measures on $(X, \mathcal{M})$ and $a_1, \dots, a_n \in [0, \infty)$, then $\sum_{1}^{n} a_j \mu_j$ is a measure on $(X, \mathcal{M})$.
\end{problem}
\begin{solution}
    $\left(\sum_i a_i\mu_i\right)(\phi)=\sum_i a_i\mu_i(\phi)=\sum_i a_i\cdot 0=0$ and for any disjoint countable collection of sets $\{E_n\}$ we find that, as $a_i\cdot\mu_i(E_n)\geq 0$ for all valid $i,n$ we can interchange the order of summation as we are working in the extended reals and deduce,
    $$\left(\sum_i a_i\mu_i\right)\left(\bigcup_{n\in\bN}E_n\right)=\sum_i a_i\mu_i\left(\bigcup_{n\in\bN}E_n\right)=\sum_i \sum_{n\in\bN} a_i\mu_i(E_n)$$
    $$=\sum_{n\in\bN}\sum_i a_i\mu_i(E_n)=\sum_{n\in\bN}\left(\sum_i a_i\mu_i\right)(E_n)$$
    proving that $\sum_i a_i\mu_i$ is a measure.
\end{solution}
\begin{problem}\label{prob:1.8}
If $(X, \mathcal{M}, \mu)$ is a measure space and $\{E_j\}_{1}^{\infty} \subset \mathcal{M}$, then $\mu(\liminf E_j) \le \liminf \mu(E_j)$. Also, $\mu(\limsup E_j) \ge \limsup \mu(E_j)$ provided that $\mu(\cup_{1}^{\infty} E_j) < \infty$.
\end{problem}
\begin{solution}
    $$\mu(\liminf E_j)=\mu\left(\bigcup_{k\geq 1}\bigcap_{n\geq k}E_n\right)=\lim_{k\to\infty}\mu\left(\bigcap_{n\geq k}E_n\right)$$
    as the sequence of sets $\{\cap_{n\geq k}E_n\}_{k\in\bN}$ is increasing. Now, $\mu(E_n)\geq\mu(\cap_{n\geq k}E_n)$ for all $k\geq n$ by monotonicity and using the definition of infimums, $\inf_{n\geq k}\mu(E_n)\geq\mu(\cap_{n\geq k}E_n)$ and taking the limits on both sides as $k$ tends to infinity we have,
    $$\liminf\mu(E_j)=\lim_{k\to\infty}\inf_{n\geq k}\mu(E_n)\geq\lim_{k\to\infty}\mu\left(\bigcap_{n\geq k}E_n\right)=\mu\left(\liminf E_j\right)$$
    Now assume that $\mu(\cup E_n)<\infty$, then similarly
    $$\mu(\limsup E_j)=\mu\left(\bigcap_{k\geq 1}\bigcup_{n\geq k}E_n\right)=\lim_{k\to\infty}\mu\left(\bigcup_{n\geq k}E_n\right)$$
    as the sequence of sets $\{\cup_{n\geq k}E_n\}_{k\in\bN}$ is decreasing and the set at index $1$ has finite measure. 
    Now, clearly $\mu(\cup_{n\geq k}E_n)\geq\sup_{n\geq k}E_n$ from the definition of supremums and taking the limits on both sides as $k\to\infty$ we have,
    $$\mu(\limsup E_j)=\lim_{k\to\infty}\mu\left(\bigcup_{n\geq k}E_n\right)\geq\lim_{k\to\infty}\sup_{n\geq k}\mu(E_n)=\limsup\mu(E_j)$$
\end{solution}
\begin{problem}\label{prob:1.9}
If $(X, \mathcal{M}, \mu)$ is a measure space and $E, F \in \mathcal{M}$, then $\mu(E) + \mu(F) = \mu(E \cup F) + \mu(E \cap F)$.
\end{problem}
\begin{solution}
    If atleast one of $\mu(E),\mu(F)$ is infinite then the equality holds trivially. 
    So assume that both of the sets $E,F$ have finit measure, in which case both $\mu(E\cap F)\leq\mu(E)$ and $\mu(E\cup F)\leq\mu(E)+\mu(F)$ also have finite measure. 
    As $(E\cup F)\setminus F=E\setminus F=E\setminus(E\cap F)$, using monotonicity we can write
    $$\mu(E\cup F)-\mu(F)=\mu((E\cup F)\setminus F)=\mu(E\setminus(E\cap F))=\mu(E)-\mu(E\cap F)$$
    rearranging(all the terms are finite so we can do this) which gives us the required equality.
\end{solution}
\begin{problem}\label{prob:1.10}
Given a measure space $(X, \mathcal{M}, \mu)$ and $E \in \mathcal{M}$, define $\mu_E(A) = \mu(A \cap E)$ for $A \in \mathcal{M}$. Then $\mu_E$ is a measure.
\end{problem}
\begin{solution}
    Clearly $\mu_E(\phi)=\mu(\phi\cap E)=\mu(\phi)=0$ and if $A_1,A_2,\ldots$ are disjoint then,
    $$\mu_E\left(\bigcup_{n\in\bN}E_n\right)=\mu\left(A\cap\bigcap_{n\in\bN}E_n\right)=\mu\left(\bigcap_{n\in\bN}(A\cap E_n)\right)=\sum_{n=1}^{\infty}\mu(A\cap E_n)=\sum_{n=1}^\infty\mu_E(A_n)$$
    as $A_n\cap E$ are also disjoint, hence proving that $\mu_E$ is a measure.
\end{solution}
\begin{problem}\label{prob:1.11}
A finitely additive measure $\mu$ is a measure iff it is continuous from below as in Theorem 1.8c. If $\mu(X) < \infty$, $\mu$ is a measure iff it is continuous from above as in Theorem 1.8d.
\end{problem}
\begin{solution}
    This is trivial.
\end{solution}
\begin{problem}\label{prob:1.12}
Let $(X, \mathcal{M}, \mu)$ be a finite measure space.
\begin{enumerate}[label=\alph*.]
    \item If $E, F \in \mathcal{M}$ and $\mu(E \triangle F) = 0$, then $\mu(E) = \mu(F)$.
    \item Say that $E \sim F$ if $\mu(E \triangle F) = 0$; then $\sim$ is an equivalence relation on $\mathcal{M}$.
    \item For $E, F \in \mathcal{M}$, define $\rho(E, F) = \mu(E \triangle F)$. Then $\rho(E, G) \le \rho(E, F) + \rho(F, G)$, and hence $\rho$ defines a metric on the space $\mathcal{M} / \sim$ of equivalence classes.
\end{enumerate}
\end{problem}
\begin{solution}
    \begin{enumerate}
        \item We have $E\Delta F=(E\cup F)\setminus (E\cap F)$.
        If $\mu(E\cap F)=\infty$ then clearly $\mu(E)=\infty=\mu(F)$ and we are done.
        Otherwise using monotonicity and rearranging we can write, $\mu(E\cup F)=\mu(E\cap F)$.
        Again, using monotonicity we have, $\mu(E)\leq\mu(E\cup F)=\mu(E\cap F)\leq\mu(E)\implies\mu(E)=\mu(E\cup F)$.
        Similarly, $\mu(F)=\mu(E\cup F)$ and we have $\mu(E)=\mu(F).$
        \item This is clearly a equivalence relation as its reflexive ($\mu(A\Delta A)=\mu(\phi)=0$), symmetric as the symmetric difference operator is symmetric in its arguements, transitive as $A\sim B,B\sim C\implies\mu(A\Delta C)\leq\mu(A\Delta B)+\mu(B\Delta C)=0+0,$ this is proved in the next part.
        \item We have the following inclusion for $A,B,C\in\mM$
        $$(A\setminus B)\cup (B\setminus C)\supseteq A\setminus C$$
        Using this we get, $\mu(A\setminus B)+\mu(B\setminus C)\geq\mu(A\setminus C).$ 
        Similarly, $\mu(B\setminus A)+\mu(C\setminus B)\geq\mu(C\setminus A)$ and adding the two gives,
        $$(\mu(A\setminus B)+\mu(B\setminus A))+(\mu(B\setminus C)+\mu(C\setminus B))\geq (\mu(A\setminus C)+\mu(C\setminus A))$$
        And as $\mu(A\Delta B)=\mu(A\setminus B)+\mu(B\setminus A)$ we get the desired inequality. 
        Thus $\rho(E,F):=\mu(E\Delta F)$ is a metric on $\mM$. If $E'\in[E]\in\mM/\sim, F'\in[F]\in\mM/\sim$ then we wish to prove that $\rho(E,F)=\rho(E',F')$. 
        This is true as $\rho(E,F)\leq\rho(E,E')+\rho(E',F)=\rho(E',F)\leq\rho(E',F')+\rho(F',F)=\rho(E',F')$ and similarly $\rho(E',F')\leq\rho(E,F)$ as well, so we have equality. 
        We showed that the choice of the representative from an equivalence class of $\sim$ does not matter so we can define $\rho$ to be a metric on $\mM/\sim$ by letting $\rho([E],[F]):=\rho(E',F')$ where $E'\in[E],F'\in[F]$ are some representatives.
    \end{enumerate}
\end{solution}
\begin{problem}\label{prob:1.13}
Every $\sigma$-finite measure is semifinite.
\end{problem}
\begin{solution}
    Let $E_1,E_2,\ldots$ be such that $X=\cup E_n$ and $\mu(E_n)<\infty$, then for any $E\in\mM$ with $\mu(E)=\infty$ we see that, $\mu(E)=\mu(\cup(E\cap E_n))\leq\sum_{n=1}^\infty \mu(E_n\cap E)$. 
    As the sum is positive, not all terms can be zero, let $\mu(E\cap E_n)>0$, we also see that $\mu(E\cap E_n)\leq\mu(E_n)<\infty$. 
    This proves that $\mu$ is semifinite as we found $E\cap E_n\subseteq E$ with positive finite measure.
\end{solution}
\begin{problem}\label{prob:1.14}
If $\mu$ is a semifinite measure and $\mu(E) = \infty$, for any $C > 0$ there exists $F \subset E$ with $C < \mu(F) < \infty$.
\end{problem}
\begin{solution}
    For the sake of contradiction assume that $s=\sup\{\mu(F)|\mM\ni F\subseteq E,\mu(F)<\infty\}<\infty$. Then there exist $F_n\subseteq E$ with $s-1/n\leq\mu(F_n)\leq s$ and thus, $K:=\cup F_n$ has measure atleast $s$. There are two cases, either $\mu(\cup F_n)$ is finite or its infinite. If its finite, as $\mu(E)=\infty,$ we must have $\mu(E\setminus K)=\infty$ and as $\mu$ is semifinite there exists some $G\subseteq E\setminus K$ with finite positive measure, but this leads to a contradiction as $\mu(E\cup K)=\mu(E)+\mu(K)$ is finite and stricly larger than $s$, here we used the fact that, by construction $G\cap K=\phi.$
    In the other case, if $\mu(K)=\infty$ then consider the increasing sequence of sets $G_n:=\cup_{k\leq n}F_k$, we thus have that $\mu(\cup G_n)=\lim\mu(G_n)=\infty$ so there exists some $n\in\bN$ such that $\mu(G_n)>S$ and also by construction we find that $\mu(G_n)\leq n\cdot s<\infty$ contradicting the maximality of $s$. So $s$ must have been $\infty$ to begin with which implies that there subsets of $E$ of arbitrarily large finite measure.
\end{solution}
\begin{problem}\label{prob:1.15}
Given a measure $\mu$ on $(X, \mathcal{M})$, define $\mu_0$ on $\mathcal{M}$ by $\mu_0(E) = \sup\{\mu(F) : F \subset E \text{ and } \mu(F) < \infty\}$.
\begin{enumerate}[label=\alph*.]
    \item $\mu_0$ is a semifinite measure. It is called the \textbf{semifinite part} of $\mu$.
    \item If $\mu$ is semifinite, then $\mu = \mu_0$. (Use Exercise 14.)
    \item There is a measure $\nu$ on $\mathcal{M}$ (in general, not unique) which assumes only the values $0$ and $\infty$ such that $\mu = \mu_0 + \nu$.
\end{enumerate}
\end{problem}
\begin{solution}
    \begin{enumerate}[label=\alph*.]
        \item Clearly, $\mu_0(\phi)=0$ and for disjoint $\{E_n\}\subseteq\mM$, 
        $$\mu_0\left(\bigcup_{n\in\bN}E_n\right):=\sup\{\mu(E)<\infty:E\in\mM,E\subseteq\bigcup_{n\in\bN}E_n\}$$
        $$=\sup\left\{\sum_{n=1}^\infty\mu(E_n\cap E)<\infty:E\in\mM,E\subseteq\bigcup_{n\in\bN}E_n\right\}$$
        $$\leq\sum_{n=1}^\infty\sup\{\mu(F)<\infty:F\in\mM,F\subseteq E_n\}=\sum_{n=1}^\infty\mu_0(E_n)$$
        as, for all $n\in\bN,\mu(E\cap E_n)\leq\mu(E)<\infty$ and $\mu(E\cap E_n)\leq\mu_0(E_n)$.
        It suffices to prove the other direction of the equality as we have already established sub additivity. 
        There are two cases:
        \begin{enumerate}[label=\Roman*.]
            \item If $\sum_{n=1}^{\infty}\mu_0(E_n)<\infty$ then for arbitrary $\epsilon>0$, for all $n\in\bN$ let $\mM\ni F_n\subseteq E_n$ be such that $\mu_0(E_n)-\epsilon 2^{-n}\leq\mu(F_n)\leq\mu_0(E_n)$ as $\mu_0(E_n)<\infty$ is a supremum. 
            Summing both sides as $n\in\bN$ we get,
            $$\sum_{n=1}^\infty\mu_0(E_n)-\epsilon\leq\sum_{n=1}^\infty\mu(F_n)\leq\sum_{n=1}^\infty\mu_0(E_n)$$
            we see that $\cup F_n$ is considered in the set $\{\mu(E)<\infty:\mM\ni E\subseteq\cup E_n\}$ as $\mu(\cup F_n)\leq\sum\mu_0(E_n)<\infty$ and thus $\sum\mu_0(E_n)-\epsilon\leq\mu_0(\cup E_n)$. 
            As $\epsilon$ was arbitrary, we must have that $\sum\mu_0(E_n)\leq\mu(\cup E_n)$ and we are done.
            \item If $\sum_{n=1}^\infty\mu_0(E_n)=\infty$, then its trivial if $\mu_0(E_n)=\infty$ for some $n\in\bN$ so assume thats not the case. 
            For any $M>0$ there must exist some $n\in\bN$ such that $\sum_{k\leq n}\mu_0(E_k)\geq M$. 
            Using arguements similar to those in the first case, for any $\epsilon>0$ we can find $\mM\ni F\subseteq\cup_{k\leq n}E_k\subseteq\cup E_k$ such that $M-\epsilon\leq\sum_{n\leq k}\mu_0(E)-\epsilon\leq\mu(F)\leq\sum_{k\leq n}\mu_0(E_k)<\infty$.
            Again using a similar line of reasoning we find that $M-\epsilon\leq\mu_0(\cup E_k)$ and as $M$ was arbitrarily large and $\epsilon$ arbitrarily small we must have $\mu_0(\cup E_n)=\infty=\sum\mu_0(E_n)$.
            we have hence proven that $\mu_0$ is a measure
        \end{enumerate}
        we have hence proven that $\mu_0$ is a measure. Now suppose $\mu_0(E)=\infty$, then there must exist $\mM\ni F\subseteq E$ with finite positive measure, and for these $F$, we trivially have $\mu_0(F)=\mu(F)$ so $\mu_0$ is semifinite.
        \item Clearly, $\mu_0\leq\mu$ always holds and we have equality when $\mu<\infty$, we will show that equality still holds when $\mu=\infty$ if $\mu$ is semifinite. 
        Say $\mu(E)=\infty,$ then there exists $\mM\ni F\subseteq E$ with arbitrarily large $\mu(F)<\infty$ by \hyperref[prob:1.14]{Problem 1.14} which implies that $\mu_0(E)=\infty$ and we are done.
        \item Let $\nu(E)$ be $0$ if $E$ is $\mu\text{-}\sigma$ finite and $\infty$ otherwise. To prove that $\nu$ is a measure it suffices to prove that if $\cup E_n$ is not $\mu\text{-}\sigma$ finite for disjoint $\{E_n\}_{n\in\bN}$ then not all $E_n$ can be $\mu\text{-}\sigma$ finite. 
        This is true because if all $E_n$ are $\mu\text{-}\sigma$ finite, we can write $E_n=\cup_j F_{n,j}$ with $\mu(F_{n,j})<\infty$ and thus, $\cup E_n=\cup_{n,j} F_{n,j}$ which is a countable union, making $\cup E_n$ $\mu\text{-}\sigma$ finite which is not true.
        This is sufficnent as, if $\{E_n\}$ are disjoint sets and $\cup E_n$ is not $\mu\text{-}\sigma$ finite then $\nu(\cup E_n)=\infty$ by definition and $\sum\nu(E_n)=\infty$ too as some $E_n$ is not $\mu\text{-}\sigma$ finite, making $\nu(E_n)=\infty$ and when $\cup E_n$ is $\mu\text{-}\sigma$ finite we see that all of $E_n$ are trivially also $\mu\text{-}\sigma$ finite, thus $\nu(\cup E_n)=0=\sum\nu(E_n)$ as $\nu(E_n)=0$ for all $n\in\bN$. As $\nu(\phi)=0$ clearly and $\nu$ has sigma additivity, its a measure.
        Now $\mu=\mu_0+\nu$ as $\mu=mu_0$ on $\mu\text{-}\sigma$ finite sets (if $E=\cup E_n$ with $\mu(E_n)<\infty$ then consider, for $n\in\bN$ $G_n:=E_n\setminus\cup_{k<n}E_k\in\mM,$ these are of finite measure too by monotonicity and partition $E$, thus $\mu(E)=\mu(\cup G_n)=\sum\mu(G_n)=\sum\mu_0(G_n)=\mu_0(\cup G_n)=\mu_0(E)$ and when $E$ is not sigma finite we clearly have $\mu(E)=\infty$ and $\mu_0(E)+\nu(E)=\mu_0(E)=\infty=\mu(E)$ so we are done.
\end{enumerate}
\end{solution}
\begin{problem}\label{prob:1.16}
Let $(X, \mathcal{M}, \mu)$ be a measure space. A set $E \subset X$ is called \textbf{locally measurable} if $E \cap A \in \mathcal{M}$ for all $A \in \mathcal{M}$ such that $\mu(A) < \infty$. Let $\tilde{\mathcal{M}}$ be the collection of all locally measurable sets. Clearly $\mathcal{M} \subset \tilde{\mathcal{M}}$; if $\mathcal{M} = \tilde{\mathcal{M}}$, then $\mu$ is called \textbf{saturated}.
\begin{enumerate}[label=\alph*.]
    \item If $\mu$ is $\sigma$-finite, then $\mu$ is saturated.
    \item $\tilde{\mathcal{M}}$ is a $\sigma$-algebra.
    \item Define $\tilde{\mu}$ on $\tilde{\mathcal{M}}$ by $\tilde{\mu}(E) = \mu(E)$ if $E \in \mathcal{M}$ and $\tilde{\mu}(E) = \infty$ otherwise. Then $\tilde{\mu}$ is a saturated measure on $\tilde{\mathcal{M}}$, called the \textbf{saturation} of $\mu$.
    \item If $\mu$ is complete, so is $\tilde{\mu}$.
    \item Suppose that $\mu$ is semifinite. For $E \in \tilde{\mathcal{M}}$, define $\underline{\mu}(E) = \sup\{\mu(A) : A \in \mathcal{M} \text{ and } A \subset E\}$. Then $\underline{\mu}$ is a saturated measure on $\tilde{\mathcal{M}}$ that extends $\mu$.
    \item Let $X_1, X_2$ be disjoint uncountable sets, $X = X_1 \cup X_2$, and $\mathcal{M}$ the $\sigma$-algebra of countable or co-countable sets in $X$. Let $\mu_0$ be counting measure on $\mathcal{P}(X_1)$, and define $\mu$ on $\mathcal{M}$ by $\mu(E) = \mu_0(E \cap X_1)$. Then $\mu$ is a measure on $\mathcal{M}$, $\tilde{\mathcal{M}} = \mathcal{P}(X)$, and in the notation of parts (c) and (e), $\tilde{\mu} \ne \underline{\mu}$.
\end{enumerate}
\end{problem}
\begin{solution}
    \begin{enumerate}[label=\alph*.]
        \item Let, $X=\cup E_n$ with $\mu(E_n)<\infty$ then, $A\in\tilde{\mM}\implies\forall n\in\bN,A\cap E_n\in\mM$ and thus $\cup(E_n\cap A)=A\cap X=A\in\mM$ proving $\tilde{\mM}\subseteq \mM$ making $\mu$ saturated.
        \item Clearly $\phi,X\in\mM$ and if $A\in\tilde{\mM}$ then for all $E\in\mu^{-1}([0,\infty)),A\cap E\in\mM$ and thus, $E\setminus (A\cap E)=E\cap A^c\in\mM\implies A^c\in\mM$ so $\tilde{\mM}$ is closed under complements. 
        If $\{A_n\}_{n\in\bN}\subseteq\tilde{\mM}$ then, $\mu(E)<\infty\implies\forall n\in\bN,A_n\cap E\in\mM\implies\cup(A_n\cap E)=E\cap(\cup A_n)\in\mM\implies\cup A_n\in\tilde{\mM}$ and hence $\tilde{\mM}$ is a $\sigma$-algebra.
        \item Clearly $\tilde{\mu}(\phi)=0$ and if $\{E_n\}_{n\in\bN}\subseteq\tilde{\mM}$ are disjoint then when $\tilde{\mu}(\cup E_n)=\infty$ there are two cases : either $\cup E_n\not\in\mM,$ in which case theres some $E_n\not\in\mM$ as well, making $\sum\tilde{\mu}(E_n)=\infty$ as well, otherwise $\cup E_n\in\mM$, now if $A=\{n\in\bN:E_n\in\mM\}$ then $\cup_{n\in\bN\setminus A}E_n\in\mM$ and thus $\tilde{\mu}(\cup E_n)=\mu(\cup_{n\in A}E_n)+\mu(\cup_{n\in\bN\setminus A}E_n)$. 
        If $\mu(\cup_{n\in A}E_n)=\infty$ then clearly we have sigma additivity for $\{E_n\}$ and otherwise, we must have $\mu(\cup_{n\in\bN\setminus A}E_n)=\infty$ and we again have sigma additivity as $\tilde{\mu}(E_k)=\infty$ for $k\in\bN\setminus A$.  
        The other case is that of $\tilde{\mu}(\cup E_n)<\infty$ in which case $\cup E_n\in\tilde{\mM}\implies E_n\cap(\cup_m E_m)=E_n\in\mM$ for all $n\in\bN$ by definition so we have sigma additivity.
        We have thus shown that $\tilde{\mu}$ is sigma additive and maps $\phi$ to $0$ making it a measure.
        To show saturation we have to show that $\tilde{\tilde{M}}\subseteq\tilde{M}$.
        If $A\in\tilde{\tilde{\mM}}$ we have that $E\in\tilde{M}$ and $\tilde{\mu}(E)<\infty\implies A\cap E\in\tilde{M}$ and as $E\in\mM$ in this case we have that $E\in\mM$ and $\tilde{\mu}(E)=\mu(E)<\infty\implies A\cap E\in\tilde{M}$ and again, $A\cap E\in\tilde{M}\implies(A\cap E)\cap E=A\cap E\in\mM\implies A\in\tilde{M}$ thus $\tilde{\tilde{\mM}}\subseteq\tilde{\mM}$ and we are done.
        \item If $A\in\tilde{\mM}$ and $\tilde{\mu}(A)=0\implies A\in\mM$ and thus $\tilde{\mu}(A)=\mu(A)=0$. 
        Now if $F\subseteq A$ then as $\mu$ is complete we have $F\in\mM\subseteq\tilde{\mM}\ni F$ and thus, $\tilde{\mu}(F)=\mu(F)=0$ so $\tilde{\mu}$ is complete as well.
        \item Clearly $\underline{\mu}\vert_{\mM}=\mu$ by monotonicity, we will prove sigma additivity now. 
        First we prove the claim : if $\mu$ is semifinite then $\underline{\mu}(E)=\sup\{\mu(F):F\in\mM,F\subseteq E,\mu(F)<\infty\}$. If $\underline{\mu}(E)=\infty$ then that means that theres some $\mM\ni F\subseteq E$ with $\mu(F)=\infty$ or, arbitrarily large $\mu(F)$, in the later case we are done but in the prior case, as $\mu$ is semifinite we can find arbitrarily large $\mu(K)$ for $\mM\ni K\subseteq F\subseteq E$ by using \hyperref[prob:1.14]{Problem 1.14} which implies $\sup\{\mu(F):F\in\mM,F\subseteq E,\mu(F)<\infty\}=\infty=\underline{\mu}(E).$
        In the other case where $\underline{\mu}(E)<\infty$ its trivial to see that our claim is true as, using similar definitions, $\mu(F)<\infty$ is forced already.
        We will use this now, let $\{E_n\}_{n\in\bN}\subseteq\tilde{\mM}$ be disjoint, then
        $$\underline{\mu}\left(\bigcup_{n\in\bN}E_n\right)=\sup\left\{\mu(E):E\in\mM,E\subseteq\bigcup_{n\in\bN}E_n,\mu(E)<\infty\right\}$$
        $$=\sup\left\{\sum_{n=1}^\infty\underbrace{\mu(E\cap E_n)}_{\leq\underline{\mu}(E_n),\text{ as }E\cap E_n\in\mM}:E\in\mM,E\subseteq\bigcup_{n\in\bN}E_n,\mu(E)<\infty\right\}\leq\sum_{n=1}^\infty\underline{\mu}(E_n)$$
        $$=\sum_{n=1}^\infty\sup\left\{\mu(F_n):F_n\in\mM,F_n\subseteq E_n,\mu(F_n)<\infty\right\}$$
        $$=\sup\left\{\sum_{n=1}^\infty\mu(F_n):n\in\bN,F_n\in\mM,F_n\subseteq E_n,\mu(F_n)<\infty\right\}$$
        $$=\sup\left\{\mu\left(\bigcup_{n\in\bN}F_n\right):\mM\ni\bigcup_{n\in\bN}F_n\subseteq\bigcup_{n\in\bN}E_n\right\}$$
        $$\leq\sup\{\mu(F):\mM\ni F\subseteq \cup E_n\}=\underline{\mu}(\cup E_n)$$
        and thus we have sigma additivity making $\underline{\mu}$ a measure. Now its saturated as $A\in\tilde{\tilde{M}}$ implies that for all $B\in\mM\subseteq\tilde{\mM}$ with $\mu(B)<\infty$ we have $\mu(B)=\underline{\mu}(B)<\infty$ and thus, $B\cap A\in\tilde{M}$ and $(B\cap A)\cap B=B\cap A\in\mM\implies A\in\mM\implies \tilde{\tilde{M}}\subseteq\tilde{M}$ and we are done. 
        \item Write subsets $S$ as $S_1\cup S_2$ where $S_i\subseteq X_i;i=1,2$ from now on. $\mu$ is a measure by \hyperref[prob:1.10]{Problem 1.10}.
        We see that a set in $\mM$ has finite measure iff it has finite intersection with $X_1$, thus these are countable or cocountable sets of form $Y_1\cup Y_2$ with $Y_1$ finite and when these are countable we clearly have that $Y_2$ is countable or otherwise we need them to be cocountable in which case $(Y_1\cup Y_2)^c=(X_1\setminus Y_1)\cup(X_2\setminus Y_2)$ is countable which is impossible as $|X_1\setminus Y_1|=\mathfrak{c}>\omega$. 
        Thus the sets of finite measure are exactly those of form $Y_1\cup Y_2$ with $Y_1$ finite and $Y_2$ countable. 
        Given any $S_1\cup S_2\subseteq X$ we see that $(S_1\cup S_2)\cap(Y_1\cup Y_2)=(S_1\cap Y_1)\cup(S_2\cap Y_2)$ is countable and as its a subset of a countable set $Y_1\cup Y_2$ making $\tilde{M}=\mathcal P(X)$. 
        Now, $\tilde{\mu}(X_2\cup\{x_1\})$ where $x_1$ is some element in $X_1$ is $\infty$ as $X_2\cup\{x_1\}$ not cocountable nor countable whereas $\underline{\mu}(X_2\cup\{x_1\})=\sup\{\mu(Y_1)=\mu(Y_1\cup Y_2):Y_1\in\{\{x_1\},\phi\},|Y_2|\leq\omega,Y_2\subseteq X_2\}=1$ and thus $\underline{\mu}\neq\tilde{\mu}.$
    \end{enumerate} 
\end{solution}
\begin{problem}\label{prob:1.17}
If $\mu^*$ is an outer measure on $X$ and $\{A_j\}_1^\infty$ is a sequence of disjoint $\mu^*$-measurable sets, then $\mu^*(E \cap (\cup_1^\infty A_j)) = \sum_1^\infty \mu^*(E \cap A_j)$ for any $E \subset X$.
\end{problem}
\begin{solution}
    As outer measures are subadditive, it suffices to show $\mu^*(E\cap(\cup A_n))\geq\sum\mu^*(E\cap A_n).$ 
    Let, $B_n=\cup_{k\leq n}A_k$ and $B=\cup A_n=\cup B_n.$ 
    Then as $A_n$ are all measurable and hence we can say,
    $$\mu^*(E\cap B)\geq\mu^*(E\cap B_n)=\mu^*(E\cap B_n\cap A_n)+\mu^*(E\cap B_n\cap A_n^c)$$
    $$=\mu^*(E\cap A_n)+\mu^*(E\cap B_{n-1})=\ldots=\sum_{m\leq n}\mu^*(E\cap A_m)$$
    by monotonicity and the definition of measurablity for outer measures. 
    As $n\to\infty$ we get our desired inequality. 
\end{solution}
\begin{problem}\label{prob:1.18}
Let $\mathcal{A} \subset \mathcal{P}(X)$ be an algebra, $\mathcal{A}_\sigma$ the collection of countable unions of sets in $\mathcal{A}$, and $\mathcal{A}_{\sigma\delta}$ the collection of countable intersections of sets in $\mathcal{A}_\sigma$. Let $\mu_0$ be a premeasure on $\mathcal{A}$ and $\mu^*$ the induced outer measure.
\begin{enumerate}[label=\alph*.]
    \item For any $E \subset X$ and $\epsilon > 0$ there exists $A \in \mathcal{A}_\sigma$ with $E \subset A$ and $\mu^*(A) \le \mu^*(E) + \epsilon$.
    \item If $\mu^*(E) < \infty$, then $E$ is $\mu^*$-measurable iff there exists $B \in \mathcal{A}_{\sigma\delta}$ with $E \subset B$ and $\mu^*(B \setminus E) = 0$.
    \item If $\mu_0$ is $\sigma$-finite, the restriction $\mu^*(E) < \infty$ in (b) is superfluous.
\end{enumerate}
\end{problem}
\begin{solution}
    \begin{enumerate}[label=\alph*.]
        \item By the definition of the outer measure as an infimum, for any $\epsilon>0$ there exists (wlog we can consider these to be disjoint) $\{K_n\}\subseteq\mA$ such that $\sum\mu_0(K_n)\leq\mu^*(E)+\epsilon$.  
        Now we see that, 
        $$\sum_{n=1}^{\infty}\mu_0(K_n)=\lim_{m\to\infty}\sum_{n<m}\mu_0(K_n)=\lim_{m\to\infty}\mu_0(\underbrace{\bigcup_{n<m}K_n}_{\in\mA})=\lim_{m\to\infty}\mu^*(\underbrace{\bigcup_{n<m}K_n}_{\in\mM(\mA)})=\mu^*(\cup K_n)$$
        using continuity from below as $\mu^*$ is a measure on the sigma algebra of $\mu^*$-measureable sets which contiains $\mA$ and thus $\mM(\mA)$ as well. As $\cup K_n\in\mA_\sigma$ we are done.
        \item Assume that $E$ is measurable. 
        By part (a) we can find $E\subseteq E_n\in A_\sigma$ with $\mu^*(E)+1/n>\mu^*(E_n)\geq\mu^*(E).$ 
        Then using continuity from above ($\mu^*(E)+1/n<\infty$ as per the statement) we get that $\mu^*(\cap E_n)=\mu^*(E)$ and $E\subseteq K:=\cap E_n\in\mA_{\sigma\delta}.$ 
        Using the definition of measurablity and finitude of $\mu^*(E)$ we get that $\mu^*(K\setminus E)=\mu^*(K)-\mu^*(E)=0$. 
        For the other direction, let $B,E$ be as in the statement, then as $\mu^*$ is complete and $\mA_{\sigma\delta}\subseteq\mM(\mA),B\setminus E$ and $B$ are measureable here, this makes their difference measureable too i.e. $B\setminus(B\setminus E)=E$ is also measurable. 
        \item Assuming $\sigma$-finiteness, let $\{H_n\}_{n\in\bN}\subseteq\mA$ be disjoint and of finite $\mu_0$-measure such that $X\subseteq\cup H_n$. 
        Assume that $E$ is measurable now. $H_n$ being measurable forces $E\cap H_n$ to be measurable as well.
        Using part (a), fixing $\epsilon>0$ we can find $B_n\in\mA_\sigma$ containing $E\cap H_n$ such that $\mu^*(B_n)\leq\mu^*(E\cap H_n)+\epsilon 2^{-n}$ which, by the finitude of $\mu^*(E\cap H_n)$ (its clearly less than $\mu_0(H_n)<\infty$) and definition of $\mu^*$-measurablity, after moving a term to another side of the equality implies $\mu^*(B_n\setminus (E\cap H_n))=\mu^*(B)-\mu^*(E\cap H_n)<\epsilon2^{-n}$. 
        Thus summing these over $n\in\bN$ and using subadditivity of outer measures we have
        $$\mu^*\left(\bigcup_{n\in\bN}(B_n\setminus (E_n\cap H_n))\right)\leq\sum_{n\in\bN}\mu^*(B_n\setminus(H_n\cap E))\leq\epsilon$$
        Now using monotonicity and the fact that $(\cup B_n)\setminus E=(\cup B_n)\setminus(\cup(H_n\cap E))\subseteq\cup(B_n\setminus(H_n\cap E))$ we see that $\mu^*((\cup B_n)\setminus E)<\epsilon$ where $E\subseteq\cup B_n\in\mA_{\sigma\sigma}=\mA_\sigma$.
        As $\epsilon$ was arbitrarily small we can find $V_n\in\mA_\sigma$ containing $E$ for which $\mu^*(V_n\setminus E)<1/n$.
        And again by monotonicity we see that $\mu^*((\cap V_n)\setminus E)=0$ with $\cap V_n\in\mA_{\sigma\delta}$. The proof for another direction does not depend on the finitude of $\mu^*(E)$ and is exactly the same to that in part (b) so we are done. 
    \end{enumerate}
\end{solution}
\begin{problem}\label{prob:1.19}
Let $\mu^*$ be an outer measure on $X$ induced from a finite premeasure $\mu_0$. If $E \subset X$, define the inner measure of $E$ to be $\mu_*(E) = \mu_0(X) - \mu^*(E^c)$. Then $E$ is $\mu^*$-measurable iff $\mu^*(E) = \mu_*(E)$. (Use Exercise 18.)
\end{problem}
\begin{solution}
    $\mu_0$ being finite cleary makes $\mu^*$ finite as well so we can rearrange terms freely.
    The only if direction is trivial, we will prove the if direction so assume that $\mu^*(E)=\mu_*(E)$, equivalently $\mu^*(X)=\mu^*(E)+\mu^*(E^c)$. 
    For any $S\subseteq X$ and any $\epsilon>0$, from \hyperref[prob:1.18]{Problem 1.18} we can find $\mu^*$-measurable $K$ containing $S$ with $\mu^*(S)+\epsilon>\mu^*(K)$. 
    As $K$ is $\mu^*$-measurable, by definition we have
    $$(\mu^*(S)+\epsilon)+\mu^*(K^c)>\mu^*(K)+\mu^*(K^c)=\mu^*(X)=\mu^*(E)+\mu^*(E^c)$$
    $$=\mu^*(E\cap K)+\mu^*(E\cap K^c)+\mu^*(E^c\cap K)+\mu^*(E^c\cap K^c)\geq\mu^*(E\cap K)+\mu^*(E^c\cap K)+\mu^*(K^c)$$
    $$\geq\mu^*(E\cap S)+\mu^*(E^c\cap S)+\mu^*(K^c)$$
    by using subadditivity and monotonicity. 
    Subtracting $\mu^*(K^c)$ we get, $\mu^*(S)+\epsilon\geq\mu^*(E\cap S)+\mu^*(E^c\cap S)$ and as $\epsilon$ was arbitrarily small we can remove it from the inequality and it will still be valid, which is enough to prove equality as we already have subadditivity by the virtue of $\mu^*$ being an outer measure.
\end{solution}
\begin{problem}\label{prob:1.20}
Let $\mu^*$ be an outer measure on $X$, $\mathcal{M}^*$ the $\sigma$-algebra of $\mu^*$-measurable sets, $\bar{\mu} = \mu^*|\mathcal{M}^*$, and $\mu^+$ the outer measure induced by $\bar{\mu}$ as in (1.12) (with $\bar{\mu}$ and $\mathcal{M}^*$ replacing $\mu_0$ and $\mathcal{A}$).
\begin{enumerate}[label=\alph*.]
    \item If $E \subset X$, we have $\mu^*(E) \le \mu^+(E)$, with equality iff there exists $A \in \mathcal{M}^*$ with $A \supset E$ and $\mu^*(A) = \mu^*(E)$.
    \item If $\mu^*$ is induced from a premeasure, then $\mu^* = \mu^+$. (Use Exercise 18a.)
    \item If $X = \{0, 1\}$, there exists an outer measure $\mu^*$ on $X$ such that $\mu^* \ne \mu^+$.
\end{enumerate}
\end{problem}
\begin{solution}
    \begin{enumerate}[label=\alph*.]
        \item $\mu^*\leq\mu^+$ follows from subadditivity, monotonicity and the definition of an infimum. 
        Now suppose we have $\mu^*(E)=\mu^+(E),$ by the definition of $\mu^+$, for all $n\in\bN$ we can find $A_n=\cup_m K_{m,n}\supseteq E$ with $K_{m,n}\in\mM^*$ such that 
        $$\mu^+(E)+1/n\geq\sum_m\mu^*(K_{m,n})\geq\mu^*(A_n)$$
        using monotonicity. As $\mM^*$ forms a $\sigma$-algebra we see that $\cap A_n\in\mM^*$ and by monotonicity, $\mu^*(\cap A_n)\leq\mu^*(E)$ and we have equality as $E\subseteq\cap A_n$. 
        The other direction is trivial.
        \item Let $\mu^*$ be induced from a premeasure $\mu_0$ on $\mA$. Then using part (a) of \hyperref[prob:1.18]{Problem 1.18}, for any $E\subseteq X$ we can find $A_n\in\mA_\sigma\subseteq\mM^*$ containing $E$ such that $\mu^*(E)+1/n\geq\mu^*(A_n)$, now as $\mM^*$ is a sigma algebra we have $E\subseteq\cap A_n\in\mM^*$ and also $\mu^*(E)\geq\mu^*(\cap A_n)$, which implies equality by monotonicity. 
        Thus we have found a set $\cap A_n\in\mM^*$ containing $E$ with $\mu^*(\cap A_n)=\mu^*(E)$ and thus by part (a) we must have $\mu^*(E)=\mu^+(E).$
        As $E$ was arbitrary we have $\mu^*=\mu^+.$
        \item Define $\mu^*:\emptyset\mapsto 0,\{0\}\mapsto 2,\{1\}\mapsto 3,\{0,1\}\mapsto 4$. 
        This meets all the criterions for being an outer measure i.e. monotonicity and $\sigma$-subadditivity and being $0$ on the emptyset.
        It can be shown that $\mM^*=\{\emptyset,\{0,1\}\}$ and that $\mu^+(\{1\})=4\neq\mu^*(\{1\})$ so we are done.
     \end{enumerate}
\end{solution}
\begin{problem}\label{prob:1.21}
Let $\mu^*$ be an outer measure induced from a premeasure and $\bar{\mu}$ the restriction of $\mu^*$ to the $\mu^*$-measurable sets. Then $\bar{\mu}$ is saturated. (Use Exercise 18.)
\end{problem}
\begin{solution}
    If $A\subseteq X$ be locally $\bar{\mu}$-measurable. 
    If $F\subseteq X$ is such that $\mu^*(F)=\infty$ then by subadditivity atleast one of $\mu^*(F\cap A),\mu^*(F\cap A^c)$ is infinite and thus we must have $\mu^*(F)=\infty=\mu^*(F\cap A)+\mu^*(F\cap A^c)$. 
    Now suppose we had $\mu^*(F)<0,$ using part (a) of \hyperref[prob:1.18]{Problem 1.18} and the arguements in part (a) of \hyperref[prob:1.20]]{Problem 1.20} we can find a $\mu^*$-measureable $C\supset F$ such that $\mu^*(F)=\mu^*(C)$.
    As $\mu^*(C)<\infty$ we must have that $C\cap F$ is measurable and thus so is $C\setminus(C\cap A)=C\cap A^c$ as the measurable sets form a $\sigma$-algebra here. 
    Thus we have,
    $$\mu^*(F)=\mu^*(C)=\bar{\mu}(C)=\bar{\mu}(C\cap A)+\bar{\mu}(C\cap A^c)=\mu^*(C\cap A)+\mu^*(C\cap A^c)$$
    $$\geq\mu^*(F\cap A)+\mu^*(F\cap A^c)$$
    using additivity of $\bar{\mu}$ and monotonicity of $\mu^*$.
    This implies, by subadditivity of $\mu^*$ that $\mu^*(F)=\mu^*(F\cap A)+\mu^*(F\cap A^c)$ and we have proven that this holds for all $F\subseteq X$, hence proving measurablity of $A$ by definition. 
    This shows that all locally measurable sets are measureable w.r.t $\bar{\mu}$, making it saturated. 
\end{solution}
\begin{problem}\label{prob:1.22}
Let $(X, \mathcal{M}, \mu)$ be a measure space, $\mu^*$ the outer measure induced by $\mu$ according to (1.12), $\mathcal{M}^*$ the $\sigma$-algebra of $\mu^*$-measurable sets, and $\bar{\mu} = \mu^*|\mathcal{M}^*$.
\begin{enumerate}[label=\alph*.]
    \item If $\mu$ is $\sigma$-finite, then $\bar{\mu}$ is the completion of $\mu$. (Use Exercise 18.)
    \item In general, $\bar{\mu}$ is the saturation of the completion of $\mu$. (See Exercises 16 and 21.)
\end{enumerate}
\end{problem}
\begin{solution}
    \begin{enumerate}[label=\alph*.]
        \item $\sigma$-algebras are algebras and measures are also premeasures, this allows us to utilise parts of \hyperref[prob:1.18]{Problem 1.18}. 
        The completion of $\mu$ is denoted by $\hat{\mu}$. 
        We will show that 
        $$\mM^*=\overline{\mM}:=\{E\cup F|E\in\mM,\exists N\in\mM:\mu(N)=0\text{ and }F\subseteq N\}$$ 
        this is sufficient as by caratheodory's theorem $\bar{\mu}$ is complete and we also know that $\mu^*|_{\mM}=\mu$ (from proposition 1.13), so it extends $\mu$. 
        As this complete extension is unique, $\bar{\mu}$ must be the completion here. 
        If $S\in\mM^*$ then using part (c) of problem 18 we can find $B\in\mM_{\sigma\delta}=\mM$ such that $S\subseteq B$ and $\mu^*(B\setminus S)=0$.
        By part (a) of problem 18, with some work we can show that there exists some $N\in\mM_\sigma=\mM$ such that $B\setminus S\subseteq N$ and $\mu^*(B\setminus S)=\mu^*(N)=\mu(N)=0$. 
        Now, to show that $E\in\mM$, it suffices to show that $E^c\in\mM$ as $\mM$ is a $\sigma$-algebra.
        Rewriting $E$ as $B\setminus(B\setminus E)$ we have $E^c=B^c\cup(B\setminus E)$ with $B^c\in\mM$ and $B\setminus E\subseteq N\in\mM$ of zero measure and thus $E^c\in\overline{\mM}$ implying $\mM^*\subseteq\overline{\mM}$.
        For the other direction of the set inclusion, any $E\cup F$ with $E\in\mM,F\subseteq N\in\mM$ with $\mu(N)=0$. 
        Also, $E\cup N\in\mM$ and $(E\cup N)\setminus(E\cup F)\subseteq N\setminus F$ so $\mu^*((E\cup N)\setminus(E\cup F))\leq\mu^*(N\setminus F)\leq\mu^*(N)=\mu(N)=0$, making $E\cup F$ measurable by part (b) of Problem 18.
        \item We need to first show that $\mM^*=\tilde{\overline{\mM}}$. 
        Let $E\in\mM^*,$ consider $K\in\overline{\mM}\subseteq\mM^*$ of finite $\mu$-measure, then we see that $E\cap K\in\mM^*$ and they have finite $\mu^*$-measure, using part (a,b) of problem 18 as we did in part (a) of this problem we can prove that $E\cap K\in\tilde{\overline{\mM}}$.
        This shows one direction of the equality $\mM^*\subseteq\overline{\mM}.$ 
        Now we show the other direction $\tilde{\overline{\mM}}\subseteq\mM^*$. 
        Suppose $G\in\tilde{\overline{\mM}}$ then for all $K\in\overline{\mM}$, of finite $\hat{\mu}$-measure we have $G\cap K\in\overline{\mM}$. 
        It suffices to show that $G\in\tilde{\mM^*}$ as $\mu^*$ is saturated. 
        If $K\in\mM^*$ has finite $\mu^*$-measure then we have shown earlier that $K\in\overline{\mM}$.
        We also have that $\mu^*=\hat{\mu}$ on $\overline{\mM}$ as any $E\cup F$ with $E\in\mM$ and $F\subseteq N\in\mu^{-1}(\{0\})$, thus $\mu^*(E\cup F)\leq\mu^*(E\cup N)\leq\mu^*(E)+\mu^*(N)=\mu(E)+\mu(N)=\mu(E)$ and we have equality by monotonicity.
        Now we see that $\mu^*(K)<\infty\implies\hat{\mu}(K)<\infty$ and thus $E\cap K\in\overline{\mM}\subseteq\mM^*$ proving $\tilde{\overline{\mM}}\subseteq\mM^*$, finally resulting in equality.
        Now to show that both of these measures agree, we already know that they agree on $\overline{\mM}$, outside which $\hat{\mu}=\infty$ and thus we need to show that $\mu^*=\infty$ too.
        However this is true as $\mu^*(E)<\infty$ for $E\in\mM^*$ implies $E\in\overline{\mM}$ so we are done.
    \end{enumerate}
\end{solution}
\begin{problem}\label{prob:1.23}
Let $\mathcal{A}$ be the collection of finite unions of sets of the form $(a, b] \cap \mathbb{Q}$ where $-\infty \le a < b \le \infty$.
\begin{enumerate}[label=\alph*.]
    \item $\mathcal{A}$ is an algebra on $\mathbb{Q}$. (Use Proposition 1.7.)
    \item The $\sigma$-algebra generated by $\mathcal{A}$ is $\mathcal{P}(\mathbb{Q})$.
    \item Define $\mu_0$ on $\mathcal{A}$ by $\mu_0(\emptyset) = 0$ and $\mu_0(A) = \infty$ for $A \ne \emptyset$. Then $\mu_0$ is a premeasure on $\mathcal{A}$, and there is more than one measure on $\mathcal{P}(\mathbb{Q})$ whose restriction to $\mathcal{A}$ is $\mu_0$.
\end{enumerate}
\end{problem}
\begin{solution}
    \begin{enumerate}[label=\alph*.]
        \item This is trivial.
        \item Given any nonempty $X\subseteq\bQ$ we see that it must be countable, thus 
        $$X=\bigcup_{x\in X}\{x\}=\bigcup_{x\in X}\bigcap_{n\in\bN}\Big(x-\frac{1}{n},x\Big]\in\mA_{\delta\sigma}\subseteq\mM(\mA)$$ 
        thus we must have equality as $\mM(\mA)\subseteq\mathcal P(\bQ)$. 
        \item All elements of $\mA$ can be easily shown to have have nonempty interior.
        Let $\mu_1=\infty$ except on $\emptyset$ and $\mu_2=\infty$ except on $\{\emptyset,\{0\}\}$.
        It can be easily shown that $\mu_1|_{\mA}=\mu_2|_{\mA}=\mu_0$ with $\mu_1\neq\mu_2$
    \end{enumerate}
\end{solution}
\begin{problem}\label{prob:1.24}
Let $\mu$ be a finite measure on $(X, \mathcal{M})$, and let $\mu^*$ be the outer measure induced by $\mu$. Suppose that $E \subset X$ satisfies $\mu^*(E) = \mu^*(X)$ (but not that $E \in \mathcal{M}$).
\begin{enumerate}[label=\alph*.]
    \item If $A, B \in \mathcal{M}$ and $A \cap E = B \cap E$, then $\mu(A) = \mu(B)$.
    \item Let $\mathcal{M}_E = \{A \cap E : A \in \mathcal{M}\}$, and define the function $\nu$ on $\mathcal{M}_E$ defined by $\nu(A \cap E) = \mu(A)$ (which makes sense by (a)). Then $\mathcal{M}_E$ is a $\sigma$-algebra on $E$ and $\nu$ is a measure on $\mathcal{M}_E$.
\end{enumerate}
\end{problem}
\begin{solution}
    \begin{enumerate}[label=\alph*.]
        \item We have $A\setminus B=((A\cap E)\setminus(B\cap E)\cup((A\cap E^c)\setminus(B\cap E^c)))=(A\cap E^c)\setminus(B\cap E^c)$ and also $(A\setminus B)^c=(B\cap E^c)\cup(A\cap E^c)^c=(B\cap E^c)\cup(A^c\cup E)\supseteq E$.
        This implies $\mu((A\setminus B)^c)=\mu^*((A\setminus B)^c)\geq\mu^*(E)=\mu(X)$ and thus we must have $\mu(A\setminus B)=\mu(X)-\mu((A\setminus B)^c)=\mu(X)-\mu(X)=0.$
        Thus, $\mu(A)=\mu(A\setminus B)+\mu(B\cap A)=\mu(B\cap A)$ by additivity and similarly $\mu(B)=\mu(A\cap B)$ too, so we have $\mu(A)=\mu(B)$.
        \item $\mM_E$ clearly contains $E,\emptyset$ and is closed under countable unions.
        If $E\cap K\in\mM_E$ with $K\in\mM$, we see that $K^c\in\mM$ and thus $E\cap K^c=E\setminus(E\cap K)$ i.e. the complement of $E\cap K$ in $E$ is also in $\mM_E$ so its closed under complements as well, making it a $\sigma$-algebra.
        Clearly $\nu(\emptyset)=0$ and if $\{A_n\cap E\}_{n\in\bN}\subseteq\mM_E$ are disjoint, then using the fact that $S\cap E=(S\cap E)\cap E$ for all $S\in\mM$ and thus $\mu(S)=\mu(S\cap E)$, we have
        $$\nu\left(\bigcup_{n\in\bN}(A_n\cap E)\right)=\nu\left(E\cap\bigcup_{n\in\bN}A_n\right)=\mu(\underbrace{\bigcup_{n\in\bN} A_n}_{\in\mM})=\mu\left(E\cap\bigcup_{n\in\bN}A_n\right)$$
        $$=\mu\left(\bigcup_{n\in\bN}(E\cap A_n)\right)=\sum_{n\in\bN}\mu(E\cap A_n)=\sum_{n\in\bN}\mu(A_n)=\sum_{n\in\bN}\nu(E\cap A_n)$$
    \end{enumerate}
\end{solution}
\begin{problem}\label{prob:1.25}
Complete the proof of Theorem 1.19.
\end{problem}
\begin{solution}
    We start with the following lemma: If $F_n\subseteq(n,n+1]$ are closed subsets for $n\in\bZ$ then so is their union. 
    For a proof, consider any sequence with terms in the union that converges. 
    If the limit is in $(n,n+1]$ then we see that the sequence is eventually in $F_n\cup F_{n+1}$ which is indeed closed and the sequence converges in $F_n\cup F_{n+1}$, and thus also in $\cup F_n$ making it closed. 
    Clearly (b,c) imply (a), we prove the other directions. 
    As the measure $\mu$ defined here is sigma finite, we can consider $\{H_n=(n,n+1]\}_{\in\bZ}\subseteq\mM_\mu$ of finite measure which partitions $\bR$. 
    Now, for every $n\in\bN$ let $U_{n,k}$ and $F_{n,k}$ be open and closed sets respectively such that $F_{n,k}\subseteq H_n\cap E\subseteq U_{n,k}$ and
    $$\mu(U_{n,k})-2^{-|n|}/3k\leq\mu(E)\leq\mu(F_{n,k})+2^{-|n|}/3k$$
    as every term is finite here we can rearrange them to get
    $$\mu(U_{n,k}\setminus E\cap H_n),\mu(E\cap H_n\setminus F_{n,k})\leq 2^{-|n|}/3k$$
    We have that $U^*_k:=\cup_n U_{n,k}$ is open, contains $E$ and $F^*_k:=\cup_n F_{n,k}$(this is a disjoint union too) is closed and contained in $E$ such that
    $$\mu\qty(U^*_k\setminus E)\leq\mu\qty(\bigcup_{n\in\bZ}(U_{n,k}\setminus H_n\cap E))\leq\sum_{n\in\bZ}\mu(U_{n,k}\setminus H_n\cap E)\leq\sum_{n\in\bZ}2^{-|n|}/3k=1/k$$
    using monotonicity as $(\cup U_{n,k})\setminus(\cup E\cap H_n)\subseteq\cup(U_{n,k}\setminus E\cap H_n)$.
    Similarly,
    $$\mu\qty(E\setminus F^*_k)\leq\mu\qty(\bigcup_{n\in\bZ}(E\cap H_n\setminus F_{n,k}))=\sum_{n\in\bZ}\mu(E\cap H_n\setminus F_{n,k})\leq\sum_{n\in\bZ}2^{-|n|}/3k=1/k$$ 
    Let, $U:=\cap U^*_k\in G_\delta$ and $F:=\cup F^*_k\in F_\sigma$ then $F\subseteq E\subseteq U$ and we have $\mu(E\setminus F)=\mu(U\setminus E)=0$. 
    Now, using additivity of $\mu$ we have $\mu(U)=\mu(U\setminus E)+\mu(E)=\mu(E)$ and $\mu(E)=\mu(F)+\mu(E\setminus F)=\mu(F)$. 
    Now we can write $E=U\setminus(U\setminus E)=F\cup(E\setminus F)$ proving that (a) implies (b,c) too and we are done.

\end{solution} 
\begin{problem}\label{prob:1.26}
Prove Proposition 1.20. (Use Theorem 1.18.)
\end{problem}
\begin{solution}
    Using theorem 1.18, given any $\epsilon>0$ we can find open $U\supset E$ with $\mu(U)\leq\mu(E)+\epsilon$ rearranging which gives $\mu(E\Delta U)<\epsilon$.
    Suppose $U=\cup I_k$ where $\{I_k\}$ are disjoint open intervals, then the sum $\sum\mu(I_k)=\mu(A)<\infty$. 
    Truncating this we can find $M\in\bN$ such that $\mu(\cup_{1\leq k\leq m}I_k)+\epsilon\geq\mu(U)$ and thus $\mu((\cup_{1\leq k\leq m}I_k)\Delta U)<\epsilon$.
    As $\mu(\cup_{1\leq k\leq m}I_k),\mu(E),\mu(U)$ are all finite, we have
    $$\mu\qty(E\Delta\qty(\bigcup_{1\leq k\leq m}I_k))\leq\mu\qty(E\Delta U)+\mu\qty(U\Delta(\bigcup_{1\leq m}I_k))<2\epsilon$$
    proving the proposition.
\end{solution}
\begin{problem}\label{prob:1.27}
Prove Proposition 1.22a. (Show that if $x, y \in C$ and $x < y$, there exists $z \notin C$ such that $x < z < y$.)
\end{problem}
\begin{solution}
    This is trivial by considering ternary expansions.
\end{solution}
\begin{problem}\label{prob:1.28}
Let $F$ be increasing and right continuous, and let $\mu_F$ be the associated measure. Then $\mu_F(\{a\}) = F(a) - F(a-)$, $\mu_F([a, b)) = F(b-) - F(a-)$, $\mu_F([a, b]) = F(b) - F(a-)$, and $\mu_F((a, b)) = F(b-) - F(a)$.
\end{problem}
\begin{solution}
    All of these follows from upper and lower continuity theorems of measures.
\end{solution}
\begin{problem}\label{prob:1.29}
Let $E$ be a Lebesgue measurable set.
\begin{enumerate}[label=\alph*.]
    \item If $E \subset N$ where $N$ is the nonmeasurable set described in \S1.1, then $m(E) = 0$.
    \item If $m(E) > 0$, then $E$ contains a nonmeasurable set. (It suffices to assume $E \subset [0, 1]$. In the notation of \S1.1, $E = \cup_{r \in R} (E \cap N_r)$.)
\end{enumerate}
\end{problem}
\begin{solution}
    \begin{enumerate}[label=\alph*.]
        \item If $E$ is a measurable subset of $N$ then we find that, for $r\in\bQ\cap[0,1),E_r:=(r+E\cap[0,1-r))\cup(-1+r+E\cap[1-r,1))\subseteq N_r$ is measurable and $m(E)=m(E_r)$.
        As $N_r$ are disjoint, $E_r$ must also be disjoint.
        Thus,
        $$\sum_{r\in\bQ\cap[0,1)}m(E)=\sum_{r\in[0,1)\cap\bQ}m(E_r)=m\qty(\bigcup_{r\in\bQ\cap[0,1)}E_r)\leq m([0,1))=1$$ 
        this is possible only when $\mu(E)=0$ so that must be the case here. 
        \item After shifting taking a finite measured subset of $E$, we can assume $E\subseteq[0,1)$.
        If all of $E\cap N_r$ are measurable, then as any measurable subset of the nonmeasurable sets $N_r$ has measure $0$ from the previous part (This can be proven using the same arguements we used to prove it for $N$ i.e. when $r=0$) we have,
        $$m(E)=m\qty(\bigcup_{r\in\bQ\cap[0,1)}N_r\cap E)=\sum_{r\in\bQ\cap[0,1)}m(E\cap N_r)=\sum_{r\in\bQ\cap[0,1)}0=0$$
        which is false as $m(E)>0$. 
    \end{enumerate}
\end{solution}
\begin{problem}\label{prob:1.30}
If $E \in \mathcal{L}$ and $m(E) > 0$, for any $\alpha < 1$ there is an open interval $I$ such that $m(E \cap I) > \alpha m(I)$.
\end{problem}
\begin{solution}
    Unless $m(E)<\infty$, we can consider some finite positive measured subset of $E$ so wlog it is safe to assume $m(E)<\infty$ from the start and we do so.
    First we define the following supremum, 
    $$C:=\sup\qty{\frac{m(E\cap I)}{m(I)}:I\text{ is a nonempty open interval}}$$
    which exists because $m(E\cap I)/m(I)\leq 1$ by monotonicity. 
    Thus we must have that $m(E\cap I)\leq Cm(I)$ for all nonempty open $I$.
    By \hyperref[prob:1.18]{Problem 1.18}, for any $\epsilon>0$ we can find open $U\supset E$ with $m(E)+\epsilon>m(U)$.
    Suppose $U=\cup I_k$ with $I_k$ being mutually disjoint open intervals. 
    Then as $E=E\cap U=\cup(E\cap I_k)$ we can write,
    $$m(E)=m\qty(\bigcup_{k}(E\cap I_k))=\sum_{k}m(E\cap I_k)\leq\sum_{k}Cm(I_k)=Cm(U)$$
    and thus, for all $\epsilon>0$ we must have $Cm(U)+\epsilon\geq m(E)+\epsilon>m(U)$ implying that 
    $$\epsilon>(1-C)m(U)\geq(1-C)m(E)$$
    Now as $m(E)>0$ this is only possible when $1-C=0$ and we are done.
\end{solution}
\begin{problem}\label{prob:1.31}
If $E \in \mathcal{L}$ and $m(E) > 0$, the set $E - E = \{x - y : x, y \in E\}$ contains an interval centered at $0$. (If $I$ is as in Exercise 30 with $\alpha > 3/4$, then $E - E$ contains $(-1/2 m(I), 1/2 m(I))$).
\end{problem}
\begin{solution}
    Consider the setup in the hint, then for $r\in(-0.5m(I),0.5m(I))$ we see that $r\in E-E$ iff $(|r|+E)\cap E$ is nonempty. 
    As it contains $\supset (|r|+E\cap I)\cap(E\cap I)$, proving this to be nonempty suffices. 
    Now by the inclusion exclusion principle for measures, for our case we can write 
    $$m((|r|+E\cap I)\cap E\cap I)=m(|r|+E\cap I)+m(E\cap I)-m((|r|+E\cap I)\cup E\cap I)$$
    $$\geq 2m(E\cap I)-m(I\cup(|r|+I))\geq 1.5m(I)-(1+|r|)m(I)=(0.5-|r|)m(I)>0$$
    so it definetly can not be empty and we are done.
\end{solution}
\begin{problem}\label{prob:1.32}
Suppose $\{\alpha_j\}_1^\infty \subset (0, 1)$.
\begin{enumerate}[label=\alph*.]
    \item $\prod_1^\infty (1 - \alpha_j) > 0$ iff $\sum_1^\infty \alpha_j < \infty$. (Compare $\sum \log(1 - \alpha_j)$ to $\sum \alpha_j$.)
    \item Given $\beta \in (0, 1)$, exhibit a sequence $\{\alpha_j\}$ such that $\prod_1^\infty (1 - \alpha_j) = \beta$.
\end{enumerate}
\end{problem}
\begin{solution}
    The first one follows from continuity of the logarithm on its domain and comparison principle for sums using the fact that $\log(1-x)=-x+O(x^2)$ for small $x$. For the second one consider $\alpha_j:=1-\beta^{2^{-j}}$.
\end{solution}
\begin{problem}\label{prob:1.33}
There exists a Borel set $A \subset [0, 1]$ such that $0 < m(A \cap I) < m(I)$ for every subinterval $I$ of $[0, 1]$. (Hint: Every subinterval of $[0, 1]$ contains Cantor-type sets of positive measure.)
\end{problem}
\begin{solution}
    By the previous problem we easily notice that there are cantor-like measurable sets $\subseteq[0,1]$ which are closed, nowhere dense and have positive measure. 
    Let $\{I_n\}$ be an ennumeration of the nonempty open (this is clearly sufficient) subintervals with rational endpoints in $[0,1]$. 
    Let $K_1\subseteq I_1$ be such a set as discussed before, as $K_1$ is nowhere dense there is an open interval in $K_1\setminus K_1$ and we can choose a positive measured nowhere dense set $C_1\subseteq I_1\setminus K_1$ after shifting and dilating sets like $K$.
    Inductively choose $K_1,C_1,\ldots,K_n,C_n$. 
    As the union of finitely many nowhere dense sets is still nowhere dense, there is an open interval in $I_{n+1}\setminus(\cup_{k\leq n}K_k\cup C_k)$ and we can choose a nowhere dense set $K_{n+1}$ in it of positive measure.
    Again we can choose a similar $C_{n+1}\subseteq I_{n+1}\setminus(K_{n+1}\cup(\cup_{k\leq n}K_k\cup C_k))$.
    Continuing ad infintum we get two collections of sets $\{K_n\}$ and $\{C_n\}$, and clearly $K:=\cup K_n$ is disjoint with $C:=\cup C_n.$
    We claim that $A=K$ works. 
    Now, given any subinterval $I$ of $[0,1]$ having positive measure we can find some $I_n\subseteq I$ and thus, because $K_n\subseteq I_n$ we see that $m(I\cap K)\geq m(I_n\cap K_n)=m(K_n)>0$ and we also have that $m(I\cap K^c)\geq m(I_n\cap C_n)=m(C_n)>0$ as $C\subseteq K^c$.
    Thus we have that $0<m(A\cap I)=m(I)-m(A^c\cap I)<m(I)$ so this $A$ satisfies the criterion,
    As all the $K_n$ were closed, $A\in F_\sigma$ so its clearly borel and we are done.
\end{solution}
\stepcounter{section}
\section*{Chapter \thesection\ : Integration}
\addcontentsline{toc}{section}{Chapter \thesection\ : Integration}
\begin{problem}\label{prob:2.1}
Let $f: X \to \overline{\bR}$ and $Y = f^{-1}(\bR)$. Then $f$ is measurable iff $f^{-1}(\{-\infty\}) \in \mM$, $f^{-1}(\{\infty\}) \in \mM$, and $f|_Y$ is measurable on $Y$.
\end{problem}
\begin{solution}
    For the only if direction, $\{\infty\},\{-\infty\},\bR=Y\in\mathcal B_{\overline{\bR}}$ making $f^{-1}$ map them into $\mM$ as $f$ is measurable and as $\mM$ is a sigma algebra we see that $f\vert_Y:Y\to\overline{\bR}$ is $\mM_Y$-measurable as $f^{-1}(E)\cap Y\in\mM_Y$ for all $E\in\mathcal B_{\overline{\bR}}$. 
    Now for the if direction, as $\mathcal B_{\overline{\bR}}=\{E\cup I:E\in\mathcal B_{\bR}, I\subseteq\{\infty,-\infty\}\}$ we see that for any such $E\cup I\in\mathcal B_{\overline{\bR}}$ we have $f^{-1}(E\cup I)=f^{-1}(E)\cup f^{-1}(I)$ and as $f^{-1}(E)\cap Y=f^{-1}(E)\in\mM_Y\subseteq\mM$ and $f^{-1}(I)\in\mM$ so $f$ is measurable on $\mM$.
\end{solution}
\begin{problem}\label{prob:2.2}
Suppose $f, g: X \to \overline{\bR}$ are measurable.
\begin{enumerate}[label=\alph*.]
    \item $fg$ is measurable (where $0 \cdot (\pm\infty) = 0$).
    \item Fix $\alpha \in \overline{\bR}$ and define $h(x) = \alpha$ if $f(x) = -g(x) = \pm\infty$ and $h(x) = f(x) + g(x)$ otherwise. Then $h$ is measurable.
\end{enumerate}
\end{problem}
\begin{solution}
    \begin{enumerate}[label=\alph*.]
        \item Consider the two maps,
        $$X\ni x\longrightarrow (f(x),g(x))\text{ and }\overline{\bR}^2\ni(x,y)\mapsto xy$$
        the first is a measurable map from $X$ (with sigma algebra $\mM$) to $R\times R$ (with sigma algebra $\otimes^2_1 \mathcal B_{\overline{\bR}}$) and the second from this to $\overline{\bR}$ (with sigma algebra $\mathcal B_{\overline{\bR}}$) and as the composition of measurable maps are measurable we are done.
        Here, the measurability of $(x,y)\to xy$ follows from its continuity (note: the metric on $\overline{\bR}$ considered throughout the chapter is usually $d(x,y):=|\arctan(x)-\arctan(y)|$ with $\arctan(\pm\infty)=\pm\pi/2$).
        \item Similarly again consider $x\mapsto(f(x),g(x))$ but this time with $(x,y)\to J(x,y)$ where,
        $$J(x,y):=\begin{cases}\alpha\text{ if }x=-y=\pm\infty\\ x+y\text{ otherwise}\end{cases}$$
        We will use \hyperref[prob:2.5]{Problem 1.5} here. 
        Let $A=\{(+\infty,-\infty),(-\infty,+\infty)\}$ and $B=\overline{\bR}^2\setminus A$, it is clear that $J$ is continuous on $B$ and hence measurable on it and as it is constant on $A$ it is trivially measurable on it so we are done.
        As compositions of measurable functions are measurable we are done.
    \end{enumerate}
\end{solution}
\begin{problem}\label{prob:2.3}
If $\{f_n\}$ is a sequence of measurable functions on $X$, then $\{x : \lim f_n(x) \text{ exists}\}$ is a measurable set.
\end{problem}
\begin{solution}
    We know that $g_1=\limsup f_n$ and $g_2=-\liminf f_n$ are measurable functions, define 
    $$h(x):=\begin{cases}2026\text{ if }g_1(x)=-g_2(x)=\pm\infty\\ g_1(x)+g_2(x)\text{ otherwise}\end{cases}$$
    $h$ is a measurable function by \hyperref[prob:2.2]{Problem 1.2} and $f_n$ converges iff $g_1=-g_2$ and we get that 
    $$\left\{x:\lim f_n(x)\text{ exists}\right\}=h^{-1}(\{0,2026\})$$
    which implies that the set infact is measurable as all the sets we performed unions and intersections on are in $\mM$.
\end{solution}
\begin{problem}\label{prob:2.4}
If $f: X \to \overline{\bR}$ and $f^{-1}((r, \infty]) \in \mM$ for each $r \in \bQ$, then $f$ is measurable.
\end{problem}
\begin{solution}
    It can easily be shown that these sets being in $\mM$ implies that $f^{-1}((k,\infty])\in\mM$ for any $k\in\overline{\bR}$ easily so we are done.
\end{solution}
\begin{problem}\label{prob:2.5}
If $X = A \cup B$ where $A, B \in \mM$, a function $f$ on $X$ is measurable iff $f|_A$ is measurable on $A$ and $f|_B$ is measurable on $B$.
\end{problem}
\begin{solution}
    If $f$ is measurable then measurability on $A,B\in\mM$ are trivial.
    For the other direction, suppsoe $f:X\to Y$ with $\mN$ being the sigma algebra on $Y$, for any $E\in\mN$ we have,
    $$f^{-1}(E)=(f^{-1}(E)\cap A)\cup(f^{-1}(E)\cap B)=(f\vert_A^{-1}(E))\cup(f\vert_B^{-1}(E))\in\mM$$
    as these two sets are in $\mM_A,\mM_B\subseteq\mM$ respectively anyways.
\end{solution}
\begin{problem}\label{prob:2.6}
The supremum of an uncountable family of measurable $\bR$-valued functions on $X$ can fail to be measurable (unless the $\sigma$-algebra $\mM$ is very special).
\end{problem}
\begin{solution}
    Suppose $X=\bR$ with some some sigma algebra containing singletons (most of the ones we care about do have them, for example the borel sigma algebra), then for any nonempty set $E\subseteq\bR$ consider the family $\{\chi_{\{r\}}\}_{r\in E}$, these are all measurable as singletons are in the sigma algebra whereas $\sup_{r\in E}\chi_{\{r\}}=\chi_E$ always being measurable would imply that the sigma algebra was the whole powerset stemming from $E$ being arbitrary.
\end{solution}
\begin{problem}\label{prob:2.7}
Suppose that for each $\alpha \in \bR$ we are given a set $E_\alpha \in \mM$ such that $E_\alpha \subset E_\beta$ whenever $\alpha < \beta$, $\bigcup_{\alpha \in \bR} E_\alpha = X$, and $\bigcap_{\alpha \in \bR} E_\alpha = \varnothing$. Then there is a measurable function $f: X \to \bR$ such that $f(x) \le \alpha$ on $E_\alpha$ and $f(x) \ge \alpha$ on $E_\alpha^c$ for every $\alpha$. (Use Exercise 4.)
\end{problem}
\begin{solution}
    We claim that the following function works
    $$f(x) := \sup\{r\in\bQ:x\notin E_r\}$$
    For $x\in E_\alpha$, from the definition we have that $x\in E_\beta$ for all $\bQ\ni\beta\geq\alpha$ and thus $f(x)=\sup\{r\in\bQ:r<\alpha, x\notin E_r\}\leq\alpha$. 
    Now, for $x\in E_\alpha^c$ its clearly $\geq\alpha$. 
    We claim that, for $r\in\bQ$ we have 
    $$f^{-1}((r,\infty))=\bigcup_{\bQ\ni\alpha> r}E_\alpha^c$$
    If $f(x)>r$, then the set of $\beta$ for which $x\in E^c_\beta$ (its an open ray unbounded below) must contain $r$, i.e its of form $(-\infty,v)$ with $v>r$ and we can find a rational $q\in(r,v)$ such that $x\in E_q$ so we have one direction of the equality, the other direction is trivial and we are done as the set on the right is in $\mM$ being a countable union of sets from $\mM$ which is sufficient to prove measurablity via \hyperref[prob:2.4]{Problem 2.4}
\end{solution}
\begin{problem}\label{prob:2.8}
If $f: \bR \to \bR$ is monotone, then $f$ is Borel measurable.
\end{problem}
\begin{solution}
    trivial.
\end{solution}
\begin{problem}\label{prob:2.9}
Let $f: [0, 1] \to [0, 1]$ be the Cantor function ($\S 1.5$), and let $g(x) = f(x) + x$.
\begin{enumerate}[label=\alph*.]
    \item $g$ is a bijection from $[0, 1]$ to $[0, 2]$, and $h = g^{-1}$ is continuous from $[0, 2]$ to $[0, 1]$.
    \item If $C$ is the Cantor set, $m(g(C)) = 1$.
    \item By Exercise 29 of Chapter 1, $g(C)$ contains a Lebesgue nonmeasurable set $A$. Let $B = g^{-1}(A)$. Then $B$ is Lebesgue measurable but not Borel.
    \item There exist a Lebesgue measurable function $F$ and a continuous function $G$ on $\bR$ such that $F \circ G$ is not Lebesgue measurable.
\end{enumerate}
\end{problem}
\begin{solution}
\begin{enumerate}[label=\alph*.]
    \item It is known that the cantor function is continuous on $[0,1]$ as well as being increasing. 
    This makes $g(x)=f(x)+x$ strictly increasing and continuous, with this we can say using the intermediate value theorem that $g$ is is a continuous bijection from $[0,1]$ to $[0,2]$ and $h=g^-1$ is also continuous. 
    \item As the cantor set is borel and $h$ is continuous, it must be borel measurable. 
    Thus, $h^{-1}(C)=g(C)$ is borel (as well as being lebesgue measurable which is weaker). 
    We can write $m(g([0,1]))=m(g(C))+m(g([0,1]\setminus C))$ as the sets here are disjoint (the second one is clearly measurable). 
    From the previous parts we see that $m(g(C))=2-m(g([0,1]\setminus C))$. 
    This set is the disjoint union of the middle third intervals removed in the process of constructing the cantor set, on which, by defintion the cantor function is constant. 
    These removed open intervals have measures summing to $1$ and as $g$ is a translation on these we see that $m(g([0,1]\setminus C))=1$ showing that $m(g(C))=1$. 
    \item As the lebesge measure is complete we see that $B$ being the subset of the cantor set which has measure zero, must be lebesgue measurable too. 
    Now, if $B$ was borel then so would be $h^-1(B)=A$ as $g^{-1}$ is continuous, so it is not borel. 
    \item Consider $\chi_{B}\circ h$, we see that $(\chi_B\circ h)^{-1}(\{1\})=h^{-1}(\chi_B^{-1}(\{1\}))=h^{-1}(B)=A$ which is not lebesgue measurable while we had $h$ be continuous (a homeomorphism even) and $\chi_B$ lebesgue measurable so we are done as $\{1\}$ is borel.

\end{enumerate}
\end{solution} 

\begin{problem}\label{prob:2.10}
Prove that the following implications are valid iff the measure $\mu$ is complete:
\begin{enumerate}[label=\alph*.]
    \item If $f$ is measurable and $f=g$ $\mu$-a.e. then $g$ is measurable.
    \item If $f_n$ is measurable for $n\in\bN$ and $f_n\longrightarrow f$ $\mu$-a.e. then $f$ is measurable. 
\end{enumerate}
\end{problem}
\begin{solution}
    \begin{enumerate}[label=\alph*.]
        \item First assume that $\mu$ is complete and that $f$ and $g$ disagree on some $N\in\mM$ with $\mu(N)=0$, then for any measurable $E$ in the range we must have that 
        $$g^{-1}(E)=(g^{-1}(E)\cap N^c)\cup(g^{-1}(E)\cap N)=(g|_{N^c}^{-1}(E))\cup(g^{-1}(E)\cap N)$$
        $$=(\underbrace{f|_{N^c(\in\mM)}^{-1}(E)}_{\in\mM})\cup(\underbrace{g^{-1}(E)\cap N}_{\subseteq N\implies\in\mM})$$
        and thus this set must be in $\mM$ itself. 
        For the other direction, suppose $\mu(N)=0$ for some $N\in\mM$ and $K\subseteq N$. 
        Let, $f=\chi_N$ and $g=\chi_K$, we see that $f=g$ on $N^c$ and thus $f=g$ $\mu$-a.e. which makes $g$ measurable and thus $g^{-1}(\{1\})=K$ must also be measurable making $\mu$ complete. 
        \item For one direction assume that $\mu$ is complete. 
        By \hyperref[prob:2.3]{Problem 2.3} we know that the set where we have convergence, is measurable so let it be $E$. 
        Further, a.e. convergence implies that $\mu(X\setminus E)=0$. 
        Let $g=\lim\chi_E f_n(x)$, we see that $g$ is measurable on $E$ as well as on $E^c$ and is thus measurable on $X$, being the limit (it always exists) of measurable functions. 
        Suppose, $N$ is the set where $f_n\not\longrightarrow f$, by our hypothesis this is of measure zero. 
        For $x\in N^c$ however $f_n(x)\longrightarrow f(x)$ and thus $f(x)=g(x)$ so $f(x)\neq g(x)$ can only happen on $N$ implying that $f=g$ a.e after which part a proves measurablity of $f$ and we are done. 
        For the other direction, suppose $E$ is a set of measure zero and $K\subseteq E$. 
        Let, $f=\chi_K$ and $f_n=0$, with this we see that $f_n\to f$ on $E^c$ which shows that $f_n\to f$ a.e. and thus $f$ is measurable so we are done. 
    \end{enumerate}
\end{solution}
\begin{problem}\label{prob:2.11}
Suppose that $f$ is a function on $\bR \times \bR^k$ such that $f(x, \cdot)$ is Borel measurable for each $x \in \bR$ and $f(\cdot, y)$ is continuous for each $y \in \bR^k$. For $n \in \bN$, define $f_n$ as follows. For $i \in \bZ$ let $a_i = i/n$, and for $a_i \le x \le a_{i+1}$ let
\[ f_n(x, y) = \frac{f(a_{i+1}, y)(x - a_i) - f(a_i, y)(x - a_{i+1})}{a_{i+1} - a_i} \]
Then $f_n$ is Borel measurable on $\bR \times \bR^k$ and $f_n \to f$ pointwise; hence $f$ is Borel measurable on $\bR \times \bR^k$. Conclude by induction that every function on $\bR^n$ that is continuous in each variable separately is Borel measurable.
\end{problem}
\begin{solution}
    We can write $\bR\times\bR^k$ as a countable union of the borel sets $\{[a_i,a_{i+1}]\times\bR^k\}_{i\in\bZ}$ and proving borel measurablity of $f_n$ on these suffices to prove borel measurablity globally. 
    It being borel on $[a_i,a_{i+1}]\times\bR^k$ follows from these maps being clearly borel 
    $$(x,y)\mapsto (x-a_i),(x,y)\mapsto (x-a_{i+1}),(x,y)\mapsto f(a_i,y),(x,y)\mapsto f(a_{i+1},y)$$ 
    as the borel measurability is preserved under products and sums. 
    To prove pointwise convergence just notice that for $(x,y)\in[a_i,a_{i+1}]\times\bR^k,f_n(x,y)$ is a convex combination of $f(a_i,y)$ and $f(a_{i+1},y)$ which converge to $f(x,y)$ as $n\to\infty$ by continuity of $f(\cdot,y)$. 
    This shows that if $f(\cdot,y)$ is continuous and $f(x,\cdot)$ borel, for all $(x,y)\in\bR\times\bR^k$ then $f$ is borel. 
    Now, suppose that $f:\bR^n\to\bR$ is continuous in each variable, suppose that the proposition holds for $n-1$, then for all $x\in\bR,f(x,\cdot):y\in\bR^{n-1}\mapsto f(x,y)$ is borel because this is continuous in each variable as well and for all $y\in\bR^{n-1},f(\cdot,y):x\in\bR\mapsto f(x,y)$ is continuous too, which implies that $f$ is borel. 
    We can induct on $n$ from the base case of $n=1$ where the proposition is true due to continuity in each variable coinciding with continuity (which is stronger than borel measurability), hence proving the proposition. 
\end{solution}
\begin{problem}\label{prob:2.12}
Prove Proposition 2.20. (See Proposition 0.20, where a special case is proved.)
\end{problem}
\begin{solution}
    If $f^{-1}(\{\infty\})$ is not null then the integral will trivially be infinite and thus it must be a null set. 
    Now, consider the sets $H_n=f^{-1}(1/n,\infty]$ and we find that $\{x:f(x)>0\}=\cup H_n$. 
    Here, we see that $f\geq n^{-1}\chi_{H_n}$ for all $n\geq 1$ and thus $\int f\geq n^{-1}\mu(H_n)$ forces all the $H_n$ to have finite measure. 
\end{solution}
\begin{problem}\label{prob:2.13}
Suppose $\{f_n\} \subset L^+$, $f_n \to f$ pointwise, and $\int f = \lim \int f_n < \infty$. Then $\int_E f = \lim \int_E f_n$ for all $E \in \mM$. However, this need not be true if $\int f = \lim \int f_n = \infty$.
\end{problem}
\begin{solution}
    For any $E\in\mM$ we can see that $f_n\chi_E\longrightarrow f\chi_E$ pointwise as well and thus, by fatou's lemma and the additivity of the integral we have, 
    $$\int f\chi_E\leq\liminf\int f_n\chi_E\text{ and }\int f\chi_{X\setminus E}\leq\liminf\int f_n\chi_{X\setminus E}$$
    $$\implies\int f=\int f\chi_E+\int f\chi_{X\setminus E}\leq\liminf\int f_n\chi E+\liminf\int f\chi_{X\setminus E}$$ 
    Using the fact that the limit infimum is subadditive we see that the above is
    $$\leq\liminf\left(\int f_n\chi_E+f_n\chi_{X\setminus E}\right)=\liminf\int f_n=\int f$$ 
    as all the terms involved in this string of inequalities are finite (this follows from the fact that $\int f<\infty$ and that $\int f_n$ is close to $\int f$ eventually) we must have that, 
    $$\int  f\chi_E=\liminf\int f_n\chi_E$$
    Assume wlog $\int f_n$ is always finite, then by the superadditivity of the limit supremum and the equality we developed earlier, we have
    $$\limsup\int f_n\chi_{E}=\limsup\left(\int f_n-\int f_n\chi_{E^c}\right)\leq\limsup\int f_n+\limsup\left(-\int f_n\chi_{E^c}\right)$$
    $$=\limsup\int f_n-\liminf\int f_n\chi_{E^c}=\int f-\int f\chi_{E^c}=\int f\chi_E$$ 
    $$\limsup\int f_n\chi_E\leq\int f\chi_E=\liminf f_n\chi_E\implies\lim\int f_n\chi_E=\int f\chi_E$$  
    For a counterexample we can consider a function on $\bR$ given by $f_n=n\chi_{(0,1/n)}+\chi_{[1,\infty)}$ which converges pointwise to $\chi_{[1,\infty)}$ and $\int f=\lim\int f_n=\infty$ but if we focus on $(0,1)$ then $\lim f_n\chi_{(0,1)}=1$ whereas $\int f\chi_{(0,1)}=0$. 
    Here the area escaped to infinity in the unit open interval and hence the contradiction. 
\end{solution}
\begin{problem}\label{prob:2.14}
If $f \in L^+$, let $\lambda(E) = \int_E f \, d\mu$ for $E \in \mM$. Then $\lambda$ is a measure on $\mM$, and for any $g \in L^+$, $\int g \, d\lambda = \int gf \, d\mu$. (First suppose that $g$ is simple.)
\end{problem}
\begin{solution}
    For any disjoint collection $\{E_n\}_{n\in\bN}\subseteq\mM$ with union $E$ then by the sigma additivity of the integral we have, 
    $$\lambda(E)=\int f\chi_Ed\mu=\int\sum_{n\in\bN}f\chi_{E_n}d\mu=\sum_{n\in\bN}\int f\chi_{E_n}d\mu=\sum_{n\in\bN}\lambda(E_n)$$ 
    Now, suppose $\phi=\sum a_j\chi_{E_j}$ is a simple function, then we see that 
    $$\int\phi d\lambda=\sum a_j\lambda(E_j)=\sum a_j\int f\chi_{E_j}d\mu=\int\sum a_jf\chi_{E_j}d\mu=\int f\phi d\mu$$ 
    Now, if $g\in L^+,$ we can find a sequence of simple functions $\{\phi_n\}$ that increase to $g$, for these we also see that $f\phi_n$ iincreases to $fg$ and by the monotone convergence theorem and the fact that $\lambda$ and $\mu$ are defined on the same sigma algebra $\mM$ we have that, 
    $$\int fg d\mu=\int\lim f\phi_n d\mu\longleftarrow\int f\phi_n d\mu=\int\phi_n\lambda\longrightarrow \int\lim\phi_n\lambda=\int g\lambda$$ 
    where we used the the monotone convergence theorem on both sides.
\end{solution}
\begin{problem}\label{prob:2.15}
If $\{f_n\} \subset L^+$, $f_n$ decreases pointwise to $f$, and $\int f_1 < \infty$, then $\int f = \lim \int f_n$.
\end{problem}
\begin{solution}
    Clearly we have $\int f\leq\lim\int f_n<\infty$. 
    For the other direction, consider the measurable set $E=\{f>0\}$. 
    Fix $\alpha>1$ and define $F_n:=\{x\in E:f_n(x)\leq\alpha f(x)\}$, then $\{F_n\}$ is an increasing sequence of measurable sets and $\cup F_n=E$. 
    From the definiton we see that $\int_{F_n}f_n\leq\alpha\int_{F_n}f$ and $\int_E f_n-\int_{F_n}f_n=\int_{E\setminus F_n}f_n\leq\int_{E\setminus F_n}f_1$. 
    As $A\mapsto\int_A f_1$ is a measure on $\mM$ and $E\setminus F_n$ is a decreasing sequence with $\int_{E\setminus F_1}f_1\leq\int f_1<\infty$ we see that 
    $$\lim_{n\to\infty}\int_{F\setminus F_n}f_1=\int_{\cap E\setminus F_n}f_1=\int_{E\setminus E}f_1=0$$
    this implies that $\lim\int_{F_n}f_n=\lim\int_E f_n$. 
    As $A\mapsto\int_Af$ is also a measure, we have 
    $$\lim\int_{F_n}f_n=\lim\int_E f_n\implies\lim\int_Ef_n\leq\lim\alpha\int_{F_n} f=\alpha\lim\int_{F_n}f=\alpha\int_{\cup F_n}f=\int_Ef$$ 
    As this is true for any $\alpha>1$ we have $\lim\int_Ef_n\leq\int_Ef$. 
    Now, fix $\alpha>0$ and consider $E^c$ on which $f$ is zero and thus $f_n$ decreases to $0$ here. 
    Let $G_n:=\{x\in E^c:f_n(x)\leq \alpha f_1(x)\}$ see that $\{G_n\}$ is an increasing sequence of measurable sets with $\cup G_n=E^c$. 
    It can also be shown that $\lim\int_{G_n}f_n=\lim\int_{E^c}f_n$ using similar arguements as we did when working over $E$ and from the definition we see that $\int_{G_n}f_n\leq\alpha\int_{G_n}f_1\leq\alpha\int f_1$ and thus, $\lim\int_{E^c}f_n=0=\int_{E^c}f$ as $\alpha>0$ was arbitrary. 
    Combining these two deductions, namely $\lim\int_Ef_n\leq\int_Ef$ and $\lim\int_{E^c}f_n=0=\int_{E^c}f$ we see that $\lim\int f_n\leq\int f$ as the integral is additive. 
    Having shown both directions of the inequality we must have $\int f=\lim\int f_n$. 
    
\end{solution}
\begin{problem}\label{prob:2.16}
If $f \in L^+$ and $\int f < \infty$, for every $\epsilon > 0$ there exists $E \in \mM$ such that $\mu(E) < \infty$ and $\int_E f > (\int f) - \epsilon$.
\end{problem}
\begin{solution}
From \hyperref[prob:2.12]{Problem 2.12} we see that the set $E=\{x:f(x)>0\}$ is sigma finite so we write it as a disjoint(standard technique) union of countably many sets of finite measure as $E=\cup F_n$, now from the additivity of the integral we have, 
$$\int f=\int f\chi_E+\underbrace{\int f\chi_{E^c}}_{=0}=\int_E f=\sum_{n\in\bN}\int_{F_n}f=\lim_{n}\sum_{k=1}^n\int_{F_n}f=\lim_n\int_{\cup_{k\leq n}F_k}f$$ 
Now, for any $\epsilon>0$ we can find a $N$ such that $\int f<\epsilon+\int_{\cup_{k\leq N}F_k}f$ and as $\mu(\cup_{k\leq N}F_k)=\sum_{k\leq N}\mu(F_k)<\infty$ we are done. 
\end{solution}
\begin{problem}\label{prob:2.17}
Assume Fatou's lemma and deduce the monotone convergence theorem from it.
\end{problem}
\begin{solution}
    Fatou's lemma implies that $\lim\int f_n=\liminf\int f_n\geq\int\liminf f_n=\int f$ while the other direction of the inequality is trivial. 
\end{solution}
\begin{problem}\label{prob:2.18}
Fatou's lemma remains valid if the hypothesis that $f_n \in L^+$ is replaced by the hypothesis that $f_n$ is measurable and $f_n \ge -g$ where $g \in L^+ \cap L^1$. What is the analogue of Fatou's lemma for nonpositive functions?
\end{problem}
\begin{solution}
Clearly $f_n+g\in L^+$ so we can apply fatou's lemma on it and get
$$\int g+\liminf\int f_n=\liminf\qty(\int f_n+\int g)=\liminf\int g+f_n$$
$$\geq\int\liminf(f_n+g)=\int g+\liminf f=\int g+\int \liminf f_n$$
and now we can just cancel out the $\int g$ terms. 
For nonpositive functions $f$ ($-f\in L^+$) we will have that
$$\limsup\int f_n\leq\int\limsup f_n$$
\end{solution}
\begin{problem}\label{prob:2.19}
Suppose $\{f_n\} \subset L^1(\mu)$ and $f_n \to f$ uniformly.
\begin{enumerate}[label=\alph*.]
    \item If $\mu(X) < \infty$, then $f \in L^1(\mu)$ and $\int f_n \to \int f$.
    \item If $\mu(X) = \infty$, the conclusions of (a) can fail. (Find examples on $\bR$ with Lebesgue measure.)
\end{enumerate}
\end{problem}
\begin{solution}
\begin{enumerate}[label=\alph*.]
\item This is trivial.
\item consider $f_n=n^{-1}\chi_{[0,n^2]}$ on $\bR$ with respect to the lebesgue measure.
\end{enumerate}
\end{solution}
\begin{problem}\label{prob:2.20}
(A generalized Dominated Convergence Theorem) If $f_n, g_n, f, g \in L^1$, $f_n \to f$ and $g_n \to g$ a.e., $|f_n| \le g_n$, and $\int g_n \to \int g$, then $\int f_n \to \int f$. (Rework the proof of the dominated convergence theorem.)
\end{problem}
\begin{solution}
We see that $g_n+f_n,g_n-f_n\in L^1$ and also as $f_n\to f,g_n\to g$ a.e. we can restrict the integrals to a set $E$ where convergence is global such that $\int=\int_E$ (as operators) and we work inside $E$ now without mention. 
By fatou's lemma
$$\int g+\liminf\int f_n=\liminf\int f_n+\lim\int g_n=\liminf\qty(\int f_n+\int g_n)=\liminf\int f_n+g_n$$
$$\geq\int\liminf(g_n+f_n)\geq\int\liminf g+\liminf f_n=\int f+g=\int f+\int g$$
cancelling out the $\int g$ terms we get $\liminf\int f_n\geq \int f$ and doing the same with $g_n-f_n$ we get $-\limsup\int f_n=\liminf\int-f_n\geq\int -f=-\int f$ which implies that $\int f\geq\limsup\int f$ and thus finally we have $\limsup\int f_n\leq\int f\leq\liminf\int f_n\implies\lim\int f_n=\int f$.
\end{solution}
\begin{problem}\label{prob:2.21}
Suppose $f_n, f \in L^1$ and $f_n \to f$ a.e. Then $\int |f_n - f| \to 0$ iff $\int |f_n| \to \int |f|$. (Use Exercise 20.)
\end{problem}
\begin{solution}
	For the if direction we can use generalised dominated convergence (\hyperref[prob:2.20]{Problem 2.20}) on $|f_n-f|\to 0$ a.e. where it is dominated by $|f_n|+|f|\to 2|f|$ a.e. as $\int |f_n|+|f|=\int|f|+\int|f_n|\to 2\int|f|=\int2|f|$ and for the only if direction we can proceed by the reverse triangle inequality as follows
\[
	\qty|{\int |f_n|-\int|f|}|\leq\int||f_n|-|f||\leq\int|f_n-f|\longrightarrow 0
\]
\end{solution}
\begin{problem}\label{prob:2.22}
Let $\mu$ be counting measure on $\bN$. Interpret Fatou's lemma and the monotone and dominated convergence theorems as statements about infinite series.
\end{problem}
\begin{solution}
    A sequence of functions $\{f_k\}$ on $\bN$ corresponds to a double sequence $a_{k,n} = f_k(n)$. Pointwise convergence $f_k \to f$ a.e. means that for each fixed $n$, $\lim_{k\to\infty} a_{k,n} = a_n$ (since the only set of measure zero for the counting measure is the empty set).
With this we can interpret the three theorems as follows: \\
\textit{Fatou's Lemma:} If $f_k \in L^+$, then $\int \liminf f_k \le \liminf \int f_k$.
In series notation: If $a_{k,n} \ge 0$ for all $k,n \in \bN$, then
$$\sum_{n=1}^\infty \left( \liminf_{k\to\infty} a_{k,n} \right) \le \liminf_{k\to\infty} \sum_{n=1}^\infty a_{k,n}$$
\textit{Monotone Convergence Theorem:} If $\{f_k\} \subset L^+$, $f_k \le f_{k+1}$, and $f_k \to f$, then $\lim \int f_k = \int f$.
In series notation: If $0 \le a_{1,n} \le a_{2,n} \le \dots$ for all $n$, and $\lim_{k\to\infty} a_{k,n} = a_n$, then
$$\lim_{k\to\infty} \sum_{n=1}^\infty a_{k,n}=\sum_{n=1}^\infty a_n$$
\textit{Dominated Convergence Theorem:} If $f_k \to f$ pointwise, and $|f_k| \le g$ where $g \in L^1(\mu)$, then $\lim \int f_k = \int f$.
In series notation: If $\lim_{k\to\infty} a_{k,n} = a_n$ for each $n$, and there exists a bounding sequence $\{b_n\}$ such that $|a_{k,n}| \le b_n$ for all $k,n$, and $\sum_{n=1}^\infty b_n < \infty$, then
$$\lim_{k\to\infty} \sum_{n=1}^\infty a_{k,n}=\sum_{n=1}^\infty a_n$$
\end{solution}
\begin{problem}\label{prob:2.23}
Given a bounded function $f: [a, b] \to \bR$, let
\[ H(x) = \lim_{\delta \to 0} \sup_{|y-x|\le\delta} f(y), \quad h(x) = \lim_{\delta \to 0} \inf_{|y-x|\le\delta} f(y). \]
Prove Theorem 2.28b by establishing the following lemmas:
\begin{enumerate}[label=\alph*.]
    \item $H(x) = h(x)$ iff $f$ is continuous at $x$.
    \item In the notation of the proof of Theorem 2.28a, $H = G$ a.e. and $h = g$ a.e. Hence $H$ and $h$ are Lebesgue measurable, and $\int_{[a, b]} H \, dm = \overline{I}_a^b(f)$ and $\int_{[a, b]} h \, dm = \underline{I}_a^b(f)$.
\end{enumerate}
\end{problem}
\begin{solution}
\begin{enumerate}[label=\alph*.]
	\item This is trivial.
	\item If it was continuous outside a null set then we can show riemann integrablity by considering partitions of form $P=\{a_1<b_1<\ldots<a_n<b_n\}\cup P_\delta$ where, by our assumptions $m(E=\cup(a_i,b_i))=\epsilon$ is small such that $f$ is continuous on the compact set $[a,b]\setminus E$ and $P_\delta$ is of small enough mesh such as to make $U(f|_{[a,b]\setminus E},P_\delta)-L(f|_{[a,b]\setminus E})<\delta$.
	Now, as $f$ is bounded there is some $M$ such that $|f|\leq M$ then
	$$U(f,P)-L(f,P)\leq 2M\epsilon+\delta$$
	which can be made arbitrarily small so $f$ must be integrable.
	Now we will prove the converse. If $f$ is integrable then there exists an increasing sequence of finite partitions $P_k$ such that $\norm{P_k}\to 0$ and $\lim\int G_{P_k}dm =\int f=\lim\int g_{P_k}dm$. It was already shown in the book that $g=G$ a.e. when $f$ is riemann integrable so if we show that $H=G$ a.e. and $h=g$ a.e. then we would be done using part (a) as this implies $H=h$ a.e. which is equivalent to continuity a.e. by part (a). 
	We will show that $H=G$ a.e. and the rest will follow similarly. Let $K$ be the set of all endpoints of some interval in $P_k$ for some $k$, this is a null set and we will show that $H=G$ outside it. 
	Say $x\in[a,b]\setminus K$ is any point then we first show that $H\geq G.$ 
First define $H_\delta(x):=\sup_{|x-y|\leq\delta}f(y)$ then we see that $H=\lim_{\delta\to 0}H_\delta$ by definition and $H_\delta\leq H_{\delta'}$ for all $0\leq\delta\leq\delta'$. 
Now, for any $\delta>0$ we can find some $k$ such that $x\in(a_k,b_k)\subseteq(x-\delta,x+\delta)$ where $(a_k,b_k)$ is an interval in $P_k$ as $\norm{P_k}\to 0$. 
Then we must have that $G_{P_k}(x)\leq H_\delta(x)$ and as $G_{P_k}$ is decreasing we have that $G(x)\leq H_\delta(x)$ and then taking the limit as $\delta\searrow 0$ we have $G(x)\leq H(x)$. 
To show the other direcrion of the equality we notice that for any $P_k$ with the interval $(a_k,b_k)$ containing $x$ within it we can find some $\delta>0$ such that $(x-\delta,x+\delta)\subseteq(a_k,b_k)$ and this implies that $G_{P_k}(x)\geq H_\delta(x)$ and as $H_\delta(x)$ is decreasing in delta, taking the limit we see that $G_{P_k}(x)\geq H(x)$ and again takinng the limit as $k\to\infty$ we see that $G\geq H$ so we are done.
\end{enumerate}
\end{solution}
\begin{problem}\label{prob:2.24}
Let $(X, \mM, \mu)$ be a measure space with $\mu(X) < \infty$, and let $(X, \overline{\mM}, \overline{\mu})$ be its completion. Suppose $f: X \to \bR$ is bounded. Then $f$ is $\overline{\mM}$-measurable (and hence in $L^1(\overline{\mu})$) iff there exist sequences $\{\phi_n\}$ and $\{\psi_n\}$ of $\mM$-measurable simple functions such that $\phi_n \le f \le  \psi_n$ and $\int (\psi_n - \phi_n) \, d\mu < n^{-1}$. In this case, $\lim \int \phi_n \, d\mu = \lim \int \psi_n \, d\mu = \int f \, d\overline{\mu}$.
\end{problem}
\begin{solution}
To show the if direction we can consider the functions $g:=\sup_n\phi_n$ and $h:=\inf_n\psi_n$then these are $\mM$-measurable and we have that $\int g-h<1/n$ for all $n$ and thus it is zero which makes $h=g$ a.e. and thus $f=g=h$ $\mu$-a.e. 
Now, If $h,g$ are $\mM$ measurable then they are clearly $\overline{\mM}$ measurable as well and as $\mM\subseteq\overline{\mM}$ we see that $\overline{\mu}$-a.e. too so we must have that $h=f=g$ $\overline{\mu}$-a.e as well. which implies that $f$ is $\overline{\mM}$ measurable. For the converse, we first deal with the case where $f\geq 0$ consider the following simple functions which are $\overline{\mM}$ measurable,
$$\phi_n^*:=\sum_{k=0}^{2^{2n}-1}\frac{k}{2^n}\cdot\chi_{f^{-1}\qty([k2^{-n},(k+1)2^{-n}))}$$
$$\psi_n^*:=\sum_{k=0}^{2^{2n}-1}\frac{k+1}{2^n}\cdot\chi_{f^{-1}\qty([k2^{-n},(k+1)2^{-n}))}$$
Then, for all large enough $n$ we have that $\psi_n^*\geq f\geq \phi_n^*$ and $0\leq f-\phi_n^*,\psi_n^*-f\leq 2^{-n}$. 
As $f$ is bounded we can find some constant $M$ with $f\leq M$. 
For brevity write
$$\phi_n^*=\sum_{k=1}^ma_k\chi_{E_k\cup F_k}\text{ where }E_i\cup F_i\text{ are disjoint }E_i\in\mM,F_i\subseteq N_i\in\mM\text{ with }\mu(N_i)=0$$
Then define, $\phi_n:=\sum_{k=1}^ma_k\chi_{E_k}$ and we have $\phi_n\leq\phi_n^*$ with equality $\overline{\mu}$ a.e. Again writing $\psi_n^*$ in the exact same notation with $N=\cup N_i\in\mM$ we see that 
$$\psi_n:=\chi_{N^c}\psi_n^*+(M+1)\chi_N=(M+1)\chi_N+\sum_{k=1}^ma_k\chi_{E_m\cap N^c}$$
is such that $\psi_n\geq\psi^*_n$ with equality with equality $\overline{\mu}$ a.e. (We use $M+1$ instead of $M$ to guarantee the inequality, we need to do this due to how $\psi_n^*\leq f+2^{-n}\leq M+1$ for large enough $n$). 
We now have $\mM$ measurable simple functions $\psi_n,\phi_n$ such that $\psi_n\geq f\geq \phi_n$ and also $\psi_n=\psi^*_n\text{ and }\phi_n=\phi_n^*$ $\overline{\mu}$-a.e. (both together), from this we can say that $\psi_n-f,f-\phi_n\leq 2^{-n}$ together $\overline{\mu}$-a.e. implying $\psi_n-\phi_n\leq 2^{-n+1}$ $\overline{\mu}$-a.e. and thus,
$$\int \psi_n-\phi_nd\overline{\mu}=\int\psi_n-\phi_nd\mu\leq\int2^{-n+1}d\mu=\frac{\mu(X)}{2^{-n+1}}\longrightarrow 0$$
and we can start from a larger index to have $\int\psi_n-\phi_n<n^{-1}.$ 
In the previous line we assumed the following: if $g$ is some $\mM$ measurable function then $\int gd\mu=\int gd\overline{\mu}$, this is easy to prove and we omit the proof. 
We also see that $\phi_n\leq f\leq \phi_n+2^{-n}$ $\overline{\mu}$-a.e. for large enough $n$ and thus for large $n$ we have $\int\phi_nd\mu\leq\int fd\overline{\mu}\leq\mu(X)\cdot2^{-n}+\int\phi_nd\mu$ using the assumption from earlier along the way which implies that $\lim\int\phi_nd\mu=\int fd\overline{\mu}$. 
The same can be proved for $\psi_n$ similarly so we are done with the proof for $f\geq 0$. 
Now, for general such $f$ we first split it into its measurable (and bounded in this case  positive and negative parts, in usual notation $f=f_+-f_-$. 
Let $\phi_n^+,\psi_n^+$ be as in the previous case, for $f_+$ and $\phi_n^-,\psi_n^-$ be those for $f_-$, then we see that $\phi_n^+-\psi_n^-\leq f_+-f_-\leq\psi_n^+-\phi_n^-$. 
As $|\phi_n^+-\psi_n^--(f_+-f_-)|\leq|\phi_n^+-f_+|+|\psi_n^--f_-|\leq 2^{-n}+2^{-n}$ $\overline{\mu}$-a.e for large enough $n$ implying 
$$\phi_n^+-\psi_n^-\leq f_+-f_-\leq\phi_n^+-\psi_n^-+2^{-n+1}$$
and we can get the corresponding bounds for $\psi_n^+-\phi_n^-$ as well and finish like in the previous case using the $\mM$ measurable simple functions functions $\zeta_n:=\phi_n^+-\psi_n^-$ and $\gamma_n:=\psi_n^+-\phi_n^-$ with $\zeta_n\leq f\leq\gamma_n$ and $\gamma_n-\zeta_n\leq 2^{-n+2}$ holding $\overline{\mu}$-a.e.

\end{solution}
\begin{problem}\label{prob:2.25}
Let $f(x) = x^{-1/2}$ if $0 < x < 1, f(x) = 0$ otherwise. Let $\{r_n\}_1^\infty$ be an enumeration of the rationals, and set $g(x) = \sum_1^\infty 2^{-n} f(x - r_n)$.
\begin{enumerate}[label=\alph*.]
    \item $g \in L^1(m)$, and in particular $g < \infty$ a.e.
    \item $g$ is discontinuous at every point and unbounded on every interval, and it remains so after any modification on a Lebesgue null set.
    \item $g^2 < \infty$ a.e., but $g^2$ is not integrable on any interval.
\end{enumerate}
\end{problem}
\begin{solution}
\begin{enumerate}[label=\alph*.]
\item It follows from the monotone convergence theorem and the fact that the lebesgue integral and riemann integral agree that $\int gdm=\sum_{n=1}^\infty\int 2^{-n}f(x-r_n)dm=\sum_{n=1}^\infty2^{-n}\cdot \int_0^1x^{-1/2}dx=2$ implying $g\in L^1(m)$ which again forces $\{g\geq\infty\}$ to be null.
\item given any real $y$, if $r_n<y<1+r_n$ then we see that $g(y)>2^{-n}(y-r_n)^{-1/2}$. Now, given any $x\in\bR$ and small enough $\delta>0$ we can find some $r_n$ such that $x-\delta<r_n<x<x+\delta<r_n+1$ and thus, for all small enough $\epsilon>0$ with $r_n<r_n+\epsilon<x$ we see that $g(r_n+\epsilon)>2^{-n}\epsilon^{-2}$ which can be made arbitrarily large, this shows that $g$ is unbounded in any open neighbourhood around $x$ from whcih the claim follows.
\item As $g<\infty$ a.e. we clearly have $g^2<\infty$ a.e. as well. With the same setup as the previous part, we see that $g(y)^2>4^{-n}(y-r_n)^{-1}$ and thus 
	\[
		\int_{x-\delta}^{x+\delta}g^2\mathrm{d}m\geq\int_{r_n}^{x} \frac{1}{2^n(y-r_n)} \mathrm{d}m(y) = \infty
	\]
thus, its not integrable on any interval of form $(x-\delta,x+\delta)$ proving the claim after noticing how all intervals have an interval of this form.
	
\end{enumerate}
\end{solution}
\begin{problem}\label{prob:2.26}
If $f \in L^1(m)$ and $F(x) = \int_{-\infty}^x f(t) \, dt$, then $F$ is continuous on $\bR$.
\end{problem}
\begin{solution}
	First we prove right continuity using the sequential criterion. Say $x_n\searrow x$, then as $A\mapsto\int\chi_A|f|$ is a finite measure on $\mathcal L$ we see that 
	\[
		\lim\qty|F(x_n)-F(x)|=\lim\qty| \int_{[x,x_n]}f\mathrm{d}m|\leq\lim\int_{[x,x_n]}|f|\mathrm{d}m=\int_{\cap[x,x_n]}|f|dm=0
	\]
	as $\cap[x,x_n]=\{x\}$ is a lebesge null set and we follow the convention of $0\cdot\infty=0$. Right continuity is proven similarly and hence we have continuity.
\end{solution}
\begin{problem}\label{prob:2.27}
Let $f_n(x) = ae^{-nax} - be^{-nbx}$ where $0 < a < b$.
\begin{enumerate}[label=\alph*.]
    \item $\sum_1^\infty \int_0^\infty |f_n(x)| \, dx = \infty$.
    \item $\sum_1^\infty \int_0^\infty f_n(x) \, dx = 0$.
    \item $\sum_1^\infty f_n \in L^1([0, \infty), m)$, and $\int_0^\infty \sum_1^\infty f_n(x) \, dx = \log(b/a)$.
\end{enumerate}
\end{problem}
\begin{solution}
    This is just a bunch of calculations. However it furnishes an example to how how $\int\sum|f_n|\notin L^1$ might result in $\int\sum f_n\neq\sum\int f_n$ even whe both are finite and $\sum f_n\in L^1$
\end{solution}

\begin{problem}\label{prob:2.28}
Compute the following limits and justify the calculations:
\begin{enumerate}[label=\alph*.]
    \item $\lim_{n \to \infty} \int_0^\infty (1 + (x/n))^{-n} \sin(x/n) \, dx$.
    \item $\lim_{n \to \infty} \int_0^1 (1 + nx^2)(1 + x^2)^{-n} \, dx$.
    \item $\lim_{n \to \infty} \int_0^\infty n \sin(x/n) [x(1 + x^2)]^{-1} \, dx$.
    \item $\lim_{n \to \infty} \int_a^\infty n(1 + n^2x^2)^{-1} \, dx$. (The answer depends on whether $a > 0$, $a = 0$, or $a < 0$. How does this accord with the various convergence theorems?)
\end{enumerate}
\end{problem}
\begin{solution}
    \begin{enumerate}[label=\alph*.]
		\item Let, $f_n(x)=(1+x/n)^{-n}\sin(x/n)$ then it converges $e^{-x}\cdot\sin(0)=0$ pointwise and it can be checked that $|f_n|\leq (1+x/2)^{-2}$ for all $n\geq 2$ and as this is $L^1$, it allows us to conclude that the limit is $0$ by DCT.
		\item Let $f_n(x)=(1+nx^2)/(1+x^2)^n$ then it converges to $\chi_{\{1\}}$ pointwise and as $|f_n|\leq 1$ in the interval $[0,1]$ we can apply DCT to conclude that the limit is $0$
		\item In a similar fashion we see that the integrands are dominated by $(1+x^2)^{-1}$ and converge pointwise to it as well after which we can conclude with DCT that the limit is $\pi/2$ using basic calculus.
		\item If $a>0$ then we see that the integrands are eventually dominated by $(1+x^2)^{-1}$ as $1/n+nx^2\geq 1+x^2\iff x^2>1/n$ which is always true in our case for large enough $n$ allowing us to apply DCT and conclude that the limit is zero as the integrands converge to $0$ pointwise. However, if $a\leq 0$ then the integrands converge to $\infty\cdot\chi_{\{0\}}$ pointwise which integrates to $0$ whereas the limit is that of the terms $\pi/2-\arctan(na)$ which has a strictly positive limit here for the $a\leq 0$ we consider. This shows us that DCT \textit{can't be applied} here, in fact it is easily proven without the DCT that any function which dominates this sequence can not be $L^1$ by considering the values on intervals of form $[1/n,1/(n-1)]$ which is left to the reader as an exercise.
    \end{enumerate}
\end{solution}
\begin{problem}\label{prob:2.29}
	Show that $\int_0^\infty x^n e^{-x} \, dx = n!$ by differentiating the equation $\int_0^\infty e^{-tx} \, dx = 1/t$. Similarly, show that $\int_{-\infty}^\infty x^{2n}e^{-x^2}\, dx = (2n)!\sqrt{\pi}/4^nn!$ by differentiating the equation\\ $\int_{-\infty}^\infty e^{-tx^2} \, dx = \sqrt{\pi/t}$ (see Proposition 2.53).
\end{problem}
\begin{solution}
	Let $f:[1,\infty)\times[0,\infty)\to\bR$ be defined as $f(t,x)=e^{-tx}$, we see that for all $t,x$ we have
	\[
		\qty|\frac{\partial}{\partial t}f(x,t)| = xe^{-tx}\leq xe^{-x}
	\]
	which is $L^1$ allowing us to use Theorem 2.27 (page 56) to deduce that
	\[
		\frac{-1}{t^2}=\frac{d}{dt}\int_{0}^{\infty}e^{-tx}\mathrm{d}x = \int_0^\infty\frac{\partial}{\partial t}f(x,t)dx=\int_0^\infty -xe^{-tx}dx
	\]
	repeating this a total of $n$ times with the dominating functions being $x^ne^{-x}$ we get the desired conclusion after setting $t=1$ at the end. 
	For the second parts we work with $t$ over the same interval but this time $x$ varies over $\bR$ with $f(t,x)=e^{-tx^2}$ and following the exact same steps with dominating functions of form $x^{2n}e^{-x^2}$ we arrive at a proof.
\end{solution}

\begin{problem}\label{prob:2.30}
Show that $\lim_{k \to \infty} \int_0^k x^n (1 - k^{-1}x)^k \, dx = n!$.
\end{problem}
\begin{solution}
	The integrands $\chi_{[0,k]}x^n(1-x/k)^k$ are dominated by $x^ne^{-x}$ which is $L^1$ and also the pointwise limit of the integrands. 
	Applying DCT and using the previosu problem we arrive at the required conclusion.
\end{solution}
\begin{problem}\label{prob:2.31}
Derive the following formulas by expanding part of the integrand into an infinite series and justifying the term-by-term integration. Exercise 29 may be useful. (Note: In (d) and (e), term-by-term integration works, and the resulting series converges, only for $a > 1$, but the formulas as stated are actually valid for all $a > 0$.)
\begin{enumerate}[label=\alph*.]
    \item For $a > 0$, $\int_0^\infty e^{-x^2} \cos ax \, dx = \sqrt{\pi}e^{-a^2/4}$.
    \item For $a > -1$, $\int_0^1 x^a (1 - x)^{-1} \log x \, dx = \sum_1^\infty (a + k)^{-2}$.
    \item For $a > 1$, $\int_0^\infty x^{a-1} (e^x - 1)^{-1} \, dx = \Gamma(a) \zeta(a)$, where $\zeta(a) = \sum_1^\infty n^{-a}$.
    \item For $a > 1$, $\int_0^\infty e^{-ax} x^{-1} \sin x \, dx = \arctan(a^{-1})$.
    \item For $a > 1$, $\int_0^\infty e^{-ax} J_0(x) \, dx = (a^2 + 1)^{-1/2}$, where $J_0(x) = \sum_0^\infty (-1)^n x^{2n} / 4^n(n!)^2$ is the Bessel function of order zero.
\end{enumerate}
\end{problem}
\begin{solution}
    All of these are very easily proven by expanding some function into its power series representation and then using DCT (with the dominating function being of form $\sum|f_n|$ for functions of form $\sum f_n$) to switch the summation and integration so we state what we expand for each.
	\begin{enumerate}[label=\alph*.]
		\item expand the cosine term and use the previous problem (we set $t=1$ at the end, here we require the result without that substitution.)
		\item expand the denominator.
		\item Rewrite the integrand as
			\[
				\frac{x^{a-1}}{e^x-1}= \frac{x^{a-1}e^{-x}}{1-e^{-x}}=\sum_{n=0}^\infty x^{a-1}e^{-x}(e^{-x})^n
			\]
		\item expand the sine term and use the previous problem.
		\item this is already given as a power series
	\end{enumerate}
	
\end{solution}

\begin{problem}\label{prob:2.32}
Suppose $\mu(X) < \infty$. If $f$ and $g$ are complex-valued measurable functions on $X$, define
\[ \rho(f, g) = \int \frac{|f - g|}{1 + |f - g|} \, d\mu. \]
Then $\rho$ is a metric on the space of measurable functions if we identify functions that are equal a.e., and $f_n \to f$ with respect to this metric iff $f_n \to f$ in measure.
\end{problem}
\begin{solution}
    For the if part, first notice that $x\mapsto x/(1+x)$ is an increasing map on the nonnegatives and thus we have for any arbitrary $\epsilon>0$ that,
\[
	\int_X \frac{|f-f_n|}{1+|f-f_n|}d\mu=\int_{|f-f_n|\geq\epsilon} \frac{|f-f_n|}{1+|f-f_n|}d\mu+\int_{|f-f_n|<\epsilon} \frac{|f-f_n|}{1+|f-f_n|}d\mu
\]
\[
	\leq 1\cdot\mu(\{|f_n-f|\geq\epsilon\})+\frac{\epsilon}{1+\epsilon}\cdot\mu(\{|f_n-f|<\epsilon\})\leq\mu(\{|f_n-f|\geq\epsilon\})+\epsilon\cdot\mu(X)
\]
And as $f_n\to f$ in measure we see that $\lim\rho(f_n,f)\leq\epsilon\mu(X)$ which can be made arbitrarily small proving this direction.
For the other direction we similarly have that
\[
	\rho(f_n,f)\geq \frac{\epsilon}{1+\epsilon}\cdot\mu(\{|f_n-f|\geq\epsilon\})
\]
As $\rho(f_n,f)\to 0$ we must have the right side of the inequality tend to $0$ proving the claim.
\end{solution}
\begin{problem}\label{prob:2.33}
If $f_n \ge 0$ and $f_n \to f$ in measure, then $\int f \le \lim \inf \int f_n$.
\end{problem}
\begin{solution}
	By the definition of the limit infimum we can find a subsequence $f_{n_k}$ such that $\liminf\int f_n=\lim\int f_{n_k}$.
	As $f_n$ converges in measure to $f$ so does $f_{n_k}$ and we can find a subsequence $f_{n_{k_m}}$ that converges a.e. to $f$. Now, we see that
\[
	\liminf_n\int f_n=\lim_k\int f_{n_k}\limsup_k\int f_{n_k}\geq\limsup_m\int f_{n_{k_m}}\geq\liminf_m\int f_{n_{k_m}}\geq\int\liminf f_{n_{k_m}}
\]
where we use fatou's lemma for the last inequality. 
As $f_{n_{k_m}}\to f$ a.e., the last expression is just $\int f$ so we are done.
\end{solution}
\begin{problem}\label{prob:2.34}
Suppose $|f_n| \le g \in L^1$ and $f_n \to f$ in measure.
\begin{enumerate}[label=\alph*.]
    \item $\int f = \lim \int f_n$.
    \item $f_n \to f$ in $L^1$.
\end{enumerate}
\end{problem}
\begin{solution}
	\begin{enumerate}[label=\alph*.]
		\item By \hyperref[prob:2.33]{Problem 1.33} we can use fatou's lemma just like we did in the proof for DCT here as well by considering the sequences $g-f_n,g+f_n$ whice are nonnegative and converge to $g-f,g+f$ respectively in measure. 
		\item By part (a) and the assumption we clearly see that $f_n,f$ are $L^1$ and as $|f_n|\to|f|$ in measure too while being dominated by $g$ we can apply part (a) again to see that $\int |f|=\lim\int |f_n|$ which is equivalent to the claim by \hyperref[prob:2.21]{Problem 1.21}.
	\end{enumerate}
	enumerate
\end{solution}
\begin{problem}\label{prob:2.35}
$f_n \to f$ in measure iff for every $\epsilon > 0$ there exists $N \in \bN$ such that $\mu(\{x : |f_n(x) - f(x)| \ge \epsilon\}) < \epsilon$ for $n \ge N$.
\end{problem}
\begin{solution}
	The only if direction follows directly from the definition and for the if direction, given $\epsilon>0$ and $\delta>0$ let, $\nu=\min\{\epsilon,\delta\}$ we see that $\{|f_n-f|\geq\epsilon\}\subseteq\{|f_n-f|\geq\nu\}$ and we can find some $N\in\bN$ such that for all $n\geq N$,
\[
	\mu(\{|f_n-f|\geq\epsilon\})\leq\mu(\{|f_n-f|\geq\nu\})<\nu\leq\delta
\]
and we are done. 
\end{solution}
\begin{problem}\label{prob:2.36}
If $\mu(E_n) < \infty$ for $n \in \bN$ and $\chi_{E_n} \to f$ in $L^1$, then $f$ is (a.e. equal to) the characteristic function of a measurable set.
\end{problem}
\begin{solution}
	As the convergence is $L^1$ we can find a subsequence $\chi_{E_{n_k}}$ that converges to $f$ outside a set of measure $0$. 
Outside this set we cleary have that $f$ is a charecteristic function and its also measurable restricted to this set being the limit of measurable functions, thus it is a.e. equal to a characteristic function of a measurable set.
\end{solution}
\begin{problem}\label{prob:2.37}
Suppose that $f_n$ and $f$ are measurable complex valued functions and $\phi:\bC\to\bC$.
\begin{enumerate}[label=\alph*.]
	\item If $\phi$ is continuous and $f_n\to f$ a.e., then $\phi\circ f_n\to\phi\circ f$ a.e.
	\item If $\phi$ is uniformly continuous and $f_n\to f$ uniformly, almost uniformly, or in measure, then $\phi\circ f_n\to\phi\circ f$ uniformly, almost uniformly, or in measure, 
respectively. 
	\item There are counter examples when the continuity assumptions on $\phi$ are not satisfied
\end{enumerate}
\end{problem}
\begin{solution}
\begin{enumerate}[label=\alph*.]
	\item This is obvious from the fact that composition by continuous functions conserve pointwise convergence which is clear from the definition of continuity.
	\item It is a fact from basic real analysis that composition with uniformly continuous functions conserves uniform convergence and in the case of almost uniform convergence we just restrict ourselves to a suitable set.
In the case of convergence in measure, first choose $\delta>0$ such that $|x-y|<\delta\implies|\phi(x)-\phi(y)|<\epsilon$, then we see that,
\[
	\{x\in X:|f_n(x)-f(x)|<\delta\}\subseteq\{x\in X:|\phi(f_n(x))-\phi(f(x))|<\epsilon\}
\]
and taking complements,
\[
	\{x\in X:|f_n(x)-f(x)|\geq\delta\}\supseteq\{x\in X:|\phi(f_n(x))-\phi(f(x))|\geq\epsilon\}
\]
Thus, we see that,
\[    
	0\longleftarrow\mu(\{x\in X:|f_n(x)-f(x)|\geq\delta\})\geq\mu(\{x\in X:|\phi(f_n(x))-\phi(f(x))|\geq\epsilon\})
\]
proving this claim. 
\item A counterexample for part (a) if continuity of $\phi$ is not assumed is: $X=\bR$ with the lebesgue measure, $\phi(x)=1-\chi_{\{0\}}$ and $f_n(x)=1/n\to 0$. For part (b) in the case of uniform and almost uniform convergence is, in the same measurable space as the previous is: $f_n(x)=x+1/n\to x$ and $\phi(x)=x^2$, this is also a counterexample for the case of convergence in measure.
\end{enumerate}
\end{solution}
\begin{problem}\label{prob:2.38}
If $f_n \to f$ in measure and $g_n \to g$ in measure, then $f_n + g_n \to f + g$ in measure, and $f_ng_n \to fg$ in measure if $\mu(X) < \infty$, but not necessarily if $\mu(X) = \infty$.
\end{problem}
\begin{solution}
The first part follows from the fact that 
\[
	\{x\in X:|f_n+g_n-(f+g)|\geq\epsilon\}\subseteq\{x\in X:|f_n-f|\geq\epsilon\}\cup\{x\in X:|g_n-g|\geq\epsilon\}
\]
from the triangle inequality. 
Forthe second part it suffices to prove that if $\psi_n\to 0$ in measure then $\psi_n \phi\to 0$ in measure for any measurable $\phi$ as,
\[
	|f_ng_n-fg|\leq |f_n-f||g_n|+|f||g_n-g|\leq|g||f_n-f|+|g-g_n||f-f_n|+|f||g_n-g|
\]
due to the middle term converging to $0$ in measure, which follows from the set inclusion below. 
For any $\epsilon>0$
\[
	\{|f_n-f||g_n-g|\geq\epsilon\}\subseteq\{|f_n-f|\geq\sqrt{\epsilon}\}\cup\{|g_n-f|\geq\sqrt{\epsilon}\}
\]
After which we can apply the previous result that convergence in measure is preserved unded addition.
Now consider $S_n:=\{x\in X:n\leq |\phi|<n+1\}$. Then $\{S_n\}_{n\geq 0}$ is a partition of $X$. We see that $$\sum_{n\geq 0}\mu(S_n)=\mu(X)<\infty$$ and thus for any $\epsilon>0$ there exists some $m_\epsilon$ such that $$\sum_{n>m_\epsilon}\mu(S_n)<\epsilon/2$$
Implying that,
\[
\mu(\{x\in X:|\phi\psi_k|\geq \epsilon\})<\sum_{n=0}^{m_\epsilon}\mu(\{x\in S_n:|\psi_k\phi|\geq\epsilon\})+\epsilon/2
\]
From construction, for $k=0,1,\ldots,m_\epsilon$ we have that
\[
	\{x\in S_n:|\phi\psi_k|\geq\epsilon\}\subseteq\{x\in X:|\psi_k|\geq\frac{\epsilon}{n+1}\}\subseteq\{|\psi_k|\geq\frac{\epsilon}{m_\epsilon+1}\}
\]
Let $N$ be such that $k\geq N\implies \mu(\{|\psi_k|\geq\epsilon/(m_\epsilon+1)
\})<\epsilon/2(m_\epsilon+1)$. 
Then from the previous work we see that
\[
	\mu(\{|\phi\psi_k|\geq\epsilon\})<\sum_{n=0}^{m_\epsilon}\mu(\{x\in S_n:|\psi_k\phi|\geq \epsilon\})+\epsilon/2
 <\frac{(m_\epsilon+1)\epsilon}{2(m_\epsilon+1)}+\frac{\epsilon}{2}=\epsilon
\]
for all $k\geq N$ implying convergence in measure (by \hyperref[prob:2.35]{Problem 1.35}) and we are done. 
For a counter example in the case of $\mu(X)=\infty$ consider $g_n(x)=x^2$ and $f_n(x)=1/n$ on $\bR$ with the lebesgue measure.
\end{solution}
\begin{problem}\label{prob:2.39}
If $f_n \to f$ almost uniformly, then $f_n \to f$ a.e. and $f_n \to f$ in measure.
\end{problem}
\begin{solution}
    Suppose $E_n$ are measurable sets such that $\mu(E_n^c)<1/n$ then we have that $\mu((\cup E_n)^c=\cap E_n^c)<1/n$ for all $n$ and hence $0$. 
	If $x\in\cup E_n$ then $x\in E_n$ for some $n$ and thus $f_n(x)\to f(x)$ as we have uniform convergence which implies pointwise convergence on $E_n$. 
	Thus, we have convergence a.e. due to $(\cup E_n)^c$ being null. 
	Given $\delta>0$, almost uniform convergence gives a set $E$ with $\mu(E^c)<\delta$ such that $f_n\to f$ uniformly on $E$. 
	Because of uniform convergence, for large enough $n$, $|f_n(x)-f(x)|<\epsilon$ for all $x\in E$. 
	This means that the set $\{|f_n-f|\geq\epsilon\}$ is entirely contained within $E^c$. 
	Therefore, its measure is less than $\delta$. 
	Since $\delta$ was arbitrary, the measure tends to $0$.
\end{solution}
\begin{problem}\label{prob:2.40}
In Egoroff's theorem, the hypothesis ``$\mu(X) < \infty$'' can be replaced by ``$|f_n| \le g$ for all $n$, where $g \in L^1(\mu)$.''
\end{problem}
\begin{solution}
We first restrict ourselves to a set witha null complement such that we have pointwise convergence on it and assume this to be the whole space here on which causes no problems due to the complement being null. It would suffice for the sets $E_{1}(k)$ to have finite measure, which they do as
\[
	\bigcup_{m\geq n}\left\{x\in X:|f_m(x)-f(x)|\geq \frac{1}{k}\right\}\subseteq\left\{x\in X:2g(x)\geq \frac{1}{k}\right\}
\]
with the set on the right being of finite measure due to $g$ being $L^1$ (more precisely, it is easily shown that it has measure atmost $2k\int_Xg)$
\end{solution}
\begin{problem}\label{prob:2.41}
If $\mu$ is $\sigma$-finite and $f_n \to f$ a.e., there exist measurable $E_1, E_2, \ldots \subset X$ such that $\mu( (\bigcup_1^\infty E_j)^c ) = 0$ and $f_n \to f$ uniformly on each $E_j$.
\end{problem}
\begin{solution}
Using sigma finiteness write $X=\cup E_n$ with $E_n$ being finite. 
We can apply egoroff's on each of the $E_n$'s to get $E_{n}^{(m)}\subseteq E_n$ on which we have uniform convergence, with $\mu(E_n\setminus E_n^{(m)})<1/m$, then we see that $\mu(E_n\setminus\cup_mE_n^{(m)})=0$. 
As,
\[
	X\setminus\bigcup_{m,n}E_n^{(m)}=\qty(\bigcup_n E_n)\setminus\bigcup_n\bigcup_m E_n^{(m)}\subseteq\bigcup_n\qty(E_n\setminus\bigcup_m E_n^{(m)})
\]
we see that the collection $\{E_n^{(m)}\}_{m,n}$ satisfies the condition in the problem statement as a countable union of null sets (which are $E_n\setminus\cup_mE_n^{(m)}$ here) is still null.
\end{solution}
\begin{problem}\label{prob:2.42}
Let $\mu$ be counting measure on $\bN$. Then $f_n \to f$ in measure iff $f_n \to f$ uniformly.
\end{problem}
\begin{solution}
	For the only if direction notice that, for any $\epsilon>0$ if $\mu(\{|f_n-f|\geq\epsilon\})<1/2$ for all $n\geq N$ then this set is empty implying that $|f_n(x)-f(x)|<\epsilon$ for all $x\in\bN$, for all $n\geq N$ proving uniform convergence.
	The only if direction is always true regardless of the domain measurable space as, if $|f_n(x)-f(x)|<\epsilon$ for all $n$ large enough, then $\{|f_n-f|\geq\epsilon\}$ is eventually always empty and thus has measure $0$ for large enough $n$ making the measure tend to $0$ as $n\to\infty$ proving convergence in measure. 
\end{solution}
\begin{problem}\label{prob:2.43}
Suppose that $\mu(X) < \infty$ and $f: X \times [0, 1] \to \bC$ is a function such that $f(\cdot, y)$ is measurable for each $y \in [0, 1]$ and $f(x, \cdot)$ is continuous for each $x \in X$.
\begin{enumerate}[label=\alph*.]
    \item If $0 < \epsilon, \delta < 1$ then $E_{\epsilon, \delta} = \{x : |f(x, y) - f(x, 0)| \le \epsilon \text{ for all } y < \delta\}$ is measurable.
    \item For any $\epsilon > 0$ there is a set $E \subset X$ such that $\mu(E) < \epsilon$ and $f(\cdot, y) \to f(\cdot, 0)$ uniformly on $E^c$ as $y \to 0$.
\end{enumerate}
\end{problem}
\begin{solution}
\begin{enumerate}[label=\alph*.]
	\item By continuity we have that,
		\[
			E_{\epsilon,\delta}=\bigcap_{r\in\bQ\cap(0,\delta)}\{x\in X:|f(x,r)-f(x,0)|\leq\epsilon\}
		\]
	as the sets we take the intersection of are measurable, so is the set itself.
	\item This is almost like a continuous version (as in, the variable we judge uniform convergence over is continuous) of egoroff's theorem. Consider the sets
		\[
			F_n(k) = \bigcup_{m\geq n}E_{1/k,1/m}^c
		\]
		\[
			= \left\{x\in X:\exists m(m\geq n)\text{ such that }|f(x,y)-f(x,0)|\geq\frac{1}{k}\text{ for some }y< \frac{1}{m}\right\}
		\]
		By continuity we know that for any fixed $x\in X$ we have $f(x,y)\to f(x,0)$ as $y\to 0$, thus for any fixed $k$, we see that $\cap_n F_n(k)$ is empty as if some $x$ were to be in it then for any $n$ we could find some $m\geq n$ such that there exists some $y$ with $0<y<1/m\leq 1/n$ and $|f(x,y)-f(x,0)|\geq 1/k$ which would contradict continuity as there is some $n$ such that $y<1/n$ implies $|f(x,y)-f(x,0)|<1/k$. 
		As we have $\mu(X)<\infty$ and $F_n(k)$ decreases in $n$ we see that $\lim_n\mu(F_n(k))=\mu(\cap_nF_n(k))=0$.
		For any $\epsilon>0$, for all $k$ we can find some $n_k$ with $\mu(F_{n_k}(k))<\epsilon 2^{-k}$.
	Thus, $\mu(\cup_k F_{n_k}(k))<\epsilon$ and as,
	\[
		\qty(\bigcup_kF_{n_k}(k))^c=\bigcap_k(F_{n_k}(k))^c
	\]
	\[
		=\left\{x\in X:\forall k,m\geq n_k\implies|f(x,y)-f(x,0)|< \frac{1}{k}\text{ for all }y< \frac{1}{m}\right\}
	\]
	Let, $E=\cup_kF_{n_k}(k)$, then for all $x\in E^c$, if $y<1/n_k$ then $|f(x,y)-f(x,0)|<1/k$ and thus we have uniform convergence on this set.
	As $\mu(E)<\epsilon$ we are done.
\end{enumerate}   
\end{solution}
\begin{problem}\label{prob:2.44}
(Lusin's Theorem) If $f: [a, b] \to \bC$ is Lebesgue measurable and $\epsilon > 0$, there is a compact set $E \subset [a, b]$ such that $m(E^c) < \epsilon$ and $f|_E$ is continuous. (Use Egoroff's theorem and Theorem 2.26.)
\end{problem}
\begin{solution}
	Fix an $\epsilon>0$, then following a a similar arguement as we did in \hyperref[prob:2.38]{Problem 1.38} we can find some $M$ such that $m(\{x\in [a,b]:|f(x)|\geq M\})<\epsilon$, let the complement of this set be $E$. 
	Then, as $f$ is bounded on $E$, it is $L^1$ and using Theorem $2.26$ we can find continuous functions arbitrarily close to this on $E$ in the $L^1$ metric. 
	We can hence choose a sequence of continuous functions that converge to $f$ on $E$ in the $L^1$ metric, due to which we can choose a subsequence of this sequence that converges in measure to our function by Corollary $2.32$, after that using Egoroff's theorem we can find some measurable $F\subseteq E$ with $m(E\setminus F)<\epsilon$ on which convergence is uniform. 
	As the limit of an uniformly convergent sequence of continuous functions is continuous we have that $f$ is continuous on $F$. 
	Now, as $F^c\subseteq(E\setminus F)\cup E^c$ we see thatm $m(F^c)=m(E^c)+m(E\setminus F)<\epsilon+\epsilon=2\epsilon$. 
	As the lebesgue measure is inner regular (Theorem $1.18$ on page $36$) we can find some compact $K\subseteq F$ such that $\mu(F\setminus K)<\epsilon$ and thus similarly, $m(K^c)=m(F\setminus K)+m(F^c)<\epsilon+2\epsilon=3\epsilon.$
	Thus, we have found a compact set $K\subseteq[a,b]$ with $m(K^c)<3\epsilon$ on which $f$ is continuous.
\end{solution}
\begin{problem}\label{prob:2.45}
If $(X_j, \mM_j)$ is a measurable space for $j = 1, 2, 3$, then $\bigotimes_1^3 \mM_j = (\mM_1 \otimes \mM_2) \otimes \mM_3$. Moreover, if $\mu_j$ is a $\sigma$-finite measure on $(X_j, \mM_j)$, then $\mu_1 \times \mu_2 \times \mu_3 = (\mu_1 \times \mu_2) \times \mu_3$.
\end{problem}
\begin{solution}
    Having identified $A\times b\times C$ with $(A\times B)\times C$ in the text, we find that by definiton we see that $\otimes^3_1\mathcal M_j=(\mathcal M_1\otimes\mathcal M_2)\otimes\mathcal M_3$ as each of these clearly contains the generators of the other.
	These two measures also clearly agree on the measurable rectangles and thus we can say that they are equal on the algebra formed by the finite disjoint unions of these rectangles, hence we can say that they are equal on the whole sigma algebra by the uniqueness criterion in caratheodory's extension theorem from chapter $1$ as the premeasures in this case are clearly $\sigma$-finite (the reasoning behind this is given for the case of two measurable spaces at the bottom of page $64$ and we thus omit the extremely similar arguement needed here)
\end{solution}
\begin{problem}\label{prob:2.46}
Let $X = Y = [0, 1]$, $\mM = \mN = \mathcal{B}_{[0, 1]}$, $\mu = \text{Lebesgue measure}$, and $\nu = \text{counting measure}$. If $D = \{(x, x) : x \in [0, 1]\}$ is the diagonal in $X \times Y$, then $\iint \chi_D \, d\mu \, d\nu$, $\iint \chi_D \, d\nu \, d\mu$, and $\int \chi_D \, d(\mu \times \nu)$ are all unequal. (To compute $\int \chi_D \, d(\mu \times \nu)$, go back to the definition of $\mu \times \nu$.)
\end{problem}
\begin{solution}
	$D$ is clearly measurable and thus $\chi_D\in L^+$.
	We have that $D_x=\{x\}$ and $D^y=\{y\}$, we use this to compute the iterated integrals combined with the fact that for any $E$ we have $(\chi_E)_x=\chi_{E_x}$ and similarly for $y$ as well.
	\[
		\int\int\chi_Dd\mu(x)d\nu(y)=\int\int\chi_{D^y}(x)d\mu(x)d\nu(y)=\int\int\chi_{\{y\}}(x)d\mu(x)d\nu(y)=\int\mu(\{y\})d\nu(y)=0
	\]
	as $\mu(\{y\})$ is always $0$. 
	\[
		\int\int\chi_Dd\nu(y)d\mu(x)=\int\nu(\{x\})d\mu(x)=\int 1d\mu(x)=\int_{[0,1]}1\mu(x)=1
	\]
	We will now show that $\int\chi_Dd(\mu\times\nu)=\mu\times\nu(D)=\infty$. 
	If $D\subseteq\cup I_n$ with $I_n\in\mathcal A$ then, we can break these $I_n$ into rectangles and have $D\subset \cup A_m\times B_m=\cup I_n$. 
	Clearly we must have $[0,1]\subseteq \cup(A_m\cap B_m)$ as $(x,x)\in A_m\times B_m$ iff $x\in A_m\cap B_m$.
	Keeping in mind that $A_m,B_m\in\mathcal B_{[0,1]}$ we thus have that $\sum\mu(A_m\cup B_m)\geq 1$ which implies that, for some $m$ we must have $\mu(A_m),\mu(B_m)\geq\mu(A_m\cap B_m)>0$ but then $B_m$ has more than countably many elements making $\nu(B_m)=\infty$ so $\mu\times\nu(A_m\times B_m)=\mu(A_m)\nu(B_m)=\infty$. 
	Thus, in the definiton of $\mu\times\nu$ the numbers we take an infimum over is the singleton $\{\infty\}$ making the measure of $D$ infinite and thus all the integrals are unequal.
\end{solution}
\begin{problem}\label{prob:2.47}
Let $X = Y$ be an uncountable linearly ordered set such that for each $x \in X$, $\{y \in X : y < x\}$ is countable. (Example: the set of countable ordinals.) Let $\mM = \mN$ be the $\sigma$-algebra of countable or co-countable sets, and let $\mu = \nu$ be defined on $\mM$ by $\mu(A) = 0$ if $A$ is countable and $\mu(A) = 1$ if $A$ is co-countable. Let $E = \{(x, y) \in X \times X : y < x\}$. Then $E_x$ and $E^y$ are measurable for all $x, y$, but $\iint \chi_E \, d\mu \, d\nu$ and $\iint \chi_E \, d\nu \, d\mu$ exist but are not equal. (If one believes in the continuum hypothesis, one can take $X = [0, 1]$ [with a nonstandard ordering] and thus obtain a set $E \subset [0, 1]^2$ such that $E_x$ is countable and $E^y$ is co-countable [in particular, Borel] for all $x, y$, but $E$ is not Lebesgue measurable.)
\end{problem}
\begin{solution}
	Clearly $E_x=\{y\in X:y<x\}$ and $E^y=\{x\in X:x>y\}$. 
	By definition we see that $E_x$ is always countable and $E^y$ always uncountable having a countable complement $E_y\cup\{y\}$, so these are always measurable.
	Thus,
	\[
	    \int\int\chi_Ed\nu d\mu=\int\nu(E_x)d\mu(x)=\int 0d\mu(x)=0
	\]
	\[
		\int\int\chi_Ed\mu d\nu=\int\mu(E^y)\nu(y)=\int 1d\nu(y)=\nu(Y)=1\text{ as }Y\text{ is }\twemoji{coconut}\text{able}
	\]
	So these integrals are unequal.	
\end{solution}
\begin{problem}\label{prob:2.48}
Let $X = Y = \bN, \mM = \mN = \mathcal{P}(\bN), \mu = \nu = \text{counting measure}$. Define $f(m, n) = 1$ if $m = n, f(m, n) = -1$ if $m = n + 1$, and $f(m, n) = 0$ otherwise. Then $\int f \, d\mu \, d\nu$ and $\int f \, d\nu \, d\mu$ exist and are unequal.
\end{problem}
\begin{solution}
    Firstly we see that $f(k,k)=1$ and $f(k,k-1)=1$. Also, $f$ is measurable (all functions are) as $\mathcal P(\bN)\otimes\mathcal P(\bN)=\mathcal M(\mathcal P(\bN)\times\mathcal P(\bN))=\mathcal P(\bN\times\bN)$. Integration with respect to these counting measures is just summation, thus we can write,
	\[
		\int\int fd\mu d\nu=\sum_{m=1}^\infty\sum_{n=1}^\infty f(m,n)=1+(1-1)+(1-1)+\ldots=1
	\]
	\[
		\int\int fd\nu d\mu=\sum_{n=1}^\infty\sum_{m=1}^\infty f(m,n)=(1-1)+(1-1)+\ldots=0
	\]
		
\end{solution}
\begin{problem}\label{prob:2.49}
Prove Theorem 2.39 by using Theorem 2.37 and Proposition 2.12 together with the following lemmas.
\begin{enumerate}[label=\alph*.]
    \item If $E \in \mM \otimes \mN$ (typo in the book) and $\mu \times \nu(E) = 0$, then $\nu(E_x) = 0$ for a.e. $x$ and $\mu(E^y) = 0$ for a.e. $y$.
    \item If $f$ is $\mathcal L$-measurable and $f = 0$ $\lambda$-a.e., then $f_x$ and $f^y$ are integrable for a.e. $x$ and $y$, and $\int f_x \, d\nu = \int f^y \, d\mu = 0$ for a.e. $x$ and $y$. (Here the completeness of $\mu$ and $\nu$ is needed.)
\end{enumerate}
\end{problem}
\begin{solution}
    We first prove the lemmas. 
	\begin{enumerate}[label=\alph*.]
		\item From the assumptions we have that
			\[
			    0=\mu\times\nu(E)=\int\mu(E^y)\nu(y)=\int\nu(E_x)\mu(x)
			\]
			thus, $\nu(E_x)=0$ for a.e. $x$ and similarly for $y$.
		\item As $f$ is $\mathcal L=\overline{\mathcal M\otimes\mathcal N}$ measurable, by proposition $2.12$, there exists a $\mathcal M\otimes\mathcal N$ measurable function $g$ such that $f=g$ $\lambda$-a.e. which clearly implies that $f=g$ $\mu\times\nu$-a.e. as well. 
			Now, if $f$ is $0$ $\lambda$ a.e. then so is $g$ and thus we also similarly have that $g=0$ $\mu\times\nu$-a.e. 
			We see that $g_x$ and $g^y$ are $\mathcal N$ and $\mathcal M$ measurable respectively. 
			Suppose we had $f=g$ outside $N$ with $\mu\times\nu(N)=0$ then by part(a) we see that $\mu(N^y)=0$ for a.e. $y$ and $\nu(N_x)=0$ for a.e $x$. 
			As, $f_x=g_x$ outside $N_x$ which is null for a.e. $x$, we have that $f_x=g_x$ a.e. for a.e. $x$. 
			In the context of the theorem, $\mathcal N$ is complete which makes $f_x$ measurable if $f_x=g_x$ a.e. 
			As $g$ is $0$ $\mu\times\nu$-a.e. we see that $g\in L^1(\mu\times\nu)$ allowing us to use Tonelli's theorem to say that,
			\[
				0=\int g\, d(\mu\times\nu)=\int\int g_x(y)\, d\nu(y) d\mu(x)\implies\int g_x\, d\nu=0\text{ for a.e. }x
			\]
			Combined with the fact that $f_x=g_x$ a.e. (implying $\int f_x=\int g_x$) for a.e. $x$ we see that $\int f_x\, d\nu=0$ for a.e. $x$. and thus $f_x$ is integrable for a.e. $x$ as well.
			The same can be shown for $y$ as well following the same procedure but with different variables.
	\end{enumerate}
	Now we prove the theorem.
	The claims concerning the measurablity of $f_x$ and $f^y$ are evident from the proof of the second lemma. 
	Now we will prove the rest assuming the same setup as used above, proving the other implications of the first case where $f\in L^+(\lambda)$ first.
	We can clearly choose $g$ such that $g\geq 0$ as well (make it $0$ on the null set wher it isn't equal to $f$, i.e. consider $g\chi_{N^c}$ which is also measurable  following the same notation as used above) and thus $g\in L^+(\mu\times\nu)$ making $x\mapsto\int g_x\,d\nu$ measurable by Tonelli's theorem. 
	From the second lemma's proof we see that $f_x=g_x$ a.e for a.e. $x$ making $\int f_x\,d\nu=\int g_x\,d\nu$ a.e. and as $\mathcal M$ is complete we must have that $x\mapsto\int f_x\,d\nu$ is measurable, being equal a.e. to a measurable function. 
	Now, as $\int g_x\,d\nu=\int f_x\,d\nu$ a.e we also see that $\int\int g\,d\nu d\mu=\int\int f\,d\nu d\mu$. 
	Using Tonelli's theorem again with the fact that the integrals w.r.t the completion and the underlying measure are equal for functions which are measurable w.r.t the underlying measure (this is easily proven) we find that,
	\[
	    \int f\,d\lambda=\int g\,d\lambda=\int g\,d(\mu\times\nu)=\int\int g\,d\nu d\mu=\int f\,d\nu d\mu
	\]
	The other order for integration can be dealt with similarly. 
	Now we prove the rest of the implications for part(b).
	If $f\in L^1(\lambda)$ then as $f=g$ a.e. we see that,
	\[
	    \int|g|\,d(\mu\times\nu)=\int|g|\,d(\lambda)=\int|f|d(\lambda)<\infty
	\]
	and thus, $g\in L^1(\mu\times\nu)$. 
	Tonelli's theorem tells us that $g_x\in L^1(\nu)$ for a.e $x$ and that the a.e. defined function $x\mapsto\int g_x\,d\nu$ is integrable (and also measurable outside a null set, this can be easily proven.) 
	As $\mathcal N$ is complete we see that it is measurable (we can ignore it on a null set, if this is hard to understand just email me.) 
	Now, clearly we have that $\int f_x=\int g_x$ for a.e. $x$ making $x\mapsto\int f_x\,d\nu$ a mesurable map.
	Thus, by Tonelli's theorem and the previous line we can say that
	\[
	    \int\int f_x(y)\,d\nu(y)d\mu(x)=\int\int g_x(y)\,d\nu(y)d\mu(x)=\int g\,d(\mu\times\nu)=\int f\,d\lambda
	\]
so we are done.
\end{solution}
\begin{problem}\label{prob:2.50}
Suppose $(X, \mM, \mu)$ is a $\sigma$-finite measure space and $f \in L^+(X)$. Let
\[ G_f = \{(x, y) \in X \times [0, \infty] : y \le f(x)\}. \]
Then $G_f$ is $\mM \times \mathcal{B}_{\bR}$-measurable and $\mu \times m(G_f) = \int f \, d\mu$; the same is also true if the inequality $y \le f(x)$ in the definition of $G_f$ is replaced by $y < f(x)$. (To show measurability of $G_f$, note that the map $(x, y) \mapsto f(x) - y$ is the composition of $(x, y) \mapsto (f(x), y)$ and $(z, y) \mapsto z - y$.) This is the definitive statement of the familiar theorem from calculus, ``the integral of a function is the area under its graph.''
\end{problem}
\begin{solution}
    $G_f$ is measurable as the composition of measurable maps is measurable, as hinted in the problem statement. 
	Now, Theorem $2.36$ tells us that,
	$\mu\times m(G_f)=\int m((G_f)_x\,d\mu(x)$
	Now, clearly $(G_f)_x=[0,f(x)]$ which has lebesge measure $f(x)$ so we are done.
	If the inequality was strict then when $f(x)=0$ the section would be an empty set, which also has $0=f(x)$ and for $f(x)>0$ the one sided open interval $[0,f(x))$ still has lebesgue measure $f(x).$
\end{solution}
\begin{problem}\label{prob:2.51}
Let $(X, \mM, \mu)$ and $(Y, \mN, \nu)$ be arbitrary measure spaces (not necessarily $\sigma$-finite).
\begin{enumerate}[label=\alph*.]
    \item If $f: X \to \bC$ is $\mM$-measurable, $g: Y \to \bC$ is $\mN$-measurable, and $h(x, y) = f(x)g(y)$, then $h$ is $\mM \otimes \mN$-measurable.
    \item If $f \in L^1(\mu)$ and $g \in L^1(\nu)$, then $h \in L^1(\mu \times \nu)$ and $\int h \, d(\mu \times \nu) = (\int f \, d\mu)(\int g \, d\nu)$.
\end{enumerate}
\end{problem}
\begin{solution}
    This is obvious for characteristic functions and we can extend it to simple functions bylinearity. 
	Part (a) is trivial as the compositon of measurable maps is measurable so we prove part (b). 
	By Theorem $2.10$ there exists measurable simple functions $\phi_n,\psi_n$,  which tend to $f,g$ such that $|\phi_n|\nearrow|f|$ and $|\psi_n|\nearrow|g|$. 
	By the monotone convergence theorem we see that,
	\[
		\int|fg|=\lim\int|\psi_n\phi_n|=\lim\qty(\int|\psi_n|\int|\phi_n|)=\int|f|\int|g|<\infty\implies fg\in L^1(\mu\times\nu)
	\]
	Now, by DCT(we use it on the extreme equalities) we have that,
	\[
		\int fg=\lim\int \phi_n\psi_n=\lim\int\phi_n\int\psi_n=\qty(\lim\int\phi_n)\qty(\lim\int\psi_n)=\int f\int g
	\]
\end{solution}
\begin{problem}\label{prob:2.52}
The Fubini-Tonelli theorem is valid when $(X, \mM, \mu)$ is an arbitrary measure space and $Y$ is a countable set, $\mN = \mathcal{P}(Y)$, and $\nu$ is counting measure on $Y$. (Cf. Theorems 2.15 and 2.25.)
\end{problem}
\begin{solution}
    By the countable additivity of the integral for any nonnegative measurable function $f$ we have that,
	\[
		\int_{X\times Y}f\,d\mu(\mu\times\nu)=\sum_{y\in Y}\int_{X\times\{y\}}f(x,y)\,d\mu(x)d\nu(y)
	\]
	Now, any mesaurable subset of $X\times\{y\}$ clearly has form $E\times\{y\}$ where $E$ is $\mathcal M$ measurable and we see that  $\mu\times\nu(E\times\{y\})=\mu(E)\nu(\{y\})=\mu(E)$. 
	From the definiton of the integral it is clear that
	\[
		\int_{X\times\{y\}}f(x,y)\,d\mu(x)d\nu(y)=\int_Xf(x,y)\,d\mu(x)
	\]
	The map $y\mapsto\int f(x,y)\,d\mu(x)$ is also trivially measurable as the sigma algebra on $Y$ is the whole power set. 
	As $\nu$ is the counting measure we thus have,
	\[
		\int\int f(x,y)\,d\mu(x)d\nu(y)=\sum_{y\in Y}\int f(x,y)\,d\mu(x)=\int f\,d(\mu\times\nu)   
	\]
	Again, by the additivity of the integral here we see that
	\[
		\int\int f(x,y)d\nu(y)d\mu(x)=\int\sum_{y\in Y}f(x,y)d\mu(x)=\sum_{y\in Y}\int f(x,y)d\mu(x)
	\]
	which is $\int f\,d(\mu\times\nu)$ as we saw before. 
	We also see that the map $x\mapsto\int f(x,y)\,d\nu(y)$ is measurable as,
	\[
		\int f(x,y)\,d\nu(y)=\sum_{y\in Y}f(x,y)
	\]
	as its the limit of a sequence of measurable functions $f^{y_1}+\ldots+f^{y_n}$ (here we enumerate $Y$ as $\{y_n:n\in\bN\}$). 
	We have proven Tonelli's theorem for this case and that of Fubini can be proven from this as in the book, or we can also reenact the same steps while using DCT instead of the additivity of the integral (which was just a corollary of MCT in this case).
\end{solution}
\begin{problem}\label{prob:2.53}
Fill in the details of the proof of Theorem 2.41.
\end{problem}
\begin{solution}
    Fix $\epsilon>0$. 
	Take a sequence of simple measurable functions $\phi_n$ such that $|\phi_n|\nearrow|f|$ and by DCT, for large enough $n$ we have $\int|\phi_n-f|<\epsilon$, we write $\phi:=\phi_n$. 	Suppose $\phi=\sum_1^ma_j\chi_{E_j}$ where $a_j\neq 0$ and $E_j$ are disjoint, then we see that
	\[
		|a_j|\mu(E_j)=\int_{E_j}|\phi|\leq\int_{E_j}|f|<\infty
	\]
	Implying that $\mu(E_j)<\infty$. 
	By theorem $2.41c$, for each $E_j$ we can find a disjoint(disjoincy is neeeded here by the way) collection $\{R_{j,i}\}_{i=1}^{k_j}$ such that $m(E_j\Delta\cup_iR_{j,i})<\epsilon$ and each $R_{j,i}$ is a product of intervals, we will call these cubes here on out. 
	Define $\psi:=\sum_{i,j}a_{j}\chi_{R_{j,i}}$, then we see that
	\[
		\int|\phi-\psi|\leq\sum_{j=1}^m|a_j|\int\cdot\bigg|\chi_{E_j}-\sum_{j=1}^{k_j}\chi_{R_{j,i}}\bigg|\leq\sum_{j=1}^m|a_j|\cdot m\qty(E_j\Delta\bigcup_{i=1}^{k_j}R_{j,i})<\epsilon\cdot\sum_{j=1}^m|a_j|
	\]
	Thus, we have that $\int|f-\psi|\leq\int|f-\phi|+\int|\phi-\psi|<\epsilon\cdot(1+\sum_1^m|a_j|)$ which can be made arbitrarily small. 
	But the fact that $\psi$ is of the form described in the problem requires proof. 
	It suffices to show that any finite union of cubes can be written as a disjont union of such. 
	We consider cubes whose sides (as defined in the book) are closed intervals, the rest follow similarly. 
	Suppose, $A=\prod[a_j,b_j]$ and $B=\prod[c_j,d_j]$ then we clearly have that $A\cap B=\prod([a_j,b_j]\cap[c_j,d_j])$ which is also a cube (it might be empty but thats a cube too.) 
	We also find that 
	\[
		A^c=\bigcup_{f\in \{1,c\}^{\{1,\ldots,n\}}}\prod_{j=1}^n[a_j,b_j]^{f(i)}
	\]
	by extreme abuse of notation. 
	Notice that this is a disjoint union. 
	Also, $[a_j,b_j]^c=(-\infty,a_j)\cup(b_j,\infty)$ and this allows us to write any of the sets in the union above as a finite disjoint union of cubes as these sets are of form (after permuting the sides so it isn't as notationally strenous to write)
	\[
		\qty((-\infty,a_1)\cup(b_1,\infty))\times\ldots\times\qty((-\infty,a_m)\cup(b_m,\infty))\times[a_{m+1},b_{m+1}]\times\ldots[a_n,b_n]
	\]
	which can be written as a disjoint union of cubes following this set theoretic identity
	\[
		\prod_{i\in A}\bigcup_{j\in B}S_{i,j}=\bigcup_{f\in B^A}\prod_{i\in A}S_{i,f(j)}
	\]
	as in our case, for fixed $i$, all the $S_{i,j}$ are disjoint. 
	This shows that the set of all cubes is an elementary family and thus by Propostion $1.7$(page $23$) we can say that that the finite disjoint union of cubes is an algebra. 
	This tells us that we can write the union of cubes $\cup_{i,j}R_{j,i}$ as a disjoint union of cubesa and using this we can write $\psi$ as a simple function where all the characteristic functions are those of cubes which are disjoint so we are done with the first claim. 
	Again, by DCT for large enough $m$ we must have that $\int|\psi\cdot\chi_{[-m,m]^n}-f|<2\cdot\epsilon$ starting with small enough $\epsilon>0$. 
	We clearly see that $\psi\chi_{[-m,m]^n}$ vainshes outside a bounded set and as its defined using charactersitc functions on disjoint cubes, we can approximate it on these cubes by a continuius function as we did in the case of the real line, while ignoring cubes of measure $0$, which on the real line are just singletons, this allows us to get $\mathcal O(\epsilon^n)$ which can be made arbitrarily small. 
	Then we can use the triangle inequality and conclude that $C_c(\bR^n)$ is dense in $L^1(\bR^n,\mathcal L^n,m)$ in the $L^1$ metric.
\end{solution}
\begin{problem}\label{prob:2.54}
How much of Theorem 2.44 remains valid if $T$ is not invertible?
\end{problem}
\begin{solution}
	Part (b) is true and we prove it below whereas part (a) is false, that is $f\circ T$ is not neccessarily lebesgue measurable, given $f$ is. 
	As $T$ is not invertible we can find some invertible linear map $S$ such that $ST(\bR^n)=\bR^k\times\{0\}^{n-k}$ where $k<n$ and this is measurable with  measure $0$. 
	Thus, $m(S^{-1}ST(\bR^n))=|\det S^{-1}|m(ST(\bR^n))=0$ showing that $T(\bR^n)$ is a lebesgue null set.
	Now, for any $E\in\mathcal L^n$(any set even) we have that $T(E)\subseteq T(\bR^n)$ proving part (b) by completeness. 
	Now to disprove part (a), say $n=2$ and consider the indicator function $\chi$ of the lebesgue measurabe set $N\times 0$ where $N$ is not lebesgue measurable on the real lina and let $T:(x,y)\mapsto(x,0)$ then, $T^{-1}(f^{-1}(0.9,1.1))=N\times\bR$. 
	Now, if $N\times\bR$ was in $\mathcal L^2$ then its indicator would be measurable and by Theorem $2.39$ we would have that $f^y$ is measurable on $\bR$ for a.e. $y$ which is false as these are all indicators of $N$ on $\bR$ so the set was not measurable to begin with.
\end{solution}
\begin{problem}\label{prob:2.55}
Let $E = [0, 1] \times [0, 1]$. Investigate the existence and equality of $\int_E f \, dm^2$, $\int_0^1 \int_0^1 f(x, y) \, dx \, dy$, and $\int_0^1 \int_0^1 f(x, y) \, dy \, dx$ for the following $f$:
\begin{enumerate}[label=\alph*.]
    \item $f(x, y) = (x^2 - y^2)(x^2 + y^2)^{-2}$.
    \item $f(x, y) = (1 - xy)^{-a}$ ($a > 0$).
    \item $f(x, y) = (x - \frac{1}{2})^{-3}$ if $0 < y < |x - \frac{1}{2}|$, $f(x, y) = 0$ otherwise.
\end{enumerate}
\end{problem}
\begin{solution}
    \begin{enumerate}[label=\alph*.]
		\item We exclude $\{(0,0)\}$ which has measure $0$ and this does not alter anything we care about. 
			As $f$ is now continuous, using Theorem $2.39$ we can work in $\mathcal B_{\bR^2}$ and the positive and negative parts evaluate to $\infty$ after some basic calculus paired with fubini's theorem and hence the integral does not exist however the iterated integrals evaluate  to $\pi/4$ and $-\pi/4$ respectively. 
		\item Removing $\{(1,1)\}$ we can say that this is again continuous and following the same arguements above, revert back to $\mathcal B_{\bR^n}$. 
			As this is nonegative, all the integrals exist and are equal.
		\item Its clearly borel, we follow the same arguements above. The first iterated integral exists and equals $0$ whereas the other does not. 
			The integral of the positive and negative parts are both infinite and thus the integral over the product measure does not exist.
    \end{enumerate}
\end{solution}
\begin{problem}\label{prob:2.56}
If $f$ is Lebesgue integrable on $(0, a)$ and $g(x) = \int_x^a t^{-1} f(t) \, dt$, then $g$ is integrable on $(0, a)$ and $\int_0^a g(x) \, dx = \int_0^a f(x) \, dx$.
\end{problem}
\begin{solution}
	Consider the function $H(t,x) := \chi_{t>x}\cdot f(t)t^{-1}$ on $(0,a)^2$. 
	$H$ is clearly measurable w.r.t $m\times m$ and by Tonelli's theorem we have that
	\[
		\int_{(0,a)^2}|H|\,d(m\times m)=\int_0^a\int_0^a|H(t,x)|\,dm(t)\,dm(x)
	\]
	\[
		=\int_0^a\int_0^t \frac{|f(t)|}{t}\,dm(x)dm(t)=\int_0^a|f(t)|\,dm(t)<\infty
	\]
	Making $H$ integrable. 
	Now, by fubini's theorem we can say that $x\mapsto\int_0^aH(x,t)\,dm(t)=\int_a^xf(t)/t\,dm(t)=g(x)$ is measurable and 
	\[
	\int_0^af(t)\,dm(t)=\int_{(0,a)^s}H\,d(m\times m)=\int_0^a\int_0^aH(t,x)\,dm(t)dm(x)=\int_0^ag(x)\,dm(x)
	\]	
\end{solution}
\begin{problem}\label{prob:2.57}
Show that $\int_0^\infty e^{-sx} x^{-1} \sin x \, dx = \arctan(s^{-1})$ for $s > 0$ by integrating $e^{-sxy} \sin x$ with respect to $x$ and $y$. (It may be useful to recall that $\tan(\frac{\pi}{2} - \theta) = (\tan \theta)^{-1}$. Cf. Exercise 31d.)
\end{problem}
\begin{solution}
	Consider $f(x,y)=se^{-sxy}\sin(x)$ on $(0,\infty)\times(1,\infty)$, this is integrable as, by tonelli's we have
	\[
		\int |f|=\int_0^\infty\int_1^\infty se^{-sxy}|\sin(x)|dydx=\int_0^\infty e^{-sx}|\sin(x)/x|dx\leq\int_0^\infty e^{-sx}=s^{-1}<\infty
	\]
	Now using fubini's theorem we see that,
	\[
		\int_0^\infty e^{-sx}x^{-1}\sin(x)dx=\int_0^\infty\int_1^\infty f(x,y)dxdy=\int_1^\infty\int_0^\infty se^{-sxy}\sin(x)dxdy
	\]
	\[
		=\int_1^\infty \frac{s}{1+(sy)^2}dy=\pi/2-\arctan(s)=\arctan(s^{-1})
	\]
\end{solution}
\begin{problem}\label{prob:2.58}
Show that $\int e^{-sx} x^{-1} \sin^2 x \, dx = \frac{1}{4} \log(1 + 4s^{-2})$ for $s > 0$ by integrating $e^{-sx} \sin 2xy$ with respect to $x$ and $y$.
\end{problem}
\begin{solution}
	Using the same setup as before and the identity $\sin(x)^2=(1-\cos(2x))/2$ we just need to evaluate integrals of form $\int e^{ax}\cos(bx)dx$ this time instead of using sine there.
\end{solution}
\begin{problem}\label{prob:2.59}
Let $f(x) = x^{-1} \sin x$.
\begin{enumerate}[label=\alph*.]
    \item Show that $\int_0^\infty |f(x)| \, dx = \infty$.
    \item Show that $\lim_{b \to \infty} \int_0^b f(x) \, dx = \frac{1}{2}\pi$ by integrating $e^{-xy} \sin x$ with respect to $x$ and $y$. (In view of part (a), some care is needed in passing to the limit as $b \to \infty$.)
\end{enumerate}
\end{problem}
\begin{solution}
    \begin{enumerate}[label=(\alph*)]
		\item it can be shown with elementary methods that 
		\[	
			\int_0^n|\sin(x)/x|dx\gg H_n
		\]
		which suffices.
	\item Consider $f(x,y)=e^{-xy}\sin(x)$. 
		For any $b>0$ we see that,
		\[
			\int_0^b\int_0^\infty| f(x,y)|dydx\leq\int_0^b\int_0^\infty xe^{-xy}dydx=b<\infty
		\]
		so its integrable on the domain used for the integral and by fubini's we have,
		\[
			\int_0^b\frac{\sin x}{x}dx=\int_0^b\int_0^\infty f(x,y)dydx=\int_0^\infty\int_0^bf(x,y)dxdy
		\]
		\[
			=\lim_{b\to\infty}\int_0^\infty- \underbrace{\frac{e^{-by}\sin(b)y+e^{-by}\cos(b)-1}{1+y^2}}_{g_b(y)}dy
		\]
		We see that $g_b(y)\longrightarrow (1+y^2)^{-1}$ as $b\to\infty$ and $|g_b(y)|\leq 3(1+y^2)^{-1}$ which is integrable allowing us to use DCT (by passing to sequences) to conclude that,
		\[
		    \int_0^\infty \frac{\sin x}{x}dx=\int_0^\infty \frac{1}{1+y^2}dy= \frac{\pi}{2}
		\]
    \end{enumerate}
\end{solution}
\begin{problem}\label{prob:2.60}
$\Gamma(x)\Gamma(y)/\Gamma(x+y) = \int_0^1 t^{x-1}(1-t)^{y-1} \, dt$ for $x, y > 0$. (Recall that $\Gamma$ was defined in $\S 2.3$. Write $\Gamma(x)\Gamma(y)$ as a double integral and use the argument of the exponential as a new variable of integration.)
\end{problem}
\begin{solution}
	Here, $x>y$ denotes the set $\{(x,y)\in\bR^2:x>y\}$. 
	Let $f(x,y)=x^{\alpha-1}y^{\beta-1}e^{-(x+y)}$. 
	Then, by \hyperref[prob:2.51]{Problem 2.51} we see that
	\[
		\Gamma(\alpha)\Gamma(\beta)=\int_{(0,\infty)^2}f\,d(m\times m)
	\]
	and by Theorem $2.47$, using the invertible linear map $G:(x,y)\mapsto(x-y,y)$ we see that,
	\[
		\int_{G(\{x>y\})=(0,\infty)^2}f\,d(m\times m)=\int_{x>y}f(x-y,y)\cdot\underbrace{|\det G|}_{1}d(m\times m)
	\]
	and by Tonelli's we have,
	\[
		\int_{x>y}f(x-y,y)=\int_0^\infty\int_0^x(x-y)^{\alpha-1}y^{\beta-1}e^{-x}dydx
	\]
	\[
	=\int_0^\infty\int_0^x\qty(1- \frac{y}{x})^{\alpha-1}\qty( \frac{y}{x})^{\beta-1}x^{\alpha+\beta-1}e^{-x}dydx=\int_0^\infty x^{\alpha+\beta-1}e^{-x}\int_0^1v^{\beta-1}(1-v)^{\alpha-1}dv
	\]
	by a change of variable $v=y/x$. 
	We thus have that
	\[
		\Gamma(\alpha)\Gamma(\beta)=\Gamma(\alpha+\beta)\int_0^1t^{\alpha-1}(1-t)^{\beta-1}dt
	\]
	proving the claim.
\end{solution}
\begin{problem}\label{prob:2.61}
If $f$ is continuous on $[0, \infty)$, for $\alpha > 0$ and $x \ge 0$ let
\[ I_\alpha f(x) = \frac{1}{\Gamma(\alpha)} \int_0^x (x - t)^{\alpha-1} f(t) \, dt. \]
$I_\alpha f$ is called the $\alpha$th \textbf{fractional integral} of $f$.
\begin{enumerate}[label=\alph*.]
    \item $I_{\alpha+\beta} f = I_\alpha (I_\beta f)$ for all $\alpha, \beta > 0$. (Use Exercise 60.)
    \item If $n \in \bN$, $I_n f$ is an $n$th-order antiderivative of $f$.
\end{enumerate}
\end{problem}
\begin{solution}
    \begin{enumerate}[label=\alph*.]
		\item By expanding we have that,
			\[
				I_\alpha(I_\beta f)(x)=\frac{1}{\Gamma(\alpha)\Gamma(\beta)}\int_0^x\int_0^t(x-t)^{\alpha-1}(t-y)^{\beta-1}f(y)dydt
			\]
			Heuristically speaking we integrate over the region $x>t>y>0$. 
			Using the invertible linear map that switches the coordinates and Theorem $2.47$ and tonelli's theorem we see that this is equal to
			\[
				\frac{1}{\Gamma(\alpha)\Gamma(\beta)}\int_0^x\int_y^x(x-t)^{\alpha-1}(t-y)^{\beta-1}f(y)dtdy
			\]
			\[
				= \frac{1}{\Gamma(\alpha)\Gamma(\beta)}\int_0^xf(y)(x-y)^{\alpha+\beta-1}\int_y^x\qty( \frac{x-t}{x-y})^{\alpha-1}\qty( \frac{t-y}{x-y})^{\beta-1} \frac{dt}{x-y}dy
			\]
			which is $I_{\alpha+\beta}f(x)$ after a change of variable $v=(t-y)/(x-y)$ and an application of the previous problem.
		\item This follows from part (a) by repeatedly applying $I_1$ which is the antiderivative.
    \end{enumerate}
\end{solution}
\begin{problem}\label{prob:2.62}
The measure $\sigma$ on $S^{n-1}$ is invariant under rotations.
\end{problem}
\begin{solution}
	Let $T$ be a rotation. 
	Given any borel set $E\subseteq S^{n-1}$, let $E_1=\Phi^{-1}((0,1]\times E)$ as in the text, then we defined $\sigma(E)=n\cdot m(E_1)$ so it suffices to show that $m(E_1)=m(T(E)_1)$. 
	For any $rx'\in E_1$ in usual notation (same as the text) we find that $T(rx')=rT(x')$ with $T(x')\in S^{n-1}$ as well due to $T$ being a rotation ( $\norm{Tv}=\norm{v}$ for all $v$), thus we have that $T(E_1)=\{rT(x'):r\in(0,1],x'\in E\}=\Phi^{-1}((0,1]\times T(E_1))$ giving us the equality $T(E_1)=T(E)_1$ and by Theorem $2.44$ we have $m(T(E_1))=|\det T|m(E_1)=m(E_1)$ as $|\det T|=1$ here.
\end{solution}
\begin{problem}\label{prob:2.63}
The technique used to prove Proposition 2.54 can also be used to integrate any polynomial over $S^{n-1}$. In fact, suppose $f(x) = \prod_1^n x_j^{\alpha_j}$ ($\alpha_j \in \bN \cup \{0\}$) is a monomial. Then $\int f \, d\sigma = 0$ if any $\alpha_j$ is odd, and if all $\alpha_j$'s are even,
\[ \int f \, d\sigma = \frac{2\Gamma(\beta_1) \dots \Gamma(\beta_n)}{\Gamma(\beta_1 + \dots + \beta_n)} \quad \text{where } \beta_j = \frac{\alpha_j + 1}{2}. \]
\end{problem}
\begin{solution}
    \[
		\int_{\bR^n}f(x)e^{-\norm{x}^2}\,dm(x)=\int_0^\infty r^{n-1}\int_{S^{n-1}}f(rx')e^{-r^2}\,d\sigma(x')\,dm(r)
    \]
	notice that $f(rx')=r^{\Sigma\alpha_i}f(x')$, this makes the above equal 
	\[
		\int_0^\infty r^{n-1+\Sigma\alpha_i}e^{-r^2}\int_{S^{n-1}} f(x')\,d\sigma(x')\,dm(r)=\int f\,d\sigma\int_0^\infty r^{n-1+\Sigma\alpha_i}e^{-r^2}dr
	\]
	Now, by a change of variable we see that,
	\[
		\int_0^\infty r^{n-1+\Sigma\alpha_i}e^{-r^2}dr= \frac{1}{2}\int_0^\infty r^{ \frac{n+\Sigma\alpha_i}{2}-1}e^{-r}= \frac{1}{2}\Gamma\qty( \frac{n+\alpha_1+\ldots+\alpha_n}{2})= \frac{1}{2}\Gamma(\beta_1+\ldots+\beta_n)
	\]
	And by fubini's theorem, unless some $\alpha_j$ is odd  we have,
	\[
	\int_{R^n}f(x)e^{-\norm{x}^2}\,dm(x)=\prod_{j=1}^n\int_{-\infty}^\infty x^{\alpha_j}e^{-x_j^2}\,dx_j=\prod_{j=1}^n2\cdot\frac{1}{2}\Gamma\qty(\frac{\alpha+1}{2})=\Gamma(\beta_1)\cdots\Gamma(\beta_n)
	\]
	now rearranging these proves the claim in this case. 
	If some $\alpha_j$ are odd then one of the integrals $\int_{-\infty}^\infty x^{\alpha_j}e^{-x_j^2}\,dx_j$ would be $0$ making $\int f\,d\sigma=0$ as $\Gamma(\Sigma\beta_i)>0$ here, this is clear from the definition of the gamma function.
\end{solution}
\begin{problem}\label{prob:2.64}
For which real values of $a$ and $b$ is $|x|^a |\log |x||^b$ integrable over $\{x \in \bR^n : |x| < \frac{1}{2}\}$? Over $\{x \in \bR^n : |x| > 2\}$?
\end{problem}
\begin{solution}
    By corollary $2.51$ we have, after attaching a suitable characteristic function, 
	\[
		\int_{B(0,1/2)}\norm{x}^a|\log(\norm{x})|^b\,dm(x)=\sigma(S^{n-1})\int_0^{1/2} r^{a+n-1}|log(r)|^bdr
	\]
	It is now trivial to check which $a,b$ work. 
	In the other case the integral is over $(2,\infty)$.
\end{solution}
\begin{problem}\label{prob:2.65}
Define $G : \bR^n \to \bR^n$ by $G(r, \phi_1, \dots, \phi_{n-2}, \theta) = (x_1, \dots, x_n)$ where
\begin{align*}
x_1 &= r \cos \phi_1, \quad x_2 = r \sin \phi_1 \cos \phi_2, \quad x_3 = r \sin \phi_1 \sin \phi_2 \cos \phi_3, \dots, \\
x_{n-1} &= r \sin \phi_1 \dots \sin \phi_{n-2} \cos \theta, \quad x_n = r \sin \phi_1 \dots \sin \phi_{n-2} \sin \theta.
\end{align*}
\begin{enumerate}[label=\alph*.]
    \item $G$ maps $\bR^n$ onto $\bR^n$, and $|G(r, \phi_1, \dots, \theta)| = |r|$.
	\item $\det D_{(r, \phi_1, \dots, \phi_{n-2}, \theta)}G = r^{n-1} \sin^{n-2} \phi_1 \sin^{n-3} \phi_2 \dots \sin \phi_{n-2}$.
    \item Let $\Omega = (0, \infty) \times (0, \pi)^{n-2} \times (0, 2\pi)$. Then $G|_\Omega$ is a diffeomorphism and $m(\bR^n \setminus G(\Omega)) = 0$.
    \item Let $F(\phi_1, \dots, \phi_{n-2}, \theta) = G(1, \phi_1, \dots, \phi_{n-2}, \theta)$ and $\Omega' = (0, \pi)^{n-2} \times (0, 2\pi)$. Then $(F|\Omega')^{-1}$ defines a coordinate system on $S^{n-1}$ except on a $\sigma$-null set, and the measure $\sigma$ is given in these coordinates by
    \[ d\sigma(\phi_1, \dots, \phi_{n-2}, \theta) = \sin^{n-2} \phi_1 \sin^{n-3} \phi_2 \dots \sin \phi_{n-2} \, d\phi_1 \dots d\phi_{n-2} \, d\theta. \]
\end{enumerate}
\end{problem}
\begin{solution}
\begin{enumerate}[label=(\alph*)]
	\item Applying the identity $\sin^2(\theta)+\cos^2(\theta)=1$ collapses $x_1^2+\ldots+x_n^2$ to $r^2$ proving the claim.
	\item Consider the transformations  $J(r,\phi_1,\ldots,\phi_{n-2},\theta)=(r\cos\phi_1,r\sin\phi_1,\ldots,\phi_2,\ldots,\phi_{n-2},\theta)$ and $G(x_1,w,\phi_2,\ldots,\phi_{n-2},\theta)=(x_1,G(w,\phi_2,\ldots,\phi_{n-2},\theta))$, then we see that $G=H\circ J$ and by the multivariate chain rule we must have $DG=DH\cdot DJ$, we can calculate $DJ$ easily and $DG$ by induction on the dimension, this proves the claim.
	\item The variables associated to the $m$-th dimension are subscripted by $m$. 
		This is trivial for $n=2$ as $G_2(\Omega_2)=\bR^2\setminus\{(x,0):x\geq 0\}$ and this is also clearly a diffeomorphism. 
		Now, we assume it to be true for $n-1$, that is $G_{n-1}$ is a diffeomorphism and $G_{n-1}(\Omega_{n-1})=U_{n-1}$ where $m(\bR^{n-1}\setminus U_{n-1})=0$. 
		We also see that $J_n(\Omega_n)=\bR\times(0,\infty)\times(0,\pi)^{n-3}\times(0,2\pi)=\bR\times\Omega_{n-1}$ and this is a diffeomorphism as we only alter the first two coordinates and this corresponds to $G_2$. 
		By the definition of $H_n$ and the assumption we see that $G_n(\Omega_n)=H_n(J_n(\Omega_n))=\bR\times U_{n-1}$ and this is also a diffeomorphism as the first coordinate is unaltered and the rest correpsonds to $G_{n-1}.$  
		Thus we have that, $\bR^n\setminus G_n(\Omega_n)=\bR^n\setminus(\bR\times U_{n-1})=\bR\times(\bR^{n-1}\setminus U_{n-1})$ and $G_n$ being a composition of diffeomorphisms, is one itself. 
		By our convention of $0\cdot\infty=0$ we see that this is of measure $0$ (one could also use the fact that this has measure $\int_\bR m_{n-1}(\bR^{n-1}\setminus u_{n-1})\,dm_1=\int 0\,dm_1=0$)
	\item By definition we see that $G_n(r,\omega)=rF_n(\omega)$, and as we had $|G_n(r,\omega)|=r$ we see that $F_n(\omega)=G_n(1,\omega)$ maps $\Omega'_n$ into $S^{n-1}$. 
		As $G_n$ was diffeomorphic on $\Omega_n=(0,\infty)\times\Omega'_n$ and acted injectively we must have $F_n$ act diffeomorphically on $\Omega'_n$ as its still $C^1$ (it shares partial derivatives with $G_n$ clearly.) 
		Thus, $F_n|_{\Omega'_n}^{-1}$ forms a valid coordinate system on $F_n(\Omega_n')\subseteq S^{n-1}$. 
		As $G_n(r,\omega)=rF_n(\omega)$ we see that all non-origin points outside a null set can be written uniquely in that form (radius times direction vector on the sphere) and as the radius can be arbitrary, the non-origin points we can not reach must be unreachable due to their direction not being in $F_n(\Omega_n')$. 
		Hence, $(\bR^n\setminus U_n)\setminus\{0\}=\{rx':r\in(0,\infty),x'\in S^{n-1}\setminus F_n(\Omega_n')\}$. 
		We thus see that, as $\chi_{(\bR^n\setminus U_n)\setminus\{0\}}(rx')=\chi_{S^{n-1}\setminus F_n(\Omega_n')}(x')$ and by polar integration, the measure of the $(\bR^n\setminus G_n(\Omega_n))\setminus\{0\}$ is
		\[
			\int_0^\infty\int_{S^n}\chi_{S^{n-1}\setminus F_n(\Omega_n)}(x')r^{n-1}\,d\sigma(x')\,dm(r)=\sigma(S^{n-1}\setminus F_n(\Omega_n))\int_0^\infty r^{n-1}\,dm(r)
		\]
		and we must have the spherical measure be zero here for the whole thing to be zero which was proved in part (c). 
		Now, we see that for borel sets $E\subseteq S^n$, 
		\[
			\int_0^1\int_{S^{n-1}}\chi_E(x')r^{n-1}\,d\sigma(x')dm(r)=\int_{\bR^n}\chi_{E_1}\,dm=\int_{G(\Omega_n)}\chi_{E_1}\,dm
		\]
		\[
		=\int_{\Omega_n}\underbrace{\chi_{E_1}(G(r,\omega))}_{\chi_{(0,1]}(r)\cdot\chi_E(F(\omega))}|\det D_{(r,\omega)}G|\,dm((r,\omega))
		\]
		\[
			=\int_0^\infty\chi_{(0,1]}(r)r^{n-1}\int_{\Omega'_n}\chi_{E}(F(\omega))\sin^{n-2}\phi_1\sin^{n-3}\phi_2\cdots\sin\phi_{n-2}\,dm_{n-1}(\omega)dm(r)
		\]
		where we used all the facts developed above freely. 
		Now, after some simplifications we have
		\[
		\sigma(E)=\int_{F^{-1}|_{\Omega_n'}^{-1}(E)}\sin^{n-2}\phi_1\sin^{n-3}\phi_2\cdots\sin\phi_{n-2}\,\mathrm{d}m((\phi_1,\phi_2,\ldots,\theta))
		\]	
		so we are done.
\end{enumerate}
\end{solution}
\setcounter{section}{2}
\stepcounter{section}
\section*{Chapter \thesection\ : Signed Measures and Differentiation}
\addcontentsline{toc}{section}{Chapter \thesection\ : Signed Measures and Differentiation}
\begin{problem}\label{prob:3.1}
Prove Proposition 3.1.
\end{problem}
\begin{solution}
Wlog suppose that $\nu$ omits the value $-\infty$, then if $E_n$ is a sequence increasing sets we have that,
\[
	\nu \qty(\bigcup_{n=1}^\infty E_n)=\nu(E_1)+ \sum_{n\geq 2}\nu(E_n\setminus E_{n-1})=\lim_{m\to\infty} \nu(E_1)+ \sum_{n\leq m}\nu(E_n\setminus E_{n-1})=\lim_{m\to\infty}\nu(E_m)
\]
And if it was decreasing with $\nu(E_1)$ finite then, $E_1\setminus E_n$ is increasing and we have
\[
	\nu(E_1)-\nu \qty(\bigcap E_n)=\nu \qty(E_1\setminus\bigcap E_n)=\nu \qty(\bigcup(E_1\setminus E_n))=\lim_{m\to\infty}\nu(E_1\setminus E_m)=\nu(E_1)- \lim_{m\to\infty}\nu(E_m)
\]
after which we can cancel out the finite term of $\nu(E_1)$.
\end{solution}
\begin{problem}\label{prob:3.2}
If $\nu$ is a signed measure, $E$ is $\nu$-null iff $|\nu|(E)=0$. Also, if $\nu$ and $\mu$ are signed measures, $\nu\perp\mu$ iff $|\nu|\perp\mu$ iff $\nu^{+}\perp\mu$ and $\nu^{-}\perp\mu$.
\end{problem}
\begin{solution}
If $|\nu|(E)=0$ then $\nu^+(E)=\nu^-(E)=0$ and as these are positive measures, by monotonicity see that $\nu^+(F)=\nu^-(F)=0$ for any measurable $F\subseteq E$ making $E$ $\nu$-null. 
Now, if $E$ was $\nu$-null then $|\nu|(E)>0$ would imply that $\nu^+(E)$ or $\nu^-(E)$ is strictly greater than $0$ and if we assume the first to be as described then $\nu(E\cap P)=\nu^+(E\cap P)>0$ which contradicts nullity of $E$ w.r.t $\nu$ as $E\cap P \subseteq E$. 
The other case leads to a contradiction similarly, thus we must have $|\nu|(E)=0$. 

As a measurable set $E$ is $|\nu|$-null iff $|\nu|(E)=0$ stemming from positivity we see that $\nu\perp\mu$ iff $|\nu|\perp\mu$ by the previous claim. 
We also see that $|\nu|\perp\mu$ iff $\nu^+,\nu^-\perp\mu$ as $|\nu|(E)=0$ iff $\nu^-(E),\nu^+(E)=0$, which proves the equivalence.
\end{solution}
\begin{problem}\label{prob:3.3}
Let $\nu$ be a signed measure on $(X, \mM)$.
\begin{enumerate}[label=\alph*.]
    \item $L^{1}(\nu)=L^{1}(|\nu|)$.
    \item If $f\in L^{1}(\nu)$, $|\int f~d\nu|\le\int|f|d|\nu|$.
    \item If $E\in\mM$, $|\nu|(E) = \sup \{ |\int_E f\, d\nu| : |f| \le 1 \}$.
\end{enumerate}
\end{problem}
\begin{solution}
\begin{enumerate}[label=\alph*.]
\item As $L^1(\nu):=L^1(\nu^+)\cap L^1(\nu^-)$ we find that, $f\in L^1(\nu)$ implies $ \int_{}^{}|f|,\mathrm{d}\nu^+, \int_{}^{}|f|\,\mathrm{d}\nu^-<\infty$ and thus, $ \int_{}^{}|f|\,\mathrm{d}|\nu|= \int_{}^{}|f|\,\mathrm{d}\nu^++ \int_{}^{}|f|\,\mathrm{d}\nu^-<\infty$ where we have used the following lemma: If $\mu_1,\mu_2$ are positive measures on the same sigma algebra then so is their sum and if $f\in L^+$ then $ \int_{}^{}f\,\mathrm{d}(\mu_1+\mu_2)= \int_{}^{}f\,\mathrm{d}\mu_1+ \int_{}^{}f\,\mathrm{d}\mu_2$. 
For a proof, one can look at simple functions first where its trivial and then routinely extend to $L^+$ by MCT and approximations using simple functions. 
Now, if we had $f\in L^1(|\nu|)$ then as $ \int_{}^{}f\,\mathrm{d}|\nu|=\int_{}^{}|f|\,\mathrm{d}\nu^++ \int_{}^{}|f|\,\mathrm{d}\nu^-<\infty$, both of these terms are individually finite which puts $f$ in $L^1(\nu^+)$ and $L^1(\nu^-)$ proving the claim.
\item \[
    \qty|\int fd\nu|=\qty|\int fd\nu^+-\int fd\nu^-|
\]
\[\leq\qty|\int fd\nu^+|+\qty|\int fd\nu^-|\leq\int|f|d\nu^++\int|f|d\nu^-=\int|f|d|\nu|
\]
 
\item Consider the hahn decomposition $P\cup N$ is usual notation leading upto the jordan decomposition, then letting $f=\chi_P-\chi_N$ we see that $ \int_{}^{}f\,\mathrm{d}\nu=\nu(P\cap E)-\nu(N\cap E)=\nu^+(E)+\nu^-(E)=|\nu|(E)$ so we have one direction of the inequality ($\leq$) and for the other direction, by part (b) we can say that $|f|\leq 1$ implies $|\int_{}^{}f\,\mathrm{d}\nu|\leq \int_{}^{}1\,\mathrm{d}|\nu|=|\nu|(E).$ so we are done.
\end{enumerate}
\end{solution}
\begin{problem}\label{prob:3.4}
If $\nu$ is a signed measure and $\lambda$, $\mu$ are positive measures such that $\nu=\lambda-\mu$, then $\lambda\ge\nu^{+}$ and $\mu\ge\nu^{-}$.
\end{problem}
\begin{solution}
Consider the hahn decomposition $P\cup N$ in usual notation, leading upto the jordan decomposition. 
Then we have $\nu^+(E)=\nu(E\cap P)=\lambda(E\cap P)-\mu(E\cap P)\leq\lambda(E\cap P)\leq\lambda(E)$ and the other inequality follows similarly. 
\end{solution}
\begin{problem}\label{prob:3.5}
If $\nu_{1}$, $\nu_{2}$ are signed measures that both omit the value $+\infty$ or $-\infty$, then $|\nu_{1}+\nu_{2}|\le|\nu_{1}|+|\nu_{2}|$. (Use Exercise 4.)
\end{problem}
\begin{solution}
It can be easily shown that in this case $\nu_1+\nu_2$ is also a signed measure. 
As all $\nu_1,\nu_2$ are finite, we can rearrange to get $\nu_1+\nu_2=\nu_1^+-\nu_1^-+\nu_2^+-\nu_2^-=(\nu_1^++\nu_2^+)-(\nu_1^-+\nu_2^-)$ which, by the previous problem implies $(\nu_1+\nu_2)^+\leq\nu_1^++\nu_2^+$ and $(\nu_1+\nu_2)^-\leq\nu_1^-+\nu_2^-$ implying $|\nu_1+\nu_2|\leq(\nu_1^++\nu_2^+)+(\nu_1^-+\nu_2^-)=(\nu_1^++\nu_1^-)+(\nu_2^++\nu_2^-)\leq|\nu_1|+|\nu_2|$.
\end{solution}
\begin{problem}\label{prob:3.6}
Suppose $\nu(E)=\int f~d\mu$ where $\mu$ is a positive measure and $f$ is an extended $\mu$-integrable function. Describe the Hahn decompositions of $\nu$ and the positive, negative, and total variations of $\nu$ in terms of $f$ and $\mu$.
\end{problem}
\begin{solution}
	Let $P= \qty{f>0}$ and $N= \qty{f\leq 0}$, then $P\cup N$ is clearly a hahn decomposition and we see that $\nu^+(E)= \int_{E\cap P}^{}f\,\mathrm{d}\mu= \int_{E}^{}f^+\,\mathrm{d}\mu$ and similarly $\nu^-(E)= \int_{E}^{}f^-\,\mathrm{d}\mu$ and thus, $|\nu|(E)= \int_{E}^{}|f|\,\mathrm{d}\mu$.
\end{solution}
\begin{problem}\label{prob:3.7}
Suppose that $\nu$ is a signed measure on $(X, \mM)$ and $E\in\mM$.
\begin{enumerate}[label=\alph*.]
    \item $\nu^{+}(E)=\sup\{\nu(F):F\in\mM,F\subset E\}$ and $\nu^{-}(E)=-\inf\{\nu(F):F\in\mM, F\subset E\}$.
    \item $|\nu|(E)=\sup\{\sum_{1}^{n}|\nu(E_{j})|:n\in\bN,E_{1},...,E_{n} \text{ are disjoint, and } \bigcup_{1}^{n}E_{j}=E\}$.
\end{enumerate}
\end{problem}
\begin{solution}
\begin{enumerate}[label=\alph*.]
	\item $F \subseteq E$ implies that $\nu(F)\leq\nu^+(F)\leq\nu^+(E)$ and $\nu(E\cap P)=\nu^+(E)$ with $E\cap P \subseteq E$ proving one part, the other follows similarly.
	\item Considering a hahn decomposition $P\cup N$ leading to a jordan decomposition, we see that $|\nu|(E)=\nu^+(E)+\nu^-(E)=|\nu(E\cap P)|+|\nu(E\cap N)|$ as $\nu^+(E)=\nu(E\cap P)\geq 0$ and $\nu^-(E)=-\nu(E\cap N)\geq 0$ proving one direction of the equality. 
	For the other, it can be easily proved that $|\nu(E)|\leq|\nu|(E)$ and thus, $\bigcup_1^nE_j=E$ (disjoint union) implies $\sum_1^n|\nu(E_j)|\leq\sum_1^n|\nu|(E_j)=|\nu|(E)$ so we are done.
\end{enumerate}
\end{solution}
\begin{problem}\label{prob:3.8}
$\nu\ll\mu$ iff $|\nu|\ll\mu$ iff $\nu^{+}\ll\mu$ and $\nu^{-}\ll\mu$.
\end{problem}
\begin{solution}
If $\nu\ll\mu$ then $\mu(E)=0\implies\nu(E)=0$ and as $\mu$ is taken positive we must have that $\mu(E\cap P),\mu(E\cap N)=0$ where $P\cup N$ is the hahn decomposition for $\nu$, this implies that $\nu(E\cap P),\nu(E\cap N)=0$ so $\nu^\pm(E)=0$ and we have that $|\nu|(E)=0$ as well proving $|\nu|\ll\mu$. 
As $|\nu|(E)\geq|\nu(E)|$ we have that $|\nu|\ll\mu$ implies $\nu\ll\mu$ and hence $\nu\ll\mu$ iff $|\nu|\ll\mu$. 
If $\nu^+,\nu^-\ll\mu$ then $\mu(E)=0$ clearly implies $|\nu|(E)=0$ and if we have $|\nu|\ll\mu$ then $\mu(E)=0$ implies $\nu^+(E)+\nu^-(E)=0$ and thus $\nu^+(E)=\nu^-(E)=0$ proving $\nu^+,\nu^-\ll\mu$ and thus all these conditions are equivalent.
\end{solution}
\begin{problem}\label{prob:3.9}
Suppose $\{\nu_{j}\}$ is a sequence of positive measures. If $\nu_{j}\perp\mu$ for all $j$, then $\sum_{1}^{\infty}\nu_{j}\perp\mu$; and if $\nu_{j}\ll\mu$ for all $j$, then $\sum_{1}^{\infty}\nu_{j}\ll\mu$. (This is also true for complex measures when $\sum_1^\infty\nu_j$ is a valid complex measure too)
\end{problem}
\begin{solution}
Say $E_j$ is $\nu_j$-null and $E_j^c$ is $\mu$-null for all $j$, then we see that, let $E=\bigcap E_j$, then we clearly see that $E$ is $\sum_1^\infty\nu_j$-null and we claim that, $E^c=\bigcup E_j^c$ is $\mu$-null. 
For any measurable $F\subseteq E^c$, we can write $F=\bigcup_{j=1}^\infty(F\cap(E_j^c\setminus\cup_{i\neq j}^\infty E_i^c))$ is a disjoint union and 
\[
    \mu(F)= \sum_{j=1}^{\infty}\mu(F\cap(E_j^c\setminus\cup_{i\neq j}^\infty E_i^c))=0
\]
as $(F\cap(E_j^c\setminus\cup_{i\neq j}^\infty E_i^c)) \subseteq E_j^c$. 
The absolute continuity result is trivial. 
\end{solution}
\begin{problem}\label{prob:3.10}
Theorem 3.5 may fail when $\nu$ is not finite. (Consider $d\nu(x)=dx/x$ and $d\mu(x)=dx$ on $(0, 1)$, or $\nu = \text{counting measure}$ and $\mu(E)=\sum_{n\in E}2^{-n}$ on $\bN$.)
\end{problem}
\begin{solution}
We choose the second counterexample. 
Clearly we have, $\nu\ll\mu$ but for any $\delta>0$ we see that there exists a $E_\delta$ with $\mu(E_\delta)<\delta$ but $\nu(E_\delta)\geq 1$ which just reads ``$E$ is nonempty'' and as the sum $\sum_{n\in\bN}2^{-n}$ converges absolutely we can find such a $E_\delta$ always.
\end{solution}
\begin{problem}\label{prob:3.11}
Let $\mu$ be a positive measure. A collection of functions $\{f_{\alpha}\}_{\alpha\in A}\subset L^{1}(\mu)$ is called uniformly integrable if for every $\epsilon>0$ there exists $\delta>0$ such that $|\int_{E}f_{\alpha}d\mu|<\epsilon$ for all $\alpha\in A$ whenever $\mu(E)<\delta$.
\begin{enumerate}[label=\alph*.]
    \item Any finite subset of $L^{1}(\mu)$ is uniformly integrable.
    \item If $\{f_{n}\}$ is a sequence in $L^{1}(\mu)$ that converges in the $L^{1}$ metric to $f\in L^{1}(\mu)$, then $\{f_{n}\}$ is uniformly integrable.
\end{enumerate}
\end{problem}
\begin{solution}
\begin{enumerate}[label=\alph*.]
	\item Consider $A=\qty{1,\ldots,n}$; for any $\epsilon>0$, as $f_jd\mu\ll\mu$ where $f_jd\mu$ are finite meausres, we can find $\delta_j>0$ such that $|\int_{E}f_jd\mu|<\epsilon$ for all $j\in A$ whenever $\mu(E)<\delta_j$ so taking $\delta=\min_j\delta_j$ suffices.
	\item As $f_n$ converges in $L^1$, it is cauchy in this metric too so for any $\epsilon>0$ we can find $N$ such that $m,n\geq N\implies|\int_Ef_m-\int_Ef_n|\leq\int|f_m-f_n|<\epsilon$ and by the previous problem, as $\{f_1,\ldots,f_N\}$ are uniformly integrable, we can find $\delta>0$ such that, $|\int_E f_j|<\epsilon<2\epsilon$ for $1\leq j\leq N$ whenever $\mu(E)<\delta$ and for $n>N,$
		\[
		    \qty|\int_Ef_nd\mu|\leq\qty|\int_Ef_Nd\mu|+\qty|\int_Ef_n-f_Nd\mu|<\epsilon+\epsilon=2\epsilon
		\]
	so we are done.
\end{enumerate}
\end{solution}
\begin{problem}\label{prob:3.12}
For $j=1,2,$ let $\mu_{j}$, $\nu_{j}$ be $\sigma$-finite measures on $(X_{j},\mM_{j})$ such that $\nu_{j}\ll\mu_{j}$. Then $\nu_{1}\times\nu_{2}\ll\mu_{1}\times\mu_{2}$ and
\[\frac{d(\nu_{1}\times\nu_{2})}{d(\mu_{1}\times\mu_{2})}(x_{1},x_{2})=\frac{d\nu_{1}}{d\mu_{1}}(x_{1})\frac{d\nu_{2}}{d\mu_{2}}(x_{2}).\]
\end{problem}
\begin{solution}
Suppose $\mu_1\times\mu_2(E)=0$ then we see that $0=\int\mu_1(E_x)d\mu_2(x)$ implying $\mu_1(E_x)=0$ (this in turn implies $\nu_1(E_x)=0$ here) for $\mu_2$-a.e. $x$ so suppose its true for all $x$ outside $N$ with $\mu_2(N)=0$ which implies $\nu_2(N)=0$ so we have that $\nu_1\times\nu_2(E)=\int\nu_1(E_x)d\nu_2(x)=0$ as $\nu_1(E_x)=0$ for $\nu_2$-a.e. $x$ thus proving $\nu_1\times\nu_2\ll\mu_1\times\mu_2.$ 
CLearly these product measures are sigma finite too, so we can find a radon nikodym derivative $f$ i.e. $d(\nu_1\times\nu_2)=fd(\mu_1\times\mu_2)$ and also $f_1,f_2$ such that $f_jd\mu_j=\nu_j$ for $j=1,2$ and by fubini's we see that,
\[
	\int\chi_E(x,y)f_1(x)f_2(y)d(\mu_1\times\mu_2(x,y))=\int\int f_1(x)\chi_{E_x}(y)f_2(y)d\mu_2(y)d\mu_1(x)
\]
\[
    =\int\nu_2(E_x)f_1(x)d\mu_1(x)=\int\nu_2(E_x)d\nu_1(x)=\nu_1\times\nu_2(E)=\int\chi_Ef(x,y)d(\mu_1\times\mu_2(x,y))
\]
which proves that $f(x,y)=f_1(x)f_2(y)$ for $\mu_1\times\mu_2$ a.e. $(x,y)$ so we have the equality treating them as equivalence classes.
\end{solution}
\begin{problem}\label{prob:3.13}
Let $X=[0,1]$, $\mM=\mathcal{B}_{[0,1]}$, $m=$ Lebesgue measure, and $\mu=$ counting measure on $\mM$.
\begin{enumerate}[label=\alph*.]
    \item $m\ll\mu$ but $dm\ne f~d\mu$ for any $f$.
    \item $\mu$ has no Lebesgue decomposition with respect to $m$.
\end{enumerate}
\end{problem}
\begin{solution}
\begin{enumerate}[label=\alph*.]
	\item As $\mu(E)=0\implies E=\emptyset$, $m\ll\mu$ holds. 
		If we had $dm=fd\mu$ then we would have $m( \qty{x})=f(x)=0$ for all $x$ which is absurd as it implies $dm=0$. 
    \item If we had $\mu=\rho+\lambda$ where $\lambda\perp m$ and $\rho\ll m$, notice how for all singletons we have $m(\{x\})=0\implies\rho(\{x\})=0$ and thus $1=\mu(\{x\})=0+\lambda(\{x\})=\lambda(\{x\})$ and we can find a partition $L\cup M$ where $\lambda(M)=m(L)=0$. 
        Now, we see that for any $E$ we have $\lambda(E)=\lambda(E\cap L)=\lambda(E\cap L)+m(E\cap L)=\mu(E\cap L)\geq 0$ as $m(E\cap L)\leq m(L)=0$ proving that $\lambda$ is positive.        Then, we must have that $\lambda(M)=0\implies\lambda(C)=0$ where $C\subseteq M$ is countably infinite, we can choose such a $C$ because $M$ is required to be uncountable ($m(M)=1$.) 
        From this we see that $\infty=\mu(C)=\rho(C)+\lambda(C)=0+0$ as $m(C)=0$ forces $\rho(C)=0$ so we have a contradiction and the hypothesis was false to begin with. 
\end{enumerate}
\end{solution}
\begin{problem}\label{prob:3.14}
If $\nu$ is an arbitrary signed measure and $\mu$ is a $\sigma$-finite measure on $(X, \mM)$ such that $\nu\ll\mu$, there exists an extended $\mu$-integrable function $f:X\rightarrow[-\infty,\infty]$ such that $d\nu=f~d\mu$. Hints:
\begin{enumerate}[label=\alph*.]
	\item It suffices to assume $\mu$ is finite and $\nu$ is positive.
	\item With these assumptions there exists $E\in\mM$ that is $\sigma$-finite for $\nu$ and $\mu(E)\geq\mu(F)$ for all sets $F$ that are $\sigma$-finite for $\nu$.
	\item The radon-nikodym theorem applies on $E$. If $F\cap E=0$, then either $\nu(F)=\mu(F)=0$ or $\mu(F)>0$ and $|\nu(F)|=\infty$.
\end{enumerate}
\end{problem}
\begin{solution}
If this is true for positive $\nu$ and finite $\mu$, then for sigma finite $\mu$ we can write $X=\bigcup_1^\infty E_n$ with $E_n$ being $\mu$-finite and mutually disjoint we have $d\nu_n=f_nd\mu_n$ (we safely set $f_n=0$ outside $E_n$) where $\mu_n(E):=\mu(E\cap E_n)$ and $\nu_n(E)=\nu(E\cap E_n)$ as $\mu_n$ is finite and clearly $\nu_n\ll\mu_n$ still holds. 
Then, with $f=\sum_1^\infty f_n$ we would have $d\nu=fd\mu$. 
If $\nu$ was signed on top of that, then using \hyperref[prob:3.8]{Problem 3.8} which states $\nu\ll\mu\implies\nu^+,\nu^-\ll\mu$ we can find $f^+,f^-$ with $d\nu^-=f^-d\mu,d\nu^+=f^+d\mu$ we have $d\nu=d\nu^+-d\nu^-=(f^+-f^-)d\mu$ and $f^+-f^-$ is clearly extended integrable as $\nu$, by definition can only assume atmost one of the values $+\infty,-\infty$. 
We now prove it for this case. 
Consider the set $ \qty{\mu(H):H\in\mM,H\text{ is }\nu\,\sigma\text{-finite}}$, let $s$ be its supremum, then $s<\infty$ as $\mu$ is finite and we can find a $\nu$ $\sigma$-finite sequence of sets $F_n$ such that $\mu(F_n)\longrightarrow s$, then we can clearly see that $\mu(\cup F_n)=s$ as $\cup F_n$ is clearly still $\nu$ $\sigma$-finite and we can show both sides of the equality easily. 
Setting $E=\cup F_n$ gives us a set as described in hint (b).
The radon nikodym theorem applies on $E$ clearly. 
If $F\cap E=\emptyset$ then if we have $\nu(F)>0$ we must have $\mu(F)>0$ as $\nu\ll\mu$which forces $\nu(F)=\infty$ as otherwise the existence of $E\cup F$ contradicts the maximality of $E$ and if we have $\mu(F)>0$ then again we similarly have $\nu(F)=\infty$ such as to not contradict the maximality of $E$ so we have proved the claim in hint (c). 
There are now two cases, first one being when $\nu(E^c)=\mu(E^c)=0$, in this case we can apply radon-nikodym to get a derivative $f$ on $E$ and extend it to the rest of the set by setting it to be zero outside this, then we see that,
\[
	\nu(K)=\nu(K\cap E)+\underbrace{\nu(K\cap E^c)}_{\leq\nu(E^c)=0}=\int_{K\cap E}f\,\mathrm{d}\mu=\int_K f\,\mathrm{d}\mu
\]
as $\nu$ was assumed to be positive and $\mu(E^c)=0$ so we have $d\nu=fd\mu$ in general so we are done.
In the other case where $\mu(E^c)>0$ and $\nu(E^c)=\infty$ we can set $f$ to be $\infty$ on $E^c$, then we see that $\nu(E\cap K)=\int_{E\cap K}fd\mu$ as usual and either $\nu(E^c\cap K)=\mu(E^c\cap K)=0$ or $\nu(E^c\cap K)=\infty$ and $\mu(E^c\cap K)>0$ so we see that $\int_{E^c\cap K}fd\mu$ and $\nu(E^c\cap K)$ are either both infinite or both $0$ by part (c) of the hint as $E^c\cap K$ is disjoint from $E$, hence we have the equality,
\[
	\nu(K)=\nu(E\cap K)+\nu(E^c\cap K)=\int_{E\cap K}fd\mu+\int_{E^c\cap K}fd\mu=\int_Kfd\mu
\]
so we are done.
\end{solution}
\begin{problem}\label{prob:3.15}
A measure $\mu$ on $(X, \mM)$ is called decomposable if there is a family $\mathfrak{F}\subset\mM$ with the following properties: (i) $\mu(F)<\infty$ for all $F\in\mathfrak{F}$; (ii) the members of $\mathfrak{F}$ are disjoint and their union is $X$; (iii) if $\mu(E)<\infty$ then $\mu(E)=\sum_{F\in\mathfrak{F}}\mu(E\cap F)$; (iv) if $E\subset X$ and $E\cap F\in\mM$ for all $F\in\mathfrak{F}$ then $E\in\mM$.
\begin{enumerate}[label=\alph*.]
    \item Every $\sigma$-finite measure is decomposable.
    \item If $\mu$ is decomposable and $\nu$ is any signed measure on $(X, \mM)$ such that $\nu\ll\mu$, there exists a measurable $f:X\rightarrow[-\infty,\infty]$ such that $\nu(E)=\int_{E}f\,d\mu$ for any $E$ that is $\sigma$-finite for $\mu$, and $|f|<\infty$ on any $F\in\mathfrak{F}$ that is $\sigma$-finite for $\nu$. (Use Exercise $14$ if $\nu$ is not $\sigma$-finite.) 
\end{enumerate}
\end{problem}
\begin{solution}
\begin{enumerate}[label=\alph*.]
	\item It can be easily shown that $\mathfrak F=\{E_n:n\in\bN\}$ works where $E_n$ partition $X$ while being $\mu$-finite.
	\item Let $\mathfrak F$ be as described in the problem. 
	As $\mu(F)<\infty$ for any $F\in\mathfrak F$ we can find a measurable $f_F$ on  $F$ such that $\nu(E\cap F)=\int_{E\cap F}f_Fd\mu$ for all $E\in\mM$ by the previous problem and if $F$ is $\sigma$-finite for $\nu$ we can clearly force $|f_F|<\infty$ (in the proof presente in the book we can take all $f_j$'s which are zero outside $A_j$ to be real valued everywhere as shown in case I which makes their sum real valued too as the $A_j$ are disjoint.) as in the proof of the radon-nikodym theorem. 
	We claim that $f:=f_F$ on $F\in\mathfrak F$ is measurable and it has the properties needed. 
	For any $E\in\mathcal B_{\overline{\bR}}$ we see that $f^{-1}(E)=\bigcup_{F\in\mathfrak F}f^{-1}(E)\cap F$ and each $f^{-1}(E)\cap F=f|_{F}^{-1}(E)\cap F=f_F^{-1}(E)\cap F\in\mM$ as $f_F$ is measurable which implies that $f^{-1}(E)$ by property (iv). 
	Now, for any $E$ that is $\sigma$-finite for $\mu$, suppose $E=\cup E_n$ (disjoint union) with $\mu(E_n)<\infty$, then we see that $\mu(E_n)=\sum_{\mathfrak F}\mu(E_n\cap F)<\infty$ and it is well known that only atmost countably many $\mu(E_n\cap F)$ are nonzero, so we can find $\{F_m\}\subseteq\mathfrak F$ such that $\mu(E_n)=\sum_{m=1}^\infty\mu(E_n\cap F_m)$ and thus we also see that $\mu(E_n\setminus(\cup_m(E_n\cap F_m)))=0$ which, by absolute continuity implies $\nu(E_n)=\nu(\cup_m(F_m\cap E_n))=\sum_{m=1}^\infty\mu(E_n\cap F_m)=\sum_{m=1}^\infty\int_{E_n\cap F_m}fd\mu=\int_{E_n}fd\mu$ and thus, $\nu(E)=\sum_1^\infty\nu(E_n)=\sum_1^\infty\int_{E_n}fd\mu=\int_Efd\mu$. 
\end{enumerate}
\end{solution}
\begin{problem}\label{prob:3.16}
Suppose that $\mu, \nu$ are measures on $(X, \mM)$ with $\nu\ll\mu$, and let $\lambda=\mu+\nu$. If $f=d\nu/d\lambda$, then $0\le f<1$ $\lambda$-a.e. and $d\nu/d\mu=f/(1-f)$.
\end{problem}
\begin{solution}
(note: we treat the derivatives as equivalence classes here.) Firstly, it is obvious that $f\geq 0$ $\lambda$-a.e We have $\int_Ed\nu=\nu(E)=\int_Efd(\mu+\nu)\leq\int_Ed\nu+\int_Ed\mu=\int_E d(\mu+\nu)$, it can then be easily shown that $f\leq 1$ $\lambda$-a.e. 
If we had $f=1$ over $E$ with $\lambda(E)=\nu(E)+\mu(E)>0$ then we see that, $\nu(E)=\int_Efd\lambda=\lambda(E)$ which implies that $\mu(E)=0$ but then as $\nu\ll\mu$ we have $\nu(E)=0$ too contradicting non-nullity, this proves that $f<1$ $\lambda$-a.e. too so we have proven the first claim. 
We clearly have $\nu\ll\mu\ll\lambda$ and $\lambda\ll\mu$.
Thus, $g=d\mu/d\lambda$ exists and by Proposition $3.9$ we see that $f=d\nu/d\lambda=(d\nu/d\mu)(d\mu/d\lambda)=g(d\nu/d\mu)$ and substituting these in we get $d\nu/d\mu=f/g=f/(1-f)$ as $f+g=d\nu/d\lambda+d\mu/d\lambda=d(\mu+\nu)/d\lambda=d\lambda/d\lambda=1.$
\end{solution}
\begin{problem}\label{prob:3.17}
	Let $(X, \mM, \mu)$ be a $\sigma$-finite measure space, $\mN$ a sub-$\sigma$-algebra of $\mM$, and $\nu=\mu|_{\mN}$ is $\sigma$-finite (This condition is missing in the book statement, without which this is false: Consider the lebesgue measure on $\bR$ as the measure space and take $\mN=\{\emptyset,\bR\}, f=\chi_{[0,1]}$.) If $f\in L^{1}(\mu)$, there exists $g\in L^{1}(\nu)$ (thus $g$ is $\mN$-measurable) such that $\int_{E}f~d\mu=\int_{E}g~d\nu$ for all $E\in\mN$; if $g^{\prime}$ is another such function then $g=g^{\prime}$ $\nu$-a.e. \\
(In probability theory, $g$ is called the conditional expectation of $f$ on $\mN$.)
\end{problem}
\begin{solution}
	As $f\in L^1$ we have that $fd\mu$ is a signed measure and it is clearly still a measure restricted to $\mN$ where we see that $(fd\mu)|_{\mN}\ll\nu=\mu|_{\mN}$ and as $\mN$ is sigma finite for $\nu$ we can apply \hyperref[prob:3.14]{Problem 3.14} to deduce existence of some $g$ such that $(fd\mu)|_{\mN}=gd\nu$ which implies $\int_Efd\mu=\int_Egd\nu$ for all $E\in\mN.$ 
    Here, $g$ is extended $\nu$ integrable by construction. 
    To show that $g\in L^1(\nu)$ consider the fact that $(|f|d\mu)|_{\mN}=|(fd\mu)|_{\mN}|=|gd\nu|=|g|d\nu$ and thus, $\int_X|g|d\nu=\int_X|f|d\mu<\infty\implies g\in L^1(\nu)$. 
    The uniqueness follows from the fact that $\int_Egd\nu=\int_Eg'd\nu$ for all $E\in\mN$ implies $g=g'$ $\nu$-a.e.
\end{solution}
\begin{problem}\label{prob:3.18}
Prove Proposition 3.13c.
\end{problem}
\begin{solution}
Let $g$ be as in the definition ($d\nu=gd\mu$ with $\mu$ being positive) of $|\nu|$ then we see that $|f|d|\nu|=|f||g|d\mu$ and we also have $fd\nu=fgd\mu$ and we have that,
\[
    \qty| \int_{}^{}f\,\mathrm{d}\nu|=\qty| \int_{}^{}fg\,\mathrm{d}\mu|\leq\int|fg|d\mu=\int|f|d|\nu|
\]
\end{solution}
\begin{problem}\label{prob:3.19}
If $\nu, \lambda$ are complex measures and $\mu$ is a positive measure, then $\nu\perp\mu$ iff $|\nu|\perp|\mu|$, and $\nu\ll\lambda$ iff $|\nu|\ll\lambda$.
\end{problem}
\begin{solution}
    First, we prove $\nu\perp\mu \iff |\nu|\perp\mu$. 
    If $\nu\perp\mu$, then by definition $\nu_r\perp\mu$ and $\nu_i\perp\mu$. By \hyperref[prob:3.2]{Problem 3.2}, this implies $|\nu_r|\perp\mu$ and $|\nu_i|\perp\mu$. Using \hyperref[prob:3.9]{Problem 3.9}, their sum is also mutually singular to $\mu$, meaning $(|\nu_r|+|\nu_i|)\perp\mu$. Since total variation satisfies $|\nu| \leq |\nu_r|+|\nu_i|$, any set that is null for $|\nu_r|+|\nu_i|$ is also null for $|\nu|$. Therefore, $|\nu|\perp\mu$.
    Conversely, if $|\nu|\perp\mu$, since $|\nu_r| \leq |\nu|$ and $|\nu_i| \leq |\nu|$, we have $|\nu_r|\perp\mu$ and $|\nu_i|\perp\mu$. By \hyperref[prob:3.2]{Problem 3.2}, this implies $\nu_r\perp\mu$ and $\nu_i\perp\mu$, meaning $\nu\perp\mu$. 
    Now, we prove $\nu\ll\lambda \iff |\nu|\ll\lambda$. 
    By definition, absolute continuity with respect to a complex measure means absolute continuity with respect to its total variation, so $\nu\ll\lambda$ means $|\lambda|(E)=0 \implies \nu(E)=0$.
    If $|\nu|\ll\lambda$, then $|\lambda|(E)=0 \implies |\nu|(E)=0$. Since $|\nu(E)| \leq |\nu|(E)$, this immediately implies $\nu(E)=0$, proving $\nu\ll\lambda$.
    Conversely, suppose $\nu\ll\lambda$ and let $|\lambda|(E)=0$. Since $|\lambda|$ is a positive measure, for any measurable subset $F \subseteq E$, we also have $|\lambda|(F)=0$, which implies $\nu(F)=0$. The total variation can be equivalently defined as $|\nu|(E) = \sup \sum_{j=1}^n |\nu(E_j)|$ over all finite measurable partitions $\{E_j\}$ of $E$ by \hyperref[prob:3.21]{Problem 3.21}. Since each $E_j \subseteq E$, we have $\nu(E_j)=0$ for all $j$. Thus, the supremum is exactly $0$, giving $|\nu|(E)=0$. This proves $|\nu|\ll\lambda$.
\end{solution}
\begin{problem}\label{prob:3.20}
If $\nu$ is a complex measure on $(X,\mM)$ and $\nu(X)=|\nu|(X)$, then $\nu=|\nu|$.
\end{problem}
\begin{solution}
We see that $\nu(X)=|\nu(X)|=|\nu(E)+\nu(X\setminus E)|\leq|\nu(E)|+|\nu(X\setminus E)|\leq|\nu|(E)+|\nu|(X\setminus E)=|\nu|(X)=\nu(X)$ so all the inequalities must be equalities, which implies that $|\nu(E)|=|\nu|(E)$ as $|\nu(E)|\leq|\nu|(E)$ for all $E\in\mM.$ 
As we have $|\nu(E)|+|\nu(X\setminus E)|=|\nu(E)+\nu(X\setminus E)|$ which is real, those two complex numbers share the same arguement (unless $0$ but that case is trivial) and as $\nu(E)+\nu(X\setminus E)=\nu(X)$ is real too, they must have no complex part making them real which implies that, $\nu(E)=|\nu(E)|=|\nu|(E)$ proving the claim.
\end{solution}
\begin{problem}\label{prob:3.21}
Let $\nu$ be a complex measure on $(X, \mM)$. If $E\in\mM$ define
\[\mu_{1}(E)=\sup\qty{\sum_{1}^{n}|\nu(E_{j})|:n\in\bN,E_{1},...,E_{n} \text{ disjoint}, E=\bigcup_{1}^{n}E_{j}},\]
\[\mu_{2}(E)=\sup\qty{\sum_{1}^{\infty}|\nu(E_{j})|:E_{1},E_{2},... \text{ disjoint}, E=\bigcup_{1}^{\infty}E_{j}},\]
\[\mu_{3}(E)=\sup\qty{\qty|\int_{E}f~d\nu|:|f|\leq 1}.\]
Then $\mu_{1}=\mu_{2}=\mu_{3}=|\nu|$. (First show that $\mu_1\leq\mu_2\leq\mu_3$· To see that $\mu_3 = |\nu|$, let $f=\overline{d\nu/d|\nu|}$ and apply Proposition $3.13.$ To see that $\nu_3\leq\nu_1$, approximate $f$ by simple functions.)
\end{problem}
\begin{solution}
$\mu_1\leq\mu_2$ is obvious and $\mu_2\leq\mu_3$ holds as for any given $E$ and any partition $\qty{E_n}$ of $E$ we can define
\[
	f=\sum_n \frac{\overline{\nu(E_n)}}{|\nu(E_n)|}\cdot\chi_{E_n}
\]
where the sum is over all $n$ for which $\nu(E_n)\neq 0$, then we clearly see that $|f|\leq 1$ and $\int_Ef\,d\nu=\sum_1^\infty\int_{E_n}f\,d\nu=\sum_1^\infty|\nu(E_n)|$ as $|z|=\overline{z}z/|z|$ for nonzero $z.$ 
Clearly $\mu_3(E)\leq\int_E d|\nu|=|\nu|(E)$ by Proposition $3.13c$ and with $f$ as in the hint (this can be made to have absolute value atmost $1$ everywhere after changing it on a null set, not affecting any of the integrals below) we see that, 
\[
	\int_E\overline{ \frac{d\nu}{d|\nu|}}d\nu=\int_E\overline{ \frac{d\nu}{d|\nu|}} \frac{d\nu}{d|\nu|}d|\nu|=\int_E|f|^2d|\nu|=\int_Ed|\nu|=|\nu|(E)
\]
as $|f|=1$ $|\nu|$-a.e. so we must have $\mu_3=|\nu|$. 
Now, we can  find a sequence of simple functions $\phi_n$ converging $f$ with $|\phi_n|\leq |f|\leq 1$ and convergence can be made uniform here as $f$ is bounded on $E$, this clearly implies that $\lim|\int_E\phi_n\,d\nu|=|\int_E f\,d\nu|$ and thus, as $|\int_E\phi_n\,d\nu|\leq\mu_1(E)$ due to linearity and the triangle inequality we must have $\lim_n|\int_E f\,d\nu|\leq\mu_1(E)$ too so we are done.
\end{solution}
\begin{problem}\label{prob:3.22}
If $f\in L^{1}(\bR^{n})$, $f\ne0,$ there exist $C, R > 0$ such that $Hf(x)\ge C|x|^{-n}$ for $|x|>R$. Hence $m(\{x:Hf(x)>\alpha\})\ge C^{\prime}/\alpha$ when $\alpha$ is small, so the estimate in the maximal theorem is essentially sharp.
\end{problem}
\begin{solution}
We are given that $f$ is not $0$ a.e. and thus, $ \int_{\bR}^{}|f|\,\mathrm{d}\mu>0$ and using DCT, we see that there exists some $r>0$ such that $K=\int_{B(0,r)}^{}|f|\,\mathrm{d}\mu>0$. 
Then, for all $x$, as $B(0,r)\subseteq B(x,|x|+r)$ we have
\[
Hf(x)\geq \frac{1}{m(B(x,|x|+r))}\int_{B(x,|x|+r)}|f|\,d\mu\geq \frac{1}{c\cdot(|x|+r)^n}\int_{B(0,r)}|f|\,d\mu\geq \frac{K}{(|x|+r)^n}
\]
where $c=m(B(0,1))$. 
Thus, we can find some $C>0$ and $R>0$ like in the statement. 
Hence, for small $\alpha>0$ we have $\qty{x:Hf(x)>\alpha}\supseteq \qty{x:C\alpha^{-1}\geq|x|^n>R^n}$. 
From this we can deduce that, $m( \qty{x:Hf(x)>\alpha})\geq c\cdot(C\alpha^{-1}-R^n)$ which is atleast $C'\alpha^{-1}$ for some $C'>0$ and small $\alpha.$
\end{solution}
\begin{problem}\label{prob:3.23}
A useful variant of the Hardy-Littlewood maximal function is $H^{*}f(x)=\sup\{\frac{1}{m(B)}\int_{B}|f(y)|dy:B \text{ is a ball and } x\in B\}$. Show that $Hf\le H^{*}f\le2^{n}Hf$.
\end{problem}
\begin{solution}
If $x\in B(p,r)=B$ then we must have that $B(p,r)\subseteq B(x,|x-p|+r)$, thus
\[
	\frac{1}{m(B)}\int_B|f|\,dm\leq \frac{c\cdot(|x-p|+r)^n}{c\cdot r^n}Hf(x)\leq 2^nHf(x)
\]
where $c=m(B(0,1))$, as $|x-p|<r$ so we must have $H^*f\leq 2^nHf$ and $H^*f\geq Hf$ is trivially true. 
\end{solution}
\begin{problem}\label{prob:3.24}
If $f\in L_{loc}^{1}$ and $f$ is continuous at $x$, then $x$ is in the Lebesgue set of $f$.
\end{problem}
\begin{solution}
	For all $\epsilon>0$ we can find $\delta>0$ such that $f(B(x,\delta)) \subseteq B(f(x),\epsilon)$, thus for all $0<\delta'<\delta$ we must have $m(B(x,\delta'))^{-1}\int_{B(x,\delta')}|f(x)-f(y)|dy\leq\epsilon$ so it must converge to $0$ implying $x\in L_f.$
\end{solution}
\begin{problem}\label{prob:3.25}
If $E$ is a Borel set in $\bR^{n}$, the density $D_{E}(x)$ of $E$ at $x$ is defined as $D_{E}(x)=\lim_{r\rightarrow0}\frac{m(E\cap B(r,x))}{m(B(r,x))}$ whenever the limit exists.
\begin{enumerate}[label=\alph*.]
    \item Show that $D_{E}(x)=1$ for a.e. $x\in E$ and $D_{E}(x)=0$ for a.e. $x\in E^{c}$.
    \item Find examples of $E$ and $x$ such that $D_{E}(x)$ is a given number $\alpha\in(0,1)$, or such that $D_{E}(x)$ does not exist.
\end{enumerate}
\end{problem}
\begin{solution}
\begin{enumerate}[label=\alph*.]
\item Consider the borel measure $\nu(F):=\mu(E\cap F)$, then we have a lebesgue decomposition $\nu=0+\chi_Edm$ and Theorem $3.22.$ we have that $D_E(x)=\chi_E(x)$ for $m$-a.e. $x.$ from which the conclusion follows.
\item Consider $E=\bigcup_{k=1}^{\infty}(B(0,2^{-2k+1})\setminus B(0,2^{-2k}))$ and $x=0$. Let $c = m(B(0,1))$. Then we see that for $m\in\bN$ we have 
    \[ m(E\cap B(0,2^{-2m+1})) = c\sum_{k=m}^{\infty}\left(2^{n(-2k+1)}-2^{-2kn}\right) = c(2^n-1)\frac{2^{-2mn}}{1-2^{-2n}}. \]
    Dividing by $m(B(0, 2^{-2m+1})) = c 2^{n(-2m+1)}$, the density along this sequence of radii is:
    \[ \frac{m(E\cap B(0,2^{-2m+1}))}{m(B(0,2^{-2m+1}))} = \frac{c(2^n-1)2^{-2mn}/(1-2^{-2n})}{c 2^{-2mn} 2^n} = \frac{2^n}{2^n+1}. \]
    Now consider the sequence of radii $r = 2^{-2m}$. The ball $B(0, 2^{-2m})$ intersects $E$ starting from index $k=m+1$, so:
    \[ m(E\cap B(0,2^{-2m})) = c\sum_{k=m+1}^{\infty}\left(2^{n(-2k+1)}-2^{-2kn}\right) = c(2^n-1)\frac{2^{-2n(m+1)}}{1-2^{-2n}}. \]
    Dividing by $m(B(0, 2^{-2m})) = c 2^{-2mn}$ yields:
    \[ \frac{m(E\cap B(0,2^{-2m}))}{m(B(0,2^{-2m}))} = \frac{c(2^n-1)2^{-2mn}2^{-2n}/(1-2^{-2n})}{c 2^{-2mn}} = \frac{1}{2^n+1}. \]
    Since $\frac{2^n}{2^n+1}\neq\frac{1}{2^n+1}$ for all $n \geq 1$, two subsequences tend to two different limits, hence $D_E(0)$ does not exist. 
    To find an example where $D_E(x) = \alpha \in (0,1)$, choose $x=0$. 
	If $n \geq 2$, define $E$ as a sector of $\bR^n$ originating at $0$. Let $A \subset S^{n-1}$ be a measurable set on the unit sphere such that its surface measure is $\alpha \sigma(S^{n-1})$ which we can do by considering the function $f:[-1,1]\to\bR$ defined by $f(t):=\sigma( \qty{x\in S^{n-1}:x_1\leq t})$ which can be showed to be continuous via lower and upper continuity, after which we can use the intermediate value theorem. Define $E = \{ y \in \bR^n \setminus \{0\} : y/|y| \in A \}$. For any $r > 0$, $m(E \cap B(0,r)) = \alpha m(B(0,r))$, meaning $D_E(0) = \alpha$ identically for all $r>0$.
    If $n = 1$, define $E = \bigcup_{k=1}^\infty \left( \left[ \frac{1}{k} - \frac{\alpha}{k(k+1)}, \frac{1}{k} \right] \cup \left[ -\frac{1}{k}, -\frac{1}{k} + \frac{\alpha}{k(k+1)} \right] \right)$. For $r = 1/k$, $m(E \cap (-r, r)) = 2 \sum_{j=k}^\infty \frac{\alpha}{j(j+1)} = \frac{2\alpha}{k} = \alpha (2r)$. Bounding the measure for $r \in (\frac{1}{k+1}, \frac{1}{k})$ shows the ratio to $\alpha$, hence $D_E(0) = \alpha$. 
\end{enumerate}
\end{solution}
\begin{problem}\label{prob:3.26}
If $\lambda$ and $\mu$ are positive, mutually singular Borel measures on $\bR^{n}$ and $\lambda+\mu$ is regular, then so are $\lambda$ and $\mu$.
\end{problem}
\begin{solution}
Define $\nu=\lambda+\mu$ and let $\lambda(A)=\mu(A^c)=0$. 
For compact $K$ we see that $\lambda(K)+\mu(K)=\nu(K)<\infty$, which implies that $\lambda(K),\mu(K)<\infty$. 
First, we show outer regularity. 
If $\lambda(E) = \infty$, the condition $\lambda(E) = \inf\{\lambda(U) : U \supseteq E, U \text{ open}\}$ is trivially satisfied by any open $U \supseteq E$. 
Assume $\lambda(E) < \infty$. 
We wish to find an open $U \supseteq E$ with $\lambda(U \setminus E) < \epsilon$. 
We split $E$ into $E \cap A^c$ and $E \cap A$. 
Notice that $\nu(E \cap A^c) = \lambda(E \cap A^c) + \mu(E \cap A^c) = \lambda(E \cap A^c) + 0 \le \lambda(E) < \infty$. 
Since $\nu(E \cap A^c) < \infty$ and $\nu$ is regular, there exists an open $U_1 \supseteq (E \cap A^c)$ with $\nu(U_1 \setminus (E \cap A^c)) < \epsilon/2$, which implies $\lambda(U_1 \setminus (E \cap A^c))\leq\nu(U_1\setminus(E\cap A^c))< \epsilon/2$. 
For the remaining piece $E \cap A$, we have $\lambda(E \cap A) \le \lambda(A) = 0$. 
We can write $A = \bigcup_{k=1}^\infty A_k$ where $A_k = A \cap B(0,k)$. 
As each $A_k$ is bounded, $\nu(A_k) < \infty$, so there exist open sets $V_k \supseteq A_k$ with $\nu(V_k \setminus A_k) < \epsilon/2^{k+1}$. 
This implies $\lambda(V_k) = \lambda(V_k \setminus A_k) + \lambda(A_k) < \epsilon/2^{k+1} + 0$. 
Taking $U_2 = \bigcup_{k=1}^\infty V_k$ gives an open set containing $A$ (and thus $E \cap A$) with $\lambda(U_2) \le \sum \lambda(V_k) < \epsilon/2$. 
Let $U = U_1 \cup U_2$. 
Then $U$ is open, $U \supseteq E$, and $U \setminus E \subseteq (U_1 \setminus (E \cap A^c)) \cup U_2$. 
Thus, $\lambda(U \setminus E) \le \lambda(U_1 \setminus (E \cap A^c)) + \lambda(U_2) < \epsilon/2 + \epsilon/2 = \epsilon$ so we have the second condition. 
The first condition is trivial.
\end{solution}
\begin{problem}\label{prob:3.27}
Verify the assertions in Examples 3.25.
\end{problem}
\begin{solution}
This is trivial.
\end{solution}
\begin{problem}\label{prob:3.28}
If $F\in NBV$, let $G(x)=|\mu_{F}|((-\infty,x])$. Prove that  $|\mu_{F}|=\mu_{T_{F}}$ by showing that $G=T_{F}$ via the following steps.
\begin{enumerate}[label=\alph*.]
	\item From the definition of $T_{F}$, $T_{F}\le G$.
	\item $|\mu_{F}(E)|\le\mu_{T_{F}}(E)$ when $E$ is an interval, and hence when $E$ is a Borel set.
    \item $|\mu_{F}|\le\mu_{T_{F}}$, and hence $G\le T_{F}$.
\end{enumerate}
\end{problem}
\begin{solution}
\begin{enumerate}[label=\alph*.]
    \item For any $x\in\bR$ and $x_0<\ldots<x_n=x$ we see that,
        \[
            \sum_{k=1}^n|F(x_k)-F(x_{k-1})|=\sum_{k=1}^n|\mu_F((x_{k-1},x_k])|\leq\sum_{k=1}^n|\mu_F|((x_{k-1},x_k])=|\mu_F|((x_0,x])
        \]
       Which is clearly less than $G(x).$ 
       This shows that $T_F\leq G$ as $T_F$ is the supremum over all sums of this form.
   \item For open (we only need these) intervals, the inequality follows from the fact that $T_F(x)-T_F(y)\geq|F(x)-F(y)|$ for $x\geq y$ as we can take limits and say that, for $a<b$
       \[
           \mu_{T_F}((a,b))=T_F(b-)-T_F(a)\geq|F(b-)-F(a)|=|\mu_F((a,b))|
       \]
       (keep in mind that all the one sided limits exist in this case.) 
       We know by construction that $\mu_F=\mu_1-\mu_2+i(\mu_3-\mu_4)$ for finite lebesgue stjeltes measures $\mu_i$ as in chapter $1.$ 
       Let $\mu_5=\mu_{T_F},$ another lebesgue stjeltes measure. 
       Given borel $E$, by  regularity we can find open sets $U_n^{(i)}\supseteq E$ such that $\mu_i(U_n^{(i)})\to\mu_i(E)$ for all $i$. 
       Notice that by the \twemoji{sandwich} theorem and monotonicty of $\mu_i$, the sequence of open sets $V_n:=\cap_i U_n^{(i)}$ have the same convergence property by monotonicity of measures and the sandwich property. 
       As $\mu_i$ are all finite we see that $\mu_F(V_n)\to\mu_F(E).$ 
       If $\cup I_n$ is a countable disjoint union of open intervals we see that $|\mu_F(\cup I_n)|\leq\mu_{T_F}(\cup I_n)$ by the triangle inequality and sigma additivity of measures.
       From this, we can conclude that $|\mu_F(V_n)|\leq\mu_{T_F}(V_n)$ for all $n$ and we have the desired inequality by taking the limit as $n\to\infty$
   \item This can be proven by  the previous part and \hyperref[prob:3.21]{Problem 3.21} directly (look at $\mu_1$) in the hyperlinked problem. 
       From this we thus have that $G\leq T_F$ by considering left open rays.
\end{enumerate}
Now that we have $G=T_F$ we must have that $\mu_G=\mu_{T_F}$ and $\mu_G$ is none other than $|\mu_F|.$
\end{solution}
\begin{problem}\label{prob:3.29}
If $F\in NBV$ is real-valued, then $\mu_{F}^{+}=\mu_{P}$ and $\mu_{F}^{-}=\mu_{N}$ where $P$ and $N$ are the positive and negative variations of $F$.
\end{problem}
\begin{solution}
It is easily seen that, the positive and negative variations are nonnegative for $NBV$ functions as $F(-\infty)=T_F(-\infty)=0.$ 
Now, we clearly have $\mu_F=\mu_{P}-\mu_{N}$ (this is true for h-intervals and the sets where this holds forms a sigma algebra due to these being finite measures) which proves that $\mu_P\geq\mu_F^+$ and $\mu_N\geq\mu_F^-$ by \hyperref[prob:3.4]{Problem 3.4}. 
Now, from page 103 we can also see that
\[
    P(x)=\sup\qty{\sum_{k=1}^n[F(x_k)-F(x_{k-1})]^+:x_0<\ldots<x_n=x} + 0
\]
as $F\in NBV$, now notice that for $a<b$, we have 
\[
    [F(b)-F(a)]^+=[\mu_F((a,b])]^+=[\mu_F^+((a,b])-\mu_F^-((a,b])]^+\leq[\mu_F^+((a,b])]^+=\mu_F^+((a,b])
\]
From this we can easily see that $\mu_P((-\infty,x])=P(x)\leq\mu_F^+((-\infty,x])$. 
But earlier we proved that $\mu_P\geq\mu_F^+$ so we must have equality above and hence $\mu_P=\mu_F^+$ can be easily shown to hold. 
\end{solution}
\begin{problem}\label{prob:3.30}
Construct an increasing function on $\bR$ whose set of discontinuities is $\bQ$.
\end{problem}
\begin{solution}
    Consider an ennumeration $\{q_n\}$ of the rationals. 
    Define $f(x):=\sum_{q_n<x}2^{-n}$ then it is obvious that for rationals $r$ we have $f(q_n+)-f(q_n-)\geq 2^{-n}$ making it discontinuous on the rationals. 
    Now, if $x$ is an irrational, let $\delta>0$ be such that $q_1,\ldots,q_m\not\in(x-\delta,x+\delta)$, we can do this becasue $x$ is irrational. 
    In this case, we see that for $y\in(x-\delta,x+\delta)$
    \[
        |f(x)-f(y)|=\sum_{q_n\in[x-\delta,x+\delta)}\frac1{2^n}\leq\sum_{n>m}\frac1{2^n}=\frac1{2^{m+1}}
    \]
    Thus we have continuity at the irrationals as this can be made arbitrarily small for larger $m$ and small enough $\delta>0.$
\end{solution}
\begin{problem}\label{prob:3.31}
Let $F(x)=x^{2}\sin(x^{-1})$ and $G(x)=x^{2}\sin(x^{-2})$ for $x\ne0,$ and $F(0)=G(0)=0$.
\begin{enumerate}[label=\alph*.]
    \item $F$ and $G$ are differentiable everywhere (including $x=0$).
    \item $F\in BV([-1,1])$, but $G\notin BV([-1,1])$.
\end{enumerate}
\end{problem}
\begin{solution}
    The first part is an exercise in calculus and the next follows from $F'$ being bounded ( $|F'|\leq3$ to be a bit more specific) and $T_G(1)-T_G(n^{-1/2})\gg\log n$ (this is not to be interpreted as $\infty-\infty$ but rather the total variation in the smaller interval $(n^{-1/2},1) )$
\end{solution}
\begin{problem}\label{prob:3.32}
If $F_{1},F_{2},...,F\in NBV$ and $F_{j}\rightarrow F$ pointwise, then $T_{F}\le \liminf T_{F_{j}}$.
\end{problem}
\begin{solution}
For any $x\in\bR$ and $\epsilon>0$ consider some $x_0<\ldots<x_n=x$ such that,
\[
    \sum_{k=1}^n|F(x_k)-F(x_{k-1})|\geq T_F(x)-\epsilon
\]
Then, by pointwise convergence we see that for large enough $j$,
\[
    -\epsilon+\sum_{k=1}^n|F(x_k)-F(x_{k-1})|\leq\sum_{k=1}^n|F_j(x_k)-F_j(x_{k-1})|\leq T_{F_j}(x)
\]
This shows that $T_{F}(x)-2\epsilon\leq T_{F_j}(x)$ for large enough $j$ and thus, $T_F(x)-2\epsilon\leq\liminf T_{F_j}(x)$. 
As $\epsilon>0$ was arbitrary we are done.
\end{solution}
\begin{problem}\label{prob:3.33}
If $F$ is increasing on $\bR$, then $F(b)-F(a)\ge\int_{a}^{b}F^{\prime}(t)dt$.
\end{problem}
\begin{solution}
As we only care about things being defined a.e. we define $F_n(x):=n(F(x+1/n)-F(x))\geq 0$ and we see that, $F_n\to F'$ a.e. thus, by fatou's lemma we have
\[
    \int_a^bF'\leq\liminf\int_a^b F_n=\liminf_nn\int_a^bF(x+1/n)-F(x)\,\mathrm{d}m(x)
\]
If we set $G(x):=\int_a^xF(t)dm(t)$ then we can pass to rieman integrablity (monotone functions are rieman integrable) and see that the last expression above is
\[
    \liminf_n\frac{G(b+1/n)-G(a+1/n)-(G(b)-G(a))}{1/n}=G'(b)-G'(a)=F(b)-F(a)
\]
As $G'(t)=F(t)$ by the fundamental theorem of calculus, when $F$ is continuous at $a,b.$ 
Now, as $F$ is continuous a.e. the set of points where it is continuous must be dense and we can thus find sequences $a_n\searrow a<b\nwarrow b_n$ where $F$ is continuous and thus, 
\[
    F(b)-F(a)\geq F(b-)-F(a+)=\lim_{n\to\infty} F(b_n)-F(a_n)\geq\lim\int_{[a_n,b_n]}F'=\int_{(a,b)}F'=\int_a^bF'
\]
where the last two equalities are due to the fact that $A\mapsto\int_AF'$ is a measure here ($(a_n,b_n)$ increases to $(a,b)$) and $m(\{a,b\})=0$.
\end{solution}
\begin{problem}\label{prob:3.34}
Suppose $F, G\in NBV$ and $-\infty<a<b<\infty.$
\begin{enumerate}[label=\alph*.]
    \item By adapting the proof of Theorem 3.36, show that
    \[\int_{[a,b]}\frac{F(x)+F(x-)}{2}dG(x)+\int_{[a,b]}\frac{G(x)+G(x-)}{2}dF(x)=F(b)G(b)-F(a-)G(a-)\]
    \item If there are no points in $[a, b]$ where $F$ and $G$ are both discontinuous, then
    \[\int_{[a,b]}F~dG+\int_{[a,b]}G~dF=F(b)G(b)-F(a-)G(a-).\]
\end{enumerate}
\end{problem}
\begin{solution}
\begin{enumerate}[label=\alph*.]
    \item Following steps similar to those the proof of theorem 3.36, taking $\Omega=\{(x,y):a\leq x\leq y\leq b\}$ we can see that,
        \[
            \int_{[a,b]}FdG-F(a-)[G(b)-G(a-)]=G(b)[F(b)-F(a-)]-\int_{[a,b]}G(x-)dF(x)
        \]
        We can also switch $F,G$ and get another equation
        \[
            \int_{[a,b]}GdF-G(a-)[F(b)-F(a-)]=F(b)[G(b)-G(a-)]-\int_{[a,b]}F(x-)dG(x)
        \]
        rearranging and adding these we get,
        \[
            \int_{[a,b]}(F(x)+F(x-))dG(x)+\int_{[a,b]}(G(x)+G(x-))dF(x)=2(F(b)G(b)-F(a-)G(a-))
        \]
    \item If $F$ is continuous at $x$ then we see that 
        \[
            \mu_F(\{x\})=\lim_n\mu_F((x+1/n,x])=\lim_n F(x)-F(x+1/n)=0
        \]
        and similarly for $G$. 
        As $F,G$ are NBV we can write them as the differnece of two increasing functions which only ever have countably many discontinuities and thus the set of discontinuities here form countable sets which we call $D_F,D_G$ for $F$ and $G$ respectively. 
        Then we can see that $\mu_F(D_G)=\sum_{x\in D_G}\mu_F(\{x\})=\sum_{x\in D_G}0=0$ because $F$ is continuous on $D_G$ and similarly $\mu_G(D_F)=0$. 
        Now, in the equation from the previous part we can write the LHS as, 
        \[
            \int_{[a,b]\setminus D_F}(F(x)+F(x-))dF(x)+\int_{[a,b]\setminus D_G}(G(x)+G(x-))dF(x)
        \]
        as the integrands are continuous over their domain of integration, we see that $F(x)+F(x-)=2F(x)$ and similarly for $G.$ 
        From this we can conclude that, 
        \[
            \int_{[a,b]}2FdG+\int_{[a,b]}2GdF=2(F(b)G(b)-F(a-)G(a-))
        \]
        from which we get the desired equality after cancelling the $2$ out.
\end{enumerate}
\end{solution}
\begin{problem}\label{prob:3.35}
If $F$ and $G$ are absolutely continuous on $[a, b]$, then so is $FG$, and
\[\int_{a}^{b}(FG^{\prime}+GF^{\prime})(x)dx=F(b)G(b)-F(a)G(a).\]
\end{problem}
\begin{solution}
It suffices to show that $F$ being absolutely continuous implies $F^2$ is absolutely continuous. 
Indeed, if $F$ is absolutely continuous then it is continuous too and thus is bounded on the compact set $[a,b]$, say $|F|\leq M$. 
If, $\epsilon>0$ is arbitrary, we can find $\delta>0$ with,
\[
    \sum_{i=1}^n(b_i-a_i)<\delta\implies\sum_{i=1}^n|F(b_i)-F(a_i)|<\frac{\epsilon}{M+1}\text{ for }a\leq a_1<b_1<\ldots<b_n\leq b
\]
Then, 
\[
    \sum_{i=1}^n|F(b_i)^2-F(a_i)^2|\leq M\sum_{i=1}^n|F(b_i)-F(a_i)|\leq\frac{\epsilon M}{M+1}<\epsilon
\]
so $F^2$ is also absolutely continuous. 
Now, it is clear that $F,G$ being absolutely continuous makes their product absolutely continuous too. 
From Theorem 3.35c we see that $(FG)'\in L^1([a,b],m)$ and the same also holds for $F',G'$. 
As they all exist a.e., they exist together a.e. as well and we can say that $(FG)'=FG'+F'G$ a.e. and again by the cited theorem,
\[
    F(b)G(b)-F(a)G(a)=\int_a^b(FG)'(x)dx=\int_a^b(FG'+F'G)(x)dx
\]
\end{solution}
\begin{problem}\label{prob:3.36}
Let $G$ be a continuous increasing (we need $G$ to be strictly increasing, otherwise consider constant functions.) function on $[a, b]$ and let $G(a)=c$, $G(b)=d.$
\begin{enumerate}[label=\alph*.]
    \item If $E\subset[c,d]$ is a Borel set, then $m(E)=\mu_{G}(G^{-1}(E)).$ (First consider the case where $E$ is an interval.)
    \item If $f$ is a Borel measurable and integrable function on $[c, d]$, then $\int_{c}^{d}f(y)dy=\int_{a}^{b}f(G(x))dG(x)$. In particular, $\int_{c}^{d}f(y)dy=\int_{a}^{b}f(G(x))G^{\prime}(x)dx$ if $G$ is absolutely continuous.
    \item The validity of (b) may fail if $G$ is merely right continuous rather than continuous.
\end{enumerate}
\end{problem}
\begin{solution}
\begin{enumerate}[label=\alph*.]
    \item The case where this is an interval follows as shown, for $c\leq x<y\leq d$ we just have $\mu_G(G^{-1}([x,y]))=\mu_G([G^{-1}(x),G^{-1}(y)])=G\circ G^{-1}(y)-G\circ G^{-1}(x)=y-x=m([x,y])$. 
        We can easily show that the collection of sets $E$ for which this works is a sigma algebra as the inverse commutes with the complement and unions, we thus have that the borel sigma algebra over $[c,d]$ is contained in this one which suffices.
    \item If $f=\chi_E$ ($E\in\mathcal B_{[c,d]}$) then, $\int_a^b\chi_E(G(x))dG(x)=\int_a^b\chi_{G^{-1}(E)}dG(x)=\mu_G(G^{-1}(E))=m(E)=\int_c^d\chi_Edm$ and we can conclude by linearity of the integral and MCT as usual. 
       If $G$ is NBV, we have that $dG=G'dm$ (can be proven from 3.35c) and by \hyperref[prob:2.14]{Problem 2.14} we are done (split the integrable function into its two parts as usual.)
   \item In that case the range of $f$ needn't be an interval in the first place. 
       For example, take $f(x):=x\chi_{[0,1/2)}(x)+(x+1)\chi_{[1/2,1]}$ which is right continuous, increasing and maps $[0,1]$ to $[0,1/2)\cup[3/2,2].$
\end{enumerate}
\end{solution}
\begin{problem}\label{prob:3.37}
Suppose $F:\bR\rightarrow\bC$. There is a constant $M$ such that $|F(x)-F(y)|\le M|x-y|$ for all $x,y\in\bR$ (that is, $F$ is Lipschitz continuous) iff $F$ is absolutely continuous and $|F^{\prime}|\le M$ a.e.
\end{problem}
\begin{solution}
    For the only if direction, absolute continuity follows directly from the definition by considering $\epsilon/{M+1}$ instead of $\epsilon$ and for the derivative to exist a.e. (it is trivially bounded as needed when it does exist) we consider intervals $(n,n+1]$ for $n\in\bZ.$ 
    Clearly $F$ is still absolutely continuous on this small interval and by Theorem 3.35c we again see that $F'$ exists a.e. in $(n,n+1]$ and thus it exists a.e. in $\bR$ too as the union of countably many disjoint null sets is still null.
    To prove the if part, for $x<y$ consider $F$ on the interval $[x,y].$ 
    By Theorem 3.35c we can write $F(x)-F(y)=\int_x^yF'(t)dt$ again and by hypothesis we have,
    \[
        |F(x)-F(y)|=\left|\int_x^yF'(t)dt\right|\leq\int_x^y|F'(t)|dt\leq M(y-x)=M|x-y|
    \]
    The case of $x>y$ is no different.
\end{solution}
\begin{problem}\label{prob:3.38}
If $f:[a,b]\rightarrow\bR$, consider the graph of $f$ as a subset of $\bC$, namely, $\{t+if(t): t\in[a,b]\}$. The length $L$ of this graph is by definition the supremum of the lengths of all inscribed polygons.
\begin{enumerate}[label=\alph*.]
    \item Let $F(t)=t+if(t);$ then $L$ is the total variation of $F$ on $[a, b]$.
    \item If $f$ is absolutely continuous, $L=\int_{a}^{b}[1+f^{\prime}(t)^{2}]^{1/2}dt$.
\end{enumerate}
\end{problem}
\begin{solution}
\begin{enumerate}[label=\alph*.]
    \item This is trivial, just remember how $|x+iy|=\sqrt{x^2+y^2}.$
    \item If $f$ is absolutely continuous then so is $F$ and by problem 3.35c it can be shown that $\mu_F= F'dm$ on this interval and thus, $\mu_{T_F}=|F'|dm$ by \hyperref[prob:3.28]{Problem 3.28}, this shows that
        \[
            L=T_F(b)-T_F(a)=\mu_{T_F}((a,b])=\int_{(a,b]}|F'(t)|dm(t)=\int_a^b[1+f'(t)^2]^{1/2}dm(t)
        \]
        as $F'(t)=1+i\cdot f'(t)$ (when it exists.)
\end{enumerate}
\end{solution}
\begin{problem}\label{prob:3.39}
    If $\{F_{j}\}$ is a sequence of nonnegative increasing functions on $[a, b]$ such that $F(x)=\sum_{1}^{\infty}F_{j}(x)<\infty$ for all $x\in[a,b]$, then $F^{\prime}(x)=\sum_{1}^{\infty}F_{j}^{\prime}(x)$ for a.e. $x\in[a,b]$. (It suffices to assume $F_j\in NBV$. onsider the meausres $\mu_{F_j}.)$
\end{problem}
\begin{solution}
    Keep in mind the following lemma : If $P_1,P_2,\ldots$ hold a.e. then $\land P_i$ does too, the proof is left as an exercsie.
    We can consider the functions $G_j(x)=F_j(x+)-F_j(a)$ and $G(x)=F(x+)-F(a)$ instead. These have the same properties as $F_j$ and $F$, and we extend them to make them $NBV$ as done before (they were trivially in $BV([a,b])$ already). Note that $G=\sum G_j$ everywhere, which follows from an application of the monotone convergence theorem with respect to the counting measure. 
    Consider the measures $\mu_G$ and $\mu_{G_j}$, for which it holds that $\mu_G=\sum_j\mu_{G_j}$ since they agree on right-closed rays. Now, consider the Lebesgue-Radon-Nikodym decompositions $dG = g \, dm + d\lambda$ and $dG_j = g_j \, dm + d\lambda_j$. 
    We have $\lambda_j \perp m$, and thus $\sum \lambda_j \perp m$ too by \hyperref[prob:3.9]{Problem 3.9} (notice that $\lambda_j,\lambda$ are positive measures as well.)
    Since $\mu_G = \sum \mu_{G_j}$, we must have $dG = (\sum_j g_j) \, dm + \sum \lambda_j$. By the uniqueness criterion in the Lebesgue-Radon-Nikodym theorem, this implies $g = \sum g_j$ a.e. and $\lambda = \sum \lambda_j$. 
    It is clear from the lebesgue differentiation theorem that $g_j = G_j'$ a.e. and $g = G'$ a.e., so the lemma at the beginning implies that $G' = \sum G_j'$ a.e.  
    If we now show that $F' = G'$ and $F_j' = G_j'$ a.e., then we would be done using the same lemmat. Suppose $p$ and $q$ are such that $p = q$ a.e. and both are differentiable a.e. We want to show that $p' = q'$ a.e. We can find a set $A \subseteq [a,b]$ where $p = q$ and both are differentiable, such that $m([a,b] \setminus A) = 0$ by the lemma. 
    Notice that $A$ must be dense in this interval, so for any $x \in A$, we can find a sequence $x_n \in A \setminus \{x\}$ that tends to $x$. By the sequential criterion for limits, we have
    \[
        p'(x) = \lim_{n \to \infty} \frac{p(x_n) - p(x)}{x_n - x} = \lim_{n \to \infty} \frac{q(x_n) - q(x)}{x_n - x} = q'(x)
    \]
    for all $x \in A$, so we are done.
\end{solution}
\begin{problem}\label{prob:3.40}
Let $F$ denote the Cantor function on $[0, 1]$ (see §1.5), and set $F(x)=0$ for $x<0$ and $F(x)=1$ for $x>1$. Let $\{[a_{n},b_{n}]\}$ be an enumeration of the closed subintervals of $[0, 1]$ with rational endpoints, and let $F_{n}(x)=F((x-a_{n})/(b_{n}-a_{n}))$. Then $G=\sum_{1}^{\infty}2^{-n}F_{n}$ is continuous and strictly increasing on $[0, 1]$, and $G^{\prime}=0$ a.e. (Use exercise 39)
\end{problem}
\begin{solution}
$F_n$ are cleary increasing, continuous and nonnegative and we also see that $G=\sum F_n\leq1<\infty$ as $F_n\leq 1$ for all $n$ by definition. 
$G$ is also continuous by weierstrass' M-test being the limit of an uniformly converging sequence of continuous functions. 
Given any $0\leq x<y\leq 1$ we can find some interval with rational endpoints $[a_n,b_n]\subset(x,y)$ and we see that $G(y)-G(x)\geq 2^{-n}(F_n(y)-F_n(x))=2^{-n}(1-0)>0$ by definition again, making $G$ strictly increasing. 
We see that $F$ is constant on intervals outside the cantor set and thus has derivative $0$ here and as the cantor set has lebesgue measure $0$ we can show that $F'=0$ a.e. by standard procedures. 
Now, by \hyperref[prob:3.39]{Problem 3.39} we see that $G'(x)=\sum 2^{-n}F_n'(x)$ a.e. and as $F_n'=0$ a.e. too (can be easily shown), we must have $G'=0$ a.e.
\end{solution}
\begin{problem}\label{prob:3.41}
Let $A\subset[0,1]$ be a Borel set such that $0<m(A\cap I)<m(I)$ for every subinterval $I$ of $[0, 1]$.
\begin{enumerate}[label=\alph*.]
    \item Let $F(x)=m([0,x]\cap A)$. Then $F$ is absolutely continuous and strictly increasing on $[0, 1]$, but $F^{\prime}=0$ on a set of positive measure.
    \item Let $G(x)=m([0,x]\cap A)-m([0,x]\setminus A).$ Then $G$ is absolutely continuous on $[0, 1]$, but $G$ is not monotone on any subinterval of $[0, 1]$.
\end{enumerate}
\end{problem}
\begin{solution}
\begin{enumerate}[label=\alph*.]
    \item Clearly $F$ is lipschitz continuous with constant $1$ and thus its absolutely continuous by \hyperref[prob:3.37]{Problem 3.37} and its strictly increasing as $F(y)-F(x)=m((x,y]\cap A)>0$ for any $x<y$ by definition. 
        As $F$ is absolutely continuous, by theorem 3.35c we see that $F'\geq 0$ exists a.e, is integrable and $F(x)=F(x)-F(0)=\int_0^xF'(t)dm(t)$. 
       $F$ can be extended to $\bR$ by setting it equal to $0=F(0)$ for $x<0$ and $m(A)=F(1)$ for $x>1$ and this is NBV so we can talk about $\mu_F.$ 
       We also see that $\mu_F((-\infty,x]))=(F'dm)((-\infty,x])$ for all $x$ and thus, $\mu_F=F'dm$. 
       If $E$ is a set of positive meausre and we have $\mu_F(E)=\int_EF'dm=0$ then it must be true that $F'=0$ a.e. in $E$ which again furnishes a set of positve meausre (same as $E$) on which $F'=0.$ 
       Notice that the measure defined by $\nu(E)=m(E\cap A)$ is the same as $\mu_F$ as they agree on right closed rays and hence we just need to find some $E$ such that $m(E)>0$ whereas $\nu(E)=0$ and the perfect (not a topological property!) candidate is $A^c$ for which $m(A^c)=1-m(A)>0$ whereas $\nu(A^c)=m(A^c\cap A)=0$ so we are done.
   \item We can easily see that $m(A^c\cap I)=m(I)-m(A\cap I)\in(0,m(I))$ as well, if we let $H(x)=m([0,x]\cap A^c)$ then we see that $H$ has the same properties as $F$ from the previous part. 
       Thus, $G$ is also absolutely continuous being the linear combination of two absolutely continuous functions. 
       Given any interval $[x,y]\subset[0,1]$ ($x<y$) suppose for the sake of contradiction that WLOG. $G$ is increasing in this interval, this implies that $G'\geq 0$ a.e. here but this can not be true as it implies $\mu_G=\mu_F-\mu_H$ has $[x,y]$ as a positive set. 
      But we clearly see that for $V=A^c\cap [x,y]\subseteq[x,y]$ we have $\mu_F(V)=m(V\cap A)=0$ whereas $\mu_H(V)=m(V\cap A^c)>0$ and thus, $\mu_G(V)<0$ which is a contradiction and thus we are done.
\end{enumerate}
\end{solution}
\begin{problem}\label{prob:3.42}
A function $F:(a,b)\rightarrow\bR(-\infty\le a<b\le\infty)$ is called convex if $F(\lambda s+(1-\lambda)t)\le\lambda F(s)+(1-\lambda)F(t)$ for all $s, t\in(a,b)$ and $\lambda\in(0,1)$.
\begin{enumerate}[label=\alph*.]
    \item $F$ is convex iff for all $s, t, s', t' \in (a, b)$ such that $s\le s^{\prime}<t^{\prime}$ and $s<t\le t^{\prime}$,
    \[\frac{F(t)-F(s)}{t-s}\le\frac{F(t^{\prime})-F(s^{\prime})}{t^{\prime}-s^{\prime}}\]
    \item $F$ is convex iff $F$ is absolutely continuous on every compact subinterval of $(a, b)$ and $F^{\prime}$ is increasing (on the set where it is defined).
    \item If $F$ is convex and $t_{0}\in(a,b)$, there exists $\beta\in\bR$ such that $F(t)-F(t_{0})\ge \beta(t-t_{0})$ for all $t\in(a,b)$.
    \item (Jensen's Inequality) If $(X, \mM, \mu)$ is a measure space with $\mu(X)=1$, $g:X\rightarrow(a,b)$ is in $L^{1}(\mu)$, and $F$ is convex on $(a, b)$, then
    \[F\left(\int g~d\mu\right)\le\int F\circ g~d\mu.\]
    (let $t_0=\int gd\mu$ and $t=g(x)$ in (c), and integrate.)
\end{enumerate}
\end{problem}
\begin{solution}
\begin{enumerate}[label=\alph*.]
    \item This is a standard result from introductory real analysis, it can be proved by just manipulating the inequalities.
    \item If $F$ is convex then we see that on any compact subinterval $[c,d]$, for $x,y\in[c,d]$ if $F'$ exists on $x,y$ ($x<y$) then, for $h>0$ small enough, by part (a)
        \[
            \frac{F(x+h)-F(x)}{(x+h)-x}\leq\frac{F(v)-F(u)}{v-u}\leq\frac{F(y+h)-F(y)}{(y+h)-y}
        \]
        where $x<x+h<y<y+h$ and as $h\to 0$ from above, as the derivatives exist at $x,y$ we get that $F'(x)\geq F'(y).$ 
        It is another standard result that convex functions are continuous and we can also see that if, $c=s'<t'$ and $c<t\leq t'$ then,
        \[
            \frac{F(t)-F(c)}{t-c}\leq \frac{F(t')-F(c)}{t'-c}
        \]
        Thus, $G(x):=(F(x)-F(c))/(x-c)$ is increasing, similarly we can show that $H(x)=(F(d)-F(x))/(d-x)$ is decreasing. 
        This, implies that $G(c+),H(d-)$ both exist and we can show that these bound all slopes using an arguement similar to the one we used to show that $F'$ is increasing on the set its defined. 
        We thus have, for all $x,y$ such that $c\leq x<y\leq d$ we have
        \[
            G(c+)\leq\frac{F(y)-F(x)}{y-x}\leq H(d-)
        \]
        and thus $F$ is lipschitz with constant atmost $\max\{|G(c+)|,|H(d-)|\}$ so it is absolutely continuous by a previous problem. 
        We now prove the if direction. 
        By hypothesis, for any such $s,t,s',t'$ as in part (a), we work in the compact subinterval $[s,t]$ where the inequality in part (a) reduces to the following by theorem 3.35c,
        \[
            \frac1{t-s}\int_s^tF'(t)dt\leq\frac1{t'-s'}\int_{s'}^{t'}F'(t)dt
        \]
        where $s\leq s'<t'$ and $s<t\leq t'.$ 
        This is trivial if $s'>t$ so we assume otherwise that $s'\leq t$, where this inequality is true as, 
        \[
            \frac1{t-s}\int_s^tF'= \frac{s'-s}{t-s}\cdot\underbrace{\frac1{s'-s}\int_s^{s'}F'}_{\leq\frac1{t-s'}\int_{s'}^tF'}+ \frac{t-s'}{t-s}\cdot\frac1{t-s'}\int_{s'}^tF'\leq\frac1{t-s'}\int_{s'}^tF'
        \]
        The inequality in the underbrace is true as all the values appearing in the integralin the left are atmost those appearing in the right and we average these. 
        It can be similarly shown that the expression on the right is atmost $\frac1{t'-s'}\int_{s'}^{t'}F'$ so we are done. 
    \item Just like we did for $c,d$ in part (b), it can be showed that the left and right derivatives exist at every point and we pick any $\beta\in[F'(t_0-),F'(t_0+)]$ (left and right derivatives and not the one sided limits of our derivative.) 
        The conclusion follows just as we did in part (b), we show part of it. 
        \[
            F'(t_0+)=\lim_{h\searrow 0}\frac{F(t_0+h)-F(t_0)}{h}\leq\frac{F(y)-F(t_0)}{y-t_0}
        \]
        for any $y>t_0$ as $t_0+h<y$ eventually. 
    \item We see that $\int_Xgd\mu\in(a,b)$ so from part (c), we can find some constant $c$ such that,
        \[
            F(t)-F\qty(\int gd\mu)\geq c\cdot\qty(t-\int gd\mu)
        \]
        for all $t\in(a,b)$ and thus we can substitute $t=g(x)$ to get,
        \[
            F\circ g(x)-F\qty(\int gd\mu)\geq c\cdot\qty(g(x)-\int gd\mu)
        \]
        for all $x\in X.$ 
        Now, integratint both sides we see that,
        \[
            \int F\circ g\,d\mu-F\qty(\int g\,d\mu)\geq c\cdot\int gd\mu-\int gd\mu=0
        \]
        so we are done.
\end{enumerate}
\end{solution}
\stepcounter{section}
\section*{Chapter \thesection\ : Point Set Topology}
\addcontentsline{toc}{section}{Chapter \thesection\ : Point Set Topology}
\begin{problem}\label{prob:4-1}
If $\operatorname{card}(X) \ge 2$, there is a topology on $X$ that is $T_0$ but not $T_1$.
\end{problem}
\begin{solution}
    For any distinct $x,y\in X$ consider the topology $\qty{\emptyset,\qty{x},\qty{x,y},X}$, then we see that there is an open set containing $x$ but not $y$ whereas there are none containing $y$ but not $x$ so this is $T_0$ and not $T_1.$
\end{solution}
\begin{problem}\label{prob:4-2}
If $X$ is an infinite set, the cofinite topology on $X$ is $T_1$ but not $T_2$, and is first countable iff $X$ is countable.
\end{problem}
\begin{solution}
If $x,y$ are distinct elements of $X$ then as $X\setminus \qty{x}$ is open (it has a finite complenent $ \qty{x}$) we see that this is $T_1$ but this is not $T_2$ as we can not have two nonempty disjoint open sets to begin with, indeed if $U,V$ are open and disjoint then we must have that $V\subseteq U^c$ which is finite ($U\neq\emptyset$ by assumption) and this in turn makes $V^c$ infinite contradicting it being open in the first place as $V\neq\emptyset$ either. 

If $X$ is first countable then consider a countable base $B_{x_0}$ of $x_0\in X$. 
For $y\neq x$, $X\setminus \qty{y}$ is a open neighborhood of $x$ and thus we can find $U\in B_{x_0}$ such that $x\in U\subset X\setminus \qty{y}$ and thus, $y\in U^c$ with $U^c$ being finite. 
From this we can write 
\[
    X\setminus \qty{x_0} = \bigcup_{U\in B_{x_0}} U^c
\]
which is countable, being a countable union of countable (here, finite) sets so we are done. 
For the other direction, consider an ennumeration $\{x_n\}_{n\in\bN}$, then we can construct a countable neighborhood base of $x_n$ as
\[
    \qty{X\setminus U: U \subseteq X\setminus \qty{x_n},U\text{ is finite}}
\]
infact, this contains exactly the open neighborhoods of $x_0$ as these must be subsets of $X$ which contain $x_0$ and have a finite complement, thus their complement has the form of $U$ above and they are of form $X\setminus U$ as described. 
Its countable because the set of countable subsets of countable sets is countable.  
 \end{solution}
\begin{problem}\label{prob:4-3}
Every metric space is normal. (If $A, B$ are closed sets in the metric space $(X,\rho)$, consider the sets of points where $\rho(x,A) < \rho(x,B)$ or $\rho(x,A) > \rho(x,B)$.)
\end{problem}
\begin{solution}
We first prove that $x\mapsto\rho(x,A)$ is continuous when $A$ is closed. 
For any $a\in A$ we see that $\rho(x,A)\leq\rho(x,a)\leq\rho(x,y)+\rho(y,a)$ and thus, $\inf_{a\in A}\rho(a,y)\geq\rho(x,A)-\rho(x,y)$, we can switch $x,y$ and get another inequality, using both of these we find that $|\rho(x,A)-\rho(y,A)|\leq\rho(x,y)$ and thus this map is indeed continuous. 
Notice that for closed sets $G$ we have that $G=\rho_G^{-1}( \qty{0})$, this can be proven easily.
Let, $\rho_A(x):=\rho(x,A)$, with this notation we see that if $A,B$ are disjoint and closed then, the disjoint open sets $U=(\rho_A-\rho_B)^{-1}((0,\infty))$ and $V=(\rho_B-\rho_A)^{-1}((0,\infty))$ suffice as $x\in A$ implies $\rho_A(x)=0$ while $\rho_B(x)>0$ hence $A\subseteq V$ and similarly $B\subseteq U.$ 
\end{solution}
\begin{problem}\label{prob:4-4}
Let $X = \bR$ and let $\mathfrak{T}$ be the family of all subsets of $\bR$ of the form $U \cup (V \cap \bQ)$ where $U, V$ are open in the usual sense. Then $\mathfrak{T}$ is a topology that is Hausdorff but not regular. (In view of Exercise 3, this shows that a topology stronger than a normal topology need not be normal or even regular.)
\end{problem}
\begin{solution}
This topology is stronger than the usual topology on $\bR$ which is hausdorff and thus this too is hausdorff. 
Before we prove that it is not regular, it is worth to mention that intuitively speaking, regularity needn't carry on to stronger topologies as new closed sets might be introduced as well. 
Now, consider the closed set $\bQ^c$ in this topology and let $x\in\bQ$, then we see that if $U_1\cup(V_1\cap\bQ)\supseteq\bQ^c,U_2\cup(V_2\cap\bQ)\ni x$ are open as described then we clearly need to have $U_1\supseteq\bQ^c$ and we claim that $U_1\cap(V_2\cap\bQ)$ is nonempty. 
We can find $a<b$ such that $x\in(a,b)\cap\bQ\subseteq(V_2\cap\bQ)$ as all the usual open sets in $\bR$ are an union of open intervals. 
Now, we can cleary see that if $q\in U_1\cap(a,b)$ then there is some $(c,d)\subseteq U_1\cap(a,b)$ containing $q$, but this interval must have nonempty intersection with $\bQ\cap(a,b)$ by the density of rationals, thus we can never find an open set containing $x$ that is disjoint to some open set containing $\bQ^c$ proving that $\mathfrak T$ is not regular and thus not normal either (in this book, normal and regular spaces need to be $T_1$ by definition and thus the singletons are closed, making $T_3\supset T_4)$. 
\end{solution}
\begin{problem}\label{prob:4-5}
Every separable metric space is second countable.
\end{problem}
\begin{solution}
For any countable dense subset $S$ we have the following countable base $ \qty{B(x,r):r\in\bQ,r>0,x\in S}$, the fact that this is a base can be easily proven. 
\end{solution}
\begin{problem}\label{prob:4-6}
Let $\mE = \{(a,b] : -\infty < a < b < \infty\}$.
\begin{enumerate}[label=\alph*.]
    \item $\mE$ is a base for a topology $\mathfrak{T}$ on $\bR$ in which the members of $\mE$ are both open and closed.
    \item $\mathfrak{T}$ is first countable but not second countable. (If $x \in \bR$, every neighborhood base at $x$ contains a set whose supremum is $x$.)
    \item $\bQ$ is dense in $\bR$ with respect to $\mathfrak{T}$. (Thus the converse of Proposition 4.5 is false.)
\end{enumerate}
\end{problem}
\begin{solution}
\begin{enumerate}[label=\alph*.]
    \item As $\mathfrak T(\mathcal E)$ contains all unions of elements in $\mathcal E$ we see that $(a,b]^c=(-\infty,a]\cup(b,\infty]$ is open in this topology as the rays can further be written as unions of the bounded h-open intervals from the base as well. 
        This shows that the complement of each element of the basis is open so they themslves are closed. 
    \item For any $x\in\bR$ we can consider $(x-1,x]\ni x$. 
        If $\mN$ is any neighborhood base then there must exist some $U\in\mN$ such that $x\in U\subseteq(x-1,x]$ which trivially implies that $\sup U=x$ and thus $\mN$ must be uncountable as $\bR$ is, so it is not second countable. 
        To show that it is first countable it suffices to prove that for $x\in\bR $ the following is a neighborhood base :
        $ \qty{(x-r,x]:r\in\bQ,r>0}$. 
        Now, gievn any $S\in\mathfrak T(\mathcal E)$ (due to proposition 4.4) is an union of finite intersections over $\mathcal E$ and thus if it contains $x$ then we can find a finite intersection over $\mathcal E$, with $x$, that is a subset of $S$. 
        Say its of form $\cap_{i=1}^n[x_i,x_i-\delta_i)$ then we see that the set $[x,x-\delta)$ for $\delta\in\bQ$ such that $0<\delta<\min_i\delta_i$ is contained in it and thus this is a neighborhood base. 
    \item We prove the equivalent statement that $\emptyset=(\overline{\bQ})^c=(\bQ^c)^\circ$ so we have to prove that that there are no open sets in $\bQ^c.$ 
        With knowledge of structure of open sets generated by some base it suffices to show that no finite intersections over this base are contained in $\bQ^c$. 
        We can use a fact used in chapter 1 that the intersection of finitely many h-intervals of this form is an union of some disjoint h-intervals so it suffices to show that no  h-intervals are in $\bQ^c$ which is trivially true by the density of $\bQ$ in $\bR$ in the usual sense. 
        Thus, every seperable space is not second countable.
\end{enumerate}
\end{solution}
\begin{problem}\label{prob:4-7}
If $X$ is a topological space, a point $x \in X$ is called a cluster point of the sequence $\{x_j\}$ if for every neighborhood $U$ of $x$, $x_j \in U$ for infinitely many $j$. If $X$ is first countable, $x$ is a cluster point of $\{x_j\}$ iff some subsequence of $\{x_j\}$ converges to $x$.
\end{problem}
\begin{solution}
    If some subsequence converges then the statement clearly holds, for the other direction first consider some neighborhood base $ \qty{ U_n}_{n\geq 0}$ and let $ \qty{\cap_{m\leq n}U_m}_{n\geq 0}$, we see that this is also a neighborhood base and $V_n$ decreases. 
    As infinitely many terms of the sequence are in any neighborhood we can inductively choose $n_1<n_2<\ldots$ such that $x_{n_k}\in V_k$. 
    Now, we see that for any neighborhood $N$ of $x$ we can find $V_k\subseteq N$ and thus, $x_{n_j}\in V_j\subseteq V_k\subseteq N$ for all $j\geq k$ so $x_{n_j}\to x$ by definition.
\end{solution}
\begin{problem}\label{prob:4-8}
If $X$ is an infinite set with the cofinite topology and $\{x_j\}$ is a sequence of distinct points in $X$, then $x_j \to x$ for every $x \in X$.
\end{problem}
\begin{solution}
If $N\ni x$ is any neighborhood then by definition $N^c$ is finite and thus for large enough $j$ we have $x_j\in N$ as the terms of the sequence are distinct, thus it must converge to $x$ by definition.
\end{solution}
\begin{problem}\label{prob:4-9}
If $X$ is a linearly ordered set, the topology $\mathfrak{T}$ generated by the sets $\{x : x < a\}$ and $\{x : x > a\}$ ($a \in X$) is called the order topology.
\begin{enumerate}[label=\alph*.]
    \item If $a, b \in X$ and $a < b$, there exist $U, V \in \mathfrak{T}$ with $a \in U$, $b \in V$, and $x < y$ for all $x \in U$ and $y \in V$. The order topology is the weakest topology with this property.
    \item If $Y \subset X$, the order topology on $Y$ is never stronger than, but may be weaker than, the relative topology on $Y$ induced by the order topology on $X$.
    \item The order topology on $\bR$ is the usual topology.
\end{enumerate}
\end{problem}
\begin{solution}
\begin{enumerate}[label=\alph*.]
    \item If there exists $c$ such that $a<c<b$ then $U= \qty{x:x<c}$ and $V= \qty{x:x>c}$ work, otherwise let $U= \qty{x:x<b}$ and $V= \qty{x:x>a}.$ 
        If some topology $J$ has this property then fix some $a\in X$, for all $x>a$ there exists $U_x,V_x$ such that $a\in U_x,x\in V_x$ such that $u<v$ for all $u\in U,v\in V$ and thus we must have that $V_x\subseteq \qty{k:k>a}$ and thus,
        \[
            \bigcup_{X\ni x>a}V_x= \qty{k:k>a}\in J
        \]
        Similarly it can be shown that $ \qty{k:k<a}\in J$ so the order topology must be contained in it. 
    \item Its clearly never strictly stronger as the rays generating the order topology on $Y$ must be in the base generating that on $X$ as well. 
        Inclusion can be strictly weaker however, for example consider $X=\bR$ with the usual order and let $Y=[0,1)\cup \qty{2}$, then we see that $ \qty{2}$ is open in the relative topology clearly but not in the order topology as all basis sets containing $2$ must be of form $ \qty{x\in Y:a<x}$ and here $a<1$ so this contains $(1-\delta,1)$ for some small $\delta>0$ too and we fail to isolate $2.$  
    \item This follows from the fact that open sets in $\bR$ in the usual sense are unions of open intervals and proposition 4.4.
\end{enumerate}
\end{solution}
\begin{problem}\label{prob:4-10}
A topological space $X$ is called disconnected if there exist nonempty open sets $U, V$ such that $U \cap V = \emptyset$ and $U \cup V = X$; otherwise $X$ is connected. When we speak of connected or disconnected subsets of $X$, we refer to the relative topology on them.
\begin{enumerate}[label=\alph*.]
    \item $X$ is connected iff $\emptyset$ and $X$ are the only subsets of $X$ that are both open and closed.
    \item If $\{E_\alpha\}_{\alpha \in A}$ is a collection of connected subsets of $X$ such that $\bigcap_{\alpha \in A} E_\alpha \neq \emptyset$, then $\bigcup_{\alpha \in A} E_\alpha$ is connected.
    \item If $A \subset X$ is connected, then $\overline{A}$ is connected.
    \item Every point $x \in X$ is contained in a unique maximal connected subset of $X$, and this subset is closed. (It is called the connected component of $x$.)
\end{enumerate}
\end{problem}
\begin{solution}
\begin{enumerate}[label=\alph*.]
\item This follows directly from the definition. 
\item For the sake of contradiction assume that $\cup E_\alpha$ is disconnected and the nonempty open sets $U,V$ separate it. 
    Then, as $E_\alpha$ are all connected we either have $E_\alpha\subseteq U$ or $E_\alpha\subseteq V$ as otherwise $E_\alpha\cap U,E_\alpha\cap V$ is a valid separation for $E_\alpha$ (in the relative topology of $E_\alpha$.) 
    But, as $U,V$ are both nonempty we must have some $E_\alpha\subseteq U$ and $E_\beta\subseteq V$ which contradicts the fact that $\cap E_\alpha\neq\emptyset.$
\item Again for the sake of contradiction assume otherwise, let $U,V$ be a separation for $\overline{A}$. 
    Then, if $A\cap U,A\cap V$ are both nonempty then this serves as a separation in the relative topology for $A$ and would imply that $A$ was disconnected as well so this can't be the case. 
    We know that $\overline{A}=A\cup acc(A)$ and if wlog $A\cap U$ was empty then we would have $U\subseteq acc(A)$ which implies that $U$ has an accumulation point of $A$ and thus by definition $U\cap A$ is nonempty which is a contradiction so both of these must be nonempty which is a contradiction, thus $\overline{A}$ must be connected to. 
\item Let $C$ be the union of all connected sets containing $x.$ 
    By part (b) we must have that $C$ is connected too making $C$ the maximal connected set containing $x$ and as its closure is connected too by part (c) we must have that, $C=\overline{C}$ making it closed. 
\end{enumerate}
\end{solution}
\begin{problem}\label{prob:4-11}
If $E_1, \dots, E_n$ are subsets of a topological space, the closure of $\bigcup_{1}^n E_j$ is $\bigcup_{1}^n \overline{E}_j$.
\end{problem}
\begin{solution}
    As $\cup E_j\subseteq\cup \overline{E_j}$ and finite unions of closed sets are closed we must have that $\overline{\cup E_j}\subseteq\cup\overline{E_j}$. 
    Now, if $C$ is a closed set containing $\cup E_j$ then it contains each $E_j$ and thus their closures as well, thus $C\supseteq\cup\overline{E_j}$. 
    Thus, by definition we have $\overline{\cup E_j}=\cap C\supseteq\cup\overline{E_j}$ where the intersection is over all closed sets containing $\cup E_j.$ 
    We thus have equality, having shown both sides of the set inclusion. 
\end{solution}
\begin{problem}\label{prob:4-12}
Let $X$ be a set. A Kuratowski closure operator on $X$ is a map $A \mapsto A^*$ from $\mathcal{P}(X)$ to itself satisfying (i) $\emptyset^* = \emptyset$, (ii) $A \subset A^*$ for all $A$, (iii) $(A^*)^* = A^*$ for all $A$, and (iv) $(A \cup B)^* = A^* \cup B^*$ for all $A, B$.
\begin{enumerate}[label=\alph*.]
    \item If $X$ is a topological space, the map $A \mapsto \overline{A}$ is a Kuratowski closure operator. (Use Exercise 11.)
    \item Conversely, given a Kuratowski closure operator, let $\mathcal{F} = \{A \subset X : A = A^*\}$ and $\mathfrak{T} = \{U \subset X : U^c \in \mathcal{F}\}$. Then $\mathfrak{T}$ is a topology, and for any set $A \subset X$, $A^*$ is its closure with respect to $\mathfrak{T}$.
\end{enumerate}
\end{problem}
\begin{solution}
\begin{enumerate}[label=\alph*.]
    \item (iv) follows from \hyperref[prob:4-11]{Problem 4.11} and the rest are trivial. 
    \item As $\emptyset^*=\emptyset$ we see that $X\in\mathfrak T$ and as $X\subseteq X^*$, we have $X=X^*$ so we similarly have $\emptyset\in\mathfrak T$ too. 
        Now, if $E_\alpha\in\mathfrak T$ then we see that $E_\alpha^c=(E_\alpha^c)^*$ for all $\alpha$ and $\cup E_\alpha\in\mathfrak T$ iff $\cap E_\alpha^c=(\cap E_\alpha^c)^*$. 
        The $\subseteq$ direction follows from definition and as $\cap (E_\alpha)^c\subseteq E_\alpha^c=(E_\alpha^c)^*$, by (iv, which can be used to deduce monotonicity) we have $(\cap (E_\alpha)^c)^*\subseteq E_\alpha^c$ for all $\alpha$ and hence the other direction $(\cap E_\alpha^c)^*\subseteq\cap E_\alpha^c$ also holds. 
        Closure under finite intersections is equivalent to closure of $\mathcal F$ under finite unions which follows from (iv) easily and we thus have that $\mathfrak T$ is a topology. 
        In this topology the closed sets are precisely the fixed points of this map and thus by monotonicty we see that,
        \[
            A^*\subseteq\bigcap_{\substack{A\subseteq K\\ K\in\mathcal F}}K\subseteq A^*
        \]
        as $A^*\in\mathcal F$ too, hence $^*$ is the closure operator with respect to $\mathfrak T.$ 
\end{enumerate}
\end{solution}
\begin{problem}\label{prob:4-13}
If $X$ is a topological space, $U$ is open in $X$, and $A$ is dense in $X$, then $\overline{U} = \overline{U \cap A}$.
\end{problem}
\begin{solution}
    It suffices to show that $U\subseteq\overline{U\cap A}.$ 
    Assume for the sake of contradiction that there exists some $x\in U$ such that $x\in(\overline{U\cap A})^c=H$ which is open. 
    Thus, $U\cap H$ is a nonempty (they both contain $x$) open set and as $A$ is dense, we must have that $A\cap(U\cap H)$ is nonempty which is absurd as $A\cap U\cap H=(A\cap U)\cap(U\setminus\overline{A\cap U})=\emptyset$ so we are done. 
\end{solution}
\begin{problem}\label{prob:4-14}
If $X$ and $Y$ are topological spaces, $f : X \to Y$ is continuous iff $f(\overline{A}) \subset \overline{f(A)}$ for all $A \subset X$ iff $\overline{f^{-1}(B)} \subset f^{-1}(\overline{B})$ for all $B \subset Y$.
\end{problem}
\begin{solution}
    If $f$ is continuous then for all closed $C$ containing $f(A)$ we must have that $A\subseteq f^{-1}(C)$ which is closed and thus $ \overline{A} \subseteq f^{-1}(C)$ and thus, $f( \overline{A} )\subseteq f(f^{-1}(C))\subseteq C$ so we must have that $f( \overline{A}) \subseteq\cap C= \overline{f(A)}$ where the union is over all such closed $C.$ 
    If the second hypothesis holds, then as
    \[
        f^{-1}( \overline{B})=\bigcap_{\substack{B \subseteq C\\ C\text{ closed}}}f^{-1}(C)
    \]
    and for all the $C$ above, $\overline{B}\subseteq C\implies \overline{f(B)} \subseteq f( \overline{B}) \subseteq f(C)$ so we see that the third hypothesis holds too. 
    Now, if the third hypothesis holds then for all closed $C$ we must have $ \overline{f^{-1}(C)} \subseteq f^{-1}(C)$ as $C=\overline{C}$ and this implies that $f^{-1}(C)= \overline{f^{-1}(C)}$ so preimages of closed sets are closed, which is equivalent to continuity so we are done. 
\end{solution}
\begin{problem}\label{prob:4-15}
If $X$ is a topological space, $A \subset X$ is closed, and $g \in C(A)$ satisfies $g = 0$ on $\partial A$, then the extension of $g$ to $X$ defined by $g(x) = 0$ for $x \in A^c$ is continuous.
\end{problem}
\begin{solution}
    We see that $X=A\cup\overline{A^c}$ and as $ \overline{A^c}=A^c\cup\partial A$, it must be the case that $g$ is continuous on both $\overline{A^c}$ (its constant here) and $A$ and hence, for any closed $C$ 
    \[
        g^{-1}(C)=(g^{-1}(C)\cap A)\cup(g^{-1}(C)\cap\overline{A^c})=g|_{A}^{-1}(C)\cup g|_{\overline{A^c}}^{-1}(C)
    \]
    and as $A$ and $ \overline{A^c}$ are both closed we must have that $g|_{A}^{-1}(C)$ and $g|_{\overline{A^c}}^{-1}(C)$ are both closed. 
    As finite unions of closed sets are closed, $g^{-1}(C)$ is closed so we are done. 
\end{solution}
\begin{problem}\label{prob:4-16}
Let $X$ be a topological space, $Y$ a Hausdorff space, and $f, g$ continuous maps from $X$ to $Y$.
\begin{enumerate}[label=\alph*.]
    \item $\{x : f(x) = g(x)\}$ is closed.
    \item If $f = g$ on a dense subset of $X$, then $f = g$ on all of $X$.
\end{enumerate}
\end{problem}
\begin{solution}
\begin{enumerate}[label=\alph*.]
    \item Consider $E=Y\times Y$ with the product topology and let $h:X\to E$ be defined by $h(x):=(f(x),g(x))$ which is continuous as its components are. 
        Here, we see that,
        \[
            \qty{x\in X : f(x) = g(x)} = h^{-1}( \qty{(y,y):y\in Y})
        \]
        and we claim that the diagonal $\Delta$ is closed. 
        It suffices to prove that its complement is open, indeed if $x\neq y$ and $U_x,U_y$ are disjoint open sets separating them, we see that $U_x\times U_y$ is open in the product topology and
        \[
            \Delta^c=\bigcup_{(x,y)\in\Delta^c}U_x\times U_y
        \]
        is open so we are done. 
    \item This follows from the fact that closed dense sets are always the whole space itself.
\end{enumerate}
\end{solution}
\begin{problem}\label{prob:4-17}
If $X$ is a set, $\mathfrak{F}$ a collection of real-valued functions on $X$, and $\mathfrak{T}$ the weak topology generated by $\mathfrak{F}$, then $\mathfrak{T}$ is Hausdorff iff for every $x, y \in X$ with $x \neq y$ there exists $f \in \mathfrak{F}$ with $f(x) \neq f(y)$.
\end{problem}
\begin{solution}
    If $x,y$ are distinct and there exists some $f$ that maps them to different values then, as $f$ is continuous over the weak topology, for disjoint open sets $U\ni f(x),V\ni f(y)$ we see that $f^{-1}(U)\ni x,f^{-1}(V)\ni y$ are disjoint open sets so $\mathfrak T$ is hausdorff. 
    Now, suppose $f(x)=f(y)=c$ for all $f$, for some  distinct $x,y$. 
    Then for all open $U\in\bR$ we see that $c\in U$ implies $f^{-1}(U)$ contains both $x,y$ for all $f$, and otherwise it contains none of $x,y$ for all $f$. 
    Thus, all the sets in the generator either have both $x,y$ or none of these and by proposition 4.4 the open sets too have this property and are unable to separate these distinct points which makes $\mathfrak T$ not hausdorff so we are done. 
\end{solution}
\begin{problem}\label{prob:4-18}
If $X$ and $Y$ are topological spaces and $y_0 \in Y$, then $X$ is homeomorphic to $X \times \{y_0\}$ where the latter has the relative topology as a subset of $X \times Y$.
\end{problem}
\begin{solution}
We can consider the bijection $x\mapsto (x,y_0)$. 
This is continuous as under the relative topology here, a generator of $X\times\{y_0\}$ is $\{U\times\{y_0\}:U\text{ open in }X\}$ and the preimages of these under this map are exactly the open sets in $X$, so we are done by proposition $4.9$. 
\end{solution}
\begin{problem}\label{prob:4-19}
If $\{X_\alpha\}$ is a family of topological spaces, $X = \prod_{\alpha} X_\alpha$ (with the product topology) is uniquely determined up to homeomorphism by the following property: There exist continuous maps $\pi_\alpha : X \to X_\alpha$ such that if $Y$ is any topological space and $f_\alpha \in C(Y,X_\alpha)$ for each $\alpha$, there is a unique $F \in C(Y,X)$ such that $f_\alpha = \pi_\alpha \circ F$. (Thus $X$ is the category-theoretic product of the $X_\alpha$ in the category of topological spaces.)
\end{problem}
\begin{solution}
We describe the situation with a commutative diagram below. 
% https://q.uiver.app/#q=WzAsNCxbMiwwLCJZIl0sWzIsMiwiXFxwcm9kIFhfXFxhbHBoYSJdLFs0LDIsIlhfXFxiZXRhIl0sWzAsMiwiWF9cXGFscGhhIl0sWzAsMSwiXFxleGlzdHMgIUYiXSxbMCwyLCJmX1xcYmV0YSJdLFsxLDIsIlxccGlfXFxiZXRhIiwyXSxbMCwzLCJmX1xcYWxwaGEiLDJdLFsxLDMsIlxccGlfXFxhbHBoYSJdXQ==
\[\begin{tikzcd}[row sep=2.25em]
	&& Y && \\
	\\
	{X_\alpha} && {\prod X_\alpha} && {X_\beta}
	\arrow["{f_\alpha}"', from=1-3, to=3-1]
	\arrow["{\exists !F}", from=1-3, to=3-3]
	\arrow["{f_\beta}", from=1-3, to=3-5]
	\arrow["{\pi_\alpha}", from=3-3, to=3-1]
	\arrow["{\pi_\beta}"', from=3-3, to=3-5]
\end{tikzcd}\]
If $f_\alpha\in C(Y,X_\alpha)$ for all $\alpha$ then we see that $f_\alpha^{-1}(U_\alpha)$ is open in $Y$ for all $\alpha$, where $U_\alpha$ is open in $X_\alpha$. 
If such a $F$ were to exist then for this diagram to commute we must have that $\pi_\alpha\circ F=f_\alpha$ and thus, $F$ is forced to be the map sending $y\mapsto (f_\alpha(y))_{\alpha\in A}$ and is also continuous as all its components are due to how the product topology is defined and so it is indeed the unique such map in $C(Y,X).$ 
Now, suppose that $Z$ had the same properties, let $\tau$ take the place of $\pi$ and if $Y=X$ and $f_\alpha=\pi_\alpha\in C(X,X_\alpha)$ (the projections are continuous) then by definition there exists a map $G\in C(X,Z)$ such that $\pi_\alpha=\tau_\alpha\circ G$ i.e this diagram commutes
% https://q.uiver.app/#q=WzAsNCxbMiwwLCJcXHByb2QgWF9cXGFscGhhIl0sWzQsMiwiWF9cXGJldGEiXSxbMCwyLCJYX1xcYWxwaGEiXSxbMiwyLCJaIl0sWzAsMSwiXFxwaV9cXGJldGEiLDJdLFswLDIsIlxccGlfXFxhbHBoYSJdLFszLDEsIlxcdGF1X1xcYmV0YSIsMl0sWzMsMiwiXFx0YXVfXFxhbHBoYSJdLFswLDMsIkciLDJdXQ==
\[\begin{tikzcd}[row sep=2.25em]
	&& {\prod X_\alpha} && \\
	\\
	{X_\alpha} && Z && {X_\beta}
	\arrow["{\pi_\alpha}", from=1-3, to=3-1]
	\arrow["G"', from=1-3, to=3-3]
	\arrow["{\pi_\beta}"', from=1-3, to=3-5]
	\arrow["{\tau_\alpha}", from=3-3, to=3-1]
	\arrow["{\tau_\beta}"', from=3-3, to=3-5]
\end{tikzcd}\]
By definition there exists $H\in C(Z,X)$ such that $\tau_\alpha=\pi_\alpha\circ H$ implying, $\pi_\alpha=\pi_\alpha\circ(H\circ G)$ and similarly, $\tau_\alpha=\tau_\alpha\circ(G\circ H)$. 
Now, by uniqueness of $F$ in the exercise statement we must have that $H\circ G=e_X$ and $G\circ H=e_Z$ as $\pi_\alpha=\pi_\alpha\circ e_X$ and $\tau_\alpha=\tau_\alpha\circ e_Z$ which tells us that $H$ is invertible and as composition of continuous maps are continuous, $X$ and $Z$ are homeomorphic through $H$ so we are done. 
\end{solution}
\begin{problem}\label{prob:4-20}
If $A$ is a countable set and $X_\alpha$ is a first (resp. second) countable space for each $\alpha \in A$, then $\prod_{\alpha \in A} X_\alpha$ is first (resp. second) countable.
\end{problem}
\begin{solution}
    Assume WLOG $A=\bN$. 
    If $X_k$ are first countable then for some $x\in\prod X_k$ consider the neighborhood bases $\mN(x_k)$ of $x_k$ in $X_k.$
    Looking at the generators of the product topology of form $\cap_{j=1}^n\pi_{k_j}^{-1}(U_{k_j})$ containing $x$ i.e. $U_{k_j}\ni x_{k_j}$ open in $X_{k_j}$ we see that there exists $V_{k_j}\in\mN(x_{k_j})$ such that $x_{k_j}\in V_{k_j} \subseteq U_{k_j}$ and thus,
    \[
        x\in\bigcap_{j=1}^n\pi^{-1}_{k_j}(V_{k_j}) \subseteq\bigcap_{j=1}^n\pi_{k_j}^{-1}(U_{k_j})
    \]
   so the following is a neighborhood base for $x$ in the product topology
   \[
       \qty{\bigcap_{j\in J}\pi_j^{-1}(N_j):J \subseteq\bN,J\text{ is finite},(\forall j\in J)N_j\in\mN(x_j)}
   \]
   We can construct an obvious injection from this to
   \[
       \bigsqcup_{J \subseteq\bN:J\text{ finite}}\underbrace{\prod_{j\in J}\bN}_{\text{countable}}
   \]
   which is countable, being a countable disjoint union of countable sets. 
   Now, if they were second countable then the exact same construction gives us a base if we replace the countable neighborhood bases with bases.
\end{solution}
\begin{problem}\label{prob:4-21}
If $X$ is an infinite set with the cofinite topology, then every $f \in C(X)$ is constant.
\end{problem}
\begin{solution}
    By continuity we see that for any closed $C\subseteq\bR$ we must have $f^{-1}(C)$ closed in $X$ and hence finite, unless its all $X.$ 
    Now, suppose $f$ is not constant and there are distinct $x,y\in f(X).$ 
    Let, $C_x,C_y$ be closed sets formed from $\bR$ by removing small enough open intervals around $x,y$ respectively such that $C_x\cup C_y=\bR.$ 
    Then we see that $X=f^{-1}(\bR)=f^{-1}(C_x)\cup f^{-1}(C_y)$ where both the preimages are closed and not the whole space as $x\neq y$ and thus they must be finite which is impossible as $X$ is infinite. 
\end{solution}
\begin{problem}\label{prob:4-22}
Let $X$ be a topological space, $(Y,\rho)$ a complete metric space, and $\{f_n\}$ a sequence in $Y^X$ such that $\sup_{x \in X} \rho(f_n(x), f_m(x)) \to 0$ as $m, n \to \infty$. There is a unique $f \in Y^X$ such that $\sup_{x \in X} \rho(f_n(x), f(x)) \to 0$ as $n \to \infty$. If each $f_n$ is continuous, so is $f$.
\end{problem}
\begin{solution}
    Clearly for any fixed $x\in X$ the sequence $ \qty{f_n(x)}_{n}$ is cauchy and thus converges to say $f(x).$ 
    For any $\epsilon>0$ choose $N$ such that $m,n\geq N$ implies $\rho(f_n(x),f_m(y))<\epsilon$ for all $x\in X$ and thus $\rho(f(x),f_n(x))\leq\rho(f_n(x),f_m(x))+\rho(f_m(x),f(x))<\epsilon+\rho(f_m(x),f(x))$ 
    As, $m\to\infty$ we find that $\rho(f_m(x),f(x))\to 0$ so we must have that $\rho(f_n(x),f(x))\leq\epsilon$ and thus $\sup_{x\in X}\rho(f(x),f_n(x))$ must converge to $0$ as $n\to\infty.$ 
    Now suppose that $g$ was another such function, then as $\rho(f(x),g(x))\leq\rho(f_n(x),f(x))+\rho(f_n(x),g(x))\to 0$ for all $x\in X$ we must have $f=g$ proving uniqueness. 
   With the metric topology in mind it suffices to prove that for all $x\in X$ and $\epsilon>0$ there exists a neighborhood $V$ of $x$ such that $f(V)\subseteq B(f(x),\epsilon)$. 
   First choose $N$ such that $\rho(f(x),f_N(x))<\epsilon/3$ for all $x$ (we can do this by the previous part), then if $V$ is a neighbornood of $x$ such that $f_N(V)\subseteq B(f_N(x),\epsilon/3)$ then, for all $y\in V$
   \[
   \rho(f(x),f(y))\leq\rho(f_N(x),f_N(y))+\rho(f_N(x),f(x))+\rho(f_N(y),f(y))< \frac{\epsilon}{3}+\frac{\epsilon}{3}+\frac{\epsilon}{3}=\epsilon 
   \]
   and thus $f(V)\subseteq B(f(x),\epsilon)$ and we are done as the balls around $f(x)$ form a neighborhood base and this easily shows that $f$ is continuous at all $x\in X$ and thus it is continuous over $X. $ 
\end{solution}
\begin{problem}\label{prob:4-23}
Give an elementary proof of the Tietze extension theorem for the case $X = \bR$.
\end{problem}
\begin{solution}
    Suppose $f\in C(A)$ with $A$ closed. 
    As open sets in $\bR$ are disjoint unions of countably many open intervals. 
    Say we have $A^c=\cup (a_n,b_n)$ where these intervals are disjoint; this might be a finite union and rays are also considered an interval. 
    If $(a_n,b_n)$ is a bounded interval then define $F$ to be the line joining $f(a_n)$ to $f(b_n)$ on this, if its a ray (suppose left open wlog) $(-\infty,b)$ then set $F$ to be $f(b)$ here and if $a\in A$ then set $F(a)=f(a).$ 
    This definetly works but going through all the cases to prove this using sequential criterion is something I will postpone for later.   
\end{solution}
\begin{problem}\label{prob:4-24}
A Hausdorff (it does not have to be) space $X$ is normal iff $X$ satisfies the conclusion of Urysohn's lemma iff $X$ satisfies the conclusion of the Tietze extension theorem.
\end{problem}
\begin{solution}
If $X$ is normal then it satisfies the conclusion of urysohn's lemma. 
If $X$ satisfies the conclusion of urysohn's lemma then it satisfies the conclusion of tietze's extension theorem as normality of the space is not used in the proof, only tietze's theorem is used (look at Theorem 4.16) 
If $X$ satisfies the conclusion of tietze's extension theorem we will prove that it is normal. 
Suppose $A,B$ are disjoint closed sets and we have $f\in C(A\cup B,[0,1])$ such that $f=1/3$ on $A$ and $f=2/3$ on $B$ and $F$ be an extension then clearly $U=F^{-1}([0,1/2)),V=F^{-1}((1/2,1])$ are open and disjoint such that $A\subseteq U$ and $B\subseteq V.$
\end{solution}
\begin{problem}\label{prob:4-25}
If $(X,\mathfrak{T})$ is completely regular, then $\mathfrak{T}$ is the weak topology generated by $C(X)$.
\end{problem}
\begin{solution}
If $A$ is a closed set and $x\in A^c$ then we can find $f_x\in C(X)$ with $f_x(x)=1$ and $f_x=0$ on $A$. 
Now consider the closed sets $f_x^{-1}( \qty{0})$, these contain $A$ and exclude $x$, thus we see that, 
\[
    A=\bigcap_{x\in A^c}f_x^{-1}( \qty{0})\quad\text{and}\quad A^c=\bigcup_{x\in A^c}f_x^{-1}(\bC\setminus \qty{0})
\]
hence we see that all open sets (these can be expressed as the complement of a closed set all ways) are in the weak topology generated by $C(X)$ as this contains all sets of form $f^{-1}(\bC\setminus \qty{0})$ with $f\in C(X).$ 
Thus, $\mathfrak T$ is contained in this weak topology (the inclusion in the other direction holds from the definition of weak topologies.)  
\end{solution}
\begin{problem}\label{prob:4-26}
Let $X$ and $Y$ be topological spaces.
\begin{enumerate}[label=\alph*.]
    \item If $X$ is connected (see Exercise 10) and $f \in C(X,Y)$, then $f(X)$ is connected.
    \item $X$ is called arcwise connected if for all $x_0, x_1 \in X$ there exists $f \in C([0,1],X)$ with $f(0) = x_0$ and $f(1) = x_1$. Every arcwise connected space is connected.
    \item Let $X = \{(s,t) \in \bR^2 : t = \sin(s^{-1})\} \cup \{(0,0)\}$, with the relative topology induced from $\bR^2$. Then $X$ is connected but not arcwise connected.
\end{enumerate}
\end{problem}
\begin{solution}
\begin{enumerate}[label=\alph*.]
    \item If $f(X)=U\cup V$ with $U,V$ nonempty, open and disjoint then we see that $X=f^{-1}(U)\cup f^{-1}(V)$ is also a valid separation for $X$ so if $X$ is connected $f(X)$ must be connected too. 
    \item Suppose $U,V$ are open, disjoint, nonempty, their union $X$ and $x_0\in U,x_1\in V$. 
        If $f\in C([0,1],X)$ is such that $f(0)=x_0$ and $f(1)=x_1$ then $f^{-1}(U),f^{-1}(V)$ are open, nonempty and partition $[0,1]$ which is a contradiction as $[0,1]$ is connected thus, arcwise connected spaces are always connected. 
    \item Let $G(x)=(x,\sin(1/x))$, then we see that $G(\bR_+)\cup G(\bR_-)\cup \qty{(0,0)}=X$ and each of these are connected by part (a). 
        It is easily provable that any open $U\ni (0,0)$ in the relative topology of $X$ has nonempty intersection with $G(\bR_+)$ and thus if $V$ is another open set in the relative topology,  disjoint to $U$ such that $U\cup V=X$ then $G(\bR_+)=(U\cap G(\bR_+))\cup(V\cap G(\bR_+))$ is a valid separation which is absurd as $G(\bR_+)$ is connected so $X$ must be connected. 
        It is however not arc conected as we can not join $(0,0)$ to $(\pi,0)$ through $X$ with a continuous function. 
        More rigourously, if $f:C([0,1],X)$ is such that $f(0)=(\pi,0)$ and $f(1)=(0,0)$ then we would need $\norm{f(x)}$ to be small enough, say less than $1/2$ for all $x\in(1-\delta,1]$ for some small $\delta>0$ so we must have that,
        \[
            f((1-\delta,1))\subseteq J=\qty{(s,\sin(1/s))\in\bR^2:1-\delta<s<1,|\sin(1/s)|<1/2}
        \]
        But $J$ is disconnected set with each connected component bounded away from $(0,0)$, i.e. they dont have this as an accumulation point. 
        Thus, we see that $f((1-\delta,1))$ must be in one of these components by part (a) and thus it does not have $f(1)$ as an accumulation point which is absurd when $f$ is continuous. 
\end{enumerate}
\end{solution}
\begin{problem}\label{prob:4-27}
If $X_\alpha$ is connected for each $\alpha \in A$ (see Exercise 10), then $X = \prod_{\alpha \in A} X_\alpha$ is connected. (Fix $x \in X$ and let $Y$ be the connected component of $x$ in $X$. Show that $Y$ includes $\{y \in X : \pi_\alpha(y) = \pi_\alpha(x) \text{ for all but finitely many } \alpha\}$ and that the latter set is dense in $X$. Use Exercises 10 and 18.)
\end{problem}
\begin{solution}
Let the set described in the question be $S$, then we find that
\[
    S=\bigcup_{\substack{J\subseteq A\\ J\text{ finite}}}\underbrace{\prod_{\alpha\in A}L_{\alpha,J}}_{\text{Let this be }P_J}\quad\text{ where }L_{\alpha,J}=X_\alpha\text{ if }\alpha\in J\text{ and }x_\alpha\text{ otherwise}
\]
Now, for any $J= \qty{\alpha_1,\ldots,\alpha_n}$ we can show using similar arguments as in \hyperref[prob:4-18]{Problem 1.18} that $\prod L_{\alpha,J}$ is homeomorphic to $X_{\alpha_1}\times\ldots\times X_{\alpha_n}$ which we show to be connected at the end, so its connectedd itslef. 
As all of these products contain $x$, thus \hyperref[prob:4-10]{Problem 4.10} implies that $S$ is continuous so $S\subseteq Y$ by definition. 
If $\cap_1^n\pi_{\alpha_i}^{-1}(U_{\alpha_i})$ is any set in the conventional base for this space then we see that there is a point in this that only differs from $x$ at the index $\alpha_1,ldots,\alpha_n$ so it has nonempty intersection with $S$ and thus all open sets do too making $S$ dense, thus we have $X= \overline{S}\subseteq \overline{Y}=Y$ as connected components are closed and we are done. 
To show that this is true for finite $A$ as used as a fact earlier we can induct by adding a coordinate on every time in a sense. 
Say $X,Y$ are connected then, the statement follows for $X\times Y$ directly using the proof above by Problem 18 as these individual spaces are connected. 
Now, we claim $(X\times Y)\times Z$ is homeomorphic to $X\times Y\times Z$ which will let us induct and prove it for finite $A.$ 
Let, $f:(X\times Y)\times Z\to X\times Y\times Z$ be defined by $((x,y),z)\mapsto (x,y,z)$ then this is continuous, looking at $\pi_1\circ f$ (the rest are proved to be continuous similarly) as, if $U,V,W$ are open in $X,Y,Z$ respectively then $(\pi_1\circ f)^{-1}(U)=f^{-1}(U\times Y\times Z)=(U\times Y)\times Z$ is open in $(X\times Y)\times Z$ due to $U\times Y$ being open in $X\times Y$ and $Z$ being open in $Z$, which is sufficient as components being continuous are equivalent to the function itself being continuous by the definition of the product topology. 
\end{solution}
\begin{problem}\label{prob:4-28}
Let $X$ be a topological space equipped with an equivalence relation, $\tilde{X}$ the set of equivalence classes, $\pi : X \to \tilde{X}$ the map taking each $x \in X$ to its equivalence class, and $\mathfrak{T} = \{U \subset \tilde{X} : \pi^{-1}(U) \text{ is open in } X\}$.
\begin{enumerate}[label=\alph*.]
    \item $\mathfrak{T}$ is a topology on $\tilde{X}$. (It is called the quotient topology.)
    \item If $Y$ is a topological space, $f : \tilde{X} \to Y$ is continuous iff $f \circ \pi$ is continuous.
    \item $\tilde{X}$ is $T_1$ iff every equivalence class is closed.
\end{enumerate}
\end{problem}
\begin{solution}
\begin{enumerate}[label=\alph*.]
    \item Clearly $\emptyset,\tilde{X}\in\mathfrak T$, the rest follows from $\pi^{-1}$ commuting with unions and intersections. 
    \item Clearly $\pi$ is continuous by definition and thus the only if direction is trivial, for the if direction it suffices to see that if $f\circ\pi$ is continuous then for all open $U \subseteq Y$, $\pi^{-1}(f^{-1}(U))$ is open, which implies that $f^{-1}(U)$ is open by definition so $f$ must be continuous. 
    \item We use the equivalence that spaces are $T_1$ iff the singletons are closed. 
        If $\{[x]\}\subseteq\tilde{X}$ were closed then $\pi^{-1}(\{[x]\})=[x]\subseteq X$ is closed too as $\pi$ is continuous and if the equivalence class are closed then for $[x]\neq[y]$ in $\tilde{X}$ we see that $[y]^c=\cup_{z\notin[y]}[z]$ is open in $X$ and omits $y$ while containing $x$ and $\cup_{z\notin[y]} \qty{[z]}$ is open in $\tilde{X}$ (its preimage under $\pi$ is $[y]^c$ which is open in $X$) and omits $[y]$ while containing $[x]$ so the quotient topology for this must be $T_1.$ 
\end{enumerate}
\end{solution}
\begin{problem}\label{prob:4-29}
If $X$ is a topological space and $G$ is a group of homeomorphisms from $X$ to itself, $G$ induces an equivalence relation on $X$, namely, $x \sim y$ iff $y = g(x)$ for some $g \in G$. Let $X = \bR^2$; describe the quotient space $\tilde{X}$ and the quotient topology on it (as in Exercise 28) for each of the following groups of invertible linear maps. In particular, show that in (a) the quotient space is homeomorphic to $[0,\infty)$; in (b) it is $T_1$ but not Hausdorff; in (c) it is $T_0$ but not $T_1$, and in (d) it is not $T_0$. (In fact, in (d) $X$ is uncountable, but there are only six open sets and there are points $p \in \tilde{X}$ such that $\overline{\{p\}} = \tilde{X}$.)
\begin{enumerate}[label=\alph*.]
    \item $\left\{ \begin{pmatrix} \cos \theta & -\sin \theta \\ \sin \theta & \cos \theta \end{pmatrix} : \theta \in \bR \right\}$
    \item $\left\{ \begin{pmatrix} 1 & a \\ 0 & 1 \end{pmatrix} : a \in \bR \right\}$
    \item $\left\{ \begin{pmatrix} a & b \\ 0 & 1 \end{pmatrix} : a > 0, b \in \bR \right\}$
    \item $\left\{ \begin{pmatrix} a & 0 \\ 0 & b \end{pmatrix} : a, b \in \bQ \setminus \{0\} \right\}$
\end{enumerate}
\end{problem}
\begin{solution}
\begin{enumerate}[label=\alph*.]
\item These are rotation matrices, each equivalence class consists of vectors of a specific length. 
    Let $f:\tilde{X}\to [0,\infty)$ be defined by $[x]\mapsto r$ where $\norm{x}=r\geq 0$, its well defined by the previous line.  
    Using \hyperref[prob:4-28]{Problem 4.28} we see that this is continuous as $f\circ\pi$ is the norm i.e. $f\circ\pi:x\mapsto\norm{x}$, which is continuous and thus so is $f.$ 
    It is clearly a bijection, hence $\tilde X$ is homeomorphic to $[0,\infty).$ 
\item These are shear matrices and if we let the one associated to $a\in\bR$ be $g_a$ then we see that $g_a(x,y)=(x+ay,y)$ from this we see that the equivalence classes are $[(x,0)]=\{(x,0)\}$ and $[(x,y)]=\bR\times\{y\}$ for $y\neq 0$. 
    By Problem 28, its $T_1$ as the equivalence classes are all closed in $\bR^2$ but its clearly not hausdorff as we can not find any open sets separating the $[(1,0)]$ and $[(2,0)]$ for example : if $U$ is open and $[(1,0)]\in U$ then $\pi^{-1}(U)$ is also open and contains $(1,0)$ which forces it to contain $(1,y)$ for all $y$ close enough to $0$ (but not $0$) and hence $[(1,y)]\in\pi(\pi^{-1}(U))\subseteq U$ and if $V$ was an open set containing $[(2,0)]$ then it too would contain all the equivalence classes $[(2,y)]=[(1,y)]$ for small enough nonzero $y$ and thus $U\cap V\neq\emptyset$, so $\tilde{X}$ is not hausdorff. 
\item Here the equivalence classes are $[(0,0)]=\{(0,0)\},[(x,0)]=\bR_+\times\{0\}$ if $x>0$ and $\bR_-\times\{0\}$ if $x<0$, $[(x,y)]=\bR\times\{y\}$ if $y\neq 0$. 
    Now, consider the points $[(0,0)]$ and $[(1,0)]$, if $U$ is any open set containing $[(0,0)]$ then by logic similar to before, it always contains $[(1,0)]$ but we the open set $V=\tilde{X}\setminus \qty{[(0,0)]}$ (it is open as $\pi^{-1}(U)=\bR^2\setminus\{(0,0)\}$) does not contain $[(0,0)]$ while containing $[(1,0)]$ so its not $T_1$. 
    From From the above we see that $[(0,0)]$ can be separated from any point so now assume that none of the two chosen classes are $[(0,0)].$ 
    If exactly one of them has points with $y$ coordinates being $0$ then either one of the open sets $U=\tilde{X}\setminus \qty{[(0,0)],[(1,0)]}$ or $V=\tilde{X}\setminus \qty{[(0,0)],[(-1,0)]}$ do the job ($\pi^{-1}(U)=\bR^2\setminus \qty{(x,0):x\leq 0}$ and $\pi^{-1}(V)=\bR^2\setminus \qty{(x,0):x\geq 0}$.) 
    Now, if both the points are equivalence classes with nonzero $y$ coordinates then they can be separated by open sets which correspond to open tubes in $\bR^2$ i.e. sets of form $\bR\times(a,b)$ so all together, this is a $T_0$ space. 
\item Let $Q=\bQ\setminus\{0\}$ for brevity and let $A= \sqrt{2}Q\times\sqrt{2}Q$ and $B=Q\times Q$ be two different equivalence classes, we show that they can not be separated even in the $T_0$ sense. 
    If $U$ is open and contains $A\in U$ then as $A\subseteq\pi^{-1}(U)$ we find that there is an open ball in $\bR^2$ around the point $(\sqrt{2},\sqrt{2})$ that is a subset of $\pi^{-1}(U)$ and this clearly contains a point from $B$ so we must have that $B\in\pi(\pi^{-1}(U))$ and the exact same argument relying on the fact that $aQ$ is dense in $\bR$ for any nonzero $a$, shows that no open set containing $B$ can exclude $A$ so the space is not $T_0.$ 
    For the extra questions at the end : it is uncountable because each equivalence class is countable whereas the real plane is uncountable, the six open sets are (here $\bR\setminus\{0\}$ is written as $R$) :
    \[
        \tilde{X},\emptyset,\pi(\bR^2\setminus\{(0,0)\}),\pi(\bR\times R),\pi(R\times\bR),\pi(R\times R),
    \]
    Now, a point $p$ for which $\overline{\{p\}}=\tilde{X}$ can be found from the above to be $p=Q\times Q$ as it has nonempty intersection with every nonempty open set in $\tilde{X}.$  
\end{enumerate} 
\end{solution}
\begin{problem}\label{prob:4-30}
If $A$ is a directed set, a subset $B$ of $A$ is called cofinal in $A$ if for each $\alpha \in A$ there exists $\beta \in B$ such that $\beta \ge \alpha$.
\begin{enumerate}[label=\alph*.]
    \item If $B$ is cofinal in $A$ and $\langle x_\alpha \rangle_{\alpha \in A}$ is a net, the inclusion map $B \to A$ makes $\langle x_\beta \rangle_{\beta \in B}$ a subnet of $\langle x_\alpha \rangle_{\alpha \in A}$.
    \item If $\langle x_\alpha \rangle_{\alpha \in A}$ is a net in a topological space, then $\langle x_\alpha \rangle$ converges to $x$ iff for every cofinal $B \subset A$ there is a cofinal $C \subset B$ such that $\langle x_\gamma \rangle_{\gamma \in C}$ converges to $x$.
\end{enumerate}
\end{problem}
\begin{solution}
\begin{enumerate}[label=\alph*.]
    \item This follows from the definition. 
    \item If $x_\alpha\to x$ then we show that all subnets (say $\langle y_\beta\rangle_{\beta\in B}$) converge to $x$ too, this clearly proves the forward direction using the previous part. 
        If $U$ is any neighborhood of $x$ then there exists $\gamma\in A$ such that $\alpha\gtrsim\gamma$ implies $x_\alpha\in U$ and by definition there is $\beta_0\in B$ such that $\beta\geq\beta_0$ implies $\alpha_\beta\geq\gamma$ and thus, $y_\beta=x_{\alpha_\beta}\in U$ for all $\beta\geq\beta_0$ proving that every subnet converges. 
        To prove the backward direction, assume that it does not converge, thus there is a neighborhood $U$ of $x$ such that $x_\alpha\in U^c$ frequently and thus by definition $B= \qty{\gamma\in A:x_\gamma\in U^c}$ is cofinal. 
        Now, for any $C$ cofinal in $B$ we find that $\langle x_\zeta\rangle_{\zeta\in C}\subseteq U^c$ so this does not converge to $x$. 
\end{enumerate}
\end{solution}
\begin{problem}\label{prob:4-31}
Let $\langle x_n \rangle_{n \in \bN}$ be a sequence.
\begin{enumerate}[label=\alph*.]
    \item If $k \mapsto n_k$ is a map from $\bN$ to itself, then $\langle x_{n_k} \rangle_{k \in \bN}$ is a subnet of $\langle x_n \rangle$ iff $n_k \to \infty$ as $k \to \infty$, and it is a subsequence (as defined in \S 0.1) iff $n_k$ is strictly increasing in $k$.
    \item There is a natural one-to-one correspondence between the subsequences of $\langle x_n \rangle$ and the subnets of $\langle x_n \rangle$ defined by cofinal sets as in Exercise 30.
\end{enumerate}
\end{problem}
\begin{solution}
\begin{enumerate}[label=\alph*.]
    \item This follows from definition.
    \item The subsequence $(x_{n_k})$ corresponds to to the cofinal set $\{n_j:j\in\bN\}$ and thus the subnet defined by it. 
\end{enumerate}
\end{solution}
\begin{problem}\label{prob:4-32}
A topological space $X$ is Hausdorff iff every net in $X$ converges to at most one point. (If $X$ is not Hausdorff, let $x$ and $y$ be distinct points with no disjoint neighborhoods, and consider the directed set $\mN_x \times \mN_y$ where $\mN_x, \mN_y$ are the families of neighborhoods of $x, y$.)
\end{problem}
\begin{solution}
If $x,y$ are distinct points and $x_\alpha\to x$ then for some neighborhoods $U\ni x,V\ni y_1$ that are disjoint we find that $x_\alpha\in U,x_\alpha\in V^c$ eventually so $x_\alpha$ is clearly not in $V^c$ eventually and thus can not converge to $y.$ 
Now, if $X$ is not hausdorff then consider $x,y$ distinct whose neighborhoods always intersect. 
Consider $\mN_x,\mN_y$ the collection of neighborhoods of $x,y$ respectively, ordered by reverse inclusion and thus consider the directed set $\mN_x\times\mN_y$ with order induced by its coordinates. 
Choose $x_{(U,V)}\in U\cap V$ which we can do for these are nonempty by hypothesis. 
Then, for any $(J,K)\in\mN_x\times\mN_y$ we see that $(U,V)\gtrsim(J,K)\implies U\subseteq J,V\subseteq K\implies U\cap V\subseteq J\cap K$ and $x_{(U,V)}\in J\cap K\subseteq J,K$ proving that the net $\langle x_{(U,V)}\rangle_{(U,V)\in \mN_x\times\mN_y}$ converges to both $x,y.$  
\end{solution}
\begin{problem}\label{prob:4-33}
Let $\langle x_\alpha \rangle_{\alpha \in A}$ be a net in a topological space, and for each $\alpha \in A$ let $E_\alpha = \{x_\beta : \beta \gtrsim \alpha\}$. Then $x$ is a cluster point of $\langle x_\alpha \rangle$ iff $x \in \bigcap_{\alpha \in A} \overline{E}_\alpha$.
\end{problem}
\begin{solution}
The only if direction follows from definition and for the if direction notice how the inclusion implies that for all neighborhoods $U$ of $x$ and $\alpha\in A$ there exists $\beta\gtrsim\alpha$ with $x_\beta\in U$, as $x\in \overline{E_\beta}$ which, by the definition of a cluster point for nets implies that $x$ is a cluster point of our net.  
\end{solution}
\begin{problem}\label{prob:4-34}
If $X$ has the weak topology generated by a family $\mathfrak{F}$ of functions, then $\langle x_\alpha \rangle$ converges to $x \in X$ iff $\langle f(x_\alpha) \rangle$ converges to $f(x)$ for all $f \in \mathfrak{F}$. (In particular, if $X = \prod_{i \in I} X_i$, then $x_\alpha \to x$ iff $\pi_i(x_\alpha) \to \pi_i(x)$ for all $i \in I$..)
\end{problem}
\begin{solution}
    If $X$ has the weak topology generated by $\mathfrak F$ then $f\in\mathfrak F$ are continuous and the conclusion follows, for the other direction notice that for any neighborhood $U$ of $x$ there exists $f_1,\ldots,f_n\in\mathfrak F$ and neighborhoods $U_j$ of $f_j(x)$ such that $x\in\cap_1^nf_j^{-1}(U_j)\subseteq U$ and as $f_j(x_\alpha)\to f_j(x)$ for these $j$ we see that there exists $\alpha_j$ such that $f_j(x_\alpha)\in U_j$ for all $\alpha\gtrsim\alpha_j$ and by definition of directed sets, we can inductively find $\beta\gtrsim\alpha_1,\ldots,\alpha_n$ and thus $f_j(x_\alpha)\in U_j$ for all $j$ when $\alpha\gtrsim\beta.$ 
    Thus, $x_\alpha\in\cap_1^nf_j^{-1}(U_j) \subseteq U$ for all $\alpha\gtrsim\beta$ so $x_\alpha\to x$ by definition. 
\end{solution}
\begin{problem}\label{prob:4-35}
Let $X$ be a set and $\mA$ the collection of all finite subsets of $X$, directed by inclusion. Let $f : X \to \bR$ be an arbitrary function, and for $A \in \mA$, let $z_A = \sum_{x \in A} f(x)$. Then the net $\langle z_A \rangle$ converges in $\bR$ iff $\{x : f(x) \neq 0\}$ is a countable set $\{x_n\}_{n \in \bN}$ and $\sum_{1}^\infty |f(x_n)| < \infty$, in which case $z_A \to \sum_{1}^\infty f(x_n)$. (Cf. Proposition 0.20.)
\end{problem}
\begin{solution}
    It suffices to assume that $f\geq 0,$ to define the sum over possibly uncountable sets we must split it into positive and negative parts anyways. 
    We prove the if direction first, let $S= \qty{f\neq 0}= \qty{x_n}_{n\in\bN}$ and $L=\sum_1^\infty f(x_n)$, then we find that for all elements in the neighborhood base $U=(L-\epsilon,L+\epsilon)$ ($\epsilon>0)$ of $L$ there exists $N$ such that $z_ \qty{x_1,\ldots,x_k}\in U$ for all $k\geq N$ and thus, if $\beta\gtrsim \qty{x_1,\ldots,x_N}$ then $L-\epsilon<z_\beta<L$ and thus $z_\beta\in U$ so the net converges to $L.$ 
    If $ \qty{f\neq 0}$ is not countable then we can easily find arbitrarily large $A$ such that $z_A=|A|/n$ where $n\in\bN$ is fixed and thus we can not have convergence in $\bR$ as, for any $B$ we have $A\cup B\gtrsim B$ and $z_{A\cup B}\geq z_A>|A|/n$ meaning that the partial sums are unbounded.  
    Now, let $S=\qty{x_n}_{n\in\bN}= \qty{f\neq 0}$ (this might be a finite set too) and, if $z_A\to L$ we see that, for any $\epsilon>0$ there exists $B$ such that $A\gtrsim B\implies z_A\in(L-\epsilon,L+\epsilon)$ and thus $\sum_1^\infty f(x_n)=z_{B\cup S}\in(L-\epsilon,L+\epsilon)$ which implies $L=\sum_1^\infty f(x_n)$ as $\epsilon$ was arbitrary. 
\end{solution}
\begin{problem}\label{prob:4-36}
Let $X$ be the set of Lebesgue measurable complex-valued functions on $[0,1]$. There is no topology $\mathfrak{T}$ on $X$ such that a sequence $\langle f_n \rangle$ converges to $f$ with respect to $\mathfrak{T}$ iff $f_n \to f$ a.e. (Use Corollary 2.32 and Exercises 30b and 31b.)
\end{problem}
\begin{solution}
    If such a topology did exist, then using the problems mentioned in the statement, we know that $\langle f_n\rangle$ converges iff for all subsequences $f_{n_k}$ of $f_n$ there is a subsequence of $f_{n_{k_j}}$ that converges. 
    There are lebesgue measurable functions $f_n,f$ such that $f_n$ converges to $f$ in $L^1$ but not pointwise a.e. and by the mentioned corollary, given any subsequence $f_{n_k}$, we can find a subsequence in the previous subsequence, say $f_{n_{k_j}}$ for which $f_{n_{k_j}}\to f$ a.e. and thus $\langle f_n\rangle$ converges to $f$ which is a contradiction as $f_n$ does not converge to $f$ a.e. (An example of such a $f_n,f$ is the jumping interval function in chapter 2's section on convergence in measure.)
\end{solution}
\begin{problem}\label{prob:4-37}
Let $0'$ denote a point that is not an element of $(-1,1)$, and let $X = (-1,1) \cup \{0'\}$. Let $\mathfrak{T}$ be the topology on $X$ generated by the sets $(-1, a)$, $(a, 1)$, $[(-1,b) \setminus \{0\}] \cup \{0'\}$, and $[(c,1) \setminus \{0\}] \cup \{0'\}$ where $-1 < a < 1$, $0 < b < 1$, and $-1 < c < 0$. (One should picture $X$ as $(-1,1)$ with the point $0$ split in two.)
\begin{enumerate}[label=\alph*.]
    \item Define $f,g : (-1,1) \to X$ by $f(x) = x$ for all $x$, $g(x) = x$ for $x \neq 0$, and $g(0) = 0'$. Then $f$ and $g$ are homeomorphisms onto their ranges.
    \item $X$ is $T_1$ but not Hausdorff, although each point of $X$ has a neighborhood that is homeomorphic to $(-1,1)$ (and hence is Hausdorff).
    \item The sets $[-\frac{1}{2}, \frac{1}{2}]$ and $([-\frac{1}{2}, \frac{1}{2}] \setminus \{0\}) \cup \{0'\}$ are compact but not closed in $X$, and their intersection is not compact.
\end{enumerate}
\end{problem}
\begin{solution}
\begin{enumerate}[label=\alph*.]
    \item Both are injective. 
    This topology agrees with the usual topology on $(-1,1)$ (which is $f$'s range) and thus $f$ being continuous in with respect to the usual topology lends us its continuity in this case as well making it a homeomorphism on its range. 
    As sets in $\mathfrak T$ are unions of finite intersections of the generators we see that any set $U$ containing $0'$ must have an interval around $0$ excluding $0$ i.e. for some $\delta>0$ we have $(-\delta,\delta)\setminus \qty{0} \subseteq U$ so we find that for any neighborhood $U$ of $g(0)$ we can find some nonempty $(-\delta,\delta)\setminus \qty{0}\subseteq U^\circ$ and we see that $V=(-\delta,\delta)$ is a neighborhood of $0$ such that $g(V) \subseteq U$ so $g$ is continuous at $0.$ 
    Continuity at other points is trivial, this makes $g$ a homeomorphism on its range too.
    \item $X$ is $T_1$ by looking at the singletons which can easily shown to be close and it is not hausdorff as we can not separate $0,0'$ by open sets, this is clear from the fact that all neighborhoods of $0'$ must contain a deleted open neighborhood around $0$. 
    \item As the ranges of $f,g$ are both open and thus these being continuous when looked at as function onto their range implies their continuity in general and we see that 
        \[
            \qty[-\frac{1}{2}, \frac{1}{2}]=f\qty(\qty[ -\frac{1}{2}, \frac{1}{2}])\quad\qty(\qty[-\frac{1}{2}, \frac{1}{2}]\setminus \qty{0})\cup \qty{0'}=g\qty(\qty[ -\frac{1}{2}, \frac{1}{2}])
        \]
        proving their continuity as images of compact sets under continuous maps is compact. 
        Their intersection $[-1/2,1/2]\setminus \qty{0}$ is clearly not compact. 
\end{enumerate}
\end{solution}
\begin{problem}\label{prob:4-38}
Suppose that $(X, \mathfrak{T})$ is a compact Hausdorff space and $\mathfrak{T}'$ is another topology on $X$. If $\mathfrak{T}'$ is strictly stronger than $\mathfrak{T}$, then $(X,\mathfrak{T}')$ is Hausdorff but not compact. If $\mathfrak{T}'$ is strictly weaker than $\mathfrak{T}$, then $(X,\mathfrak{T}')$ is compact but not Hausdorff.
\end{problem}
\begin{solution}
If $\mathfrak T'$ is coarser then it is clearly compact as all open covers in it are also open covers in $\mathfrak T$ and if its finer then it is hausdorff as all neighborhoods in $\mathfrak T$ are also neighborhoods in $\mathfrak T'$. 
Now, if $\mathfrak T'$ is finer then the identity map from $\mathfrak T'$ to $\mathfrak T$ is continuous. 
If $\mathfrak T'$ is compact then its closed sets are also compact and thus get mapped to compact sets under the identity map, which are closed in $\mathfrak T$ as it is hausdorff hence implying $\mathfrak T' \subseteq \mathfrak T$ so $\mathfrak T'$ can not be strictly finer. 
Now, if $\mathfrak T'$ was strictly coarser then by what we just proved we see that this can not be hausdorff.
\end{solution}
\begin{problem}\label{prob:4-39}
Every sequentially compact space is countably compact.
\end{problem}
\begin{solution}
    If it isn't countably compact then we can find a countably infinite open subcover $ \qty{U_n}_{n\in\bN}$ with no finite subcover. 
    Removing $U_k$ for which $U_k\subseteq\cup_{n<k}U_n$ and reindexing we get another countablly infinite (otherwise, a finite subcover existed to begin with) open cover with no finite subcover from which we can pick $x_n\in U_n\setminus\cup_{k<n}U_k.$ 
   Now, for any $x$ and $U_j\ni x$ we see that $x_n\notin U_j$ for all $n>j$ and thus it can not have $x$ as a cluster point so no convergent subsequences exist either (these are also subnets, treating the main sequence as a net so we can't have these by Proposition 4.20, it is also easy to show that the existence of a convergent subsequence implies the existence of cluster points straight away.)  
\end{solution}
\begin{problem}\label{prob:4-40}
If $X$ is countably compact, then every sequence in $X$ has a cluster point. If $X$ is also first countable, then $X$ is sequentially compact.
\end{problem}
\begin{solution}
If the space is countably compact then analogous to general compactness we have another variation of proposition 4.21 : If $ \qty{F_n}$ is a countable collection of closed sets with FIP then $\cap F_n$ is nonempty. 
Now, if we set $F_n:=\overline{\qty{x_m:m\geq n}}$ then this collection has FIP and thus has nonempty intersection. 
The conclusion now follows from \hyperref[prob:4-33]{Problem 4.33} treating the sequence as a net in the usual manner. 
For the second conclusion, first pick a cluster point $x$ and a neighborhood base $ \qty{U_n}$ that is countable and decreasing (this is described in the book). 
Inductively pick integers $n_k>n_{k-1}>\ldots n_1$ such that $x_{n_k}\in U_{n_k}$ (we can do this because $x$ is a cluster point) and we clearly have $x_{n_k}\to x.$
\end{solution}
\begin{problem}\label{prob:4-41}
A $T_1$ space $X$ is countably compact iff every infinite subset of $X$ has an accumulation point.
\end{problem}
\begin{solution}
    We first show that cluster points of a sequence and accumulation points treating the sequence as a set are the same in $T_1$ spaces. 
    If $x$ is an accumulation point of $(x_n)$ then we can find neighborhoods $V_j$ containing $x$ but not $x_j$ and thus for any $n$ and neighourhood $U$ of $x$ we see that $U\cap V_1\cap\ldots\cap V_n \subseteq U$ must contain some $x_m$ for $m>n$ making $x$ a cluster point. 
    The cluster points are clearly accumulation points. 
From the previous problem we can say that any infinite set has an accumulation point by choosing a distinct sequence of points in it and the converse follows from the fact that if $X$ is not countably compact then we can find a sequence $(x_n)$ as in \hyperref[prob:4-39]{Problem 4.39} that has no cluster points and $ \qty{x_n:n\in\bN}$ is an infinite set with no accumulation points (these are the same as cluster points of this sequence in this case as proved in the start.)
\end{solution}
\begin{problem}\label{prob:4-42}
The set of countable ordinals $\Omega$ (\S 0.4) with the order topology (Exercise 9) is sequentially compact and first countable but not compact. (To prove sequential compactness, use Proposition 0.19.)
\end{problem}
\begin{solution}
For any $p\in\Omega$, let $q$ be strictly greater than $p$ (this must exist by the definition of $\Omega$.) 
Then the sets $S(a,b):=\qty{x:a<x<b}$ for $a\leq p\leq b\leq q$ forms a neighborhood base as all generator sets containing $p$ contain a set of this form : if $p\in \qty{x:x>t}$ then $p>t$ so $p\in S(t,q) \subseteq \qty{x:x>t}$ and $p\in \qty{x:x<t}$ means $p<t$ so $p\in S(\inf\Omega,\min(t,q)) \subseteq \qty{x:x<t}$. 
This is a countable because $I_q\cup \qty{q}$ is countable (by the definition of $\Omega$.)  
Now, consider a sequence $(x_n)$ in $\Omega$, we see that $\sup_{n>k}x_n$ exists in $\Omega$ for all $k$ by Proposition 0.19 and as $\Omega$ is well ordered by definition, $s=\inf_{n\in\bN}\sup_{n>k}x_n\in\Omega$. 
By construction we see that if $s<t$ then $x_n<t$ for all large enough $n$ and if $s>t$ then
$x_n>t$ infinitely often. 
If $x_n=s$ infinitely often then we can easily construct the converging subsequence, otherwise we see that $s$ is an accumulation point of $ \qty{x_n:n\in\bN}$. 
This makes $s$ a cluster point of the sequence by the first part of \hyperref[prob:4-39]{Problem 4.39} as the order topology is normal (and hence $T_1$) by \hyperref[prob:4-9]{Problem 4.9}a. 
As this is first countable, we can choose a subsequence converging to $s$ as done in the second part of \hyperref[prob:4-40]{Problem 4.40} (all neighborhoods of $s$ contain some terms from the sequence) and thus $\Omega$ is sequentially compact. 
The open cover $ \qty{I_x}_{x\in\Omega}$ (this is a cover as $\Omega$ has no maximal elements) does not admit finite subcovers as if it did then $\Omega$ would be countable, being a union of a finite number of countable sets. 
\end{solution}
\begin{problem}\label{prob:4-43}
For $x \in [0,1)$, let $\sum a_n(x) 2^{-n}$ ($a_n(x) = 0$ or $1$) be the base-2 decimal expansion of $x$. (If $x$ is a dyadic rational, choose the expansion such that $a_n(x) = 0$ for $n$ large.) Then the sequence $\langle a_n \rangle$ in $\{0,1\}^{[0,1)}$ has no pointwise convergent subsequence. (Hence $\{0,1\}^{[0,1)}$, with the product topology arising from the discrete topology on $\{0,1\}$, is not sequentially compact. It is, however, compact, as we shall show in \S 4.6.)
\end{problem}
\begin{solution}
    We know that convergence in the product topology is equivalent to pointwise convergence so for the subsequence $a_{n_k}$ to converge to say $a$ we would need to have $a_{n_k}(x)\to a(x)$ for all $x\in[0,1).$ 
    We show that for all such $a_{n_k}$ there exists $x\in[0,1)$ such that $a_{n_k}(x)$ does not converge. 
    For example, take $x=\sum_{k\in\bN}2^{-n_{2k}}$ (this is clearly not dyadic, each $1$ in the binary expansion must be $2$ zeroes apart from another.) and we have $a_{n_k}(x)=1$ for even $k$ and $0$ for odd $k$ so $a_{n_k}(x)$ does not converge and we are done. 
\end{solution}
\begin{problem}\label{prob:4-44}
If $X$ is countably compact and $f : X \to Y$ is continuous, then $f(X)$ is countably compact.
\end{problem}
\begin{solution}
The proof is the same as the case where we have normal continuity, see Proposition 4.26.
\end{solution}
\begin{problem}\label{prob:4-45}
If $X$ is normal, then $X$ is countably compact iff $C(X) = BC(X)$. (Use Exercises 40 and 44. If $\langle x_n \rangle$ is a sequence in $X$ with no cluster point, then $\{x_n : n \in \bN\}$ is closed, and Corollary 4.17 applies.)
\end{problem}
\begin{solution}
    If $X$ is countably compact then we find that $f(X) \subseteq\bC$ is countably compact too by the previous problem and thus we can find a finite cover of the countable open cover $ \qty{B(0,n)}_{n\in\bN}$ proving that $f$ is bounded. 
    For the other direction, say its not countably compact, then we can find a sequence $(x_n)$ of distinct points with no cluster points as in \hyperref[prob:4-39]{Problem 4.39} and as normal spaces are $T_1$ by definition, by using the first part of \hyperref[prob:4-41]{Problem 4.41} we can say that $A=\qty{x_n:n\in\bN}$ has no accumulation points and it is thus closed. 
    Further, as this has no accumulation points, there exists open neighborhoods of $x_m$ that omits all other points in $A$ and thus, in the relative topology of $A$ we find that $ \qty{x_m}\cap A= \qty{x_m}$ is open and thus $A$ inherits a discrete topology as its relative topology, this means that $C(A)=\bC^A$. 
    Let $f\in C(A)$ be such that $f(x_m)=m$, we can extend $f$ to $F\in C(X)$ on $X$ such that $F(x_m)=m$ too by Corollary 4.17 and as $F$ is continuous and unbounded, $C(X)\neq BC(X)$ so we are done. 
\end{solution}
\begin{problem}\label{prob:4-46}
Prove Theorem 4.34.
\end{problem}
\begin{solution}
    We first prove this for $f\in C(K,[0,1])$. 
    We can find some $U$ is precompact and $K \subseteq U \subseteq \overline{U} \subseteq X$ by Proposition 4.31. 
    The sets $f^{-1}([2/3,1])=F,f^{-1}([0,1/3])=G$ are closed in $K$ and thus compact (we work in a hausdorff space). 
    As $K$ is closed too, $G$ is closed in $X$ making $U\cap G^c$ a open set containing the compact set  $F$ and thus we can find $g_1\in C(X,[0,1])$ such that $g_1=1$ on $F$ and $\supp(g) \subseteq U\cap G^c \subseteq U$ by Theorem 4.32 so $0\leq f-g_1/3\leq 2/3$. 
    Now replacing $f$ with $(3/2)\cdot(f-g_1/3)\in C(K,[0,1])$ we can find $g_2\in C(X,[0,1])$ such that $\supp(g_2) \subseteq U$ and $0\leq (3/2)\cdot(f-g_1/3)-g_2/3\leq 2/3\implies 0\leq f-g_1/3-2g_2/3^2\leq (2/3)^2$. 
    We can continue like this and have a sequence $ \qty{g_n}_{n\in\bN} \subseteq C(X,[0,1])$ with $\supp(g_n) \subseteq U$ such that for all $m\in\bN$, $0\leq f-F_m\leq (2/3)^m$ on $K.$ where $F_m=(1/2)\cdot\sum_{k=1}^m (2/3)^kg_k$. 
    By the weierstrass M-test and proposition 4.13 we see that $F=\lim F_m$ is continuous and extends $f$ and as $\supp(g_n) \subseteq U$ we see that $ \qty{F\neq 0} \subseteq\cup \qty{g_n\neq 0}\implies \qty{F\neq 0} \subseteq U$ and thus $\supp(F) \subseteq \overline{U}$ so $F$ vanishes outside a compact set.  
    Now we prove this for $f\in C(K)$ (it can be easily proved from the previous result that it works on $C(K,[a,b])$ for real $a<b$.)  
    Let $g=f/(1+|f|)\in C(K,[-1,1])$ and choose $F\in C_c(X,[-1,1]))$ that extends $g$ and vanishes outside a compact set, say $L$. 
    Now, we define $h$ on the compact (it is a closed subset of the compact set $L\cup K$) set $F^{-1}( \qty{1,-1})\cup K$ to be $1$ on $K$ and $0$ on $F^{-1}( \qty{-1,1})$ (we can do this because $L,K$ are disjoint as $|g|=|G|<1$ on $K$ and $h$ is continuous on their union as its continuous on each of them and they are closed by the gluing lemma used previously) and extend to $H\in C_c(X,[0,1]).$ 
    Then we see that $|HG|<1$ and $\Phi=HG/(1-|HG|)\in C_c(X)$ (this inclusion is true as $HG\in C_c(X)$ already.) 
    It can be easily proven that $\Phi|_K=f.$
\end{solution}
\begin{problem}\label{prob:4-47}
Prove Proposition 4.36. Also, show that if $X$ is Hausdorff but not locally compact, Proposition 4.36 remains valid except that $X^*$ is not Hausdorff.
\end{problem}
\begin{solution}
Let the topology on $X$ be $\mathfrak T$ and that on $X^*$ be $\mathfrak T^*.$ 
If we have an open cover of $X^*$ then some of these sets must contain $\infty$ and thus are of form $K^c\cup \qty{\infty}$ for some compact $K$ in $(X,\mathfrak T)$, thus the rest of the sets  must cover $K$. 
These open sets in $\mathfrak T^*$ that cover $K$ are of either open in $\mathfrak T$ or the union of an open set in $\mathfrak T$ with $\infty$, so we can extract a finite subcover easily by ignoring $\infty$ as $\infty\not\in K$ giving us a finite subcover. 
If $p,q$ are distinct points in $X^*$ which are not $\infty$ then we can find open neighborhoods separating them as $\mathfrak T \subseteq \mathfrak T^*$ and if exactly one of them is $\infty$, say $q$ then as $X$ is LC we can find compact $N \subseteq X$ with respect to $\mathfrak T$ such that $p\in N^\circ$ which is open in $\mathfrak T$ and thus $\mathfrak T^*$, then $N^c\cup \qty{\infty}\in \mathfrak T^*$ is an open set containing $\infty$ and as $N^\circ\cap(N^c\cup\infty)=\emptyset$ as $\infty\notin N$ proving hausdorffness. 
The inclusion map is an embedding as the relative topology of $X$ in $X^*$ is that of $X$ (for compact (and hence closed) $K$ in $X$ we have  $X\cap(K^c\cup \qty{\infty})=K^c$ which is open in $X.)$ 
Now, if $C(X)\ni f=g+c$ for a constant $c$ and some $g\in C_0(X)$ then we show that $f$ can be extended to $C(X^*)$ by defining $f(\infty)=c.$ 
We first show continuity at $\infty.$ 
For any $\epsilon$ we see that $U=\qty{x\in X:|g(x)|\geq\epsilon}^c\cup\infty$ is such that $f(U) \subseteq B_{\bC}(f(\infty),\epsilon)$ and as $U$ is open in $\mathfrak \mathfrak T^*$ (as $g\in C_0(X)$) we have continuity at $\infty$ as the balls generate the topology on $\bC.$ 
Now, for any $x\in X$, as $f\in C(X)$ we readily have continuity at $x$ as $\mathfrak T \subseteq \mathfrak T^*$so $f$ is continuous on every point in $X^*$ and thus $f\in C(X^*).$ 
For the converse, suppose that $f$ can be extended continuously to $C(X^*)$ and we have $f(\infty)=c$, then we claim that $f-c\in C_0(X).$ 
For any $\epsilon>0$ we must have a neighborhood $B(f(\infty),\epsilon)$, we must have $f^{-1}(B(c,\epsilon))= \qty{\infty}\cup \qty{x\in X:|f(x)-c|\geq\epsilon}^c\in \mathfrak T^*$ which, by definition means that $ \qty{|f-c|\geq\epsilon}$ is compact in $X$ so, on $X$ we can write $f=(f-c)+c$ where $f-c\in C_0(X)$ and $c$ is some constant. 
Now, if $X$ was not locally compact then choose some point $p\in X$ without any compact neighborhoods and try to separate it from $\infty.$ 
If $p\in K^c\cup \qty{\infty}\in\mathfrak T^*$ for $K$ compact in $K$ then if $p\in U\in\mathfrak T$ we must have that $U\subseteq K$ which makes $K$ a compact neighborhood of $p$ and if otherwise $p\in N^c\cup \qty{\infty}$ in the same notation then this would mean that $N^c \subseteq K$ which again makes $K$ a compact neighborhood as $N^c$ is open (compact sets are closed in hausdorff spaces.) 
This proves that $X^*$ can't be hausdorff, but the rest of the proposition remains valid as we did not use local comapctness anywhere else in the proof. 
\end{solution}
\begin{problem}\label{prob:4-48}
Complete the proof of Proposition 4.40.
\end{problem}
\begin{solution}
    Given any compact $K$ as we have $K \subseteq\cup U_n$ and thus we can find a finite subcover $U_{n_1},\ldots,U_{n_k}$ and as $ \overline{ U_n} \subseteq U_{n+1}$ we see that $K \subseteq \overline{U_{\max n_k}}$, this shows that 
    \[
        \qty{g\in\bC^X:\sup_{x\in \overline{U_{n+1}}}|f(x)-g(x)|< \frac{1}{m}}\subseteq\qty{g\in\bC^X:\sup_{x\in K}|f(x)-g(x)|< \frac{1}{m}}
    \]
    so we have that its a neighborhood base for $f$ as claimed. 
    Now as this is a countable neighborhood base, we can get a decreasing neighborhood base from this as usual, the rest directly follows using this new base.  
\end{solution}
\begin{problem}\label{prob:4-49}
Let $X$ be a compact Hausdorff space and $E \subset X$.
\begin{enumerate}[label=\alph*.]
    \item If $E$ is open, then $E$ is locally compact in the relative topology.
    \item If $E$ is dense in $X$ and locally compact in the relative topology, then $E$ is open. (Use Exercise 13.)
    \item $E$ is locally compact in the relative topology iff $E$ is relatively open in $\overline{E}$.
\end{enumerate}
\end{problem}
\begin{solution}
    If $A \subseteq E$ is compact then for an open (relative to $E$) cover $ \qty{U_\alpha\cap E}$ of $A\cap E=A$, where $U_\alpha$ are open in $X$, we see that $\qty{U_\alpha}$ forms an open cover of $A$ in $X$ and we can extract a finite subcover $U_{\alpha_1},\ldots,U_{\alpha_n}$ and further get a finite open cover $ \qty{E\cap U_{\alpha_{k}}}$ for $A$ relative to $E$ too. 
    If it was the case $A \subseteq E$ is relatively compact in $E$ instead then for any open cover $ \qty{U_\alpha}$ of $A$ in $X$ we see that $ \qty{U_\alpha\cap E} $ is a open cover of $A$ in $E$ and thus we can find a finite subcover $ \qty{U_{\alpha_k}\cap E}$ which implies that $A \subseteq\cup U_{\alpha_k}$ which proves compactness in $X.$ 
    We have proved that for any $E$, $Y \subseteq E \subseteq X$ is compact relative to $X$ iff its compact relative to $E.$
\begin{enumerate}[label=\alph*.]
    \item 
    By normality of CH spaces, for $p\in E$ we can find open $U,V$ which are disjoint such that $p\in U$ and $V \subseteq E^c$ implying $p\in U \subseteq V^c \subseteq E$ and as $V^c$ is a closed subset of a compact space, its compact too so we can conclude that its a compact neighborhood of $p$ in $E$ as $p\in A\cap E \subseteq V^c$ using the proof at the start.
    \item
    For any $p\in E$ we can by hypothesis find a compact $K \subseteq E$ and open $U$ such that $p\in E\cap U \subseteq K$. 
    Using the mentioned exercise and the proof in the beginning we see that $ \overline{E\cap U}= \overline{U} \subseteq K$ as $K$ is compact (we work in a hausdorff space.) 
    This proves that $p\in U \subseteq K \subseteq E$ so $E$ must be open. 
    \item
    $ \overline{E}$ is compact being a closed subset of a compact space and if $E$ is relatively open in $ \overline{E}$ then $E$ is locally compact in the relative topology by part (a). 
    If $E$ is locally compact in the relative topology, then as $E$ is dense in $ \overline{E}$ in the relative topology of $ \overline{E}$ then by part (b) we see that it is relatively open in $ \overline{E}.$
\end{enumerate}
\end{solution}
\begin{problem}\label{prob:4-50}
Let $U$ be an open subset of a compact Hausdorff space $X$ and $U^*$ its one-point compactification (see Exercise 49a). If $\phi : X \to U^*$ is defined by $\phi(x) = x$ if $x \in U$ and $\phi(x) = \infty$ if $x \in U^c$, then $\phi$ is continuous.
\end{problem}
\begin{solution}
By the previous problem's part (a) we have that $U$ is locally compact in the relative topology and thus $U^*$ is compact and hausdorff ($U$ is hausdorff as well clearly.) 
$\phi$ is clearly continuous on $U$ and for $x\in U^c$, for any neighborhood $K^c\cup \qty{\infty}$ of $\phi(x)$ for $K$ compact in $U$ we see that $\phi^{-1}(K^c\cup \qty{\infty})=\phi^{-1}(K^c)\cup U^c=(U\cap K^c)\cup U^c=K^c$ is a neighborhood containing $x$ as its open in $X$ (this is because $K$ being compact in $U$ implies that it is compact in $X$ and thus closed too as $X$ is hausdorff, making its complement open.) 
So we have that $\phi$ is continuous everywhere making it continuous on $X.$ 
\end{solution}
\begin{problem}\label{prob:4-51}
If $X$ and $Y$ are topological spaces, $\phi \in C(X,Y)$ is called proper if $\phi^{-1}(K)$ is compact in $X$ for every compact $K \subset Y$. Suppose that $X$ and $Y$ are LCH spaces and $X^*$ and $Y^*$ are their one-point compactifications. If $\phi \in C(X,Y)$, then $\phi$ is proper iff $\phi$ extends continuously to a map from $X^*$ to $Y^*$ by setting $\phi(\infty_X) = \infty_Y$.
\end{problem}
\begin{solution}
    Suppose it is extensible as such, then for any open neighborhood $K^c\cup \qty{\infty}$ containing $\infty_Y$ ($K$ is an arbitrary compact set in $Y$), by continuity $\phi^{-1}(K^c\cup\infty_Y)=\phi^{-1}(K)^c\cup \qty{\infty_X}$ must be open in $X^*$ which forces $\phi^{-1}(K)$ to be compact in $X$ so $\phi$ is proper. 
    Now, if $\phi$ is proper and continuous, to prove that setting $\phi(\infty_X)=\infty_Y$ is a valid continuous extension it suffices to prove continuity at $\infty_X$ and we indeed see that for all open neighborhoods $K^c\cup \qty{\infty_Y}$ ($K$ is some arbitrary compact set in $Y$) of $\phi(\infty_X)$ we have $\phi^{-1}(K^c\cup \qty{\infty})=\phi^{-1}(K)^c\cup \qty{\infty_X}$ which is open in $X^*$ when $\phi$ is proper so we are done. 
\end{solution}
\begin{problem}\label{prob:4-52}
The one-point compactification of $\bR^n$ is homeomorphic to the $n$-sphere $\{x \in \bR^{n+1} : |x| = 1\}$.
\end{problem}
\begin{solution}
It can be shown using the stereographic projection $j$ that $\bR^n$ is homeomorphic to $S^n\setminus \qty{p}$ for some $p\in S^n$. 
Also, its clear that all homeomorphisms are proper so by the previous problem it suffices to show that $(S^n\setminus \qty{p})^*\cong S^n$ which follows directly from \hyperref[prob:4-50]{Problem 4.50} as the $\phi$ there is a homeomorphism in this case (we set $X=S^n,U=S^n\setminus \qty{p}$) as it maps $p$ to $\infty$ in the compactification. 
\end{solution}
\begin{problem}\label{prob:4-53}
Lemma 4.37 remains true if the assumption that $X$ is locally compact is replaced by the assumption that $X$ is first countable.
\end{problem}
\begin{solution}
    The forward direction works the same but for the if direction, if we have $x\in \overline{E}\setminus E$ then we can choose $(x_n) \subseteq E$ such that $x_n\to x$ by Proposition 4.6. Now, we claim that $K:= \qty{x_n:n\in\bN}\cup \qty{x}$ is compact, indeed from any open cover, the open set containing $x$ contains all but finitely many terms of the sequence so an open cover can be extracted. 
    Now consider $E\cap K= \qty{x_n:n\in\bN}$, this is not closed as it does not contain one of its accumulation points $x$ so we are done. 
\end{solution}
\begin{problem}\label{prob:4-54}
Let $\bQ$ have the relative topology induced from $\bR$.
\begin{enumerate}[label=\alph*.]
    \item $\bQ$ is not locally compact.
    \item $\bQ$ is $\sigma$-compact (it is a countable union of singleton sets), but uniform convergence on singletons (i.e., pointwise convergence) does not imply uniform convergence on compact subsets of $\bQ$.
\end{enumerate}
\end{problem}
\begin{solution}
\begin{enumerate}[label=\alph*.]
    \item
    We know from the beginning of \hyperref[prob:4-49]{Problem 4.49} that $K \subseteq Q$ is compact in the relative topology iff $K$ is compact in $\bR.$ 
    No such compact sets $K$ can contain any open sets $\bQ\cap U$ in $\bQ$ ($U$ is open in $\bR$)as by the density of rationals and hausdorffness (making compact sets closed), it would imply that $K\supseteq \overline{\bQ\cap U}=U$ and no nonempty open sets in $\bR$ can be subsets of $\bR$ as the irrationals are also dense in the reals. 
    If there are no such compact sets containing open sets in $\bQ$ to begin with, it can not be LC. 
    \item
        Its $\sigma$-compact by the hint already and a counterexample to prove the other claim is: $f_m(r):=\chi_{ \qty{1/n:n\in\bN,n\geq m}}(r)$, this does not converge uniformly to its limit, the constant zero function uniformly on the compact set $[0,1]$.
\end{enumerate}
\end{solution}
\begin{problem}\label{prob:4-55}
Every open set in a second countable LCH space is $\sigma$-compact.
\end{problem}
\begin{solution}
Suppose $ \qty{U_n}$ is a countable base. 
Given any open $U$, for $p$ suppose $K_p$ is a compact neighborhood of $p$ in $U$, we can choose such a set by Proposition 4.30.  
We can find some function $f:U\to\bN$ such that $p\in U_{f(p)} \subseteq K_p.$ 
Now, we see that $U \subseteq\cup_{n\in f(U)} U_n$ by construction and this is also a countable union as $f$ maps into the naturals. 
For $n\in f(U)$ let $x_n$ be such that $f(x_n)=n$ and thus, $U_n \subseteq K_{x_n}$ hence proving $U \subseteq\cup_{n\in f(U)}K_{x_n}$ proving that $U$ is sigma compact.
\end{solution}
\begin{problem}\label{prob:4-56}
Define $\Phi : [0,\infty] \to [0,1]$ by $\Phi(t) = \frac{t}{t+1}$ for $t \in [0,\infty)$ and $\Phi(\infty) = 1$.
\begin{enumerate}[label=\alph*.]
    \item $\Phi$ is strictly increasing and $\Phi(t+s) \le \Phi(t) + \Phi(s)$.
    \item If $(Y,\rho)$ is a metric space, then $\Phi \circ \rho$ is a bounded metric on $Y$ that defines the same topology as $\rho$.
    \item If $X$ is a topological space, the function $\rho(f,g) = \Phi(\sup_{x \in X} |f(x) - g(x)|)$ is a metric on $\bC^X$ whose associated topology is the topology of uniform convergence.
    \item If $X$ is a $\sigma$-compact LCH space and $\{U_n\}_{1}^\infty$ is as in Proposition 4.39, the function
        \[ \rho(f,g) = \sum_{1}^\infty 2^{-n} \Phi\qty(\sup_{x \in \overline{U_n}} |f(x) - g(x)|) \]
    is a metric on $\bC^X$ whose associated topology is the topology of uniform convergence on compact sets.
\end{enumerate}
\end{problem}
\begin{solution}
    We use the following lemma throughout the problem: If $A,B \subseteq \mathcal P(X)$ generate the topologies $P,Q$ respectively and if for all $U\in A$ and $x\in U$ we can find $E\in Q$ such that $x\in E \subseteq U$ then $A \subseteq Q$ and hence $P \subseteq Q$. 
    If this works both ways then $P=Q$. 
    The proof is simple. 
\begin{enumerate}[label=\alph*.]
    \item It is clearly strictly increasing, the inequality holds trivially if at least one of $t,s$ is infinity, so assume they aren't. 
        Then, 
        \[
            \Phi(t+s) = \frac{t+s}{t+s+1} = \frac{t}{t+s+1} + \frac{s}{t+s+1} \leq \frac{t}{t+1} + \frac{s}{s+1} = \Phi(t) + \Phi(s)
        \]
    \item As $\Phi(x)=0$ iff $x=0$, using part (a) and the triangle inequality for $\rho$, we can easily verify that $\Phi\circ\rho$ is a metric, and it is bounded as $\Phi \leq 1$. 
        Since $\Phi(t) \le t$, we immediately have $\Phi\circ\rho \leq \rho$, so $B_\rho(f, \epsilon) \subseteq B_{\Phi \circ \rho}(f, \epsilon)$. 
        For the other direction, since $\Phi$ is strictly increasing, $B_{\Phi \circ \rho}(f, \Phi(\epsilon)) \subseteq B_\rho(f, \epsilon)$. 
        As the metric balls generate the metric topology, this proves that they generate the same topology (we can employ the lemma as for any element in a ball, we can find another ball centered at this element that is contained in the larger one.) 
    \item For brevity we write $\rho_S(f,g)$ for $\Phi(\sup_{x\in S}|f(x)-g(x)|)$ and let $d_u$ be the standard uniform metric (we call it a metric because it assuming infinity is not an issue here). 
        By the exact same ball logic from part (b), $B_{d_u}(f, \epsilon) \subseteq B_{\rho}(f, \epsilon)$ and $B_{\rho}(f, \Phi(\epsilon)) \subseteq B_{d_u}(f, \epsilon)$.
        Thus, the topology associated with $\rho$ is clearly the topology of uniform convergence.
    \item For any compact $K$ and $m\in\bN$, we want to find $\epsilon>0$ such that (keep in mind that the inequality towards the right is equivalent to $d_u(f,g)<1/m$) 
        \[
            f\in B_\rho(f,\epsilon) \subseteq\qty{g\in\bC^X:\rho_K(f,g)< \Phi\qty(\frac{1}{m})}
        \]
        Notice that $\rho_{ \overline{U_n}}$ is increasing in $n$ as $\Phi$ is strictly increasing and $ \overline{U_n} \subseteq \overline{U_{n+1}}$. 
        We also have $K \subseteq \overline{U_n}$ and thus $\rho_K\leq\rho_{ \overline{U_n}}$ for large enough $n$ as proved in \hyperref[prob:4-48]{Problem 4.48}. 
        Thus,
        \[
            2^{-n}\rho_K(f,g)\leq 2^{-n}\rho_{ \overline{U_n}}(f,g)\leq\rho(f,g)
        \]
        implying that $\epsilon=2^{-n}\Phi(m^{-1})$ works. 
        Now, for any $\epsilon>0$ we want to find $K$ compact and $m\in\bN$ such that
        \[
            f\in\qty{g\in\bC^X:\rho_K(f,g)<\Phi(1/m)} \subseteq B_\rho(f,\epsilon)
        \]
        If $K= \overline{U_n}$, then we see that 
        \[
            \rho(f,g)\leq\sum_{k\leq n}2^{-k}\rho_{ \overline{U_k}}(f,g)+\sum_{k>n}2^{-k}\leq\rho_{ K}(f,g)+2^{-n}
        \]
        If $m,n$ are large enough such that $2^{-n}<\epsilon/2$ and $1/m<\Phi^{-1}(\epsilon/2)$, then $\rho_{K}(f,g) < \Phi(1/m)$ implies $\rho(f,g) < \epsilon$. 
        Thus, the topology defined by $\rho$ is exactly the topology of uniform convergence on compact sets, this can be proved using the same ball logic we have been using (fix the compact sets we take supremum over when necessary.)  
\end{enumerate}
\end{solution}
\begin{problem}\label{prob:4-57}
An open cover $\mathcal{U}$ of a topological space $X$ is called locally finite if each $x \in X$ has a neighborhood that intersects only finitely many members of $\mathcal{U}$. If $\mathcal{U}, \mathcal{V}$ are open covers of $X$, $\mathcal{V}$ is a refinement of $\mathcal{U}$ if for each $V \in \mathcal{V}$ there exists $U \in \mathcal{U}$ with $V \subset U$. $X$ is called paracompact if every open cover of $X$ has a locally finite refinement.
\begin{enumerate}[label=\alph*.]
    \item If $X$ is a $\sigma$-compact LCH space, then $X$ is paracompact. In fact, every open cover $\mathcal{U}$ has locally finite refinements $\{V_{\alpha}\}$, $\{W_{\alpha}\}$ such that $\overline{V}_{\alpha}$ is compact and $\overline{W_{\alpha}} \subset V_{\alpha}$ for all $\alpha$. (Let $\{U_n\}_{1}^\infty$ be as in Proposition 4.39. For each $n$, $\{E \cap (U_{n+2} \setminus \overline{U_{n-1}}) : E \in \mathcal{U}\}$ is an open cover of $\overline{U}_{n+1} \setminus U_n$. Choose a finite subcover to obtain $\{V_{\alpha}\}$ and mimic the beginning of the proof of Proposition 4.41 to obtain $\{W_{\alpha}\}$.)
    \item If $X$ is a $\sigma$-compact LCH space, for any open cover $\mathcal{U}$ of $X$ there is a partition of unity on $X$ subordinate to $\mathcal{U}$ and consisting of compactly supported functions.
\end{enumerate}
\end{problem}
\begin{solution}
\begin{enumerate}[label=\alph*.]
    \item As $U_{n+2}\setminus \overline{U_{n-1}} \supseteq \overline{U_{n+1}}\setminus U_n$ and $U_{n+2}\setminus \overline{U_{n-1}}$ is open (compact sets are closed here due to hausdorff property) we see that the cover mentioned in the statement is indeed open and as $ \overline{U_{n+1}}\setminus U_n$ is a closed subset of a compact set $ \overline{U_{n+1}}$ we see that its compact as well. 
        Thus we can extract a finite open cover for each $n$ and get a new open cover $ \qty{V_\alpha}$ (this is also countable) as $X=\cup_{n\geq 2}(U_{n+2}\setminus \overline{U_{n-1}})$ clearly. 
        $ \overline{V_\alpha}$ is compact, being a closed subset of the compact set $ \overline{U_{n+2}}$
        Now, for any $x\in X$ we see that its in atmost a finite number of sets of form $U_{n+2}\setminus\overline{U_{n-1}}$ and as each of these is covered by a finite number of $V_\alpha$'s, this must be a locally finite cover, thus proving paracompactness. 
        Now, for each $x\in X$ we can find a $V_\alpha$ containing it and then choose a compact neighborhood of $x$ contained in it by proposition 4.30 and as their interiors form a open cover of $U_{n+2}\setminus \overline{U_{n-1}}$ for each $n$ we can choose finite covers for these, say $N_1^\circ,\ldots,N_m^\circ$ (thus, the sets themselves also form a cover) and then let $W_\alpha$ be the union of all these $N_i^\circ$ where $N_i \subseteq V_\alpha$ (keep in mind that $n$ was fixed here) then we see that $ \overline{W_\alpha} = \cup N_i^\circ \subseteq\cup \overline{N_i}=\cup N_i \subseteq V_\alpha$ as $N_i$ are closed too. 
        $ \qty{W_\alpha}$ is clearly a refinement too and it is locally finite by the same logic we used for $ \qty{V_\alpha}$. 
    \item For each $\alpha$ let $g_\alpha\in C(X,[0,1])$ such that $g_\alpha=1$ on $ \overline{W_\alpha}$ and $\supp(g_\alpha) \subseteq V_\alpha$ by Urysohn's lemma as $ \overline{W_\alpha}$ is compact being a closed subset of the compact set $ \overline{V_\alpha}$. 
        We also see that as $\qty{ \overline{W_\alpha}}$ forms a locally finite cover (follows from local finiteness (only finitely many nonzero terms exist due to this) of $ \qty{V_\alpha}$), that $\sum_\alpha g_\alpha(x)$ is both always positive and always finite. 
        Let, $h_\alpha:=g_\alpha\cdot(\sum_\alpha g_\alpha)^{-1}$, this is well defined and $\supp(h_\alpha)=\supp(g_\alpha) \subseteq V_\alpha \subseteq \overline{V_\alpha}$ which is compact so $h_\alpha\in C_c(X)$ and for any $x\in X$ if $g_{\alpha_1},\ldots,g_{\alpha_k}$ are exactly the $g_\alpha$ that are nonzero at $x$ then, 
        \[
            \sum_\alpha h_\alpha(x)=\sum_{j=1}^kh_{\alpha_j}(x)=\sum_{j=1}^k \frac{g_{\alpha_j}(x)}{\sum_{i=1}^kg_{\alpha_i}(x)}=1
        \]
        as $\sum_\alpha g_\alpha(x)=\sum_{i=1}^kg_{\alpha_i}(x)$ here. 
        $ \qty{h_\alpha}$ is thus a partition of unity subordinate to $\mathcal U$ as $\qty{V_\alpha}$ was a refinement and it consists of compactly supported functions so we are done. 
\end{enumerate}
\end{solution}
\begin{problem}\label{prob:4-58}
If $\{X_\alpha\}_{\alpha \in A}$ is a family of topological spaces of which infinitely many are noncompact, then every closed compact subset of $\prod_{\alpha \in A} X_\alpha$ is nowhere dense.
\end{problem}
\begin{solution}
    Let $C$ be a compact closed set that is not nowhere dense, thus it has nonempty interior. ( $C$ is nowhere dense iff $(\overline{C})^\circ$ is empty and here, $\overline{C}=C$), we must have some open $U \subseteq C$ of form $U=\prod U_\alpha$ where $U_\alpha=X_\alpha$ for all but finitely many $\alpha.$ 
    As projections are continuous, $\pi_\alpha(C)$ is compact for all $\alpha$ and as $\pi_\alpha(U)=X_\alpha$ for all but finitely many $\alpha$ we have $X_\alpha=\pi_\alpha(U) \subseteq\pi_\alpha(C) \subseteq X_\alpha$ for all but finitely many $\alpha$ implying that $X_\alpha$ are compact for all but finitely many $\alpha$ which is a contradiction so we are done. 
\end{solution}
\begin{problem}\label{prob:4-59}
The product of finitely many locally compact spaces is locally compact.
\end{problem}
\begin{solution}
Say $X_1,\ldots,X_n$ are LC, then for any $x\in\prod X_i$ we can find compact neighborhoods $K_i\ni\pi_i(x)$ and open $U_i$ such that $\pi_i(x)\in U_i \subseteq K_i$, then by tychonoff's theorem $K=\prod K_i$ is compact in its relative topology which means that it is compact as we have proved before. 
Moreover, $U=\prod U_i$ is an open set and it contains $x$ so we see that $x\in U \subseteq K^\circ$ proving that $K$ is a compact neighborhood proving LC of the product as $x$ was arbitrary. 
\end{solution}
\begin{problem}\label{prob:4-60}
The product of countably many sequentially compact spaces is sequentially compact. (Use the ``diagonal trick'' as in the proof of Theorem 4.44.)
\end{problem}
\begin{solution}
    Suppose $ \qty{X_k}_{k\in\bN}$ are the sequentially compact spaces, then for a sequence $(x_n)$ in their product, we see that $\pi_1(x_n)$ must have a converging subsequence, let this be $\pi_1({x_{n,1}})$, then again $\pi_2(x_{n,1})$ must have a converging subsequence too which we call $\pi_2(x_{n,2})$ where $x_{n,2}$ is a subsequence of $x_{n,1}$, we keep doing this to get sequences $ \qty{ (x_{n,k})}_{k}$ where $\pi_k(x_{n,k})$ converges. 
    Now consider $y_k:=x_{k,k}$ which is also a subsequence of $(x_n)$ (it is easier to see this by writing some of the sequences out, you can also find a similar construction in Rudin's PMA) as $y_k$ is clearly a subsequence of $(x_{n,m})$ for $k\geq m$ we have that $\pi_m(y_k)$ converges. 
    As $(x_n)$ was arbitrary to begin with and we constructed a converging subsequence (all coordinates converging is equivalent to the sequence converging \hyperref[prob:4-34]{Problem 4.34}, treating the sequence as a net), $\prod X_k$ must be sequentially compact too.  
\end{solution}
\begin{problem}\label{prob:4-61}
Theorem 4.43 remains valid for maps from a compact Hausdorff space $X$ into a complete metric space $Y$ provided the hypothesis of pointwise boundedness is replaced by pointwise total boundedness. (Make this statement precise and then prove it.)
\end{problem}
\begin{solution}
The precise statement should be : Let $X$ be a compact Hausdorff space. 
If $\mathcal F$ is an equicontinuous, pointwise totally bounded (i.e. for all $x\in X$ the set $ \qty{f(x):f\in\mathcal F}$ is totally bounded in $Y$) subset of $C(X,Y)$, then $\mathcal F$ is totally bounded in the uniform metric, and the closure of $\mathcal F$ in $C(X,Y)$ is compact. \\
This is proved in the exact same way with these minor changes : After picking the $x_j$ as we did in the book, we can pick the $z_k$ as such because a finite union of totally bounded sets $ \qty{f(x_j):f\in\mathcal F,1\leq j\leq n}=\cup_{i=1}^n \qty{f(x_i):f\in\mathcal F}$ is totally bounded and $C(X,Y)$ is complete by \hyperref[prob:4-22]{Problem 4.22}.  
\end{solution}
\begin{problem}\label{prob:4-62}
Rephrase Theorem 4.44 in a form similar to Theorem 4.43 by using the metric in Exercise 56d.
\end{problem}
\begin{solution}
Let $X$ be a $\sigma$-compact LCH space. 
If $\mathcal F \subseteq C(X)$ is pointwise bounded and equicontniuous then it is precompact with respect to the metric in \hyperref[prob:4-56]{Problem 4.56}d. 
\end{solution}
\begin{problem}\label{prob:4-63}
Let $K \in C([0,1] \times [0,1])$. For $f \in C([0,1])$, let 
\[
    Tf(x) = \int_0^1 K(x,y)f(y)\,\mathrm{d}y
\] 
Then $Tf \in C([0,1])$, and $\{Tf : \|f\|_u \le 1\}$ is precompact in $C([0,1])$.
\end{problem}
\begin{solution}
    As $K$ is continuous on a compact set it is uniformly continuous, if  $\norm{u-v}_{\bR^2}<\delta$ implies $|K(u)-K(v)|<\epsilon,$ then we see that for $|x-y|<\delta$ 
\[
    |Tf(x)-Tf(y)|\leq\norm{f}_u\cdot\qty(\int_0^1|K(x,t)-K(y,t)|dt)\leq\epsilon\norm{f}_u
\]
as $\norm{(x,t)-(y,t)}_{\bR^2}<\delta$ too and $f$ is bounded, being continuous on a compact set.  
This proves that $Tf\in C([0,1]).$ 
Now, as $[0,1]$ is compact hausdorff, $ \qty{f\in C([0,1]):\norm{f}_u\leq 1}$ is equicontinuous as $|Tf(x)-Tf(y)|\leq\epsilon\norm{f}_u\leq\epsilon$ for $|x-y|<\delta$ as shown above and pointwise bounded due to $K$ being bounded (it is continuous on a compact set) say $|K|\leq M$, then $|Tf(x)|\leq M\norm{f}_u\leq M.$ 
Thus, by Theorem 4.43 we see that it is precompact. 
\end{solution}
\begin{problem}\label{prob:4-64}
Let $(X,\rho)$ be a metric space. A function $f \in C(X)$ is called H\"older continuous of exponent $\alpha$ ($\alpha > 0$) if the quantity
\[ N_\alpha(f) = \sup_{x \neq y} \frac{|f(x) - f(y)|}{\rho(x,y)^\alpha} \]
is finite. If $X$ is compact, $\{f \in C(X) : \|f\|_u \le 1 \text{ and } N_\alpha(f) \le 1\}$ is compact in $C(X)$.
\end{problem}
\begin{solution}
    This set is clearly closed as if we have $N_\alpha(f_n)\leq 1$ where $f_n\in C(X)$ and $\norm{f_n}_u\leq 1$ then $f_n\to f$ implies $\norm{f}\leq 1$ trivially and as $|f_n(x)-f_n(y)|\leq N_\alpha(f_n)\rho(x,y)^\alpha\leq\rho(x,y)^\alpha$ we have $|f(x)-f(y)|\leq\rho(x,y)^\alpha$ so $N_\alpha(f)\leq 1$ by passing to limits. 
    Now, it is equicontinuous as all functions $f$ in it follow $|f(x)-f(y)|\leq\rho(x,y)^\alpha$ and pointwise bounded due to each member having uniform norm atmost $1.$ 
    As metric spaces are inherently hausdorff we see that the closure of this set is compact by Theorem 4.43, which is exactly the set itself as it was closed already. 
\end{solution}
\begin{problem}\label{prob:4-65}
Let $U$ be an open subset of $\bC$, and let $\{f_n\}$ be a sequence of holomorphic functions on $U$. If $\{f_n\}$ is uniformly bounded on compact subsets of $U$, there is a subsequence that converges uniformly to a holomorphic function on compact subsets of $U$. (Use the Cauchy integral formula to obtain equicontinuity.)
\end{problem}
\begin{solution}
Lemma 2.43 shows $\sigma$-compactness and $\bC$ is clearly LCH. 
Now we will show equicontinuity of $ \qty{f_n}$ at each $x\in U.$ 
We can find a nonempty compact disc $D=\qty{y:|y-x|\leq r} \subseteq U$ and cauchy's integral formula tells us that for all $y\in D^\circ$, (the boundary is assumed to be oriented counterclockwise)
\[
    f_n(y)=\oint_{\partial D} \frac{f_n(\zeta)}{\zeta-y}d\zeta
\]
By hypothesis, $|f_n|\leq M$ on $D$ for some $M.$ 
Then, we have 
\[
    |f_n(y)-f_n(x)|\leq M\oint_{\partial D} \frac{|x-y|}{|(\zeta-x)(\zeta-y)|}d\zeta\leq\frac{2\pi r M|x-y|}{r(r-|x-y|)}
\]
proving equicontinuity (one can also use cauchy's estimates to do this.) 
Now we can use Theorem 4.44 to find a uniformly converging (on compact sets) subsequence say $f_{n_k}\to f.$ 
Holomorphicity of $f$ can be easily shown using morera's theorem and the fact that we can swap the limit and integration by the virtue of uniform convergence. 

\end{solution}
\begin{problem}\label{prob:4-66}
Let $1 - \sum_{1}^\infty c_n t^n$ be the Maclaurin series for $(1-t)^{1/2}$.
\begin{enumerate}[label=\alph*.]
    \item The series converges absolutely and uniformly on compact subsets of $(-1,1)$, as does the termwise differentiated series $-\sum_{1}^\infty n c_n t^{n-1}$. Thus, if $f(t) = 1 - \sum_{1}^\infty c_n t^n$ then $f'(t) = -\sum_{1}^\infty n c_n t^{n-1}$.
    \item By explicit calculation, $f(t) = -2(1-t)f'(t)$, from which it follows that $(1-t)^{-1/2}f(t)$ is constant. Since $f(0) = 1$, $f(t) = (1-t)^{1/2}$.
\end{enumerate}
\end{problem}
\begin{solution}
\begin{enumerate}[label=\alph*.]
    \item $c_n=(-1/n!)\cdot\prod_{k=1}^n(2k-3)/2$ and thus $|c_{n+1}/c_n|=(2n-1)/2(n+1)<1$ so the series converges absolutely and uniformly on compact subsets by the weierstrass M-test (the compact subsets $K$ are contained in $[-R,R]$ for some $|R|<1$, on which we have absolute and uniform convergence.) 
        It is a standard theorem in analysis that the same holds for the termwise differentiated series and this is actually the derivative. 
    \item The calculation is easily done, and if this is true then we see that $(1-t)^{-1/2}f(t)$ has derivative $(1-t)^{-1/2}f'(t)+(1-t)^{-3/2}f(t)/2=0$ from the equality in the statement, this shows that $f(t)=C\cdot(1-t)^{-1/2}$ for some $C$ and as $f(0)=1$ we must have that $C=1$ proving the claim. 
\end{enumerate}
\end{solution}
\begin{problem}\label{prob:4-67}
Prove Theorem 4.52. (If there exists $x_0 \in X$ such that $f(x_0) = 0$ for all $f \in \mA$, let $Y$ be the one-point compactification of $X \setminus \{x_0\}$; otherwise let $Y$ be the one-point compactification of $X$. Apply Proposition 4.36 and the Stone-Weierstrass theorem on $Y$.)
\end{problem}
\begin{solution}
We first prove a lemma: If $f\in C_0(X,\bR)$ then $f|_{X\setminus \qty{x}}\in C_0(X\setminus \qty{x},\bR)$ iff $f(x)=0.$ 
If $f$ is continuous and $|f(x)|>0$, say $|f(x)|>\epsilon>0$ then in any neighborhood of $x$ in $X$ we can find a point where $|f|>\epsilon/2$ using continuity, so $x$ is a cluster point of $ \qty{y\in X\setminus \qty{x}:|f|>\epsilon/2}$ that is not in it and thus it is not closed in $X$ and thus not compact as $X$ is hausdorff and thus it can not be compact in the open set $X\setminus \qty{x}$ as well proving one direction. 
The other direction is trivial. \\
If there exists a $x_0\in X$ such that $f(x_0)=0$ for all $f\in\mA$, as $X\setminus \qty{x_0}$ (this is open) is still LCH, let $Y$ be as in the hint. 
Proposition 4.36 tells us that we can extend all $f|_{X\setminus \qty{x_0}}\in\mA|_{X\setminus \qty{x_0}} \subseteq C_0(X\setminus \qty{x_0},\bR)$ (we can do this by the lemma) continuously to functions in $C(Y,\bR)$ by setting them to be $0$ at $\infty$ and the new algebra, call it $\mA|_{X\setminus \qty{x_0}}^*$ we get still separates points, is closed under the uniform norm (can be checked easily) and by stone weierstrass we see that $\mA|_{X\setminus x_0}^* = \qty{f\in C(Y,\bR):f(\infty)=0}$ so before the extension we must have had $\mA|_{X\setminus \qty{x_0}} = C_0(X\setminus \qty{x_0},\bR)$. 
Again by the lemma we must have that $\mA = \qty{f\in C_0(X,\bR):f(x_0)=0}$ so we are done. 
Now if otherwise $\mA$ did not vanish at any point then again after extending like we did here we would get that $\mA^*=C(X\cup \qty{\infty},\bR)$ which implies that $\mA= C_0(X,\bR)$ proving the theorem. 
\end{solution}
\begin{problem}\label{prob:4-68}
Let $X$ and $Y$ be compact Hausdorff spaces. The algebra generated by functions of the form $f(x,y) = g(x)h(y)$, where $g \in C(X)$ and $h \in C(Y)$, is dense in $C(X \times Y)$.
\end{problem}
\begin{solution}
    First, to see that $f(x,y)=h(x)g(y)\in C(X\times Y)$ for $h\in C(X),g\in C(Y)$ it suffices to see that $f=\Psi\circ\lambda$ where $\Psi:\bC\times\bC\to\bC$ is the continuous map $(x,y)\mapsto xy$ and $\lambda:X\times Y\to\bC\times\bC$ is continuous the map $(x,y)\mapsto(f(x),g(y))$, this is continuous because its components $i:(x,y)\mapsto f(x)$ and $j:(x,y)\mapsto g(y)$ are too (for open $U,$ $g^{-1}(U)$ is open and $i^{-1}(U)=g^{-1}(U)\times Y$ is open in the product topology so this is continuous and so is the other component.) 
This is clearly an algebra and for any distinct $(a,b),(c,d)$ in $X\times Y$, some coordinate must differ, say wlog $a\neq c$, then by Urysohn's lemma we can find $g\in C(X)$ such that $0=g(a)\neq g(c)=1$ and the function $f(x,y)=g(x)h(y)$ where $h$ is any constant non zero function separates these two points. 
Clearly this also contains the constant function so we must have it being dense in $C(X\times Y)$ by stone weierstrass. 
\end{solution}
\begin{problem}\label{prob:4-69}
Let $A$ be a nonempty set, and let $X = [0,1]^A$. The algebra generated by the coordinate maps $\pi_\alpha : X \to [0,1]$ ($\alpha \in A$) and the constant function 1 is dense in $C(X)$.
\end{problem}
\begin{solution}
This separates points as for any unique $f,g\in X$ there is some $a\in A$ such that $f(a)\neq g(a)$ and then $\pi_a(f)\neq\pi_a(g)$, as it already contains the constant function, the algebra generated by it must also have these properties, making it dense in $C(X)$ by the stone weierstrass theorem. 
\end{solution}
\begin{problem}\label{prob:4-70}
Let $X$ be a compact Hausdorff space. An ideal in $C(X,\bR)$ is a subalgebra $I$ of $C(X,\bR)$ such that if $f \in I$ and $g \in C(X,\bR)$ then $fg \in I$.
\begin{enumerate}[label=\alph*.]
    \item If $I$ is an ideal in $C(X,\bR)$, let $h(I) = \{x \in X : f(x) = 0 \text{ for all } f \in I\}$. Then $h(I)$ is a closed subset of $X$, called the hull of $I$.
    \item If $E \subset X$, let $k(E) = \{f \in C(X,\bR) : f(x) = 0 \text{ for all } x \in E\}$. Then $k(E)$ is a closed ideal in $C(X,\bR)$, called the kernel of $E$.
    \item If $E \subset X$, then $h(k(E)) = \overline{E}$.
    \item If $I$ is an ideal in $C(X,\bR)$, then $k(h(I)) = \overline{I}$. (Hint: $k(h(I))$ may be identified with a subalgebra of $C_0(U,\bR)$ where $U = X \setminus h(I)$.)
    \item The closed subsets of $X$ are in one-to-one correspondence with the closed ideals of $C(X,\bR)$.
\end{enumerate}
\end{problem}
\begin{solution}
\begin{enumerate}[label=\alph*.]
    \item As $I \subseteq C(X,\bR)$ we know that $f^{-1}( \qty{0})$ is closed for all $f\in I$ and thus, 
    \[
        h(I)=\bigcap_{f\in I} f^{-1}( \qty{0})
    \]
    is closed too as arbitrary intersections of closed sets are closed too. 
    \item This is clearly an ideal and its closed because if $f_n(x)=0$ for all $n$ and $f_n$ converge to $f$ in the uniform metric then $f(x)=0$ too (this now works for all $x\in E$.) 
    \item If $f\in k(E)$ then $f(E)= \qty{0}$ is dense in $f( \overline{E})$ by continuity, implying $f( \overline{E})= \qty{0}$ too so $k(E) \subseteq k( \overline{E})$ and we have inclusion clearly so $h(k(E))=h(k( \overline{E}))$ and it suffices to assume that $E$ is closed. 
    Now, if $E$ is closed and $x\notin E$ then by normality of CH spaces and urysohn's lemma we can find $f\in C(X,\bR)$ such that $f\in k(E)$ but $f(x)\neq 0$ so $x\notin h( \qty{f})$ and as $h(k(E)) \subseteq h( \qty{f})$ we have $x\notin h(k(E))$ so $h(k(E)) \subseteq E$. 
    Clearly $E \subseteq h(k(E))$ and we are done. 
    \item Clearly $h(I)\supseteq h( \overline{I})$ and if $x\in h(I)$ then, for $f_n\in I$ that converge to $f$ in the uniform metric, we have $f(x)=0$ too so $h( \overline{I}) \subseteq h(I)$ as well, thus it similarly suffices to assume $I$ to be closed. 
    By part (a) we know that $h(I)$ is closed and thus $U$ as in the hint, is open and LCH in the relative topology. 
    It can be proved that if $f\in C(X,\bR)$ then $f|_{U}\in C_0(U,\bR)$ if $f(h(I))= \qty{0}$ i.e. $f\in k(h(I)).$ by just checking the definition. 
    Thus we see that $k(h(I))$ can be identified with a subalgebra (a closed one infact) of $C_0(U,\bR)$. 
    We clearly have $I \subseteq k(h(I))$ so, as $I|_{U} \subseteq k(h(I))|_U \subseteq C_0(U,\bR)$ proving $k(h(I))=I$ is equivalent to proving $C_0(U,\bR) = I|_U$. 
    $I$ was a closed subalgebra and thus $I|_U$ is one too. 
    Given any distinct $x,y\in U$, by construction we can find $f\in I$ such that $f(x)\neq 0$ and we can find $g\in C(X,[0,1])$ such that $g(x)=1,g(y)=0,g(h(I))= \qty{0}$ by tietze's lemma and as $I$ was an ideal, $fg\in I$ and $fg$ separates $x,y$ as well as being $0$ on $h(I)$ as $g\in h(I)$ so $I|_U$ clearly separates points and from the above we also see that it vanishes at no points in $U$ so by Theorem 4.52 (\hyperref[prob:4-67]{Problem 4.67}) we have $I|_U=C_0(U,\bR)$ and we have $I=k(h(I))$ so we are done. 
    \item The map $E\mapsto k(E)$ from the closed sets in $X$ to the closed ideals in $C(X,\bR)$ has $h$ as an inverse by part (c,d) so it must be a bijection.  
\end{enumerate}
\end{solution}
\begin{problem}\label{prob:4-71}
(This is a variation on the theme of Exercise 70; it does not use the Stone-Weierstrass theorem.) Let $X$ be a compact Hausdorff space, and let $M$ be the set of all nonzero algebra homomorphisms from $C(X,\bR)$ to $\bR$. Each $x \in X$ defines an element $\hat{x}$ of $M$ by $\hat{x}(f) = f(x)$.
\begin{enumerate}[label=\alph*.]
    \item If $\phi \in M$ then $\{f \in C(X,\bR) : \phi(f) = 0\}$ is a maximal proper ideal in $C(X,\bR)$.
    \item If $I$ is a proper ideal in $C(X,\bR)$, there exists $x_0 \in X$ such that $f(x_0) = 0$ for all $f \in I$. (Suppose not; construct an $f \in I$ with $f > 0$ everywhere and conclude that $1 \in I$. This requires no deep theorems.)
    \item The map $x \to \hat{x}$ is a bijection from $X$ to $M$.
    \item If $M$ is equipped with the topology of pointwise convergence, then the map $x \to \hat{x}$ is a homeomorphism from $X$ to $M$. (Since $M$ is defined purely algebraically, it follows that the topological structure of $X$ is completely determined by the algebraic structure of $C(X,\bR)$.)
\end{enumerate}
\end{problem}
\begin{solution}
\begin{enumerate}[label=\alph*.]
    \item This is clearly a proper ideal as $\phi$ is a nonzero algebra homomorphism. 
        Suppose $J$ is an ideal containing $\ker\phi$ strictly, then we can find $f\in J$ with $\phi(f)\neq 0$ and defining $g=f/\phi(f)^2\in J$ we have $1-fg\in\ker\phi \subseteq J$ so $(1-fg)+fg=1\in J$ implying that $J$ is the full ring, here $J=C(X,\bR)$ proving maximality of $\ker\phi.$ 
    \item If that wasn't the case then we can find $f_x\in I$ such that $f_x^2\ni f_x(x)^2>0$ for all $x\in X$ and by continuity, there are open neighborhoods $U_x$ of $x$ such that $f_x^2>0$ on $U_x.$ 
        The $U_x$ form an open cover so by compactness we can find a subcover $U_{x_1},\ldots,U_{x_m}$ which implies that $f:=\sum_{i=1}^mf_{x_i}^2\in I$ is a continuous function that is never $0$ and thus $1/f\in C(X,\bR)$ so as $I$ is an ideal we have $f\cdot(1/f)=1\in I$ implying $I= C(X,\bR)$ (for all $f\in C(X,\bR),1\cdot f=f\in I$.) 
    \item This is clearly an injection as for unique $x,y$ we can find $f\in C(X,\bR)$ such that $\hat{x}(f)=f(x)\neq f(y)=\hat{y}(f)$, by tietze's extension theorem due to normality of CH spaces. 
        Now, given any $\phi\in M$ as $\ker\phi$ is a maximal proper ideal by part (a), and by part (b) there exists some point $x_0$ where it $\ker\phi$ vanishes, but to preserve maximality we must have $\ker\phi= \qty{f\in C(X,\bR):f(x_0)=0}$ and thus as $(x\mapsto f(x)-f(x_0))\in\ker\phi$ we see that $\phi(f-f(x_0)+f(x_0))=\phi(f(x_0))=f(x_0)\phi(1)=\hat{x_0}(f)$ as $\phi(1)=1$ (otherwise it would be zero everywhere.) so we see that $x\mapsto\hat{x}$ is also a surjection, proving that it is indeed a bijection.  
    \item Call this map $e$. 
        Then we see that the open sets in $M$ are generated from the base
        \[
            \qty{\hat{x}\in M:\hat{x}(f_j)=f_j(x)\in U_j\text{ for }j=1,\ldots,n}\text{ where } n\in\bN,f_j\in C(X,\bR),U_j\text{ open}
        \]
        After moving some symbols around we can see that $M$ has a topology generated by sets of form $ \qty{\hat{x}\in M:x\in f^{-1}(U)}=\widehat{f^{-1}(U)}$ where $f\in C(X,\bR)$ and $U$ is open in $X$ so $M$ is homeomorphic to $(X,w)$ where $w$ is the weak topology on $X$ generated by $C(X,\bR)$ (this is determined by the previous topology, not this one, we are not making any cyclic argument.) 
        Now, for any such $\widehat{f^{-1}(U)}$ we see that $e^{-1}(\widehat{f^{-1}(U)})=f^{-1}(U)$ is open in $X$ and as those sets generated $M$ we can say that $e$ is continuous. 
        As $X$ is compact and $M$ hausdorff and $e$ a continuous bijection we must have it being a homeomorphism as for all open $U$ in $X$, $e(U)=e(U^c)^c$ which is open as $U^c$ is closed and thus compact in $X$ implying that $e(U^c)$ is compact in $M$ and thus closed as it was hausdorff, making its complement open and thus proving continuity of $e^{-1}$ so we can say that $e$ is a homeomorphism. 
\end{enumerate}
\end{solution}
\begin{problem}\label{prob:4-72}
Every subset of a completely regular space is completely regular in the relative topology.
\end{problem}
\begin{solution}
    If $X$ is completely regular and $Y \subseteq X$ then for any closed $K=A\cap Y$ in the relative topology ($A$ is closed in $X$) and $p\in Y\cap K^c$, we see that $p\notin A$ so by hypothesis we can find $f\in C(X,[0,1])$ such that $f(p)=1$ and $f=0$ on $A$ and we see that the continuous restriction $f|_{A}$ separates $p$ and $A\cap Y$ in $Y$ so $Y$ is completely regular in the relative topology too. \\
    This proves the if direcition as $I^{C(X,I)}$ is compact hausdorff and thus it is normal so its completely regular and so is any subset with the relatitve topology and if $X$ is homeomorphic to say $Y \subseteq I^{C(X,I)}$ via $\phi\in C(X,Y)$ (also write $h=\phi^{-1}\in C(Y,X)$) then for any closed $A \subseteq Y$ and $p\notin Y$, we can find $f\in C(Y,I)$ such that $f=0$ on the closed set $h^{-1}(A)=\phi(A)$ and $f(h^{-1}(p))=f(\phi(p))=1$ and thus, $f\circ\phi\in C(X,I)$ separates $p$ and $A$ proving $X$ to be completely regular too as $p,A$ were arbitrary. 
\end{solution}
\begin{problem}\label{prob:4-73}
If $X$ is a completely regular space, a subalgebra $\mA$ of $BC(X)$ is called completely regular if (i) it is closed and contains the constant functions, and (ii) $\mA \cap C(X,I)$ separates points and closed sets.
\begin{enumerate}[label=\alph*.]
    \item If $(Y, e)$ is a Hausdorff compactification of $X$, $\mA_Y = \{f \circ e : f \in C(Y)\}$ is a completely regular subalgebra of $BC(X)$.
    \item If $(Y, e)$ and $(Y',e')$ are Hausdorff compactifications of $X$ such that $\mA_Y = \mA_{Y'}$, there is a homeomorphism $\phi : Y \to Y'$ such that $\phi \circ e = e'$. (Adapt the proof of Theorem 4.57, which deals with the case $Y = \beta X$.)
    \item If $(Y, e)$ is the compactification of $X$ associated to $\mathfrak{F} \subset C(X,I)$, then $\mA_Y$ is the smallest closed subalgebra (it should be completely regular subalgebra) of $BC(X)$ that contains $\mathfrak{F}$. (Use Exercise 69.)
    \item The Hausdorff compactifications of $X$ (upto homeomorphism) are in one-to-one correspondence with the completely regular subalgebras of $BC(X)$.
\end{enumerate}
\end{problem}
\begin{solution}
\begin{enumerate}[label=\alph*.]
    \item The constant functions are always continuous so they are clearly in $\mA_Y$ and by the previous problem we see that $e(X)$ is completely regular due to $Y$ being completely regular (it is normal, being CH) so $\mA_Y$ is clearly a completely regular subalgebra of $BC(X)$ (all the maps are bounded because $Y$ is compact.) 
    \item Let $F=C(Y,I),F'=C(Y',I)$, consider the usual embeddings (these embed as $Y,Y'$ are CH) $i,i'$ of $Y,Y'$ into the cubes $I^F,I^{F'}$ respectively.
        For all $g\in F'$ there exists $f_g\in F$ such that $g\circ e'=f_g\circ e$ and this $f_g$ is unique as for any similar $f_g'$ we see that $f_g'=f_g$ on $e(X)$ which is dense in $Y$ and thus $f_g'=f_g$, thus $g\mapsto f_g$ is actually a bijection.  
        Now, we define a continuous map $\Phi:I^F\to I^{F'}$ by $\pi_g(\Phi(p))=\pi_{f_g}(p)$ for $g\in F'$, with this we see that for all $x\in X$, 
        \[
            \pi_g(\Phi(i(e(x))))=\pi_{f_g}(i(e(x))=f_g(e(x))=g(e'(x))=\pi_g(i'(e'(x))
        \]
        and thus $\Phi\circ i\circ e=i'\circ e'.$ 
        As $\Phi(i(e(X))=i'(e'(X)) \subseteq\beta Y'$ and $\Phi(\beta Y)\supseteq i'(e'(X))$ we see that 
        \[
        \Phi(\beta Y) \supseteq \overline{i'(e'(X))}\supseteq i(\overline{e'(X)})=i'(Y')=\beta Y'
        \] 
        as $\Phi(\beta Y)$ is compact and hence closed in $Y'$ as $Y'$ is hausdorff. 
        W thus see that $\Phi$ maps $\beta Y$ surjectively on $\beta Y'$. 
        This can be summarised in the following commutative diagram, 
       \[
        \begin{tikzcd}
	    & Y && {\beta Y} && {I^F} \\
	    X \\
	    & {Y'} && {\beta Y'} && {I^{F'}}
	    \arrow["i", from=1-2, to=1-4]
	    \arrow[hook, from=1-4, to=1-6]
	    \arrow["{\Phi|_{\beta Y}}"', from=1-4, to=3-4]
	    \arrow["\Phi"', from=1-6, to=3-6]
    	\arrow["e", from=2-1, to=1-2]
    	\arrow["{e'}"', from=2-1, to=3-2]
    	\arrow["{i'}"', from=3-2, to=3-4]
    	\arrow[hook, from=3-4, to=3-6]
        \end{tikzcd}
        \]
        Now, if we define $\Psi:I^{F'}\to I^F$ by $\pi_{f_g}(\Psi(q))=\pi_g(q)$ for $g\in F'$ (this is well defined as $g\mapsto f_g$ is a bijection.) then symmetrically we see that $\Psi\circ i'\circ e'=i\circ e$ and $\Psi(\beta Y') \subseteq Y$. 
        We also see that $\pi_g(\Phi(\Psi(q)))=\pi_{f_g}(\Psi(q))=\pi_g(q)$ and thus $\Phi\circ\Psi=Id_{I^{F'}}$, similarly we have $\Psi\circ\Phi=Id_{I^{F}}$ so $\Psi$ is the continuous inverse of $\Phi,$ making $\Phi$ a homeomorphism. 
        Thus, $\Phi|_{\beta Y}$ is also a homeomorphism into $\beta Y'$ and we can define a homeomorphism $\phi=i'^{-1}\circ\Phi|_{\beta Y}\circ i$ from $Y$ to $Y'$ for which $\phi\circ e=i'^{-1}\circ\Phi|_{\beta Y}\circ i\circ e=i^{-1}\circ i'\circ e'=e'$ and we are done.
    \item By the mentioned problem, the algebra generated by the maps $\pi_f:I^{\mathfrak F}\to I$ for $f\in\mathfrak F$ and constant $1$ is dense in $C(I^{\mathfrak F})$ and we can restrict this to $Y$. 
        Thus, the algebra generated by $\pi_f\circ e=f$ and $1|_Y\circ e=1|_X$ ( $\mathfrak F$ needn't contain a constant function but as completely regular algebras always do, its safe to assume we have the constants) is dense in $\mA_Y$ as right composition by continuous functions (here it is $e$) preserves uniform convergence and thus the smallest completely regular subalgebra of $BC(X)$ containing $\mathfrak F$ is $A_Y$ itself. 
    \item Consider the map sending a hausdorff compactification $(Y,e)$ to $\mA_Y$, this is injective upto homeomorphism by part (b) and surjective because for any such subalgebra $\mA$ we can consider the compactification $Y$ associated to $\mA$ and by part (c), we see that $\mA_Y \subseteq A$ by minimality and $A \subseteq \mA_Y$ by definition so we have equality, this proves the existence of an one-to-one correspondence. 
\end{enumerate}
\end{solution}
\begin{problem}\label{prob:4-74}
Consider $\bN$ (with the discrete topology) as a subset of its Stone-\v{C}ech compactification $\beta\bN$.
\begin{enumerate}[label=\alph*.]
    \item If $A$ and $B$ are disjoint subsets of $\bN$, their closures in $\beta\bN$ are disjoint. (Hint: $\chi_A \in C(\bN,I)$.)
    \item No sequence in $\bN$ converges in $\beta\bN$ unless it is eventually constant (so $\beta\bN$ is emphatically not sequentially compact).
\end{enumerate}
\end{problem}
\begin{solution}
 We identify $\bN$ with $e(\bN)$ here. 
\begin{enumerate}[label=\alph*.]
    \item Following the hint, we see that $\pi_{\chi_A}:I^{C(\bN,I)}\to I$ is a continuous map and thus so is the restriction $\pi_{\chi_A}|_{\beta\bN}$, making $\pi_{\chi_A}|_{\beta\bN}^{-1}( \qty{1})$ is a closed set that contains $e(A)$ as $\pi_{\chi_A}(e(x))=\chi_A(x)=1$ if $x\in A$ and $0$ otherwise. 
        This also implies that $e(B) \subseteq \pi_{\chi_A}|_{\beta\bN}^{-1}( \qty{0})$, so as $A,B$ are in disjoint closed sets, they have disjoint closure. 
    \item If $(x_n)$ is a sequence that is not eventually constant but converges to say $x$, then every subsequence also converges to $x$. 
    Now, as it is not eventually constant, there two cases: either there are finitely many distinct terms and thus there are atleast two $a,b\in\bN$ such that $x_n=a$ for infinitely many $n$ and $x_n=b$ for infinitely many $n$ so we can choose two subsequences $(x_{n_k}),(x_{m_k})$ of $(x_n)$ which are all $a$'s and all $b$'s and by part (a), $ \overline{ \qty{a}}$ and $ \overline{ \qty{b}}$ are disjoint, but by hypothesis they contain a common point $x$ which is a contradiction, and in the other case there are infinitely many terms so we can again partition this infinite set into two infinite sets and get disjoint subsequences and by the same arguement we see that no sequence in $\bN$ that is not eventually constant can converge in $\beta\bN.$ 
\end{enumerate}
\end{solution}
\begin{problem}\label{prob:4-75}
Suppose $X$ is a completely regular space. The set $M$ of nonzero algebra homomorphisms from $BC(X,\bR)$ to $\bR$, equipped with the topology of pointwise convergence, is homeomorphic to $\beta X$. (See Exercise 71. This realization of $\beta X$ is the natural one from the point of view of Banach algebra theory.)
\end{problem}
\begin{solution}
    By the mentione problem we see that $\beta X$ is homeomorphic to $M_\beta$, the set of nonzero algebra homomorphisms from $C(\beta X,\bR)=BC(\beta X,\bR)$ (due to compactness) to $\bR$ under the topology of pointwise convergence. 
    It thus suffices to prove that $M_\beta$ is homeomorphic to $M,$ the same but with $BC(X,\bR).$ 
    First of all, by the universal property of stone cech compactification, the map $\phi:f\mapsto f\circ e$ is an algebra isomorphism from $BC(\beta X,\bR)$ to $BC(X,\bR).$ 
    Now, define $\Phi:M\to M_\beta$ by $u\mapsto u\circ \phi$ and $\Psi:M_\beta\to M$ by $v\to\phi^{-1}\circ v$, then we see that $\Psi$ is the inverse of $\Phi$ and thus both are bijections. 
    To prove continuity of $\Phi$, consider any net $\langle h_\alpha\rangle_{\alpha\in A}$ converging to $h$ in $M$, equivalently by \hyperref[prob:4-34]{Problem 4.34} $\pi_f(h_\alpha)=h_\alpha(f)\to h(f)=\pi_f(h)$ for all $f\in BC(X,\bR)$ and thus for any $g\in BC(\beta X,\bR)$ we have $\Phi(h_\alpha)(g)=h_\alpha(g\circ e)\to h(g\circ e)=\Phi(h)(g)$ as $g\circ e\in BC(X,\bR)$ proving that $\Phi$ is continuous and similarly if $k_\alpha\to k$ in $M_\beta$ then for all $g\in M_\beta$ we must have $k_\alpha(g)\to k(g)$ and thus for all $f\in BC(X,\bR)$ we have $\Psi(k_\alpha)(f)=k_\alpha(\phi^{-1}(f))\to\phi^{-1}(f)=\Psi(k)(f)$ as $\phi^{-1}(f)\in BC(\beta X,\bR)$ proving continuity if $\Psi=\Phi^{-1}$ hence proving that $M_{\beta X}\cong M$ so we are done. 
\end{solution}
\begin{problem}\label{prob:4-76}
If $X$ is normal and second countable, there is a countable family $\mathfrak{F} \subset C(X,I)$ that separates points and closed sets. (Let $\mathcal{B}$ be a countable base for the topology. Consider the set of pairs $(U,V) \in \mathcal{B} \times \mathcal{B}$ such that $\overline{U} \subset V$ and use Urysohn's lemma.)
\end{problem}
\begin{solution}
    Following the hint, we can find a family of functions $f_{U,V}\in C(X,[0,1])$ such that $f=1$ on $ \overline{U}$ and $0$ on $V^c$. 
    Now, if $A$ is some closed set and $p\notin A$ then we can find some $U\in\mathcal B$ such that $p\in U \subseteq A^c$ and as the space is normal we can find disjoint open sets $P,Q$ such that $p\in P$ and $U^c \subseteq Q$ so we can again find $V\in\mathcal B$ such that $p\in V \subseteq P$ and thus $V \subseteq\overline{V} \subseteq Q^c \subseteq U$ and we see that $f_{V,U}$ separates $p,A$ as $f(p)=1$ and $f(A)=0$ from the fact that $p\in V, A \subseteq U^c$ by construction. 
\end{solution}
\begin{problem}\label{prob:4-77}
Let $\{(X_n,\rho_n)\}_{1}^\infty$ be a countable family of metric spaces whose metrics take values in $[0,1]$. (The latter restriction can always be satisfied; see Exercise 56b.) Let $X = \prod_{n=1}^\infty X_n$. If $x, y \in X$, say $x = (x_1, x_2, \dots)$ and $y = (y_1, y_2, \dots)$, define $\rho(x,y) = \sum_{1}^\infty 2^{-n} \rho_n(x_n, y_n)$. Then $\rho$ is a metric that defines the product topology on $X$.
\end{problem}
\begin{solution}
    It is easy to see that $\rho$ is a valid metric as symmetry and the triangle inequality follow term-wise from $\rho_n$, and $\rho(x,y)=0$ iff $\rho_n(x_n,y_n)=0$ for all $n$ iff $x=y$.
    The series converges as $\rho_n\leq 1$ means it is bounded by $\sum_1^\infty 2^{-n}=1$.
    We use the same ball logic as in \hyperref[prob:4-56]{Problem 4.56} to show the topologies are equivalent.
    If $U$ is a basic open set in the product topology containing $x$, it restricts only finitely many coordinates, say $n_1,\ldots,n_k$, so there is some $\epsilon>0$ such that $y\in U$ if $\rho_{n_i}(x_{n_i},y_{n_i})<\epsilon$ for all $1\leq i\leq k$.
    If we choose $\delta=2^{-\max n_i}\epsilon$, then for any $y\in B_\rho(x,\delta)$ we see that $2^{-n_i}\rho_{n_i}(x_{n_i},y_{n_i})\leq\rho(x,y)<\delta$ which implies $\rho_{n_i}(x_{n_i},y_{n_i})<\epsilon$ for all the restricting coordinates, thus $y\in U$ proving $B_\rho(x,\delta) \subseteq U$.
    For the other direction, given any metric ball $B_\rho(x,\epsilon)$, choose $N$ large enough such that $2^{-N}<\epsilon/2$ (since $\sum_{n>N}2^{-n}=2^{-N}$).
    Consider the basic open set $V=\qty{y\in X:\rho_n(x_n,y_n)<\epsilon/2\text{ for }1\leq n\leq N}$, which clearly contains $x$.
    For any $y\in V$ we see that
    \[
        \rho(x,y)=\sum_{n=1}^N 2^{-n}\rho_n(x_n,y_n)+\sum_{n=N+1}^\infty 2^{-n}\rho_n(x_n,y_n)< \frac{\epsilon}{2}\sum_{n=1}^N 2^{-n} + 2^{-N} < \frac{\epsilon}{2}+\frac{\epsilon}{2}=\epsilon
    \]
    so $V \subseteq B_\rho(x,\epsilon)$ and we are done.
\end{solution}
\end{document}
